\chapter{Methods}
\section{Inference via Message Passing}
\label{sec:1}

In this section I present a method for deriving message passing algorithms that enable inference. Inference is the mechanism that allows behavioral agents (presented in later sections) to process observations, infer unobservable environmental states, perform planning, and infer useful actions. The underlying message passing algorithm rests upon a factor graph representation of the agents' generative models where messages are sent along the edges.

To focus on the core aspects of this mechanism I give the derivations including a numerical example for a simple agent that can receive observations and infer hidden environmental causes but it cannot plan nor act.

Concretely, the simple agent is concerned with its cat (yes, the agent has a pet). He is \(90\,\%\) certain that its cat is a sweet, loving angel. Unfortunately, the agent cannot look inside the cat's brain to really find out but can only observe her behavior. Sometimes, the cat walks straight up to a full glas of water, makes eye contact with the agent and pushes the glas from the table (breaks glas). The agent believes: If the cat were a little chaos demon this is exactly the kind of behavior I expect (with \(99\,\%\) certainty). On the other hand, if she really was a sweet loving angel then she would rarely do such cruel things (maybe \(10\,\%\) chance).

The agent's model of this situation can formally be described via a (generative) model \(p(s,o)=p(o|s)p(s)\) consisting of a prior probability \(p(s)\) about hidden causal states \(s\) (is the cat an angel or a daemon?) and a likelihood distribution \(p(o|s)\) that relates the hidden states to observations \(o\) (does the cat break glas or not?) as
\begin{alignat*}{3}
     & p(s=\text{"angel"})             &  &                    &  & = 0.9  \\
     & p(s=\text{"demon"})             &  &                    &  & = 0.1  \\
     & p(o=\text{"breaks glas"}        &  & |s=\text{"demon"}) &  & = 0.99 \\
     & p(o=\text{"doesn't break glas"} &  & |s=\text{"demon"}) &  & = 0.01 \\
     & p(o=\text{"breaks glas"}        &  & |s=\text{"angel"}) &  & = 0.1  \\
     & p(o=\text{"doesn't break glas"} &  & |s=\text{"angel"}) &  & = 0.9.
\end{alignat*}
With this probabilistic model the agent can then infer the true character of its cat upon making an observation. This inference can be computed via Bayes theorem. For example when the agent observes \(o=\text{"breaks glas"}\) it can calculate the posterior probability
\begin{equation*}
    p(s|o=\text{"breaks glas"}) = \frac{p(o=\text{"breaks glas"}|s)p(s)}{\sum_s p(o=\text{"breaks glas"}|s)p(s)},
\end{equation*}
which concretely results in
\begin{alignat*}{3}
     & p(s=\text{"angel"}|o=\text{"breaks glas"})  &  & \approx 0.48  \\
     & p(s=\text{"daemon"}|o=\text{"breaks glas"}) &  & \approx 0.52,
\end{alignat*}
i.e, the agent infers that the cat's character is quite ambiguous and he is actually slightly more certain that she is a demon (\(52\,\%\)).

This inference can be formalized as a message passing algorihm on a factor graph. Towards this, I will first give the simple agent's generative model as its (Forney) factor graph representation.

\subsection{The Forney Factor Graph Representation}
A Forney factor graph (FFG) is a neat tool to represent a factorized probabilistic model where probabilistic factors (e.g., \(p(s)\)) become nodes and variables (e.g., \(s\)) become edges. The generative model of the simple agent has two variables \(s\), and \(o\) and consists of two factors \(p(s)\) and \(p(o|s)\). This leads to an FFG that equally has two edges and two nodes as shown in Fig.~\ref{fig:1}.
\begin{figure}[htbp]
    \centering

    \begin{tikzpicture}[
            factor/.style={draw=black, thick, rectangle, minimum size=0.8cm, font=\normalsize},
            edge/.style={thick},
            varlabel/.style={font=\normalsize}
        ]

        % Knoten
        \node[factor] (a) {$p(s)$};
        \node[factor, right=1.5cm of a] (b) {$p(o|s)$};

        % Kanten
        \draw[edge] (a.east) -- (b.west) node[midway, above, varlabel] {$s$};
        \draw[edge] (b.east) -- ++(1.5cm, 0) node[midway, above, varlabel] {$o$};

    \end{tikzpicture}

    \caption{FFG of the model \(p(s,o)\).}
    \label{fig:1}
\end{figure}

In this work we utilize a version of an FFG that requires every edge to be terminated by a node on both ends, resulting in a (terminated) Forney factor graph. To this end, the model \(p(s,o)\) can be extended by a factor of 1 resulting in \(p(s,o)=p(s)p(o|s)\cdot1\) with the (terminated) FFG shown in Fig.~\ref{fig:2}
\begin{figure}[htbp]
    \centering

    \begin{tikzpicture}[
            factor/.style={draw=black, thick, rectangle, minimum size=0.8cm, font=\normalsize},
            edge/.style={thick},
            varlabel/.style={font=\normalsize}
        ]

        % Knoten
        \node[factor] (a) {$p(s)$};
        \node[factor, right=1.5cm of a] (b) {$p(o|s)$};
        \node[factor, right=1.5cm of b] (c) {$1$};

        % Kanten
        \draw[edge] (a.east) -- (b.west) node[midway, above, varlabel] {$s$};
        \draw[edge] (b.east) -- (c.west) node[midway, above, varlabel] {$o$};

    \end{tikzpicture}

    \caption{FFG of the model \(p(s,o)\) with terminated edges.}
    \label{fig:2}
\end{figure}

To facilitate future derivations we rename node factors with functions \(f\) that have indices as unique identifiers. Following this shift, the three nodes become \(p(s)=f_a(s)\), \(p(o|s)=f_b(s,o)\), and \(1=f_c(o)\) resulting in the FFG in Fig.~\ref{fig:3}. Note that \(f_c\) has an argument \(o\) because it is connected to the edge \(o\) but always results into \(1\).
\begin{figure}[htbp]
    \centering

    \begin{tikzpicture}[
            factor/.style={draw=black, thick, rectangle, minimum size=0.8cm, font=\normalsize},
            edge/.style={thick},
            varlabel/.style={font=\normalsize}
        ]

        % Knoten
        \node[factor] (a) {$f_a(s)$};
        \node[factor, right=1.5cm of a] (b) {$f_b(s,o)$};
        \node[factor, right=1.5cm of b] (c) {$f_c(o)$};

        % Kanten
        \draw[edge] (a.east) -- (b.west) node[midway, above, varlabel] {$s$};
        \draw[edge] (b.east) -- (c.west) node[midway, above, varlabel] {$o$};

    \end{tikzpicture}

    \caption{FFG of the model \(p(s,o)\) with renamed factors (\(f_a(s)\), \(f_b(s,o)\), \(f_c(o)\)).}
    \label{fig:3}
\end{figure}

\subsection{Belief Distributions}

In the numerical example I conditioned the generative model on a specific observation ("breaks glas") and computed the posterior distribution(s) of the remaining variables (only \(s\)). This conditioning on specific events is the "driving force" during inference that effectively "pushes" the probability mass of the remaining variables to its respective posterior forms. These form-flexible (variational) distributions exist for all variables (and factors) and effectively represent the agent's beliefs about these variables (and factors). Concretely, these belief distributions are named \(q\) and exist for all edges and nodes (for all variables and factors) in the graph (\cref{fig:4}).
\begin{figure}[htbp]
    \centering

    \begin{tikzpicture}[
            factor/.style={draw=black, thick, rectangle, minimum size=0.8cm, font=\normalsize},
            edge/.style={thick},
            varlabel/.style={font=\normalsize}
        ]

        % Knoten
        \node[factor, label={[varlabel]below:$q_a(s)$}] (a) {$f_a(s)$};
        \node[factor, label={[varlabel]below:$q_b(s,o)$}, right=1.5cm of a] (b) {$f_b(s,o)$};
        \node[factor, label={[varlabel]below:$q_c(o)$}, right=1.5cm of b] (c) {$f_c(p)$};

        % Kanten
        \draw[edge, ] (a.east) -- (b.west) node[midway, above, varlabel] {$s$} node[midway, below, varlabel] {$q_i(s)$};
        \draw[edge] (b.east) -- (c.west) node[midway, above, varlabel] {$o$} node[midway, below, varlabel] {$q_j(o)$};

    \end{tikzpicture}

    \caption{FFG of the model \(p(s,o)\) with belief distributions \(q\).}
    \label{fig:4}
\end{figure}
Relating to the numerical example, the conditioning of the generative model on the observation ("breaks glas") is realized by "clamping" the respective belief distribution (\(q_j(o)\)) to this event, i.e., it is set to delta distribution that has probability \(1\) for the event "breaks glas". This clamping exerts a kind of "force" on all other belief distributions and "pushes" them to their posterior forms. For example, after inference the belief \(q_i(s)\) will assume the posterior \(p(s|o=\text{"breaks glas"})\). In summary, we can now represent the agent's generative model as a factor graph and understand that it maintains belief distributions \(q\) that are form-flexible and "react" to new evidence (e.g., an observation) anywhere in the graph to form posteriors

\subsection{The (Bethe) Free Energy Functional}

As a next step towards the message passing procedere that drives inference (which forms the \(q\) distributions), it is important to realize that the metaphorical force that exists between the belief distributions (and the factors of the generative model) can exactly be defined as a variational free energy \(F[q]\). It is a functional of the belief distributions (indicated by \(q\)) and is defined as the Kullback-Leibler divergence between the joint belief distribution \(q(s,o)\) and the generative model \(p(s,o)\) with
\begin{equation}
    F[q] = D_{\mathrm{KL}}(q(s,o) \parallel p(s,o)) = \sum_{s,o} q(s,o) \log \frac{q(s,o)}{p(s,o)},
\end{equation}
where the joint belief distribution is the product of all node beliefs divided by the product of all edge beliefs, i.e., \(q(s,o)=\frac{q_a(s)q_b(s,o)q_c(o)}{q_i(s)q_j(o)}\). In this case \(F\) is equal to the Bethe free energy. Note that we assumed discrete distributions (e.g., categorical), leading to a sum instead of an integral.

For the message passing algorithm we must identify local messages that are sent on the edges of the factor graph. These local messages are driven by local free energies, which are part of the global free energy \(F\). To facilitate the derivations of these local messages, it is convenient to partition \(F\) into a sum of these local free energies. Following this, \(F\) can be rearranged to
\begin{equation}
    \begin{split}
        F[q] = & \underbrace{\sum_s q_a(s) \log \frac{q_a(s)}{f_a(s)} + \sum_{s,o} q_b(s,o) \log \frac{q_b(s,o)}{f_b(s,o)} + \sum_o q_c(o) \log \frac{q_c(o)}{f_c(o)}}_{\text{Local Free Energy Terms}} \\
               & + \underbrace{\sum_s q_i(s) \log \frac{1}{q_i(s)} + \sum_o q_j(o) \log \frac{1}{q_j(o)}}_{\text{Correcting Entropy Terms}}.
    \end{split}
\end{equation}
This yields three node-local free energy terms that measure the Kullback-Leibler divergence between the local belief distribution \(q\) and the local factor \(f\). Importantly, these local free energies contain entropy terms that correspond to the edges' entropies they are connected to. As an example consider the first local free energy term \(\sum_s q_a(s) \log \frac{q_a(s)}{f_a(s)}\) which can be decomposed into \(- (-\sum_s q_a(s) \log q_a(s)) - \sum_s q_a(s) \log f_a(s)\). The term \(- (-\sum_s q_a(s) \log q_a(s))\) is the entropy of edge \(i\), i.e., it is equivalent to \(- (-\sum_s q_i(s) \log q_i(s))\). Similarly node \(b\) also contains this entropy of edge \(i\) as part of \(- (-\sum_{s,o} q_b(s,o) \log q_b(s,o))\). Consequently, having two nodes for every edge results into overcounting of edge entropies, i.e., it is subtracted twice. To correct for this, one entropy term is added for each edge in \(F\), e.g., \(\sum_s q_i(s) \log \frac{1}{q_i(s)}\). This corrective measure is not a deliberative choice but a natural consequence of defining belief distributions over edges and nodes in the graph. This circumstance makes the variational free energy \(F\) equal to the Bethe free energy.

\subsection{The Lagrangian of the Free Energy}

Now, that I have defined the "force" that couples the belief distributions with the factors of the generative model, i.e., \(F\), and how it can be decomposed to a sum of local terms, it next becomes necessary to understand that the belief distributions \(q\) are freely optimizable functions that do not automatically satisfy the requirements of probability distributions. Concretely, first they must meet the normalization constraints

\noindent
\begin{minipage}[t]{0.5\textwidth}
    \setlength{\abovedisplayskip}{0pt}
    \setlength{\belowdisplayskip}{0pt}
    \begin{align}
        \sum_s q_a(s)       & \overset{!}{=} 1 \\[0.1em]
        \sum_{s,o} q_b(s,o) & \overset{!}{=} 1 \\
        \sum_o q_c(o)       & \overset{!}{=} 1
    \end{align}
\end{minipage}\hfill
\begin{minipage}[t]{0.5\textwidth}
    \setlength{\abovedisplayskip}{0pt}
    \setlength{\belowdisplayskip}{0pt}
    \begin{align}
        \sum_s q_i(s) & \overset{!}{=} 1  \\
        \sum_o q_j(o) & \overset{!}{=} 1,
    \end{align}
\end{minipage}
and second they must meet marginalization constraints with respect to neighboring belief distributions in the factor graph with

\noindent
\begin{minipage}[t]{0.5\textwidth}
    \setlength{\abovedisplayskip}{0pt}
    \setlength{\belowdisplayskip}{0pt}
    \begin{align}
        q_i(s) & \overset{!}{=} q_a(s)          \\[0.8em]
        q_i(s) & \overset{!}{=} \sum_y q_b(s,o)
    \end{align}
\end{minipage}\hfill
\begin{minipage}[t]{0.5\textwidth}
    \setlength{\abovedisplayskip}{0pt}
    \setlength{\belowdisplayskip}{0pt}
    \begin{align}
        q_j(o) & \overset{!}{=} \sum_x q_b(s,o) \\
        q_j(o) & \overset{!}{=} q_c(o).
    \end{align}
\end{minipage}

These constraints are then added to \(F\) and lead to the Lagrangian
\begin{equation}
    \begin{split}
        L = F & + \sum_s \lambda_{ia}(s) (q_i(s) - q_a(s)) + \sum_s \lambda_{ib}(s) (q_i(s) - \sum_y q_b(s,o)) \\
              & + \sum_o \lambda_{jb}(o) (q_j(o) - \sum_s q_b(s,o)) - \sum_o \lambda_{jc}(o) (q_j(o) - q_c(o)) \\
              & + \psi_a (\sum_s q_a(s) - 1) + \psi_b (\sum_{s,o} q_b(s,o) - 1) + \psi_c (\sum_o q_c(o) - 1)   \\
              & + \psi_i (\sum_s q_i(s) - 1) + \psi_j (\sum_o q_j(o) - 1),
    \end{split}
    \label{eq:36}
\end{equation}
where \(\lambda\) und \(\psi\) are Lagrange multipliers that enforce marginalization and normalization respectively.

\subsection{Messages and Belief Updates}

Having defined the Lagrangian of the factor graph it is now possible to derive the messages that are necessary for the final belief updates. Recall, that the message passing algorithm aims at finding belief distributions \(q\) by minimizing \(F\) (which is part of \(L\)), i.e. the beliefs "relax" to their solutions driven by the metaphorical force \(F\). This "relaxation" process is realized by the message passing in the graph. Following this, \(L\) is partially differentiated with respect to each belief (variational) distribution \(q\), then variational derivatives are set to \(0\), and solved for \(q^*\):
\begin{alignat}{3}
     & \frac{\delta L}{\delta q_a(s)} &  & = \log q_a(s) + 1 - \log f_a(s) - \lambda_{ia}(s) + \psi_a             & \overset{!}{=} 0 \\
     &                                &  & \Rightarrow q_a^*(s) = f_a(s) \exp(\lambda_{ia}(s)) \exp(-\psi_a - 1). & \label{eq:6}
\end{alignat}
By a similar argument it follows that
\begin{alignat}{3}
     & \Rightarrow q_b^*(s,o) &  & = f_b(s,o) \exp(\lambda_{ib}(s)) \exp(\lambda_{jb}(o)) \exp(-\psi_b - 1)                & \label{eq:7} \\
     & \Rightarrow q_c^*(o)   &  & = f_c(o) \exp(\lambda_{jc}(o)) \exp(-\psi_c - 1)                           \label{eq:8}                \\
     & \Rightarrow q_i^*(s)   &  & = \exp(\lambda_{ia}(s)) \exp(\lambda_{ib}(s)) \exp(\psi_i - 1)             \label{eq:9}                \\
     & \Rightarrow q_j^*(o)   &  & = \exp(\lambda_{jb}(o)) \exp(\lambda_{jc}(o)) \exp(\psi_j - 1).\label{eq:10}
\end{alignat}
Second, we partially solve the system of equations (\cref{eq:6,eq:7,eq:8,eq:9,eq:10}) with respect to the normalization constraints with
\begin{alignat}{2}
     & \sum_s q_a^*(s)                                       &  & \overset{!}{=} 1                                                             \\
     & \sum_s f_a(s) \exp(\lambda_{ia}(s)) \exp(-\psi_a - 1) &  & = 1                                                                          \\
     & \exp(-\psi_a - 1) \sum_x f_a(s) \exp(\lambda_{ia}(s)) &  & = 1                                                                          \\
     & \exp(-\psi_a - 1)                                     &  & = \frac{1}{\sum_s f_a(s) \exp(\lambda_{ia}(s))}                              \\
     & \Rightarrow q_a^*(s)                                  &  & = f_a(s) \exp(\lambda_{ia}(s)) \frac{1}{\sum_s f_a(s) \exp(\lambda_{ia}(s))} \\
     & \Rightarrow q_a^*(s)                                  &  & \propto f_a(s) \exp(\lambda_{ia}(s)). \label{eq:1}
\end{alignat}
The term \(\exp(-\psi_a - 1)\) effectively becomes the normalization constant for \(q_a^*(s)\). Instead of doing this derivation repeatedly in subsequent sections, I will directly use the proportional sign \(\propto\) when defining belief distributions. More generally, in some cases I will also use "\(\propto\)" even if the relation allows for the stricter "\(=\)" sign.
Analogously we obtain
\begin{alignat}{2}
     & \Rightarrow q_b^*(s,o) &  & \propto f_b(s,o) \exp(\lambda_{ib}(s)) \exp(\lambda_{jb}(o)) \label{eq:2} \\
     & \Rightarrow q_c^*(o)   &  & \propto f_c(o) \exp(\lambda_{jc}(o))                         \label{eq:3} \\
     & \Rightarrow q_i^*(s)   &  & \propto \exp(\lambda_{ia}(s)) \exp(\lambda_{ib}(s))          \label{eq:4} \\
     & \Rightarrow q_j^*(o)   &  & \propto \exp(\lambda_{jb}(o)) \exp(\lambda_{jc}(o)).\label{eq:5}
\end{alignat}
Third, we solve the remaining system of equations (\cref{eq:1,eq:2,eq:3,eq:4,eq:5}) with respect to the marginalization constraints with
\begin{alignat}{2}
     & q_a^*(s)                          &  & \overset{!}{=} q_i(s)                               \\
     & f_a(s) \exp(\lambda_{ia}(s))      &  & \propto \exp(\lambda_{ia}(s)) \exp(\lambda_{ib}(s)) \\
     & \Rightarrow \exp(\lambda_{ib}(s)) &  & \propto f_a(s). \label{eq:26}
\end{alignat}
Similarly, it can be shown that
\begin{alignat}{2}
     & \Rightarrow \exp(\lambda_{ia}(s)) &  & \propto \sum_o f_b(s,o) \exp(\lambda_{jb}(o)) \label{eq:27}    \\
     & \Rightarrow \exp(\lambda_{jc}(o)) &  & \propto \sum_s f_b(s,o) \exp(\lambda_{ib}(s))    \label{eq:28} \\
     & \Rightarrow \exp(\lambda_{jb}(o)) &  & \propto f_c(o). \label{eq:29}
\end{alignat}
Crucially, the derivations do not only yield the messages but also how they can be used to finally compute the belief distributions \(q\). Concretely, the final beliefs are given in \cref{eq:2,eq:3,eq:4,eq:5} and thus depend on the messages \(\exp(\lambda_{ib}(s))\), \(\exp(\lambda_{jb}(o))\), \(\exp(\lambda_{jc}(o))\), and \(\exp(\lambda_{ia}(s))\). Moreover, the definitions of these four messages are given in \cref{eq:26,eq:27,eq:28,eq:29}.

\subsection{Messages in the Factor Graph}
Having derived the four messages above, we can now view them inside the factor graph.
Their names give a clear interpretation of their position and direction of information flow in the graph. For example, the message \(\exp(\lambda_{ia}(s))\) has indices \(i\) and \(a\), meaning that its location in the FFG is on the edge \(i\) and points towards node \(a\). The resulting graph with all messages is given in Fig.~\ref{fig:6}.
\begin{figure}[htbp]
    \centering

    \begin{tikzpicture}[
            factor/.style={draw=black, thick, rectangle, minimum size=0.8cm, font=\normalsize},
            edge/.style={thick},
            varlabel/.style={font=\normalsize},
            greycircle/.style={draw=black, thick, circle, fill=gray!30, minimum size=0.5cm, font=\normalsize}
        ]

        % Knoten
        \node[factor, label={[varlabel]below:$q_a(s)$}] (a) {$f_a(s)$};
        \node[factor, label={[varlabel]below:$q_b(s,o)$}, right=6cm of a] (b) {$f_b(s,o)$};
        \node[factor, label={[varlabel]below:$q_c(o)$}, right=6cm of b] (c) {$f_c(o)$};

        % Kanten
        \draw[edge] (a.east) -- (b.west) node[midway, above, varlabel] {$s$} node[midway, below, varlabel] {$q_i(s)$} node[pos=0.2, below, varlabel] {$\exp(\lambda_{ia}(s))$}  node[pos=0.8, below, varlabel] {$\exp(\lambda_{ib}(s))$};
        \draw[edge] (b.east) -- (c.west) node[midway, above, varlabel] {$o$} node[midway, below, varlabel] {$q_j(o)$} node[pos=0.2, below, varlabel] {$\exp(\lambda_{jb}(o))$}  node[pos=0.8, below, varlabel] {$\exp(\lambda_{jc}(o))$};

    \end{tikzpicture}

    \caption{FFG of the model \(p(s,o)\) with messages.}
    \label{fig:6}
\end{figure}

Using the full message names renders the FFG visually overloaded.
Instead we use abbreviations (\(\textcolor{black}{\bm{\overleftarrow{\mu}}},\textcolor{black}{\bm{\overrightarrow{\mu}}}\)) with arrows clearly indicating the direction of information flow with \(\textcolor{black}{\bm{\overleftarrow{\mu}(}s\bm{)}} = \exp(\lambda_{ia}(s))\), \(\textcolor{black}{\bm{\overrightarrow{\mu}(}s\bm{)}} = \exp(\lambda_{ib}(s))\), \(\textcolor{black}{\bm{\overleftarrow{\mu}(}o\bm{)}} = \exp(\lambda_{jb}(o))\), and \(\textcolor{black}{\bm{\overrightarrow{\mu}(}o\bm{)}} = \exp(\lambda_{jc}(o))\).
The resulting FFG is given in Fig.~\ref{fig:7}.
\begin{figure}[htbp]
    \centering

    \begin{tikzpicture}[
            factor/.style={draw=black, thick, rectangle, minimum size=0.8cm, font=\normalsize},
            edge/.style={thick},
            varlabel/.style={font=\normalsize},
            greycircle/.style={draw=black, thick, circle, fill=gray!30, minimum size=0.5cm, font=\normalsize}
        ]

        % Knoten
        \node[factor, label={[varlabel]below:$q_a(s)$}] (a) {$f_a(s)$};
        \node[factor, label={[varlabel]below:$q_b(s,o)$}, right=4cm of a] (b) {$f_b(s,o)$};
        \node[factor, label={[varlabel]below:$q_c(o)$}, right=4cm of b] (c) {$f_c(o)$};

        % Kanten
        \draw[edge] (a.east) -- (b.west) node[midway, above, varlabel] {$s$} node[midway, below, varlabel] {$q_i(s)$} node[pos=0.15, below, varlabel] {$\textcolor{black}{\bm{\overleftarrow{\mu}(}s\bm{)}}$}  node[pos=0.85, below, varlabel] {$\textcolor{black}{\bm{\overrightarrow{\mu}(}s\bm{)}}$};
        \draw[edge] (b.east) -- (c.west) node[midway, above, varlabel] {$o$} node[midway, below, varlabel] {$q_j(o)$} node[pos=0.15, below, varlabel] {$\textcolor{black}{\bm{\overleftarrow{\mu}(}o\bm{)}}$}  node[pos=0.85, below, varlabel] {$\textcolor{black}{\bm{\overrightarrow{\mu}(}o\bm{)}}$};

    \end{tikzpicture}

    \caption{FFG of the model \(p(s,o)\) with \(\bm{\mu}\) messages.}
    \label{fig:7}
\end{figure}
In principle, the message passing algorithm is now complete and messages can be computed to finally determine the belief distributions. Still missing is the explicit modeling of the observation in the factor graph. Without an observation, inference has no basis for computing beliefs that relate to posteriors. Next I show the integration of an observation in the factor graph and how that locally leads to alternative messages.

\subsection{Observation (Delta) Constraints}

Having a concrete observation \(\hat o\) for a variable \(o\) renders its corresponding (edge) belief distribution \(q_j(o)\) a delta distribution. This effectively assigns all probability mass to this single observation while other events have probability \(0\). In this work, distributions are discrete such that delta distributions are given as Kronecker deltas that evaluate to \(1\) for the respective observation event. Relating to the cat example, the delta distribution is given as \(\delta_{o,\hat o}\) with \(\hat o = \text{"breaks glas"}\), i.e., only the event "breaks glas" evaluates to \(1\), and thus defines the edge's belief distribution as \(q_j(o) = \delta_{o,\hat o}\). This delta constraint is visualized in the factor graph (\cref{fig:29}) by a gray bead at the respective edge.
\begin{figure}[htbp]
    \centering

    \begin{tikzpicture}[
            factor/.style={draw=black, thick, rectangle, minimum size=0.8cm, font=\normalsize},
            edge/.style={thick},
            varlabel/.style={font=\normalsize},
            greycircle/.style={draw=black, thick, circle, fill=gray!30, minimum size=0.5cm, font=\normalsize}
        ]

        % Knoten
        \node[factor, label={[varlabel]below:$q_a(s)$}] (a) {$f_a(s)$};
        \node[factor, label={[varlabel]below:$q_b(s,o)$}, right=4cm of a] (b) {$f_b(s,o)$};
        \node[factor, label={[varlabel]below:$q_c(o)$}, right=4cm of b] (c) {$f_c(o)$};

        % Kanten
        \draw[edge] (a.east) -- (b.west) node[midway, above, varlabel] {$s$} node[midway, below, varlabel] {$q_i(s)$} node[pos=0.15, below, varlabel] {$\textcolor{black}{\bm{\overleftarrow{\mu}(}s\bm{)}}$}  node[pos=0.85, below, varlabel] {$\textcolor{black}{\bm{\overrightarrow{\mu}(}s\bm{)}}$};
        \draw[edge] (b.east) -- (c.west) node[midway, greycircle] {$\delta$} node[midway, above=0.3, varlabel] {$o$} node[midway, below=0.3, varlabel] {$q_j(o)$} node[pos=0.15, below, varlabel] {$\textcolor{black}{\bm{\overleftarrow{\mu}(}o\bm{)}}$}  node[pos=0.85, below, varlabel] {$\textcolor{black}{\bm{\overrightarrow{\mu}(}o\bm{)}}$};

    \end{tikzpicture}

    \caption{FFG of the model \(p(s,o)\) with delta constraint.}
    \label{fig:29}
\end{figure}
Importantly, introducing the constraint \(q_j(o) = \delta_{o,\hat o}\) leads to alternative messages \(\textcolor{black}{\bm{\overleftarrow{\mu}(}o\bm{)}}\) and \(\textcolor{black}{\bm{\overrightarrow{\mu}(}o\bm{)}}\). These messages are derived in the following.

Starting point for this derivation are the local stationary solutions from \cref{eq:1,eq:2,eq:3,eq:4,eq:5}. During the following resolution of marginalization constraints that depend on \(q_j(o)\) (i.e., \cref{eq:28,eq:29} in prior derivations) the belief \(q_j(o)\) is substituted by \(\delta_{o,\hat o}\), which leads to:
\begin{alignat}{2}
     & \sum_s q_b^*(s,o)                                                                           &  & \overset{!}{=} q_j^*(o)                                                               \\
     & \sum_s f_b(s,o) \exp(\lambda_{ib}(s)) \exp(\lambda_{jb}(o))                                 &  & \propto \delta_{o,\hat o}                                                             \\
     & \exp(\lambda_{jb}(o)) \sum_s f_b(s,o) \exp(\lambda_{ib}(s))                                 &  & \propto \delta_{o,\hat o}                                                             \\
     & \exp(\lambda_{jb}(o))                                                                       &  & \propto \frac{\delta_{o,\hat o}}{\sum_s f_b(s,o) \exp(\lambda_{ib}(s))} \label{eq:23} \\
     & \Rightarrow \exp(\lambda_{jb}(o))              \propto \delta_{o,\hat o}                                                                                                               \\
     & \Rightarrow \textcolor{black}{\bm{\overleftarrow{\mu}(}o\bm{)}}  \propto \delta_{o,\hat o}.
\end{alignat}
Note that \(\frac{\delta_{o,\hat o}}{\sum_s f_b(s,o) \exp(\lambda_{ib}(s))}\) in \cref{eq:23} evaluates to \(\delta_{o,\hat o}\). To understand this, consider the following two cases. First, if \(o \neq \hat o\) then both \(\frac{\delta_{o,\hat o}}{\sum_s f_b(s,o) \exp(\lambda_{ib}(s))}\) and \(\delta_{o,\hat o}\) yield 0 as the result. Second, if \(o = \hat o\) then \(\delta_{o,\hat o}\) will result into 1 and \(\frac{\delta_{o,\hat o}}{\sum_s f_b(s,o) \exp(\lambda_{ib}(s))}\) will result into some positive number (not 0). For the the latter, normalization ensures that this arbitrary number will be 1 because all other outcomes already have probability 0 as outlined in the first case.

Similarly, using the constraint \(q_c^*(o) \overset{!}{=} q_j^*(o)\) we obtain
\begin{alignat}{2}
     & \Rightarrow \textcolor{black}{\bm{\overrightarrow{\mu}(}o\bm{)}}  \propto \delta_{o,\hat o}.
\end{alignat}

Having derived the messages \(\textcolor{black}{\bm{\overleftarrow{\mu}(}o\bm{)}}\) and \(\textcolor{black}{\bm{\overrightarrow{\mu}(}o\bm{)}}\) that result from the observation \(\hat o\) we can now show how this message passing algorithm achieves the inference from the cat example. The algorithm is formally given in the following section.

\subsection{Message Passing Algorithm}

As outlined before, the algorithm consists of first computing all messages and second the edges' belief distributions and is given in \cref{alg:1}.
\begin{algorithm}[htbp]
    \caption{Message Passing Algorithm of the Simple Agent}
    \label{alg:1}
    \begin{algorithmic}[1]
        \State $\textcolor{black}{\bm{\overleftarrow{\mu}(}o\bm{)}}$,\quad$\textcolor{black}{\bm{\overrightarrow{\mu}(}o\bm{)}}$,\quad$\textcolor{black}{\bm{\overrightarrow{\mu}(}s\bm{)}}$
        \State $\textcolor{black}{\bm{\overleftarrow{\mu}(}s\bm{)}}$
        \State $q_i(s)$,\quad$q_j(o)$
    \end{algorithmic}
\end{algorithm}
Note that some messages depend on the result of other messages which leads to a certain order of execution, e.g., \(\textcolor{black}{\bm{\overleftarrow{\mu}(}s\bm{)}} \propto \sum_o f_b(s,o) \textcolor{black}{\bm{\overleftarrow{\mu}(}o\bm{)}}\) depends on \(\textcolor{black}{\bm{\overleftarrow{\mu}(}o\bm{)}}\).

In the following I present the algorithm's execution using the numerical cat example.
First, the messages \(\textcolor{black}{\bm{\overleftarrow{\mu}(}o\bm{)}}\), \(\textcolor{black}{\bm{\overrightarrow{\mu}(}o\bm{)}}\), and \(\textcolor{black}{\bm{\overrightarrow{\mu}(}s\bm{)}}\) are evaluated (there is no order between the three messages because they do not depend on each other):
\begin{alignat}{2}
     & \Rightarrow \textcolor{black}{\bm{\overleftarrow{\mu}(}o\bm{)}}  &  & \propto \delta_{o,\hat o}                     \\
     & \Rightarrow \textcolor{black}{\bm{\overrightarrow{\mu}(}o\bm{)}} &  & \propto \delta_{o,\hat o}                     \\
     & \Rightarrow \textcolor{black}{\bm{\overrightarrow{\mu}(}s\bm{)}} &  & = \exp(\lambda_{ib}(s)) \propto f_a(s) = p(s)
\end{alignat}
with \(\hat o = \text{"breaks glas"}\).

Second, \(\textcolor{black}{\bm{\overleftarrow{\mu}(}s\bm{)}}\) is evaluated:
\begin{alignat}{2}
     & \Rightarrow \textcolor{black}{\bm{\overleftarrow{\mu}(}s\bm{)}} &  & \propto \sum_o f_b(s,o) \textcolor{black}{\bm{\overleftarrow{\mu}(}o\bm{)}} = \sum_o p(o|s) \delta_{o,\hat o} = p(\hat o|s)
\end{alignat}

Third, the edges' belief distributions are computed:
\begin{alignat}{2}
     & \Rightarrow q_i(s) &  & \propto \textcolor{black}{\bm{\overleftarrow{\mu}(}s\bm{)}} \textcolor{black}{\bm{\overrightarrow{\mu}(}s\bm{)}} = p(\hat o|s) p(s)                                        \\
     & \Rightarrow q_j(o) &  & \propto \textcolor{black}{\bm{\overleftarrow{\mu}(}o\bm{)}} \textcolor{black}{\bm{\overrightarrow{\mu}(}o\bm{)}} = \delta_{o,\hat o} \delta_{o,\hat o} = \delta_{o,\hat o}
\end{alignat}
In principle, the evaluation of \(q_j(o)\) is unnecessary because it will trivially yield the start condition (delta constraint).
Finally, we can insert the probabilities from the cat example:
\begin{alignat}{2}
    q_i(s=\text{"angel"}) & \propto p(o=\text{"breaks glas"}|s=\text{"angel"}) p(s=\text{"angel"}) = 0.1 \cdot 0.9 = 0.09   \\
    q_i(s=\text{"demon"}) & \propto p(o=\text{"breaks glas"}|s=\text{"demon"}) p(s=\text{"demon"}) = 0.99 \cdot 0.1 = 0.099
\end{alignat}
After normalization, the agent's belief \(q_i(s)\) yields:
\begin{alignat}{2}
    q_i(s=\text{"angel"}) & = \frac{0.09}{0.09+0.099} \approx 0.48  \\
    q_i(s=\text{"demon"}) & = \frac{0.099}{0.09+0.099} \approx 0.52
\end{alignat}
Furthermore, the observation belief \(q_j(o)\) trivially results into
\begin{alignat}{2}
    q_j(o=\text{"breaks glas"})        & = 1  \\
    q_j(o=\text{"doesn't break glas"}) & = 0.
\end{alignat}
because it is constrained by a real observation.

In conclusion, the message passing algorithm achieves Bayesian inference and the resulting belief distribution \(q_i(s)\) is equal to the posterior \(p(s|o=\text{"breaks glas"})\) that has been analytically shown before. Furthermore, the algorithm that results from the Bethe free energy as defined in this section is also known as the sum-product algorithm or belief propagation.


\section{Agents}

In this section I present the four agents that are developed in this work. First, I detail what agents entail.

\subsubsection{The Perception-Action Loop}

In general, agents are entities that are in constant exchange with an external environment. They receive observations from and act upon the environment. In this work observations and actions occur at discrete timesteps and are signified via the variables \(o_t\) and \(a_t\) respectively at the current timestep \(t\) (\cref{fig:23}).
\begin{figure}[htbp]
    \centering
    \begin{tikzpicture}[
            kreis/.style={
                    draw,
                    circle,
                    minimum size=3.5cm,
                    fill=white,
                    text centered,
                    font=\large\bfseries
                },
            pfeil/.style={
                    draw,
                    ->,
                    >={triangle 45},
                    dashed,
                    line width=3.5,
                    color=darkgray!50,
                    shorten >=3pt,
                    shorten <=3pt
                },
            label/.style={
                    midway,
                    font=\Large,
                    color=black,
                    text depth=0.5ex
                }
        ]

        \node[kreis] (agent) {Agent};
        \node[kreis, right=3cm of agent] (env) {Environment};

        \draw[pfeil] (env.north) to[bend right=40] node[label, above] {$o_t$} (agent.north);
        \draw[pfeil] (agent.south) to[bend right=40] node[label, below] {$a_t$} (env.south);

    \end{tikzpicture}
    \caption{Perception-Action Loop between the agent and the environment. The environment sends observations \(o_t\) to the agent. Vice versa, the agent sends actions \(a_t\) to the environment. The current timestep is \(t\).}
    \label{fig:23}
\end{figure}
The environment generates observations in possibly arbitrary ways but the agent is assumed to have certain mechanisms that process these observations and trigger actions. Specifically, agents are assumed to have a preference for certain observations that underwrite their ecological niche in the environment. Moreover, they are equipped with the means to understand (simulate) the chain of causality from hidden environmental causes to its sensors, i.e., they have of model how the state of the environment leads to certain observations. Even more, they also have a model of the causal chain of these unobservable (hidden) environmental states through time and how they can be influenced via actions. In total, this allows them to plan a sequence of actions that leads to preferred observations. Additionally, it allows them to keep past observations as a memory and enables them to retrospectively draw conclusions about past events in the light of new observations from the present.

\subsubsection{The Generative Model}

This mechanism rests upon a generative model that probabilistically relates observations, actions, and the (discrete) succession of hidden states. Formally, it is given by a likelihood \(p(o|s)\) that states how observations \(o\) are generated by hidden states \(s\) and by a transition model \(p(s_{t+1}|s_t,a_t)\) that defines how hidden states \(s\) change from timestep \(t\) to \(t+1\) given an action \(a_t\). Actions possess an action prior \(p(a)\) and preferences for observations are encoded by the goal prior \(\tilde{p}(o)\) that assigns higher probabilities to preferred observations.

\subsubsection{The Temporal Generative Model (POMDP)}

The four distributions (likelihood, transition model, action prior, and goal prior) are building blocks for a temporal rollout (from timestep \(0\) to \(T\)), which formally leads to a Partially Observable Markov Decision Process (POMDP). In later simulations there will be three timesteps, i.e., \(T=2\), and hidden states at \(t=0\) are initially distributed with \(p(s_0)\). The resulting POMDP is given by
\begin{equation}
    \begin{split}
        p(o_0, o_1, o_2, s_0, s_1, s_2, a_0, a_1) = \; & \tilde{p}(o_0) p(o_0|s_0) p(s_0)                \\
                                                       & \tilde{p}(o_1) p(o_1|s_1) p(s_1|s_0,a_0) p(a_0) \\
                                                       & \tilde{p}(o_2) p(o_2|s_2) p(s_2|s_1,a_1) p(a_1)
        \label{eq:15}
    \end{split}
\end{equation}
and has a factor graph representation (\cref{fig:8}) (including the variables' belief distributions \(q\) on the edges). Note, that in the POMDP (\cref{eq:15}) some variables occur in more than two factors (e.g., \(s_0\) in \(p(s_0)\), \(p(o_0|s_0)\) and \(p(s_1|s_0,a_0)\)). When representing the POMDP as a factor graph a special situation arises in which these variables (modeled by edges) should be connected to more than two factor nodes. Crucially, factor graphs are allowed to have exactly two nodes for each edge which requires an artificial duplication of variables to make them available to more than two factor nodes (e.g., \(s_0\) is duplicated to \(s^{\scriptscriptstyle I}_0\) and \(s^{\scriptscriptstyle II}_0\)). Duplicated variables are forced to be equivalent via the special equality node (indicated by the \(=\) sign) they are all connected to. Importantly, this situation is only a formal requirement of factor graphs and has no consequences for the results of inference processes presented in following sections.

\begin{figure}[htbp]
    \centering

    \resizebox{1.0\textwidth}{!}{%
        \begin{tikzpicture}[
                factor/.style={draw=black, thick, rectangle, minimum size=0.8cm, font=\normalsize},
                eqfactor/.style={draw=black, thick, rectangle, minimum size=0.5cm, font=\normalsize},
                edge/.style={thick},
                varlabel/.style={font=\normalsize},
                greycircle/.style={draw=black, thick, circle, fill=gray!30, minimum size=0.5cm, font=\normalsize}
            ]

            % Knoten
            \node[factor] (d) {$p(s_0)$};
            \node[eqfactor, right=1.8cm of d] (eq0) {$=$};
            \node[factor, below=1.5cm of eq0] (A0) {$p(o_0|s^{\scriptscriptstyle II}_0)$};
            \node[factor, below=1.5cm of A0] (c0) {$\tilde{p}(o_0)$};
            \node[factor, right=1.8cm of eq0] (B0) {$p(s_1|s^{\scriptscriptstyle I}_0,a_0)$};
            \node[eqfactor, right=1.8cm of B0] (eq1) {$=$};
            \node[factor, below=1.5cm of eq1] (A1) {$p(o_1|s^{\scriptscriptstyle II}_1)$};
            \node[factor, below=1.5cm of A1] (c1) {$\tilde{p}(o_1)$};
            \node[factor, right=1.8cm of eq1] (B1) {$p(s_2|s^{\scriptscriptstyle I}_1,a_1)$};
            \node[factor] (A2) at ([xshift=2cm] B1.east |- A1) {$p(o_2|s_2)$};
            \node[factor, below=1.5cm of A2] (c2) {$\tilde{p}(o_2)$};
            \node[factor, above=1.5cm of B0] (u0) {$p(a_0)$};
            \node[factor, above=1.5cm of B1] (u1) {$p(a_1)$};

            % Kanten zwischen Knoten
            \draw[edge] (d.east) -- (eq0.west) node[midway, above, varlabel] {$q(s_0)$};

            \draw[edge] (eq0.south) -- (A0.north) node[midway, left, varlabel] {$q(s^{\scriptscriptstyle II}_0)$};

            \draw[edge] (A0.south) -- (c0.north) node[midway, left, varlabel] {$q(o_0)$};

            \draw[edge] (eq0.east) -- (B0.west) node[midway, above, varlabel] {$q(s^{\scriptscriptstyle I}_0)$};

            \draw[edge] (B0.north) -- (u0.south) node[midway, left, varlabel] {$q(a_0)$};

            \draw[edge] (B0.east) -- (eq1.west) node[midway, above, varlabel] {$q(s_1)$};

            \draw[edge] (eq1.south) -- (A1.north) node[midway, left, varlabel] {$q(s^{\scriptscriptstyle II}_1)$};

            \draw[edge] (A1.south) -- (c1.north) node[midway, left, varlabel] {$q(o_1)$};

            \draw[edge] (eq1.east) -- (B1.west) node[midway, above, varlabel] {$q(s^{\scriptscriptstyle I}_1)$};

            \draw[edge] (B1.east) -| (A2.north) node[pos=0.25, above, varlabel] {$q(s_2)$};

            \draw[edge] (A2.south) -- (c2.north) node[midway, left, varlabel] {$q(o_2)$};

            \draw[edge] (B1.north) -- (u1.south) node[midway, left, varlabel] {$q(a_1)$};

            % Kanten an Knoten

        \end{tikzpicture}
    }

    \caption{Foundational Factor Graph of the two non-Context Agents.}
    \label{fig:8}
\end{figure}


\subsubsection{Lifecycle of an Agent}

The lifecycle of agents consists of three recurring stages (\cref{fig:30}): First, it receives an observation, which leads to a delta constraint at the respective belief distribution. For example, when the current timestep is \(t=0\) and it receives an observation \(\hat {o_0}\) then \(q(o_0)=\delta_{o_0,\hat {o_0}}\).

Second, the agent performs inference over the whole factor graph, i.e., the corresponding message passing algorithm is executed. In this way it computes the posterior distributions of all variables in the graph based on the observation \(\hat {o_0}\). These posteriors are then given by the respective belief distributions \(q\). In other words, during inference the agent forms beliefs about all variables in its graph. For example, it infers the hidden cause \(s_0\) of the observation \(\hat {o_0}\). Also, it infers a likely future trajectory of hidden states (via \(q(s_1)\) and \(q(s_2)\)). Moreover, it also forms beliefs about future observations. Crucially, these future observation beliefs are not solely driven by future hidden states but also by the preferences (i.e., \(\tilde{p}(o_1)\) and \(\tilde{p}(o_2)\)) on observations. The additional influence from the preferences couples back to other parts of graph such that the agent infers a future trajectory of hidden states that also leads to preferred observation. Moreover, it will also form appropriate action beliefs that fulfill these preferences.

In the third stage, the agent samples an action from the current action belief (e.g., \(q(a_0)\)) and "sends" it to the environment. In this work, the action belief will also be constrained by a delta distribution which enables the agent to "remember" which action it has executed.

After executing an action the cycle starts over again, i.e., the agent receives the next observation (\(\hat {o_1}\)). In this way, the agent collects observations from each timestep which remain as delta constraints in the factor graph. This enables the agent to remember past observations. Moreover, past beliefs about hidden states can still change when receiving new observations. This way, the agent can change its beliefs retrospectively. The other way around, past observations also influence the future trajectory of hidden states beliefs, i.e, the agent utilizes its memory of past events to make better predictions about the future.

% --- 1. GRAFIKEN ISOLIERT LADEN UND SKALIEREN ---
% Falls du die external-library in der Masterarbeit nutzt, hier kurz pausieren
\tikzexternaldisable

% Wir nutzen \scalebox, da es im Gegensatz zu TikZ-scale AUCH Label-Abstände schrumpft!
\newsavebox{\boxTL}
\sbox{\boxTL}{\scalebox{0.75}{\begin{tikzpicture}[
    factor/.style={draw=black, thick, rectangle, minimum size=0.8cm, font=\normalsize, fill=white},
    eqfactor/.style={draw=black, thick, rectangle, minimum size=0.5cm, font=\normalsize, fill=white},
    edge/.style={thick},
    varlabel/.style={font=\normalsize},
    greycircle/.style={draw=black, thick, circle, fill=gray!30, minimum size=0.5cm, font=\normalsize}
]

% Knoten der linken Hälfte (t=0)
\node[factor] (d) {$p(s_0)$};
\node[eqfactor, right=1.8cm of d] (eq0) {$=$};
\node[factor, below=1.5cm of eq0] (A0) {$p(o_0|s^{\scriptscriptstyle II}_0)$};
\node[factor, below=1.5cm of A0] (c0) {$\tilde{p}(o_0)$};
\node[factor, right=1.8cm of eq0] (B0) {$p(s_1|s^{\scriptscriptstyle I}_0,a_0)$};
\node[factor, above=1.5cm of B0] (u0) {$p(a_0)$};

% Unsichtbarer Endpunkt für die ausgehende Kante
\coordinate[right=1.2cm of B0] (out_s1);

% Kanten
\draw[edge] (d.east) -- (eq0.west) node[midway, above, varlabel] {$q(s_0)$};
\draw[edge] (eq0.south) -- (A0.north) node[midway, left, varlabel] {$q(s^{\scriptscriptstyle II}_0)$};
\draw[edge] (A0.south) -- (c0.north) 
    node[midway, left=0.3cm, varlabel] {{\Large \textcolor{accorange}{$\hat{o}_0$}} $\rightarrow q(o_0)$} 
    node[midway, greycircle] {$\delta$};\draw[edge] (eq0.east) -- (B0.west) node[midway, above, varlabel] {$q(s^{\scriptscriptstyle I}_0)$};
\draw[edge] (B0.north) -- (u0.south) node[midway, left, varlabel] {$q(a_0)$};

% Abgebrochene Kante, die anzeigt, dass der Graph weitergeht (optional gestrichelt)
\draw[edge, densely dashed] (B0.east) -- (out_s1) node[midway, above, varlabel] {$q(s_1)$};

\end{tikzpicture}}}

\newsavebox{\boxTR}
\sbox{\boxTR}{\scalebox{0.75}{\begin{tikzpicture}[
    factor/.style={draw=black, thick, rectangle, minimum size=0.8cm, font=\normalsize, fill=white},
    eqfactor/.style={draw=black, thick, rectangle, minimum size=0.5cm, font=\normalsize, fill=white},
    edge/.style={thick},
    varlabel/.style={font=\normalsize},
    greycircle/.style={draw=black, thick, circle, fill=gray!30, minimum size=0.5cm, font=\normalsize}
]

% --- KNOTEN ---
% Wir starten mit eq0 als "Anker" für diesen Ausschnitt
\node[eqfactor] (eq0) {$=$};

% Knoten t=0 (unterhalb eq0)
\node[factor, below=1.5cm of eq0] (A0) {$p(o_0|s^{\scriptscriptstyle II}_0)$};
\node[factor, below=1.5cm of A0] (c0) {$\tilde{p}(o_0)$};

% Übergang & Knoten t=1
\node[factor, right=1.8cm of eq0] (B0) {$p(s_1|s^{\scriptscriptstyle I}_0,a_0)$};
\node[factor, above=1.5cm of B0] (u0) {$p(a_0)$};

\node[eqfactor, right=1.8cm of B0] (eq1) {$=$};
\node[factor, below=1.5cm of eq1] (A1) {$p(o_1|s^{\scriptscriptstyle II}_1)$};
\node[factor, below=1.5cm of A1] (c1) {$\tilde{p}(o_1)$};

% Unsichtbare Koordinaten für die gestrichelten Enden links und rechts
\coordinate[left=1.2cm of eq0] (in_s0);
\coordinate[right=1.2cm of eq1] (out_sI1);


% --- KANTEN ---

% 1. Gestrichelte Eingangskante q(s_0)
\draw[edge, densely dashed] (in_s0) -- (eq0.west) node[midway, above, varlabel] {$q^*(s_0)$};

% 2. Durchgezogene innere Kanten
\draw[edge] (eq0.south) -- (A0.north) node[midway, left, varlabel] {$q^*(s^{\scriptscriptstyle II}_0)$};
\draw[edge] (A0.south) -- (c0.north) node[midway, left=0.3, varlabel] {$q(o_0)$} node[midway, greycircle] {$\delta$};

\draw[edge] (eq0.east) -- (B0.west) node[midway, above, varlabel] {$q^*(s^{\scriptscriptstyle I}_0)$};
\draw[edge] (B0.north) -- (u0.south) node[midway, left=0.3cm, varlabel] {{\Large \textcolor{accorange}{$\hat{a}_0$}} $\sim q^*(a_0)$} node[midway, greycircle] {$\delta$};
\draw[edge] (B0.east) -- (eq1.west) node[midway, above, varlabel] {$q^*(s_1)$};

\draw[edge] (eq1.south) -- (A1.north) node[midway, left, varlabel] {$q^*(s^{\scriptscriptstyle II}_1)$};
\draw[edge] (A1.south) -- (c1.north) node[midway, left, varlabel] {$q^*(o_1)$};

% 3. Gestrichelte Ausgangskante q(s^I_1)
\draw[edge, densely dashed] (eq1.east) -- (out_sI1) node[midway, above, varlabel] {$q^*(s^{\scriptscriptstyle I}_1)$};

\end{tikzpicture}}}

\newsavebox{\boxBOT}
\sbox{\boxBOT}{\scalebox{0.75}{\begin{tikzpicture}[
    factor/.style={draw=black, thick, rectangle, minimum size=0.8cm, font=\normalsize, fill=white},
    eqfactor/.style={draw=black, thick, rectangle, minimum size=0.5cm, font=\normalsize, fill=white},
    edge/.style={thick},
    varlabel/.style={font=\normalsize},
    greycircle/.style={draw=black, thick, circle, fill=gray!30, minimum size=0.5cm, font=\normalsize}
]

% Knoten
\node[factor] (d) {$p(s_0)$};
\node[eqfactor, right=1.8cm of d] (eq0) {$=$};
\node[factor, below=1.5cm of eq0] (A0) {$p(o_0|s^{\scriptscriptstyle II}_0)$};
\node[factor, below=1.5cm of A0] (c0) {$\tilde{p}(o_0)$};
\node[factor, right=1.8cm of eq0] (B0) {$p(s_1|s^{\scriptscriptstyle I}_0,a_0)$};
\node[eqfactor, right=1.8cm of B0] (eq1) {$=$};
\node[factor, below=1.5cm of eq1] (A1) {$p(o_1|s^{\scriptscriptstyle II}_1)$};
\node[factor, below=1.5cm of A1] (c1) {$\tilde{p}(o_1)$};
\node[factor, right=1.8cm of eq1] (B1) {$p(s_2|s^{\scriptscriptstyle I}_1,a_1)$};
\node[factor] (A2) at ([xshift=2cm] B1.east |- A1) {$p(o_2|s_2)$};
\node[factor, below=1.5cm of A2] (c2) {$\tilde{p}(o_2)$};
\node[factor, above=1.5cm of B0] (u0) {$p(a_0)$};
\node[factor, above=1.5cm of B1] (u1) {$p(a_1)$};

% Kanten zwischen Knoten
\draw[edge] (d.east) -- (eq0.west) node[midway, above, varlabel] {$q(s_0)$} node[near end, below, varlabel] {$\textcolor{accorange}{\bm{\overrightarrow{\mu}}}$} node[near start, below, varlabel] {$\textcolor{accorange}{\bm{\overleftarrow{\mu}}}$};
\draw[edge] (eq0.south) -- (A0.north) node[midway, left, varlabel] {$q(s^{\scriptscriptstyle II}_0)$} node[near start, right, varlabel] {$\textcolor{accorange}{\bm{{\uparrow}{\mu}}}$} node[near end, right, varlabel] {$\textcolor{accorange}{\bm{{\downarrow}{\mu}}}$};
\draw[edge] (A0.south) -- (c0.north) node[midway, left=0.3, varlabel] {$q(o_0)$} node[midway, greycircle] {$\delta$} node[near start, right=0.2, varlabel] {$\textcolor{accorange}{\bm{{\uparrow}{\mu}}}$} node[near end, right=0.2, varlabel] {$\textcolor{accorange}{\bm{{\downarrow}{\mu}}}$};
\draw[edge] (eq0.east) -- (B0.west) node[midway, above, varlabel] {$q(s^{\scriptscriptstyle I}_0)$} node[near end, below, varlabel] {$\textcolor{accorange}{\bm{\overrightarrow{\mu}}}$} node[near start, below, varlabel] {$\textcolor{accorange}{\bm{\overleftarrow{\mu}}}$};
\draw[edge] (B0.north) -- (u0.south) node[midway, left, varlabel] {$q(a_0)$} node[near end, right, varlabel] {$\textcolor{accorange}{\bm{{\uparrow}{\mu}}}$} node[near start, right, varlabel] {$\textcolor{accorange}{\bm{{\downarrow}{\mu}}}$};
\draw[edge] (B0.east) -- (eq1.west) node[midway, above, varlabel] {$q(s_1)$} node[near end, below, varlabel] {$\textcolor{accorange}{\bm{\overrightarrow{\mu}}}$} node[near start, below, varlabel] {$\textcolor{accorange}{\bm{\overleftarrow{\mu}}}$};
\draw[edge] (eq1.south) -- (A1.north) node[midway, left, varlabel] {$q(s^{\scriptscriptstyle II}_1)$} node[near start, right, varlabel] {$\textcolor{accorange}{\bm{{\uparrow}{\mu}}}$} node[near end, right, varlabel] {$\textcolor{accorange}{\bm{{\downarrow}{\mu}}}$};
\draw[edge] (A1.south) -- (c1.north) node[midway, left, varlabel] {$q(o_1)$} node[near start, right, varlabel] {$\textcolor{accorange}{\bm{{\uparrow}{\mu}}}$} node[near end, right, varlabel] {$\textcolor{accorange}{\bm{{\downarrow}{\mu}}}$};
\draw[edge] (eq1.east) -- (B1.west) node[midway, above, varlabel] {$q(s^{\scriptscriptstyle I}_1)$} node[near end, below, varlabel] {$\textcolor{accorange}{\bm{\overrightarrow{\mu}}}$} node[near start, below, varlabel] {$\textcolor{accorange}{\bm{\overleftarrow{\mu}}}$};
\draw[edge] (B1.east) -| (A2.north) node[pos=0.25, above, varlabel] {$q(s_2)$} node[pos=0.11, below, varlabel] {$\textcolor{accorange}{\bm{\overleftarrow{\mu}}}$} node[pos=0.89, right, varlabel] {$\textcolor{accorange}{\bm{{\downarrow}{\mu}}}$};
\draw[edge] (A2.south) -- (c2.north) node[midway, left, varlabel] {$q(o_2)$} node[near start, right, varlabel] {$\textcolor{accorange}{\bm{{\uparrow}{\mu}}}$} node[near end, right, varlabel] {$\textcolor{accorange}{\bm{{\downarrow}{\mu}}}$};
\draw[edge] (B1.north) -- (u1.south) node[midway, left, varlabel] {$q(a_1)$} node[near end, right, varlabel] {$\textcolor{accorange}{\bm{{\uparrow}{\mu}}}$} node[near start, right, varlabel] {$\textcolor{accorange}{\bm{{\downarrow}{\mu}}}$};

\end{tikzpicture}}}


% --- 2. DAS LAYOUT ZUSAMMENSETZEN ---
\begin{figure}[htbp]
    \centering
    \resizebox{1.0\textwidth}{!}{%
        \begin{tikzpicture}[
                >=Stealth,
                % Hier definieren wir die blauen Boxen. Die runden Ecken hier 
                % können jetzt NICHT mehr in deine Faktorgraphen durchbluten.
                panel/.style={
                        draw=gray!80, thick, fill=blue!5,
                        rounded corners=4pt, inner sep=8pt, align=center
                    },
                cycle arrow/.style={
                        ->, line width=2.0pt, color=gray!80!black
                    }
            ]

            % Die vorbereiteten, perfekt skalierten Boxen einfügen
            \node[panel] (TL) at (0,0) {\usebox{\boxTL}};

            \node[panel, right=0.5cm of TL] (TR) {\usebox{\boxTR}};

            \path (TL.south west) -- (TR.south east) coordinate[midway] (centerBottom);
            \node[panel, below=0.5cm of centerBottom, anchor=north] (BOT) {\usebox{\boxBOT}};

            % --- NEU: Freischwebender Text IN den Boxen (oben rechts) ---
            % xshift und yshift rücken den Text etwas vom Rand weg nach innen.
            \node[anchor=north west, font=\sffamily\bfseries\normalsize, text=gray!80!black]
            at ([xshift=8pt, yshift=-8pt]TL.north west) {Perception};

            \node[anchor=north west, font=\sffamily\bfseries\normalsize, text=gray!80!black]
            at ([xshift=8pt, yshift=-8pt]TR.north west) {Action};

            \node[anchor=north west, font=\sffamily\bfseries\normalsize, text=gray!80!black]
            at ([xshift=8pt, yshift=-8pt]BOT.north west) {Inference};
            % -------------------------------------------------------------

            % Die Pfeile außen herum
            \draw[cycle arrow, dashed] (TR.north) to[bend right=20]
            node[midway, above=4pt, font=\sffamily\bfseries] {} (TL.north);

            \draw[cycle arrow] (TL.west) to[bend right=35]
            node[midway, left=4pt, font=\sffamily\bfseries] {} (BOT.west);

            \draw[cycle arrow] (BOT.east) to[bend right=35]
            node[midway, right=4pt, font=\sffamily\bfseries] {} (TR.east);

        \end{tikzpicture}
    }

    \tikzexternalenable

    \caption{Agents' lifecycle with the recurring stages perception, inference, and action. During perception the observation belief that corresponds to the current timestep (e.g., \(q(o_0\)) becomes a delta distribution with all probability mass at the actual observation \(\hat{o_0}\). When performing inference the messages are sent on the edges (message passing) and all belief distributions are updated (indicated by \(q^*\)) (expect the delta constrained distributions). During action, an action \(\hat{a_0}\) is sampled from the corresponding belief distribution \(q(a_0)\) that has been infered during inference. Additionally, this belief also becomes a delta distribution with all probability mass on the sampled action \(\hat{a_0}\). After an action, the cycle starts over again and the next observation (e.g., \(\hat{o_1}\)) is perceived.}
    \label{fig:30}
\end{figure}

The POMDP in \cref{eq:15} and the corresponding factor graph in \cref{fig:8} will be the basis for the BFE- and GFE agents that are introduced in following sections.


\subsection{Bethe Free Energy (BFE) Agent}

In this section I present the message passing algorithm of the BFE agent. This agent serves as a baseline without having an additional context layer nor does it do information seeking. All messages can be derived from the Bethe free energy of its factor graph (compare \cref{sec:1}). Messages from equality nodes also result from the Bethe free energy (see \cref{app:1}). The agent's generative model is given by \cref{eq:15}.

\subsubsection{Message Passing Algorithm}

The factor graph that results from the POMDP is given in \cref{fig:31} (including variable beliefs and messages) and the corresponding message passing algorithm that enables inference is given in \cref{alg:2}. Note that in the graph the agent has already received an observation \(\hat {o_0}\) to exemplify the alternative messages resulting from the corresponding delta constraint.

\begin{figure}[htbp]
    \centering

    \resizebox{1.0\textwidth}{!}{%
        \begin{tikzpicture}[
                factor/.style={draw=black, thick, rectangle, minimum size=0.8cm, font=\normalsize},
                eqfactor/.style={draw=black, thick, rectangle, minimum size=0.5cm, font=\normalsize},
                edge/.style={thick},
                varlabel/.style={font=\normalsize},
                greycircle/.style={draw=black, thick, circle, fill=gray!30, minimum size=0.5cm, font=\normalsize}
            ]

            % Knoten
            \node[factor] (d) {$p(s_0)$};
            \node[eqfactor, right=1.8cm of d] (eq0) {$=$};
            \node[factor, below=1.5cm of eq0] (A0) {$p(o_0|s^{\scriptscriptstyle II}_0)$};
            \node[factor, below=1.5cm of A0] (c0) {$\tilde{p}(o_0)$};
            \node[factor, right=1.8cm of eq0] (B0) {$p(s_1|s^{\scriptscriptstyle I}_0,a_0)$};
            \node[eqfactor, right=1.8cm of B0] (eq1) {$=$};
            \node[factor, below=1.5cm of eq1] (A1) {$p(o_1|s^{\scriptscriptstyle II}_1)$};
            \node[factor, below=1.5cm of A1] (c1) {$\tilde{p}(o_1)$};
            \node[factor, right=1.8cm of eq1] (B1) {$p(s_2|s^{\scriptscriptstyle I}_1,a_1)$};
            \node[factor] (A2) at ([xshift=2cm] B1.east |- A1) {$p(o_2|s_2)$};
            \node[factor, below=1.5cm of A2] (c2) {$\tilde{p}(o_2)$};
            \node[factor, above=1.5cm of B0] (u0) {$p(a_0)$};
            \node[factor, above=1.5cm of B1] (u1) {$p(a_1)$};

            % Kanten zwischen Knoten
            \draw[edge] (d.east) -- (eq0.west) node[midway, above, varlabel] {$\textcolor{accblue}{q(s_0)}$} node[near end, below, varlabel, scale=0.7] {${\large \textcolor{accorange}{\bm{\overrightarrow{\mu}(}s_0\bm{)}}}$} node[near start, below, varlabel, scale=0.7] {${\large \textcolor{accorange}{\bm{\overleftarrow{\mu}(}s_0\bm{)}}}$};
            \draw[edge] (eq0.south) -- (A0.north) node[midway, left, varlabel] {$\textcolor{accblue}{q(s^{\scriptscriptstyle II}_0)}$} node[pos=0.25, right=0.05, varlabel, scale=0.7] {${\large \textcolor{accorange}{\bm{{\uparrow}{\mu}(}s^{\scriptscriptstyle II}_0\bm{)}}}$} node[pos=0.75, right=0.05, varlabel, scale=0.7] {${\large \textcolor{accorange}{\bm{{\downarrow}{\mu}(}s^{\scriptscriptstyle II}_0\bm{)}}}$};
            \draw[edge] (A0.south) -- (c0.north) node[midway, greycircle] {$\delta$} node[midway, left=0.3, varlabel] {$\textcolor{accblue}{q(o_0)}$} node[pos=0.25, right=0.25, varlabel, scale=0.7] {${\large \textcolor{accorange}{\bm{{\uparrow}{\mu}(}o_0\bm{)}}}$} node[pos=0.8, right=0.25, varlabel, scale=0.7] {${\large \textcolor{accorange}{\bm{{\downarrow}{\mu}(}o_0\bm{)}}}$};
            \draw[edge] (eq0.east) -- (B0.west) node[midway, above, varlabel] {$\textcolor{accblue}{q(s^{\scriptscriptstyle I}_0)}$} node[near start, below, varlabel, scale=0.7] {${\large \textcolor{accorange}{\bm{\overleftarrow{\mu}(}s^{\scriptscriptstyle I}_0\bm{)}}}$} node[near end, below, varlabel, scale=0.7] {${\large \textcolor{accorange}{\bm{\overrightarrow{\mu}(}s^{\scriptscriptstyle I}_0\bm{)}}}$};
            \draw[edge] (B0.north) -- (u0.south) node[midway, left, varlabel] {$\textcolor{accblue}{q(a_0)}$} node[pos=0.25, right=0.05, varlabel, scale=0.7] {${\large \textcolor{accorange}{\bm{{\downarrow}{\mu}(}a_0\bm{)}}}$} node[pos=0.8, right=0.05, varlabel, scale=0.7] {${\large \textcolor{accorange}{\bm{{\uparrow}{\mu}(}a_0\bm{)}}}$};
            \draw[edge] (B0.east) -- (eq1.west) node[midway, above, varlabel] {$\textcolor{accblue}{q(s_1)}$} node[near end, below, varlabel, scale=0.7] {${\large \textcolor{accorange}{\bm{\overrightarrow{\mu}(}s_1\bm{)}}}$} node[near start, below, varlabel, scale=0.7] {${\large \textcolor{accorange}{\bm{\overleftarrow{\mu}(}s_1\bm{)}}}$};
            \draw[edge] (eq1.south) -- (A1.north) node[midway, left, varlabel] {$\textcolor{accblue}{q(s^{\scriptscriptstyle II}_1)}$} node[pos=0.25, right=0.05, varlabel, scale=0.7] {${\large \textcolor{orange!75!black}{\bm{{\uparrow}{\mu}(}s^{\scriptscriptstyle II}_1\bm{)}}}$} node[pos=0.75, right=0.05, varlabel, scale=0.7] {${\large \textcolor{orange!75!black}{\bm{{\downarrow}{\mu}(}s^{\scriptscriptstyle II}_1\bm{)}}}$};
            \draw[edge] (A1.south) -- (c1.north) node[midway, left, varlabel] {$\textcolor{accblue}{q(o_1)}$} node[pos=0.25, right=0.05, varlabel, scale=0.7] {${\large \textcolor{orange!75!black}{\bm{{\uparrow}{\mu}(}o_1\bm{)}}}$} node[pos=0.8, right=0.05, varlabel, scale=0.7] {${\large \textcolor{accorange}{\bm{{\downarrow}{\mu}(}o_1\bm{)}}}$};
            \draw[edge] (eq1.east) -- (B1.west) node[midway, above, varlabel] {$\textcolor{accblue}{q(s^{\scriptscriptstyle I}_1)}$} node[near start, below, varlabel, scale=0.7] {${\large \textcolor{accorange}{\bm{\overleftarrow{\mu}(}s^{\scriptscriptstyle I}_1\bm{)}}}$} node[near end, below, varlabel, scale=0.7] {${\large \textcolor{accorange}{\bm{\overrightarrow{\mu}(}s^{\scriptscriptstyle I}_1\bm{)}}}$};
            \draw[edge] (B1.east) -| (A2.north) node[pos=0.25, above, varlabel] {$\textcolor{accblue}{q(s_2)}$} node[pos=0.89, right=0.05, varlabel, scale=0.7] {${\large \textcolor{accorange}{\bm{{\downarrow}{\mu}(}s_2\bm{)}}}$} node[pos=0.11, below, varlabel, scale=0.7] {${\large \textcolor{accorange}{\bm{\overleftarrow{\mu}(}s_2\bm{)}}}$};
            \draw[edge] (A2.south) -- (c2.north) node[midway, left, varlabel] {$\textcolor{accblue}{q(o_2)}$} node[pos=0.25, right=0.05, varlabel, scale=0.7] {${\large \textcolor{accorange}{\bm{{\uparrow}{\mu}(}o_2\bm{)}}}$} node[pos=0.8, right=0.05, varlabel, scale=0.7] {${\large \textcolor{accorange}{\bm{{\downarrow}{\mu}(}o_2\bm{)}}}$};
            \draw[edge] (B1.north) -- (u1.south) node[midway, left, varlabel] {$\textcolor{accblue}{q(a_1)}$} node[pos=0.25, right=0.05, varlabel, scale=0.7] {${\large \textcolor{accorange}{\bm{{\downarrow}{\mu}(}a_1\bm{)}}}$} node[pos=0.8, right=0.05, varlabel, scale=0.7] {${\large \textcolor{accorange}{\bm{{\uparrow}{\mu}(}a_1\bm{)}}}$};

        \end{tikzpicture}
    }

    \caption{Factor graph of the BFE agent with variables' belief distributions \(\textcolor{accblue}{q}\) and messages \(\textcolor{accorange}{\mu}\).}
    \label{fig:31}
\end{figure}

\begin{algorithm}[htbp]
    \caption{Message Passing Algorithm of the BFE Agent}
    \label{alg:2}
    \small
    \begin{algorithmic}[1] % Die [1] sorgt für automatische Zeilennummern!
        \linespread{1.4}\selectfont
        \State $\textcolor{accorange}{\bm{\overrightarrow{\mu}(}s_0\bm{)}} \propto p(s_0)$,\quad$\textcolor{accorange}{\bm{{\uparrow}{\mu}(}o_0\bm{)}} \propto \delta_{o_0,\hat{o_0}}$,\quad$\textcolor{accorange}{\bm{{\downarrow}{\mu}(}a_0\bm{)}} \propto p(a_0)$,\quad$\textcolor{accorange}{\bm{{\uparrow}{\mu}(}o_1\bm{)}} \propto \tilde{p}(o_1)$,\quad$\textcolor{accorange}{\bm{{\downarrow}{\mu}(}a_1\bm{)}} \propto p(a_1)$,
        \Statex $\textcolor{accorange}{\bm{{\uparrow}{\mu}(}o_2\bm{)}} \propto \tilde{p}(o_2)$,\quad$\textcolor{accorange}{\bm{{\downarrow}{\mu}(}o_0\bm{)}} \propto \delta_{o_0,\hat{o_0}}$
        \State $\textcolor{accorange}{\bm{{\uparrow}{\mu}(}s^{\scriptscriptstyle II}_0\bm{)}} \propto \sum_{o_0} p(o_0|s^{\scriptscriptstyle II}_0) \textcolor{black}{\bm{{\uparrow}{\mu}(}o_0\bm{)}}$,\quad$\textcolor{accorange}{\bm{{\uparrow}{\mu}(}s^{\scriptscriptstyle II}_1\bm{)}} \propto \sum_{o_1} p(o_1|s^{\scriptscriptstyle II}_1) \textcolor{black}{\bm{{\uparrow}{\mu}(}o_1\bm{)}}$,
        \Statex $\textcolor{accorange}{\bm{\overleftarrow{\mu}(}s_2\bm{)}} \propto \sum_{o_2} p(o_2|s_2) \textcolor{black}{\bm{{\uparrow}{\mu}(}o_2\bm{)}}$
        \State $\textcolor{accorange}{\bm{\overrightarrow{\mu}(}s^{\scriptscriptstyle I}_0\bm{)}} \propto \sum_{s_0,s^{\scriptscriptstyle II}_0} \delta_{s_0,s^{\scriptscriptstyle I}_0}\delta_{s_0,s^{\scriptscriptstyle II}_0} \textcolor{black}{\bm{\overrightarrow{\mu}(}s_0\bm{)}} \textcolor{black}{\bm{{\uparrow}{\mu}(}s^{\scriptscriptstyle II}_0\bm{)}}$,\quad$\textcolor{accorange}{\bm{\overleftarrow{\mu}(}s^{\scriptscriptstyle I}_1\bm{)}} \propto \sum_{a_1,s_2} p(s_2|s^{\scriptscriptstyle I}_1,a_1) \textcolor{black}{\bm{{\downarrow}{\mu}(}a_1\bm{)}} \textcolor{black}{\bm{\overleftarrow{\mu}(}s_2\bm{)}}$
        \State $\textcolor{accorange}{\bm{\overrightarrow{\mu}(}s_1\bm{)}} \propto \sum_{s^{\scriptscriptstyle I}_0,a_0} p(s_1|s^{\scriptscriptstyle I}_0,a_0) \textcolor{black}{\bm{\overrightarrow{\mu}(}s^{\scriptscriptstyle I}_0\bm{)}} \textcolor{black}{\bm{{\downarrow}{\mu}(}a_0\bm{)}}$,\quad$\textcolor{accorange}{\bm{\overleftarrow{\mu}(}s_1\bm{)}} \propto \sum_{s^{\scriptscriptstyle I}_1,s^{\scriptscriptstyle II}_1} \delta_{s_1,s^{\scriptscriptstyle I}_1}\delta_{s_1,s^{\scriptscriptstyle II}_1} \textcolor{black}{\bm{\overleftarrow{\mu}(}s^{\scriptscriptstyle I}_1\bm{)}} \textcolor{black}{\bm{{\uparrow}{\mu}(}s^{\scriptscriptstyle II}_1\bm{)}}$
        \State $\textcolor{accorange}{\bm{\overrightarrow{\mu}(}s^{\scriptscriptstyle I}_1\bm{)}} \propto \sum_{s_1,s^{\scriptscriptstyle II}_1} \delta_{s_1,s^{\scriptscriptstyle I}_1}\delta_{s_1,s^{\scriptscriptstyle II}_1} \textcolor{black}{\bm{\overrightarrow{\mu}(}s_1\bm{)}} \textcolor{black}{\bm{{\uparrow}{\mu}(}s^{\scriptscriptstyle II}_1\bm{)}}$,\quad$\textcolor{accorange}{\bm{{\downarrow}{\mu}(}s^{\scriptscriptstyle II}_1\bm{)}} \propto \sum_{s_1,s^{\scriptscriptstyle I}_1} \delta_{s_1,s^{\scriptscriptstyle I}_1}\delta_{s_1,s^{\scriptscriptstyle II}_1} \textcolor{black}{\bm{\overrightarrow{\mu}(}s_1\bm{)}} \textcolor{black}{\bm{\overleftarrow{\mu}(}s^{\scriptscriptstyle I}_1\bm{)}}$,
        \Statex $\textcolor{accorange}{\bm{{\uparrow}{\mu}(}a_0\bm{)}} \propto \sum_{s_1,s^{\scriptscriptstyle I}_0} p(s_1|s^{\scriptscriptstyle I}_0,a_0) \textcolor{black}{\bm{\overrightarrow{\mu}(}s^{\scriptscriptstyle I}_0\bm{)}} \textcolor{black}{\bm{\overleftarrow{\mu}(}s_1\bm{)}}$,\quad$\textcolor{accorange}{\bm{\overleftarrow{\mu}(}s^{\scriptscriptstyle I}_0\bm{)}} \propto \sum_{s_1,a_0} p(s_1|s^{\scriptscriptstyle I}_0,a_0) \textcolor{black}{\bm{{\downarrow}{\mu}(}a_0\bm{)}} \textcolor{black}{\bm{\overleftarrow{\mu}(}s_1\bm{)}}$
        \State $\textcolor{accorange}{\bm{{\uparrow}{\mu}(}a_1\bm{)}} \propto \sum_{s^{\scriptscriptstyle I}_1,s_2} p(s_2|s^{\scriptscriptstyle I}_1,a_1) \textcolor{black}{\bm{\overrightarrow{\mu}(}s^{\scriptscriptstyle I}_1\bm{)}} \textcolor{black}{\bm{\overleftarrow{\mu}(}s_2\bm{)}}$,\quad$\textcolor{accorange}{\bm{{\downarrow}{\mu}(}s_2\bm{)}} \propto \sum_{s^{\scriptscriptstyle I}_1,a_1} p(s_2|s^{\scriptscriptstyle I}_1,a_1) \textcolor{black}{\bm{\overrightarrow{\mu}(}s^{\scriptscriptstyle I}_1\bm{)}} \textcolor{black}{\bm{{\downarrow}{\mu}(}a_1\bm{)}}$,
        \Statex $\textcolor{accorange}{\bm{\overleftarrow{\mu}(}s_0\bm{)}} \propto \sum_{s^{\scriptscriptstyle I}_0,s^{\scriptscriptstyle II}_0} \delta_{s_0,s^{\scriptscriptstyle I}_0}\delta_{s_0,s^{\scriptscriptstyle II}_0} \textcolor{black}{\bm{\overleftarrow{\mu}(}s^{\scriptscriptstyle I}_0\bm{)}} \textcolor{black}{\bm{{\uparrow}{\mu}(}s^{\scriptscriptstyle II}_0\bm{)}}$,\quad$\textcolor{accorange}{\bm{{\downarrow}{\mu}(}s^{\scriptscriptstyle II}_0\bm{)}} \propto \sum_{s_0,s^{\scriptscriptstyle I}_0} \delta_{s_0,s^{\scriptscriptstyle I}_0}\delta_{s_0,s^{\scriptscriptstyle II}_0} \textcolor{black}{\bm{\overrightarrow{\mu}(}s_0\bm{)}} \textcolor{black}{\bm{\overleftarrow{\mu}(}s^{\scriptscriptstyle I}_0\bm{)}}$,
        \Statex $\textcolor{accorange}{\bm{{\downarrow}{\mu}(}o_1\bm{)}} \propto \sum_{s^{\scriptscriptstyle II}_1} p(o_1|s^{\scriptscriptstyle II}_1) \textcolor{black}{\bm{{\downarrow}{\mu}(}s^{\scriptscriptstyle II}_1\bm{)}}$
        \State $\textcolor{accorange}{\bm{{\downarrow}{\mu}(}o_2\bm{)}} \propto \sum_{s_2} p(o_2|s_2) \textcolor{black}{\bm{{\downarrow}{\mu}(}s_2\bm{)}}$
        \State $\textcolor{accblue}{q(s_0)} \propto \textcolor{black}{\bm{\overleftarrow{\mu}(}s_0\bm{)}} \textcolor{black}{\bm{\overrightarrow{\mu}(}s_0\bm{)}}$,\quad$\textcolor{accblue}{q(s^{\scriptscriptstyle I}_0)} \propto \textcolor{black}{\bm{\overleftarrow{\mu}(}s^{\scriptscriptstyle I}_0\bm{)}} \textcolor{black}{\bm{\overrightarrow{\mu}(}s^{\scriptscriptstyle I}_0\bm{)}}$,\quad$\textcolor{accblue}{q(s^{\scriptscriptstyle II}_0)} \propto \textcolor{black}{\bm{{\uparrow}{\mu}(}s^{\scriptscriptstyle II}_0\bm{)}} \textcolor{black}{\bm{{\downarrow}{\mu}(}s^{\scriptscriptstyle II}_0\bm{)}}$,
        \Statex $\textcolor{accblue}{q(s_1)} \propto \textcolor{black}{\bm{\overleftarrow{\mu}(}s_1\bm{)}} \textcolor{black}{\bm{\overrightarrow{\mu}(}s_1\bm{)}}$,\quad$\textcolor{accblue}{q(s^{\scriptscriptstyle I}_1)} \propto \textcolor{black}{\bm{\overleftarrow{\mu}(}s^{\scriptscriptstyle I}_1\bm{)}} \textcolor{black}{\bm{\overrightarrow{\mu}(}s^{\scriptscriptstyle I}_1\bm{)}}$,\quad$\textcolor{accblue}{q(s^{\scriptscriptstyle II}_1)} \propto \textcolor{black}{\bm{{\uparrow}{\mu}(}s^{\scriptscriptstyle II}_1\bm{)}} \textcolor{black}{\bm{{\downarrow}{\mu}(}s^{\scriptscriptstyle II}_1\bm{)}}$,
        \Statex $\textcolor{accblue}{q(s_2)} \propto \textcolor{black}{\bm{\overleftarrow{\mu}(}s_2\bm{)}} \textcolor{black}{\bm{{\downarrow}{\mu}(}s_2\bm{)}}$,\quad$\textcolor{accblue}{q(a_0)} \propto \textcolor{black}{\bm{{\uparrow}{\mu}(}a_0\bm{)}} \textcolor{black}{\bm{{\downarrow}{\mu}(}a_0\bm{)}}$,\quad$\textcolor{accblue}{q(a_1)} \propto \textcolor{black}{\bm{{\uparrow}{\mu}(}a_1\bm{)}} \textcolor{black}{\bm{{\downarrow}{\mu}(}a_1\bm{)}}$,
        \Statex $\textcolor{accblue}{q(o_0)} \propto \textcolor{black}{\bm{{\uparrow}{\mu}(}o_0\bm{)}} \textcolor{black}{\bm{{\downarrow}{\mu}(}o_0\bm{)}}$,\quad$\textcolor{accblue}{q(o_1)} \propto \textcolor{black}{\bm{{\uparrow}{\mu}(}a_1\bm{)}} \textcolor{black}{\bm{{\downarrow}{\mu}(}a_1\bm{)}}$,\quad$\textcolor{accblue}{q(o_2)} \propto \textcolor{black}{\bm{{\uparrow}{\mu}(}a_2\bm{)}} \textcolor{black}{\bm{{\downarrow}{\mu}(}a_2\bm{)}}$
    \end{algorithmic}
\end{algorithm}



\subsection{Generalized Free Energy (GFE) Agent}
\label{sec:gfe}

In this section I introduce the GFE agent. It can be considered a real Active Inference agent because it does not solely minimize the Bethe free energy but also respects additional motives that lead to information seeking. This makes up for the generalized free energy that Active Inference agents are mandated to minimize. The GFE agent is indeed similar to the BFE agent and is also based on the same generative model (\cref{eq:15}) but instead of only realizing preferred observations it also seeks information in the environment, i.e., it is curious. In other words, during inference it infers actions that do not only lead to preferred observation \(o\) but which are also informative regarding the hidden states \(s\) of its generative model. The amount of information that a certain observation yields with respect to its hidden state (that the agent assumes to generate this observation via its generative model) can be quantified via the mutual information between these two variables. Formally, the GFE agent is driven to maximize the mutual information between future hidden state beliefs \(q(s)\) and observation beliefs \(q(o)\). This mechanism enables it to sample actions at present time that maximize this quantity. The mechanism is encoded in the free energy functional (Lagrangian) of the agent and leads to alternative messages at the edges in the corresponding factor graph that resemble future hidden states. Moreover, we will see that messages at edges of future observation beliefs disappear under this motive. Importantly, the mutual information mechanism is only applied to future parts of the factor graph, i.e., it does not apply to parts of the graph that resemble the present or past. Intuitively, the agent is only interested in informative observations that are yet to be observed, i.e., that are in the future. This circumstance comes into play during the lifecycle of the agent, which leads to alternative messages at edges that resemble hidden states and observations depending on the timestep (future timestep or present/past timestep). Moreover, messages resulting from likelihood nodes that resemble the present or past, i.e., where an observation is available, will also take on an alternative form because they will be based on a local mean-field approximation of belief distributions at these nodes. In the case of the GFE agent this will not make a difference in comparison to using messages derived from the Bethe free energy (without mean-field assumption) from \cref{sec:1}. For the later presented Context-GFE agent this will be important to break cyclic dependencies introduced by the context state (see \cref{sec:context-bfe} for more explanation). To allow for comparison between the GFE- and Context-GFE agent both include this mean-field approximation for belief distributions of present and past likelihood nodes. The resulting messages from this assumption are variational messages (presented below).

\subsubsection{GFE-based Messages}

In this section I derive the alternative (GFE-based) message (\(\textcolor{accorange}{\bm{{\uparrow}{\mu}(}s\bm{)}}\) in \cref{fig:32.2}) that is sent by future likelihood nodes upwards to transition nodes and enables the agent with information seeking.
\begin{figure}[htbp]
    \centering
    \begin{subfigure}[t]{0.48\textwidth}
        \centering
        \begin{tikzpicture}[
                factor/.style={draw=black, thick, rectangle, minimum size=0.8cm, font=\normalsize},
                eqfactor/.style={draw=black, thick, rectangle, minimum size=0.5cm, font=\normalsize},
                edge/.style={thick},
                varlabel/.style={font=\normalsize},
                greycircle/.style={draw=black, thick, circle, fill=gray!30, minimum size=0.5cm, font=\normalsize}
            ]

            % Knoten
            \node[factor, label={[varlabel]left:$q_a(s,o)$}] (A1) {$p(o|s)$};
            \node[factor, below=1.0cm of A1, label={[varlabel]left:$q_b(o)$}] (c1) {$\tilde{p}(o)$};

            % Kanten zwischen Knoten
            \draw[edge] ([yshift=1.5cm]A1.north) -- (A1.north) node[midway, left, varlabel] {$q_i(s)$};
            \draw[edge] (A1.south) -- (c1.north) node[midway, left, varlabel] {$q_j(o)$};

        \end{tikzpicture}
        \caption{Perception-related factor graph sector with Bethe free energy.}
        \label{fig:32.1}
    \end{subfigure}
    \hfill
    \begin{subfigure}[t]{0.48\textwidth}
        \centering
        \begin{tikzpicture}[
                factor/.style={draw=black, thick, rectangle, minimum size=0.8cm, font=\normalsize},
                eqfactor/.style={draw=black, thick, rectangle, minimum size=0.5cm, font=\normalsize},
                edge/.style={thick},
                varlabel/.style={font=\normalsize},
                greycircle/.style={draw=black, thick, circle, fill=gray!30, minimum size=0.5cm, font=\normalsize}
            ]

            % Knoten
            \node[factor, label={[varlabel]left:$q_a(s,o)$}] (A1) {$p(o|s)$};
            \node[factor, below=1.0cm of A1, label={[varlabel]left:$q_b(o)$}] (c1) {$\tilde{p}(o)$};

            % Kanten zwischen Knoten
            \draw[edge] ([yshift=1.5cm]A1.north) -- (A1.north) node[midway, left, varlabel] {$q_i(s)$} node[pos=0.25, right, varlabel, scale=1.0] {$\textcolor{accorange}{\bm{{\uparrow}{\mu}(}s\bm{)}}$} node[pos=0.75, right, varlabel, scale=1.0] {$\bm{{\downarrow}{\mu}(}s\bm{)}$} node[pos=1, circle, draw=black, fill=white, inner sep=2pt] {};
            \draw[edge] (A1.south) -- (c1.north) node[midway, left, varlabel] {$q_j(o)$} node[pos=0, rectangle, draw=black, fill=white, inner sep=2.5pt] {};

        \end{tikzpicture}
        \caption{Perception-related factor graph sector with generalized free energy and corresponding GFE-based message (\(\textcolor{accorange}{\bm{{\uparrow}{\mu}(}s\bm{)}}\)).}
        \label{fig:32.2}
    \end{subfigure}
    \caption{Perception-related factor graph sectors.}
    \label{fig:32}
\end{figure}

Starting point for this derivation is the Lagrangian \(L\) of this perception-related factor graph sector with Bethe free energy (\cref{fig:32.1}) which is also used by the BFE agent:
\begin{equation}
    \begin{split}
        L = & \sum_{o,s} q_a(o,s) \ln \frac{q_a(o,s)}{p(o|s)} + \sum_o q_b(o) \ln \frac{q_b(o)}{\tilde{p}(o)} + \sum_o q_j(o) \ln \frac{1}{q_j(o)} \\
            & + \sum_s \lambda_{ia}(s) \left( q_i(s) - \sum_o q_a(o,s) \right) + \sum_o \lambda_{ja}(o) \left( q_j(o) - \sum_s q_a(o,s) \right)    \\
            & + \sum_o \lambda_{jb}(o) \left( q_j(o) - q_b(o) \right)
    \end{split}
\end{equation}
Note, that I refrain from including the normalization constraints here define following distributions via the proportional (\(\propto\)) sign.
Subtracting the entropy of \(q_i(s)\) is not necessary because it appears just once, i.e., in the local free energy of \(q_a\) (see ...).

To derive the GFE-based message there are two changes necessary in the Lagrangian \(L\). First, we add a negative mutual information term. Note, that the resulting message passing algorithm achieves a minimization of the Lagrangian which means the mutual information must have a negative sign to enable the agent to maximize the mutual information. The resulting intermediate Lagrangian \(L_{\text{MI}}\) is given in the following:
\begin{equation}
    L_{\text{MI}} = L + \sum_{o,s} q_a(o,s) \ln \frac{q_a(o)q_a(s)}{q_a(o,s)}
\end{equation}
Before applying the second change we first rearrange \(L_{\text{MI}}\) to a shorter form:
\begin{align}
    L_{\text{MI}} = & \sum_{o,s} q_a(o,s) \ln \frac{q_a(o)q_a(s)}{p(o|s)} + \sum_o q_b(o) \ln \frac{q_b(o)}{\tilde{p}(o)} + \sum_o q_j(o) \ln \frac{1}{q_j(o)} \nonumber \\
                    & + \sum_s \lambda_{ia}(s) \left( q_i(s) - \sum_o q_a(o,s) \right) + \sum_o \lambda_{ja}(o) \left( q_j(o) - \sum_s q_a(o,s) \right) \nonumber        \\
                    & + \sum_o \lambda_{jb}(o) \left( q_j(o) - q_b(o) \right)       \label{eq:37}                                                                        \\
    =               & \sum_{o,s} q_a(o,s) \ln \frac{q_a(o)q_a(s)}{p(o|s)} + \sum_o q_a(o) \ln \frac{q_a(o)}{\tilde{p}(o)} + \sum_o q_a(o) \ln \frac{1}{q_a(o)} \nonumber \\
                    & + \sum_s \lambda_{ia}(s) \left( q_i(s) - \sum_o q_a(o,s) \right) \label{eq:38}                                                                     \\
    =               & \sum_{o,s} q_a(o|s)q_a(s) \ln \frac{q_a(o)q_a(s)}{p(o|s)\tilde{p}(o)} + \sum_s \lambda_{ia}(s) \left( q_i(s) - q_a(s) \right) \label{eq:39}
\end{align}
In \cref{eq:37} I merged the local free energy of node \(a\) with the mutual information term. In \cref{eq:38} I utilized the equivalence of \(q_b(o)\), \(q_j(o)\), and \(q_a(o)\) as given by the respective marginalization constraints to substitute all \(q_b(o)\) and \(q_j(o)\) with \(q_a(o)\). In \cref{eq:39} I merged the first three terms from \cref{eq:38}, factorized \(q_a(s,o)\) in the expectation of the resulting term, and solved the sum in the remaining marginalization constraint term.

The second change consists of a "p-substitution" in which \(q_a(o|s)\) is substituted by the real likelihood \(p(o|s)\) in \cref{eq:39}. This leads to the Lagrangian \(L_{\text{GFE}}\) that entails the generalized free energy of the corresponding perception-related sector in the factor graph:
\begin{align}
    L_{\text{GFE}} = \sum_{o,s} p(o|s)q_a(s) \ln \frac{q_a(o)q_a(s)}{p(o|s)\tilde{p}(o)} + \sum_s \lambda_{ia}(s) \left( q_i(s) - q_a(s) \right)
\end{align}

The two changes to the Lagrangian are indicated in corresponding factor graph (\cref{fig:32.2}) by a small bead and a small square where the two edges terminate at node \(a\). In general, having beads at both ends of the edges means that the node local belief distribution \(q_a(o,s)\) assumes a mean-field approximation, i.e, \(q_a(o,s) = q_a(o)q_a(s)\), which is indeed the case in \(\sum_{o,s} q_a(o,s) \ln \frac{q_a(o)q_a(s)}{p(o|s)}\). Moreover, in the case when one of the variables' belief distributions (e.g., \(q_a(o)\)) is substituted by the real likelihood (i.e, \(p(o|s)\)) in the expection of the local free energy (\(\sum_{o,s} p(o|s)q_a(s) \ln \frac{q_a(o)q_a(s)}{p(o|s)}\)), then the corresponding bead of the edge that relates to this variable (edge \(j\) relates to variable \(o\)) is replaced by a rectangle.

Next, \(F_{\text{GFE}}\) is rearranged to a form that facilitates following derivations:
\begin{align}
    L_{\text{GFE}} = \sum_{o} p(o|s) \ln \frac{q_a(o)}{p(o|s)\tilde{p}(o)} + \sum_s q_a(s) \ln q_a(s) + \sum_s \lambda_{ia}(s) \left( q_i(s) - q_a(s) \right)
\end{align}

Next, the Lagrangian is differentiated with respect to \(q_a(s)\), set to \(0\), and finally solved for \(q^*_a(s)\):
\begin{alignat}{3}
     & \frac{\delta L_{\text{GFE}}}{\delta q_a(s)} &  & = \sum_{o} p(o|s) \ln \frac{q_a(o)}{p(o|s)\tilde{p}(o)} + \ln q_a(s) + 1 - \lambda_{ia}(s) \overset{!}{=} 0                                \\
     &                                             &  & \Rightarrow q^*_a(s) \propto \exp \left( \sum_{o} p(o|s) \ln \frac{p(o|s)\tilde{p}(o)}{q_a(o)} \right) \exp \left( \lambda_{ia}(s) \right)
\end{alignat}

Finally, we use the result of \cref{sec:1} for the local stationary point of \(q_i(s)\), which is:
\begin{equation}
    q^*_i(s) \propto \exp \left( \lambda_{ia}(s) \right) \exp \left( \lambda_{i\bullet}(s) \right)
\end{equation}
Here, \(\lambda_{i\bullet}\) is the Lagrange multiplier that enforces marginalization between the belief \(q_i(s)\) and the belief of the other node (besided node \(a\)) that the edge \(i\) is connected to (node not shown in \cref{fig:32}), e.g., the transition node. In the next step we will see that \(\exp \left( \lambda_{i\bullet}(s) \right)\) is the GFE-based message, i.e., the central result of this whole derivation.

To this end, we use the marginalization constraint \(q_i(s) \overset{!}{=} \sum_o q_a(o,s)\), i.e., \(q_i(s) \overset{!}{=} q_a(s)\) to relate both stationary solutions \(q^*_i(s)\) and \(q^*_a(s)\):
\begin{alignat}{2}
     & q^*_i(s)                                                                      &  & \overset{!}{=} q^*_a(s)                                                                                               \\
     & \exp \left( \lambda_{ia}(s) \right) \exp \left( \lambda_{i\bullet}(s) \right) &  & \propto \exp \left( \sum_{o} p(o|s) \ln \frac{p(o|s)\tilde{p}(o)}{q_a(o)} \right) \exp \left( \lambda_{ia}(s) \right) \\
     & \exp \left( \lambda_{i\bullet}(s) \right)                                     &  & \propto \exp \left( \sum_{o} p(o|s) \ln \frac{p(o|s)\tilde{p}(o)}{q_a(o)} \right)                                     \\
     & \textcolor{accorange}{\bm{{\uparrow}{\mu}(}s\bm{)}}                                                  &  & \propto \exp \left( \sum_{o} p(o|s) \ln \frac{p(o|s)\tilde{p}(o)}{q_a(o)} \right) \label{eq:40}
\end{alignat}

Finally, the GFE-based message that is sent by the likehood node \(a\) upwards in \cref{fig:32.2} is given in \cref{eq:40}.

Although \cref{eq:40} yields the necessary message that enables the agent to seek information in the environment, in practice the message can lead to convergence issues during inference (message passing). To accomodate for this problem we derive an alternative message that leads to a more robust message passing algorithm. In general this is achieved by developing an inference procedere that is local to the perception-related factor graph sector. During this procedere \(q_i(s)\) and \(q_j(o)\) are optimized (inferred) until convergence and only then we compute an outgoing message (analog to \(\textcolor{accorange}{\bm{{\uparrow}{\mu}(}s\bm{)}}\) in \cref{fig:32.2}) that is sent to the rest of the graph. The local inference loop consists of three recursive steps. First, we use the current (at loop step \(l-1\)) \(q^{(l-1)}_i(s)\) to compute \(q^{(l)}_j(o)\) via:
\begin{align}
    q^{(l)}_j(o) \propto \sum_s p(o|s) q^{(l-1)}_i(s) \label{eq:41}
\end{align}

Second, this result is used to compute the message from \cref{eq:40}:
\begin{align}
    \textcolor{accorange}{\bm{{\uparrow}{\mu}}^{(l)}\bm{(}s\bm{)}} \propto \exp \left( \sum_{o} p(o|s) \ln \frac{p(o|s)\tilde{p}(o)}{q_a^{(k)}(o)} \right) \label{eq:42}
\end{align}
Here, \(q_a^{(l)}(o) = q^{(l)}_j(o)\).

Third, this upward message together with the downward message \(\bm{{\downarrow}{\mu}(}s\bm{)}\) coming from the rest of the graph leads to an updated state belief:
\begin{align}
    q^{(l)}_i(s) \propto \textcolor{accorange}{\bm{{\uparrow}{\mu}}^{(l)}\bm{(}s\bm{)}} \bm{{\downarrow}{\mu}(}s\bm{)} \label{eq:43}
\end{align}

\cref{eq:41,eq:42,eq:43} are iterated until convergence of \(q_i(s)\). As a last step, the converged \(q^*_i(s)\) (signified by \(^*\)) is used to compute a final upward message (analog to \cref{eq:40}) that leads to a more stable message passing algorithm:
\begin{align}
    \textcolor{accorange}{\bm{{\uparrow}{\mu}(}s\bm{)}} \propto \frac{q^*_i(s)}{\bm{{\downarrow}{\mu}(}s\bm{)}}
\end{align}

\subsubsection{Variational Messages}

In this section I derive the variational message (\(\textcolor{accorange}{\bm{{\uparrow}{\mu}(}s\bm{)}}\)) that results from a local mean-field approximation \(q_a(o,s) = q_a(o)q_a(s)\) at the node \(a\) in \cref{fig:37.2}. This message will be sent by likelihood nodes upwards to the rest of the graph that resemble present and past timesteps of the GFE- and Context-GFE agents. In principle, also the message that is sent downwards to the observation edge will take on the same form but in this work there will always be a delta constraint (from the observation) that effectively blocks this downward message.
\begin{figure}[htbp]
    \centering
    \begin{subfigure}[t]{0.48\textwidth}
        \centering
        \begin{tikzpicture}[
                factor/.style={draw=black, thick, rectangle, minimum size=0.8cm, font=\normalsize},
                eqfactor/.style={draw=black, thick, rectangle, minimum size=0.5cm, font=\normalsize},
                edge/.style={thick},
                varlabel/.style={font=\normalsize},
                greycircle/.style={draw=black, thick, circle, fill=gray!30, minimum size=0.5cm, font=\normalsize}
            ]

            % Knoten
            \node[factor, label={[varlabel]left:$q_a(o,s)$}] (A1) {$p(o|s)$};

            % Kanten zwischen Knoten
            \draw[edge] ([yshift=1.5cm]A1.north) -- (A1.north) node[midway, left, varlabel] {$q_i(s)$};
            \draw[edge] (A1.south) -- ([yshift=-1.5cm]A1.south) node[midway, left, varlabel] {$q_j(o)$};

        \end{tikzpicture}
        \caption{Likelihood node with Bethe free energy.}
        \label{fig:37.1}
    \end{subfigure}
    \hfill
    \begin{subfigure}[t]{0.48\textwidth}
        \centering
        \begin{tikzpicture}[
                factor/.style={draw=black, thick, rectangle, minimum size=0.8cm, font=\normalsize},
                eqfactor/.style={draw=black, thick, rectangle, minimum size=0.5cm, font=\normalsize},
                edge/.style={thick},
                varlabel/.style={font=\normalsize},
                greycircle/.style={draw=black, thick, circle, fill=gray!30, minimum size=0.5cm, font=\normalsize}
            ]

            % Knoten
            \node[factor, label={[varlabel]left:$q_a(o)q_a(s)$}] (A1) {$p(o|s)$};

            % Kanten zwischen Knoten
            \draw[edge] ([yshift=1.5cm]A1.north) -- (A1.north) node[midway, left, varlabel] {$q_i(s)$} node[pos=0.25, right, varlabel, scale=1.0] {$\textcolor{accorange}{\bm{{\uparrow}{\mu}(}s\bm{)}}$} node[pos=1, circle, draw=black, fill=white, inner sep=2pt] {};
            \draw[edge] (A1.south) -- ([yshift=-1.5cm]A1.south) node[midway, left, varlabel] {$q_j(o)$} node[pos=0, circle, draw=black, fill=white, inner sep=2pt] {};

        \end{tikzpicture}
        \caption{Likelihood node with mean-field approximation between \(s\) and \(o\) (indicated by the two beads) and corresponding variational message (\(\textcolor{accorange}{\bm{{\uparrow}{\mu}(}s\bm{)}}\)).}
        \label{fig:37.2}
    \end{subfigure}
    \caption{Likelihood nodes for the derivation of variational messages.}
    \label{fig:37}
\end{figure}

Starting point for this derivation is the Bethe-based Lagrangian of the factor graph in \cref{fig:37.1} (excluding the marginalization constraint for \(q_j(o)\) and normalization constraints):
\begin{alignat}{2}
    L = \sum_{o,s} q_a(o,s) \ln \frac{q_a(o,s)}{p(o|s)} + \sum_s \lambda_{ia}(s) \left( q_i(s) - \sum_o q_a(o,s) \right)
\end{alignat}

Next, we introduce the mean-field approximation, i.e., \(q_a(o,s) = q_a(o)q_a(s)\), which leads to a Lagrangian that supports variational message passing (VMP):
\begin{alignat}{2}
    L_{\text{VMP}} = \sum_{o,s} q_a(o)q_a(s) \ln \frac{q_a(o)q_a(s)}{p(o|s)} + \sum_s \lambda_{ia}(s) \left( q_i(s) - \sum_o q_a(o)q_a(s) \right)
\end{alignat}

Next, we differentiate \(L_{\text{VMP}}\) with respect to \(q_a(o)\) and \(q_a(s)\), set the results to \(0\), and rearrange for \(q^*_a(o)\) and \(q^*_a(s)\):
\begin{alignat}{3}
    & \frac{\delta L_{\text{VMP}}}{\delta q_a(s)} &  & = \ln q_a(s) + 1 - \sum_o q_a(o) \ln p(o|s) - \lambda_{ia}(s) \overset{!}{=} 0                                \\
    &                                             &  & \Rightarrow q^*_a(s) \propto \exp \left( \sum_o q_a(o) \ln p(o|s) \right) \exp \left( \lambda_{ia}(s) \right)
\end{alignat}

Finally, we use the marginalization constraint (\(q_i(s) = q_a(s)\)) and the result of \cref{sec:1} for \(q_i(s)\) to derive the message \(\textcolor{accorange}{\bm{{\uparrow}{\mu}(}s\bm{)}}\):
\begin{alignat}{2}
    & q^*_i(s)                                                                      &  & \overset{!}{=} q^*_a(s)                                                                                               \\
    & \exp \left( \lambda_{ia}(s) \right) \exp \left( \lambda_{i\bullet}(s) \right) &  & \propto \exp \left( \sum_o q_a(o) \ln p(o|s) \right) \exp \left( \lambda_{ia}(s) \right) \\
    & \exp \left( \lambda_{i\bullet}(s) \right)                                     &  & \propto \exp \left( \sum_o q_a(o) \ln p(o|s) \right)         \\
    & \textcolor{accorange}{\bm{{\uparrow}{\mu}(}s\bm{)}}                                                  &  & \propto \exp \left( \sum_o q_a(o) \ln p(o|s) \right)
\end{alignat}


\subsubsection{Message Passing Algorithm}

The resulting message passing algorithm that enables the agent to perform inference is based on the whole factor graph of its temporal generative model. The factor graph is very similar to the graph from the BFE agent but additionally shows small beads and sqares at future likelihood nodes that resemble the generalized free energy which is used in these parts of the graph and leads to the information seeking behavior of the agent. Again, the graph shows an exemplary integration of an observation \(\hat{o_0}\) as a delta constraint which leads to a variational upward message there. Crucially, using GFE-based- and variational messages leads to a message passing algorithm which necessitates a recursive execution to let all belief distribution converge to a stable solution. In contrast to the algorithm of the BFE agent, this leads to an algorithm that must be run repeatedly (signified by the arrow on the left side in \cref{alg:3}). Additionally, the recursive computation of the GFE-based messages is also shown with an arrow at the left side to display the loopy execution. For this, the messages \(\textcolor{orange!75!black}{\bm{{\downarrow}{\mu}(}s^{\scriptscriptstyle II}_1\bm{)}}\), \(\textcolor{orange!75!black}{\bm{{\downarrow}{\mu}(}s_2\bm{)}}\) and beliefs \(\textcolor{accblue}{q(s^{\scriptscriptstyle II}_1)}\), \(\textcolor{accblue}{q(s_2)}\) must be initialized (e.g., uniformly) before the execution is started.

\begin{figure}[htbp]
    \centering

    \resizebox{1.0\textwidth}{!}{%
        \begin{tikzpicture}[
                factor/.style={draw=black, thick, rectangle, minimum size=0.8cm, font=\normalsize},
                eqfactor/.style={draw=black, thick, rectangle, minimum size=0.5cm, font=\normalsize},
                edge/.style={thick},
                varlabel/.style={font=\normalsize},
                greycircle/.style={draw=black, thick, circle, fill=gray!30, minimum size=0.5cm, font=\normalsize}
            ]

            % Knoten
            \node[factor] (d) {$p(s_0)$};
            \node[eqfactor, right=1.8cm of d] (eq0) {$=$};
            \node[factor, below=1.5cm of eq0] (A0) {$p(o_0|s^{\scriptscriptstyle II}_0)$};
            \node[factor, below=1.5cm of A0] (c0) {$\tilde{p}(o_0)$};
            \node[factor, right=1.8cm of eq0] (B0) {$p(s_1|s^{\scriptscriptstyle I}_0,a_0)$};
            \node[eqfactor, right=1.8cm of B0] (eq1) {$=$};
            \node[factor, below=1.5cm of eq1] (A1) {$p(o_1|s^{\scriptscriptstyle II}_1)$};
            \node[factor, below=1.5cm of A1] (c1) {$\tilde{p}(o_1)$};
            \node[factor, right=1.8cm of eq1] (B1) {$p(s_2|s^{\scriptscriptstyle I}_1,a_1)$};
            \node[factor] (A2) at ([xshift=2cm] B1.east |- A1) {$p(o_2|s_2)$};
            \node[factor, below=1.5cm of A2] (c2) {$\tilde{p}(o_2)$};
            \node[factor, above=1.5cm of B0] (u0) {$p(a_0)$};
            \node[factor, above=1.5cm of B1] (u1) {$p(a_1)$};

            % Kanten zwischen Knoten
            \draw[edge] (d.east) -- (eq0.west) node[midway, above, varlabel] {$\textcolor{accblue}{q(s_0)}$} node[near end, below, varlabel, scale=0.7] {${\large \textcolor{accorange}{\bm{\overrightarrow{\mu}(}s_0\bm{)}}}$} node[near start, below, varlabel, scale=0.7] {${\large \textcolor{accorange}{\bm{\overleftarrow{\mu}(}s_0\bm{)}}}$};
            \draw[edge] (eq0.south) -- (A0.north) node[midway, left, varlabel] {$\textcolor{accblue}{q(s^{\scriptscriptstyle II}_0)}$} node[pos=0.25, right, varlabel, scale=0.7] {${\large \textcolor{accorange}{\bm{{\uparrow}{\mu}(}s^{\scriptscriptstyle II}_0\bm{)}}}$} node[pos=0.75, right, varlabel, scale=0.7] {${\large \textcolor{accorange}{\bm{{\downarrow}{\mu}(}s^{\scriptscriptstyle II}_0\bm{)}}}$}  node[pos=1, circle, draw=black, fill=white, inner sep=2pt] {};
            \draw[edge] (A0.south) -- (c0.north) node[midway, greycircle] {$\delta$} node[midway, left=0.3, varlabel] {$\textcolor{accblue}{q(o_0)}$} node[pos=0.25, right=0.25, varlabel, scale=0.7] {${\large \textcolor{accorange}{\bm{{\uparrow}{\mu}(}o_0\bm{)}}}$} node[pos=0.75, right=0.25, varlabel, scale=0.7] {${\large \textcolor{accorange}{\bm{{\downarrow}{\mu}(}o_0\bm{)}}}$}  node[pos=0, circle, draw=black, fill=white, inner sep=2pt] {};
            \draw[edge] (eq0.east) -- (B0.west) node[midway, above, varlabel] {$\textcolor{accblue}{q(s^{\scriptscriptstyle I}_0)}$} node[near start, below, varlabel, scale=0.7] {${\large \textcolor{accorange}{\bm{\overleftarrow{\mu}(}s^{\scriptscriptstyle I}_0\bm{)}}}$} node[near end, below, varlabel, scale=0.7] {${\large \textcolor{accorange}{\bm{\overrightarrow{\mu}(}s^{\scriptscriptstyle I}_0\bm{)}}}$};
            \draw[edge] (B0.north) -- (u0.south) node[midway, left, varlabel] {$\textcolor{accblue}{q(a_0)}$} node[pos=0.25, right, varlabel, scale=0.7] {${\large \textcolor{accorange}{\bm{{\downarrow}{\mu}(}a_0\bm{)}}}$} node[pos=0.75, right, varlabel, scale=0.7] {${\large \textcolor{accorange}{\bm{{\uparrow}{\mu}(}a_0\bm{)}}}$};
            \draw[edge] (B0.east) -- (eq1.west) node[midway, above, varlabel] {$\textcolor{accblue}{q(s_1)}$} node[near end, below, varlabel, scale=0.7] {${\large \textcolor{accorange}{\bm{\overrightarrow{\mu}(}s_1\bm{)}}}$} node[near start, below, varlabel, scale=0.7] {${\large \textcolor{accorange}{\bm{\overleftarrow{\mu}(}s_1\bm{)}}}$};
            \draw[edge] (eq1.south) -- (A1.north) node[midway, left, varlabel] {$\textcolor{accblue}{q(s^{\scriptscriptstyle II}_1)}$} node[pos=0.25, right, varlabel, scale=0.7] {${\large \textcolor{orange!75!black}{\bm{{\uparrow}{\mu}(}s^{\scriptscriptstyle II}_1\bm{)}}}$} node[pos=0.75, right, varlabel, scale=0.7] {${\large \textcolor{orange!75!black}{\bm{{\downarrow}{\mu}(}s^{\scriptscriptstyle II}_1\bm{)}}}$} node[pos=1, circle, draw=black, fill=white, inner sep=2pt] {};
            \draw[edge] (A1.south) -- (c1.north) node[midway, left, varlabel] {$\textcolor{accblue}{q(o_1)}$} node[pos=0, rectangle, draw=black, fill=white, inner sep=2.5pt] {};
            \draw[edge] (eq1.east) -- (B1.west) node[midway, above, varlabel] {$\textcolor{accblue}{q(s^{\scriptscriptstyle I}_1)}$} node[near start, below, varlabel, scale=0.7] {${\large \textcolor{accorange}{\bm{\overleftarrow{\mu}(}s^{\scriptscriptstyle I}_1\bm{)}}}$} node[near end, below, varlabel, scale=0.7] {${\large \textcolor{accorange}{\bm{\overrightarrow{\mu}(}s^{\scriptscriptstyle I}_1\bm{)}}}$};
            \draw[edge] (B1.east) -| (A2.north) node[pos=0.25, above, varlabel] {$\textcolor{accblue}{q(s_2)}$} node[pos=0.89, right, varlabel, scale=0.7] {${\large \textcolor{accorange}{\bm{{\downarrow}{\mu}(}s_2\bm{)}}}$} node[pos=0.11, below, varlabel, scale=0.7] {${\large \textcolor{accorange}{\bm{\overleftarrow{\mu}(}s_2\bm{)}}}$} node[pos=1, circle, draw=black, fill=white, inner sep=2pt] {};
            \draw[edge] (A2.south) -- (c2.north) node[midway, left, varlabel] {$\textcolor{accblue}{q(o_2)}$} node[pos=0, rectangle, draw=black, fill=white, inner sep=2.5pt] {};
            \draw[edge] (B1.north) -- (u1.south) node[midway, left, varlabel] {$\textcolor{accblue}{q(a_1)}$} node[pos=0.25, right, varlabel, scale=0.7] {${\large \textcolor{accorange}{\bm{{\downarrow}{\mu}(}a_1\bm{)}}}$} node[pos=0.75, right, varlabel, scale=0.7] {${\large \textcolor{accorange}{\bm{{\uparrow}{\mu}(}a_1\bm{)}}}$};

        \end{tikzpicture}
    }

    \caption{Factor graph of the GFE agent with variables' belief distributions \(\textcolor{accblue}{q}\) and messages \(\textcolor{accorange}{\mu}\).}
    \label{fig:33}
\end{figure}

\begin{algorithm}[htbp]
    \caption{Message Passing Algorithm of the GFE Agent}
    \label{alg:3}
    \small
    \begin{algorithmic}[1] % Die [1] sorgt für automatische Zeilennummern!
        \linespread{1.4}\selectfont
        \setlength{\leftskip}{0.3cm}
        \State \tikzmark{algtop3}$\textcolor{accorange}{\bm{\overrightarrow{\mu}(}s_0\bm{)}} \propto p(s_0)$,\quad$\textcolor{accorange}{\bm{{\uparrow}{\mu}(}o_0\bm{)}} \propto \delta_{o_0,\hat{o_0}}$,\quad$\textcolor{accorange}{\bm{{\downarrow}{\mu}(}a_0\bm{)}} \propto p(a_0)$,\quad$\textcolor{accorange}{\bm{{\downarrow}{\mu}(}a_1\bm{)}} \propto p(a_1)$,\quad$\textcolor{accorange}{\bm{{\downarrow}{\mu}(}o_0\bm{)}} \propto \delta_{o_0,\hat{o_0}}$
        \State $\textcolor{accorange}{\bm{{\uparrow}{\mu}(}s^{\scriptscriptstyle II}_0\bm{)}} \propto \exp \left( \sum_{o_0} q(o_0) p(o_0|s^{\scriptscriptstyle II}_0) \right)$,
        \Statex \quad\quad\tikzmark{onetop3}$q^{(l)}(o_1) \propto \sum_{s^{\scriptscriptstyle II}_1} p(o_1|s^{\scriptscriptstyle II}_1) q^{(l-1)}(s^{\scriptscriptstyle II}_1)$
        \Statex \quad\quad$\bm{{\uparrow}{\mu}}^{(l)}\bm{(}s^{\scriptscriptstyle II}_1\bm{)} \propto \exp \left( \sum_{o_1} p(o_1|s^{\scriptscriptstyle II}_1) \ln \frac{p(o_1|s^{\scriptscriptstyle II}_1)\tilde{p}(o_1)}{q^{(l)}(o_1)} \right)$
        \Statex \quad\quad\tikzmark{onebottom3}$q^{(l)}(s^{\scriptscriptstyle II}_1) \propto \bm{{\uparrow}{\mu}}^{(l)}\bm{(}s^{\scriptscriptstyle II}_1\bm{)} \bm{{\downarrow}{\mu}(}s^{\scriptscriptstyle II}_1\bm{)}$
        \Statex $\textcolor{accorange}{\bm{{\uparrow}{\mu}(}s^{\scriptscriptstyle II}_1\bm{)}} \propto \frac{q^*(s^{\scriptscriptstyle II}_1)}{\bm{{\downarrow}{\mu}(}s^{\scriptscriptstyle II}_1\bm{)}}$,
        \Statex \quad\quad\tikzmark{twotop3}$q^{(l)}(o_2) \propto \sum_{s_2} p(o_2|s_2) q^{(l-1)}(s_2)$
        \Statex \quad\quad$\bm{{\uparrow}{\mu}}^{(l)}\bm{(}s_2\bm{)} \propto \exp \left( \sum_{o_2} p(o_2|s_2) \ln \frac{p(o_2|s_2)\tilde{p}(o_2)}{q^{(l)}(o_2)} \right)$
        \Statex \quad\quad\tikzmark{twobottom3}$q^{(l)}(s_2) \propto \bm{{\uparrow}{\mu}}^{(l)}\bm{(}s_2\bm{)} \bm{{\downarrow}{\mu}(}s_2\bm{)}$
        \Statex $\textcolor{accorange}{\bm{\overleftarrow{\mu}(}s_2\bm{)}} \propto \frac{q^*(s_2)}{\bm{{\downarrow}{\mu}(}s_2\bm{)}}$
        \State $\textcolor{accorange}{\bm{\overrightarrow{\mu}(}s^{\scriptscriptstyle I}_0\bm{)}} \propto \sum_{s_0,s^{\scriptscriptstyle II}_0} \delta_{s_0,s^{\scriptscriptstyle I}_0}\delta_{s_0,s^{\scriptscriptstyle II}_0} \textcolor{black}{\bm{\overrightarrow{\mu}(}s_0\bm{)}} \textcolor{black}{\bm{{\uparrow}{\mu}(}s^{\scriptscriptstyle II}_0\bm{)}}$,\quad$\textcolor{accorange}{\bm{\overleftarrow{\mu}(}s^{\scriptscriptstyle I}_1\bm{)}} \propto \sum_{a_1,s_2} p(s_2|s^{\scriptscriptstyle I}_1,a_1) \textcolor{black}{\bm{{\downarrow}{\mu}(}a_1\bm{)}} \textcolor{black}{\bm{\overleftarrow{\mu}(}s_2\bm{)}}$,
        \State $\textcolor{accorange}{\bm{\overrightarrow{\mu}(}s_1\bm{)}} \propto \sum_{s^{\scriptscriptstyle I}_0,a_0} p(s_1|s^{\scriptscriptstyle I}_0,a_0) \textcolor{black}{\bm{\overrightarrow{\mu}(}s^{\scriptscriptstyle I}_0\bm{)}} \textcolor{black}{\bm{{\downarrow}{\mu}(}a_0\bm{)}}$,\quad$\textcolor{accorange}{\bm{\overleftarrow{\mu}(}s_1\bm{)}} \propto \sum_{s^{\scriptscriptstyle I}_1,s^{\scriptscriptstyle II}_1} \delta_{s_1,s^{\scriptscriptstyle I}_1}\delta_{s_1,s^{\scriptscriptstyle II}_1} \textcolor{black}{\bm{\overleftarrow{\mu}(}s^{\scriptscriptstyle I}_1\bm{)}} \textcolor{black}{\bm{{\uparrow}{\mu}(}s^{\scriptscriptstyle II}_1\bm{)}}$
        \State $\textcolor{accorange}{\bm{\overrightarrow{\mu}(}s^{\scriptscriptstyle I}_1\bm{)}} \propto \sum_{s_1,s^{\scriptscriptstyle II}_1} \delta_{s_1,s^{\scriptscriptstyle I}_1}\delta_{s_1,s^{\scriptscriptstyle II}_1} \textcolor{black}{\bm{\overrightarrow{\mu}(}s_1\bm{)}} \textcolor{black}{\bm{{\uparrow}{\mu}(}s^{\scriptscriptstyle II}_1\bm{)}}$,\quad$\textcolor{accorange}{\bm{{\downarrow}{\mu}(}s^{\scriptscriptstyle II}_1\bm{)}} \propto \sum_{s_1,s^{\scriptscriptstyle I}_1} \delta_{s_1,s^{\scriptscriptstyle I}_1}\delta_{s_1,s^{\scriptscriptstyle II}_1} \textcolor{black}{\bm{\overrightarrow{\mu}(}s_1\bm{)}} \textcolor{black}{\bm{\overleftarrow{\mu}(}s^{\scriptscriptstyle I}_1\bm{)}}$,
        \Statex $\textcolor{accorange}{\bm{{\uparrow}{\mu}(}a_0\bm{)}} \propto \sum_{s_1,s^{\scriptscriptstyle I}_0} p(s_1|s^{\scriptscriptstyle I}_0,a_0) \textcolor{black}{\bm{\overrightarrow{\mu}(}s^{\scriptscriptstyle I}_0\bm{)}} \textcolor{black}{\bm{\overleftarrow{\mu}(}s_1\bm{)}}$,\quad$\textcolor{accorange}{\bm{\overleftarrow{\mu}(}s^{\scriptscriptstyle I}_0\bm{)}} \propto \sum_{s_1,a_0} p(s_1|s^{\scriptscriptstyle I}_0,a_0) \textcolor{black}{\bm{{\downarrow}{\mu}(}a_0\bm{)}} \textcolor{black}{\bm{\overleftarrow{\mu}(}s_1\bm{)}}$
        \State $\textcolor{accorange}{\bm{{\uparrow}{\mu}(}a_1\bm{)}} \propto \sum_{s^{\scriptscriptstyle I}_1,s_2} p(s_2|s^{\scriptscriptstyle I}_1,a_1) \textcolor{black}{\bm{\overrightarrow{\mu}(}s^{\scriptscriptstyle I}_1\bm{)}} \textcolor{black}{\bm{\overleftarrow{\mu}(}s_2\bm{)}}$,\quad$\textcolor{accorange}{\bm{{\downarrow}{\mu}(}s_2\bm{)}} \propto \sum_{s^{\scriptscriptstyle I}_1,a_1} p(s_2|s^{\scriptscriptstyle I}_1,a_1) \textcolor{black}{\bm{\overrightarrow{\mu}(}s^{\scriptscriptstyle I}_1\bm{)}} \textcolor{black}{\bm{{\downarrow}{\mu}(}a_1\bm{)}}$,
        \Statex $\textcolor{accorange}{\bm{\overleftarrow{\mu}(}s_0\bm{)}} \propto \sum_{s^{\scriptscriptstyle I}_0,s^{\scriptscriptstyle II}_0} \delta_{s_0,s^{\scriptscriptstyle I}_0}\delta_{s_0,s^{\scriptscriptstyle II}_0} \textcolor{black}{\bm{\overleftarrow{\mu}(}s^{\scriptscriptstyle I}_0\bm{)}} \textcolor{black}{\bm{{\uparrow}{\mu}(}s^{\scriptscriptstyle II}_0\bm{)}}$,\quad$\textcolor{accorange}{\bm{{\downarrow}{\mu}(}s^{\scriptscriptstyle II}_0\bm{)}} \propto \sum_{s_0,s^{\scriptscriptstyle I}_0} \delta_{s_0,s^{\scriptscriptstyle I}_0}\delta_{s_0,s^{\scriptscriptstyle II}_0} \textcolor{black}{\bm{\overrightarrow{\mu}(}s_0\bm{)}} \textcolor{black}{\bm{\overleftarrow{\mu}(}s^{\scriptscriptstyle I}_0\bm{)}}$
        \State \tikzmark{algbottom3}$\textcolor{accblue}{q(s_0)} \propto \textcolor{black}{\bm{\overleftarrow{\mu}(}s_0\bm{)}} \textcolor{black}{\bm{\overrightarrow{\mu}(}s_0\bm{)}}$,\quad$\textcolor{accblue}{q(s^{\scriptscriptstyle I}_0)} \propto \textcolor{black}{\bm{\overleftarrow{\mu}(}s^{\scriptscriptstyle I}_0\bm{)}} \textcolor{black}{\bm{\overrightarrow{\mu}(}s^{\scriptscriptstyle I}_0\bm{)}}$,\quad$\textcolor{accblue}{q(s^{\scriptscriptstyle II}_0)} \propto \textcolor{black}{\bm{{\uparrow}{\mu}(}s^{\scriptscriptstyle II}_0\bm{)}} \textcolor{black}{\bm{{\downarrow}{\mu}(}s^{\scriptscriptstyle II}_0\bm{)}}$,
        \Statex $\textcolor{accblue}{q(s_1)} \propto \textcolor{black}{\bm{\overleftarrow{\mu}(}s_1\bm{)}} \textcolor{black}{\bm{\overrightarrow{\mu}(}s_1\bm{)}}$,\quad$\textcolor{accblue}{q(s^{\scriptscriptstyle I}_1)} \propto \textcolor{black}{\bm{\overleftarrow{\mu}(}s^{\scriptscriptstyle I}_1\bm{)}} \textcolor{black}{\bm{\overrightarrow{\mu}(}s^{\scriptscriptstyle I}_1\bm{)}}$,\quad$\textcolor{accblue}{q(s^{\scriptscriptstyle II}_1)} \propto \textcolor{black}{\bm{{\uparrow}{\mu}(}s^{\scriptscriptstyle II}_1\bm{)}} \textcolor{black}{\bm{{\downarrow}{\mu}(}s^{\scriptscriptstyle II}_1\bm{)}}$,
        \Statex $\textcolor{accblue}{q(s_2)} \propto \textcolor{black}{\bm{\overleftarrow{\mu}(}s_2\bm{)}} \textcolor{black}{\bm{{\downarrow}{\mu}(}s_2\bm{)}}$,\quad$\textcolor{accblue}{q(a_0)} \propto \textcolor{black}{\bm{{\uparrow}{\mu}(}a_0\bm{)}} \textcolor{black}{\bm{{\downarrow}{\mu}(}a_0\bm{)}}$,\quad$\textcolor{accblue}{q(a_1)} \propto \textcolor{black}{\bm{{\uparrow}{\mu}(}a_1\bm{)}} \textcolor{black}{\bm{{\downarrow}{\mu}(}a_1\bm{)}}$,
        \Statex $\textcolor{accblue}{q(o_0)} \propto \textcolor{black}{\bm{{\uparrow}{\mu}(}o_0\bm{)}} \textcolor{black}{\bm{{\downarrow}{\mu}(}o_0\bm{)}}$,\quad$\textcolor{accblue}{q(o_1)} \propto \sum_{s^{\scriptscriptstyle II}_1} p(o_1|s^{\scriptscriptstyle II}_1) q(s^{\scriptscriptstyle II}_1)$,\quad$\textcolor{accblue}{q(o_2)} \propto \sum_{s_2} p(o_2|s_2) q(s_2)$
    \end{algorithmic}

    \begin{tikzpicture}[remember picture, overlay]
        \draw[-{Stealth[length=2.2mm, width=1.6mm]}, thick, gray!75!black, rounded corners=5pt]
        ([yshift=0.5ex, xshift=-2.5ex]pic cs:algbottom3) -- ++(-0.4cm, 0) |- ([xshift=-2.5ex, yshift=0.5ex]pic cs:algtop3);
    \end{tikzpicture}
    \begin{tikzpicture}[remember picture, overlay]
        \draw[-{Stealth[length=2.2mm, width=1.6mm]}, thick, gray!75!black, rounded corners=5pt]
        ([yshift=0.5ex, xshift=-0.5ex]pic cs:onebottom3) -- ++(-0.4cm, 0) |- ([xshift=-0.5ex, yshift=0.5ex]pic cs:onetop3);
    \end{tikzpicture}
    \begin{tikzpicture}[remember picture, overlay]
        \draw[-{Stealth[length=2.2mm, width=1.6mm]}, thick, gray!75!black, rounded corners=5pt]
        ([yshift=0.5ex, xshift=-0.5ex]pic cs:twobottom3) -- ++(-0.4cm, 0) |- ([xshift=-0.5ex, yshift=0.5ex]pic cs:twotop3);
    \end{tikzpicture}

\end{algorithm}


\subsection{Context-BFE Agent}
\label{sec:context-bfe}

The Context-BFE agent is similar to the BFE agent but entails another hidden context state as part of its generative model. Specifically, the agent does not only assume that observations are generated by hidden states \(s\) but additionaly assumes a hidden context state \(c\) that also influences observations. In contrast to states \(s\), the context state does not change for each timestep and thus is assumed to be persistent during multiple timesteps. In later simulations, there will only be a single context state during the succession of three regular hidden states \(s\). The corresponding POMDP is given by
\begin{equation}
    \begin{split}
        p(o_0, o_1, o_2, s_0, s_1, s_2, c, a_0, a_1) = \; & \tilde{p}(o_0) p(o_0|s_0,c) p(s_0) p(c)           \\
                                                          & \tilde{p}(o_1) p(o_1|s_1,c) p(s_1|s_0,a_0) p(a_0) \\
                                                          & \tilde{p}(o_2) p(o_2|s_2,c) p(s_2|s_1,a_1) p(a_1)
        \label{eq:35}
    \end{split}
\end{equation}
and entails a prior over the context \(p(c)\) similarly to the initial prior over the hidden state \(p(s_0)\). The corresponding factor graph of \cref{eq:35} is given in \cref{fig:28}.

\begin{figure}[htbp]
    \centering

    \resizebox{1.0\textwidth}{!}{%
        \begin{tikzpicture}[
                factor/.style={draw=black, thick, rectangle, minimum size=0.8cm, font=\normalsize},
                eqfactor/.style={draw=black, thick, rectangle, minimum size=0.5cm, font=\normalsize},
                edge/.style={thick},
                varlabel/.style={font=\normalsize},
                greycircle/.style={draw=black, thick, circle, fill=gray!30, minimum size=0.5cm, font=\normalsize}
            ]

            % Knoten
            \node[factor] (d) {$p(s_0)$};
            \node[eqfactor, right=1.8cm of d] (eq0) {$=$};
            \node[factor] (A0) at ([xshift=1.5cm, yshift=-2.15cm] d.east) {$p(o_0|s^{\scriptscriptstyle II}_0,c^{\scriptscriptstyle II})$};
            \node[factor, below=1.5cm of A0] (c0) {$\tilde{p}(o_0)$};
            \node[factor, right=1.8cm of eq0] (B0) {$p(s_1|s^{\scriptscriptstyle I}_0,a_0)$};
            \node[eqfactor, right=1.8cm of B0] (eq1) {$=$};
            \node[factor] (A1) at ([xshift=1.5cm, yshift=-2.15cm] B0.east) {$p(o_1|s^{\scriptscriptstyle II}_1,c^{\scriptscriptstyle IV})$};
            \node[factor, below=1.5cm of A1] (c1) {$\tilde{p}(o_1)$};
            \node[factor, right=1.8cm of eq1] (B1) {$p(s_2|s^{\scriptscriptstyle I}_1,a_1)$};
            \node[factor] (A2) at ([xshift=1.5cm] B1.east |- A1) {$p(o_2|s_2,c^{\scriptscriptstyle III})$};
            \node[factor, below=1.5cm of A2] (c2) {$\tilde{p}(o_2)$};
            \node[factor, above=1.5cm of B0] (u0) {$p(a_0)$};
            \node[factor, above=1.5cm of B1] (u1) {$p(a_1)$};
            \node[eqfactor, above right=3cm and 0.65cm of d] (eq3) {$=$};
            \node[eqfactor, above right=3cm and 0.65cm of B0] (eq4) {$=$};
            \node[factor, left=1.8cm of eq3] (dc) {$p(c)$};

            % Kanten zwischen Knoten
            \draw[edge] (d.east) -- (eq0.west) node[pos=0.75, above, varlabel] {$s_0$};
            \draw[edge] (eq0.south) -- (eq0.south |- A0.north) node[midway, left, varlabel] {$s^{\scriptscriptstyle II}_0$};
            \draw[edge] (A0.south) -- (c0.north) node[midway, left, varlabel] {$o_0$};
            \draw[edge] (eq0.east) -- (B0.west) node[midway, above, varlabel] {$s^{\scriptscriptstyle I}_0$};
            \draw[edge] (B0.north) -- (u0.south) node[midway, left, varlabel] {$a_0$};
            \draw[edge] (B0.east) -- (eq1.west) node[pos=0.75, above, varlabel] {$s_1$};
            \draw[edge] (eq1.south) -- (eq1.south |- A1.north) node[midway, left, varlabel] {$s^{\scriptscriptstyle II}_1$};
            \draw[edge] (A1.south) -- (c1.north) node[midway, left, varlabel] {$o_1$};
            \draw[edge] (eq1.east) -- (B1.west) node[midway, above, varlabel] {$s^{\scriptscriptstyle I}_1$};
            \draw[edge] (B1.east) -| ([xshift=0.57cm]A2.north) node[pos=0.37, above, varlabel] {$s_2$};
            \draw[edge] (A2.south) -- (c2.north) node[midway, left, varlabel] {$o_2$};
            \draw[edge] (B1.north) -- (u1.south) node[midway, left, varlabel] {$a_1$};
            \draw[edge] (dc.east) -- (eq3.west) node[midway, above, varlabel] {$c$};
            \draw[edge] (eq3.east) -- (eq4.west) node[midway, above, varlabel] {$c^{\scriptscriptstyle I}$};
            \draw[edge] (eq4.east) -| ([xshift=-0.5cm]A2.north) node[pos=0.25, above, varlabel] {$c^{\scriptscriptstyle III}$};
            \draw[edge] (eq3.south) -- (eq3.south |- A0.north) node[midway, left, varlabel] {$c^{\scriptscriptstyle II}$};
            \draw[edge] (eq4.south) -- (eq4.south |- A1.north) node[midway, left, varlabel] {$c^{\scriptscriptstyle IV}$};

        \end{tikzpicture}
    }
    \caption{Foundational Factor Graph of the two Context Agents.}
    \label{fig:28}
\end{figure}

In principle, the derivation of messages based on the factor graph's Bethe free energy is straightforward (compare \cref{sec:1}). In practice, the addition of a hidden context state leads to factor graph that contains loops (e.g., \(c^{\scriptscriptstyle II}_0 \rightarrow c^{\scriptscriptstyle I}_0 \rightarrow c^{\scriptscriptstyle IV}_0 \rightarrow s^{\scriptscriptstyle II}_1 \rightarrow s_1 \rightarrow s^{\scriptscriptstyle I}_0 \rightarrow s^{\scriptscriptstyle II}_0 \rightarrow c^{\scriptscriptstyle II}_0\)). Given this situation, messages derived from the Bethe free energy might lead to a message passing algorithm which might not yield stable solutions for the belief distributions, i.e, the algorithm might not converge. The problem lies at the nodes' belief distributions that explicitly respect dependencies between variables of the node, e.g., the belief distribution \(q(o_0,s^{\scriptscriptstyle II}_0,c^{\scriptscriptstyle II}_0)\) at the node \(p(o_0|s^{\scriptscriptstyle II}_0,c^{\scriptscriptstyle II}_0)\) respects the dependencies between the variables \(o_0\), \(s^{\scriptscriptstyle II}_0\), and \(c^{\scriptscriptstyle II}_0\). The solution to the convergence issues of pure Bethe free energy related message passing algorithm in graphs with loops lies in breaking up these dependencies to break the loop. Following this, it is then sufficient to define this node's belief distribution as \(q(o_0,s^{\scriptscriptstyle II}_0,c^{\scriptscriptstyle II}_0) = q(o_0,s^{\scriptscriptstyle II}_0)q(c^{\scriptscriptstyle II}_0)\). In this way, the dependencies between \(c^{\scriptscriptstyle II}_0\) and the other two variables are broken up in the belief distribution which eventually breaks the prior exemplary loop and the resulting message passing algorithm yields stable solutions, i.e., it converges. Breaking up these dependencies in nodes' belief distributions leads to structured variational messages that are sent from the respective nodes. These messages are derived in the following section.

\subsubsection{Structured Variational Messages}

Starting point for this derivation is a factor node with three edges (\cref{fig:34.1}), i.e, three variables.

\begin{figure}[htbp]
    \centering
    \begin{subfigure}[t]{0.48\textwidth}
        \centering
        \begin{tikzpicture}[
                factor/.style={draw=black, thick, rectangle, minimum size=0.8cm, minimum width=2.0cm, font=\normalsize},
                eqfactor/.style={draw=black, thick, rectangle, minimum size=0.5cm, font=\normalsize},
                edge/.style={thick},
                varlabel/.style={font=\normalsize},
                greycircle/.style={draw=black, thick, circle, fill=gray!30, minimum size=0.5cm, font=\normalsize}
            ]

            % Knoten
            \node[factor, label={[varlabel]left:$q_a(o,s,c)$}] (A1) {$p(o|s,c)$};

            % Kanten zwischen Knoten
            \draw[edge] ([yshift=1.5cm,xshift=0.7cm]A1.north) -- ([xshift=0.7cm]A1.north) node[midway, left, varlabel] {$q_i(s)$};
            \draw[edge] (A1.south) -- ([yshift=-1.5cm]A1.south) node[midway, left, varlabel] {$q_j(o)$};
            \draw[edge] ([yshift=1.5cm,xshift=-0.7cm]A1.north) -- ([xshift=-0.7cm]A1.north) node[midway, left, varlabel] {$q_k(c)$};

        \end{tikzpicture}
        \caption{Factor node with node belief that respects variables' dependencies.}
        \label{fig:34.1}
    \end{subfigure}
    \hfill
    \begin{subfigure}[t]{0.48\textwidth}
        \centering
        \begin{tikzpicture}[
                factor/.style={draw=black, thick, rectangle, minimum size=0.8cm, minimum width=3.0cm, font=\normalsize},
                eqfactor/.style={draw=black, thick, rectangle, minimum size=0.5cm, font=\normalsize},
                edge/.style={thick},
                varlabel/.style={font=\normalsize},
                greycircle/.style={draw=black, thick, circle, fill=gray!30, minimum size=0.5cm, font=\normalsize}
            ]

            % Knoten
            \node[factor, label={[varlabel]left:$q_a(o,s)q_a(c)$}] (A1) {$p(o|s,c)$};

            % Kanten zwischen Knoten
            \draw[edge] ([yshift=1.5cm,xshift=1.2cm]A1.north) -- ([xshift=1.2cm]A1.north) node[midway, left, varlabel] {$q_i(s)$} node[pos=0.25, right, varlabel] {$\textcolor{accorange}{\bm{{\uparrow}{\mu}(}s\bm{)}}$} node[pos=0.75, right, varlabel] {$\bm{{\downarrow}{\mu}(}s\bm{)}$} {};
            \draw[edge] (A1.south) -- ([yshift=-1.5cm]A1.south) node[midway, left, varlabel] {$q_j(o)$} node[pos=0.25, right, varlabel] {$\bm{{\uparrow}{\mu}(}o\bm{)}$} node[pos=0.75, right, varlabel] {$\textcolor{accorange}{\bm{{\downarrow}{\mu}(}o\bm{)}}$} {};
            \draw[edge] ([yshift=1.5cm,xshift=-1.2cm]A1.north) -- ([xshift=-1.2cm]A1.north) node[midway, left, varlabel] {$q_k(c)$} node[pos=0.25, right, varlabel] {$\textcolor{accorange}{\bm{{\uparrow}{\mu}(}c\bm{)}}$} node[pos=0.75, right, varlabel] {$\bm{{\downarrow}{\mu}(}c\bm{)}$} node[pos=1, circle, draw=black, fill=white, inner sep=2pt] {};
            \draw[edge, rounded corners=3.5pt]
            ([xshift=1.2cm]A1.north)
            |- ([yshift=0.15cm]A1.south)
            -- (A1.south);

        \end{tikzpicture}
        \caption{Factor node with node belief that assumes independence between \(c\) and the variables \(s\) and \(o\). This situation is visualized in graph, where both variables (\(s\) and \(o\)) whose dependencies are modeled in \(q_a\) are connected by a line through the node. The independent variable \(c\) (with respect to \(q_a\)) has its edge terminated by a bead.}
        \label{fig:34.2}
    \end{subfigure}
    \caption{Exemplary factor nodes that are used for the derivation of structured variational messages (\(\bm{{\uparrow}{\mu}(}c\bm{)}\), \(\bm{{\uparrow}{\mu}(}s\bm{)}\), and \(\bm{{\downarrow}{\mu}(}o\bm{)}\)).}
    \label{fig:34}
\end{figure}

The corresponding Lagrangian (compare \cref{sec:1}) (excluding normalization constraints) is given by:
\begin{equation}
    \begin{split}
        L = & \sum_{o,s,c} q_a(o,s,c) \ln \frac{q_a(o,s,c)}{p(o|s,c)} + \sum_c \lambda_{ka}(c) \left( q_k(c) - \sum_{o,s} q_a(o,s,c) \right)                \\
            & + \sum_s \lambda_{ia}(s) \left( q_i(s) - \sum_{o,c} q_a(o,s,c) \right) + \sum_o \lambda_{ja}(o) \left( q_j(o) - \sum_{s,c} q_a(o,s,c) \right)
    \end{split}
\end{equation}

In the following I show how the constraint \(q_a(o,s,c) = q_a(o,s)q_a(c)\) leads to alternative messages (structured variational messages) for \(\bm{{\uparrow}{\mu}(}c\bm{)}\), \(\bm{{\uparrow}{\mu}(}s\bm{)}\), and \(\bm{{\downarrow}{\mu}(}o\bm{)}\) in \cref{fig:34.2}. Following this, we introduce the constraint which leads to a Lagrangian \(L_{\text{SVMP}}\) that supports structured variational message passing (SVMP):
\begin{equation}
    \begin{split}
        L_{\text{SVMP}} = & \sum_{o,s,c} q_a(o,s)q_a(c) \ln \frac{q_a(o,s)q_a(c)}{p(o|s,c)} + \sum_c \lambda_{ka}(c) \left( q_k(c) - \sum_{o,s} q_a(o,s)q_a(c) \right)            \\
                          & + \sum_s \lambda_{ia}(s) \left( q_i(s) - \sum_{o,c} q_a(o,s)q_a(c) \right) + \sum_o \lambda_{ja}(o) \left( q_j(o) - \sum_{s,c} q_a(o,s)q_a(c) \right)
    \end{split}
\end{equation}

Next, the Lagrangian is differentiated with respect to \(q_a(o,s)\) and \(q_a(c)\), both are set to \(0\), and solved for \(q^*_a(o,s)\) and \(q^*_a(c)\) respectively:
\begin{alignat}{3}
     & \frac{\delta L_{\text{SVMP}}}{\delta q_a(o,s)} &  & =           &  & \ln q_a(o,s) + 1 + \sum_c q_a(c) \ln q_a(c) \nonumber                                                                             \\
     &                                                &  &             &  & - \sum_c q_a(c) \ln p(o|s,c) - \lambda_{ja}(o) - \lambda_{ia}(s) - \sum_c \lambda_{ka}(c) q_a(c) \overset{!}{=} 0                 \\
     &                                                &  & \Rightarrow &  & \; q^*_a(o,s) \propto \exp \left( \sum_{c} q_a(c) \ln p(o|s,c) \right) \exp \left( \lambda_{ja}(o) \right) \exp \left( \lambda_{ia}(s) \right) \label{eq:44} \\
     & \frac{\delta L_{\text{SVMP}}}{\delta q_a(c)}   &  & =           &  & \sum_{o,s} q_a(o,s) \ln q_a(o,s) + \ln q_a(c) + 1 - \sum_{o,s} q_a(o,s) \ln p(o|s,c) \nonumber \\
     &                                                &  &             &  & - \sum_{o,s} \lambda_{ja}(o) q_a(o,s) - \sum_{o,s} \lambda_{ia}(s) q_a(o,s) - \lambda_{ka}(c) \sum_{s,o} q_a(o,s) \overset{!}{=} 0  \\
     &                                                &  & \Rightarrow &  & \;q^*_a(c) \propto \exp \left( \sum_{o,s} q_a(o,s) \ln p(o|s,c) \right) \exp \left( \lambda_{ka}(c) \right) \label{eq:45}
\end{alignat}

Note, that in \cref{eq:44} some terms are removed because they are constant with respect to \(o\) and \(s\). The same is true for \cref{eq:45}, where terms are removed that are constant with respect to \(c\).
Finally, the marginalization constraints are used to derive the three outgoing messages \(\textcolor{accorange}{\bm{{\uparrow}{\mu}(}c\bm{)}}\), \(\textcolor{accorange}{\bm{{\uparrow}{\mu}(}s\bm{)}}\), and \(\textcolor{accorange}{\bm{{\downarrow}{\mu}(}o\bm{)}}\) in \cref{fig:34.2} (here the results for edge beliefs, e.g., for \(q_i(s)\), are used from \cref{sec:1}):
\begin{alignat}{3}
    &q_i(s) &&\overset{!}{=} \sum_o q_a(o,s) \\
    &\exp \left( \lambda_{ia}(s) \right) \exp \left( \lambda_{i\bullet}(s) \right) &&\propto \sum_o \exp \left( \sum_c q_a(c) \ln p(o|s,c) \right) \exp \left( \lambda_{ja}(o) \right) \exp \left( \lambda_{ia}(s) \right) \\
    &\mathrlap{\Rightarrow \exp \left( \lambda_{i\bullet}(s) \right) \propto \sum_o \exp \left( \sum_c q_a(c) \ln p(o|s,c) \right) \exp \left( \lambda_{ja}(o) \right)} \\
    &\mathrlap{\Rightarrow \textcolor{accorange}{\bm{{\uparrow}{\mu}(}s\bm{)}} \propto \sum_o \exp \left( \sum_c q_a(c) \ln p(o|s,c) \right) \bm{{\uparrow}{\mu}(}o\bm{)}}
\end{alignat}

By a similar argument it follows:
\begin{alignat}{1}
    &q_j(o) \overset{!}{=} \sum_s q_a(o,s) \\
    &\Rightarrow \textcolor{accorange}{\bm{{\downarrow}{\mu}(}o\bm{)}} \propto \sum_s \exp \left( \sum_c q_a(c) \ln p(o|s,c) \right) \bm{{\downarrow}{\mu}(}s\bm{)}
\end{alignat}

Lastly, \(\bm{{\uparrow}{\mu}(}c\bm{)}\) is given:
\begin{alignat}{3}
    &q_k(c) &&\overset{!}{=} q_a(c) \\
    &\exp \left( \lambda_{ka}(c) \right) \exp \left( \lambda_{k\bullet}(c) \right) &&\propto \exp \left( \sum_{o,s} q_a(o,s) \ln p(o|s,c) \right) \exp \left( \lambda_{ka}(c) \right) \\
    &\mathrlap{\Rightarrow \exp \left( \lambda_{k\bullet}(s) \right) \propto \exp \left( \sum_{o,s} q_a(o,s) \ln p(o|s,c) \right)} \\
    &\mathrlap{\Rightarrow \textcolor{accorange}{\bm{{\uparrow}{\mu}(}c\bm{)}} \propto \exp \left( \sum_{o,s} q_a(o,s) \ln p(o|s,c) \right)}
\end{alignat}

The "\(\bullet\)" signs in the lambdas indicate the other nodes (not shown in \cref{fig:34}) the respective edge is connected to. The terms that contain these lambdas (e.g., \(\exp \left( \lambda_{i\bullet}(s) \right)\)) are the outgoing messages to these other nodes.

\subsubsection{Message Passing Algorithm}

The factor graph of the context-BFE agent with messages is given in \cref{fig:35}. Again, the graph shows an exemplary integration of an observation \(\hat{o_0}\) as a delta constraint.
\begin{figure}[htbp]
    \centering

    \resizebox{1.0\textwidth}{!}{%
        \begin{tikzpicture}[
                factor/.style={draw=black, thick, rectangle, minimum size=0.8cm, font=\normalsize},
                eqfactor/.style={draw=black, thick, rectangle, minimum size=0.5cm, font=\normalsize},
                edge/.style={thick},
                varlabel/.style={font=\normalsize},
                greycircle/.style={draw=black, thick, circle, fill=gray!30, minimum size=0.5cm, font=\normalsize}
            ]

            % Knoten
            \node[factor] (d) {$p(s_0)$};
            \node[eqfactor, right=1.8cm of d] (eq0) {$=$};
            \node[factor] (A0) at ([xshift=1.5cm, yshift=-2.25cm] d.east) {$p(o_0|s^{\scriptscriptstyle II}_0,c^{\scriptscriptstyle II})$};
            \node[factor, below=1.6cm of A0] (c0) {$\tilde{p}(o_0)$};
            \node[factor, right=1.8cm of eq0] (B0) {$p(s_1|s^{\scriptscriptstyle I}_0,a_0)$};
            \node[eqfactor, right=1.8cm of B0] (eq1) {$=$};
            \node[factor] (A1) at ([xshift=1.5cm, yshift=-2.25cm] B0.east) {$p(o_1|s^{\scriptscriptstyle II}_1,c^{\scriptscriptstyle IV})$};
            \node[factor, below=1.6cm of A1] (c1) {$\tilde{p}(o_1)$};
            \node[factor, right=1.8cm of eq1] (B1) {$p(s_2|s^{\scriptscriptstyle I}_1,a_1)$};
            \node[factor] (A2) at ([xshift=1.5cm] B1.east |- A1) {$p(o_2|s_2,c^{\scriptscriptstyle III})$};
            \node[factor, below=1.6cm of A2] (c2) {$\tilde{p}(o_2)$};
            \node[factor, above=1.6cm of B0] (u0) {$p(a_0)$};
            \node[factor, above=1.6cm of B1] (u1) {$p(a_1)$};
            \node[eqfactor, above right=3cm and 0.65cm of d] (eq3) {$=$};
            \node[eqfactor, above right=3cm and 0.65cm of B0] (eq4) {$=$};
            \node[factor, left=1.8cm of eq3] (dc) {$p(c)$};

            % Kanten zwischen Knoten
            \draw[edge] (d.east) -- (eq0.west) node[pos=0.77, above, varlabel] {$\textcolor{accblue}{q(s_0)}$} node[near start, below, varlabel, scale=0.7] {$\textcolor{accorange}{\bm{\overleftarrow{\mu}(}s_0\bm{)}}$} node[near end, below, varlabel, scale=0.7] {$\textcolor{accorange}{\bm{\overrightarrow{\mu}(}s_0\bm{)}}$};
            \draw[edge] (eq0.south) -- (eq0.south |- A0.north) node[midway, left=-0.05, varlabel] {$\textcolor{accblue}{q(s^{\scriptscriptstyle II}_0)}$} node[pos=0.25, right, varlabel, scale=0.7] {$\textcolor{accorange}{\bm{{\uparrow}{\mu}(}s^{\scriptscriptstyle II}_0\bm{)}}$} node[pos=0.75, right, varlabel, scale=0.7] {$\textcolor{accorange}{\bm{{\downarrow}{\mu}(}s^{\scriptscriptstyle II}_0\bm{)}}$};
            \draw[edge] (A0.south) -- (c0.north) node[midway, greycircle] {$\delta$} node[midway, left=0.3, varlabel] {$\textcolor{accblue}{q(o_0)}$} node[pos=0.25, right=0.25, varlabel, scale=0.7] {$\textcolor{accorange}{\bm{{\uparrow}{\mu}(}o_0\bm{)}}$} node[pos=0.75, right=0.25, varlabel, scale=0.7] {$\textcolor{accorange}{\bm{{\downarrow}{\mu}(}o_0\bm{)}}$};
            \draw[edge] (eq0.east) -- (B0.west) node[midway, above, varlabel] {$\textcolor{accblue}{q(s^{\scriptscriptstyle I}_0)}$} node[near start, below, varlabel, scale=0.7] {$\textcolor{accorange}{\bm{\overleftarrow{\mu}(}s^{\scriptscriptstyle I}_0\bm{)}}$} node[near end, below, varlabel, scale=0.7] {$\textcolor{accorange}{\bm{\overrightarrow{\mu}(}s^{\scriptscriptstyle I}_0\bm{)}}$};
            \draw[edge] (B0.north) -- (u0.south) node[midway, left, varlabel] {$\textcolor{accblue}{q(a_0)}$} node[pos=0.75, right, varlabel, scale=0.7] {$\textcolor{accorange}{\bm{{\uparrow}{\mu}(}a_0\bm{)}}$} node[pos=0.25, right, varlabel, scale=0.7] {$\textcolor{accorange}{\bm{{\downarrow}{\mu}(}a_0\bm{)}}$};
            \draw[edge] (B0.east) -- (eq1.west) node[pos=0.77, above, varlabel] {$\textcolor{accblue}{q(s_1)}$} node[near start, below, varlabel, scale=0.7] {$\textcolor{accorange}{\bm{\overleftarrow{\mu}(}s_1\bm{)}}$} node[near end, below, varlabel, scale=0.7] {$\textcolor{accorange}{\bm{\overrightarrow{\mu}(}s_1\bm{)}}$};
            \draw[edge] (eq1.south) -- (eq1.south |- A1.north) node[midway, left, varlabel] {$\textcolor{accblue}{q(s^{\scriptscriptstyle II}_1)}$} node[pos=0.25, right, varlabel, scale=0.7] {$\textcolor{accorange}{\bm{{\uparrow}{\mu}(}s^{\scriptscriptstyle II}_1\bm{)}}$} node[pos=0.75, right, varlabel, scale=0.7] {$\textcolor{accorange}{\bm{{\downarrow}{\mu}(}s^{\scriptscriptstyle II}_1\bm{)}}$};
            \draw[edge] (A1.south) -- (c1.north) node[midway, left, varlabel] {$\textcolor{accblue}{q(o_1)}$} node[pos=0.25, right, varlabel, scale=0.7] {$\textcolor{accorange}{\bm{{\uparrow}{\mu}(}o_1\bm{)}}$} node[pos=0.75, right, varlabel, scale=0.7] {$\textcolor{accorange}{\bm{{\downarrow}{\mu}(}o_1\bm{)}}$};
            \draw[edge] (eq1.east) -- (B1.west) node[midway, above, varlabel] {$\textcolor{accblue}{q(s^{\scriptscriptstyle I}_1)}$} node[near start, below, varlabel, scale=0.7] {$\textcolor{accorange}{\bm{\overleftarrow{\mu}(}s^{\scriptscriptstyle I}_1\bm{)}}$} node[near end, below, varlabel, scale=0.7] {$\textcolor{accorange}{\bm{\overrightarrow{\mu}(}s^{\scriptscriptstyle I}_1\bm{)}}$};
            \draw[edge] (B1.east) -| ([xshift=0.57cm]A2.north) node[pos=0.37, above, varlabel] {$\textcolor{accblue}{q(s_2)}$} node[pos=0.11, below, varlabel, scale=0.7] {$\textcolor{accorange}{\bm{\overleftarrow{\mu}(}s_2\bm{)}}$} node[pos=0.89, right, varlabel, scale=0.7] {$\textcolor{accorange}{\bm{{\downarrow}{\mu}(}s_2\bm{)}}$};
            \draw[edge] (A2.south) -- (c2.north) node[midway, left, varlabel] {$\textcolor{accblue}{q(o_2)}$} node[pos=0.25, right, varlabel, scale=0.7] {$\textcolor{accorange}{\bm{{\uparrow}{\mu}(}o_2\bm{)}}$} node[pos=0.75, right, varlabel, scale=0.7] {$\textcolor{accorange}{\bm{{\downarrow}{\mu}(}o_2\bm{)}}$};
            \draw[edge] (B1.north) -- (u1.south) node[midway, left, varlabel] {$\textcolor{accblue}{q(a_1)}$} node[pos=0.75, right, varlabel, scale=0.7] {$\textcolor{accorange}{\bm{{\uparrow}{\mu}(}a_1\bm{)}}$} node[pos=0.25, right, varlabel, scale=0.7] {$\textcolor{accorange}{\bm{{\downarrow}{\mu}(}a_1\bm{)}}$};
            \draw[edge] (dc.east) -- (eq3.west) node[midway, above, varlabel] {$\textcolor{accblue}{q(c)}$} node[near start, below, varlabel, scale=0.7] {$\textcolor{accorange}{\bm{\overleftarrow{\mu}(}c\bm{)}}$} node[near end, below, varlabel, scale=0.7] {$\textcolor{accorange}{\bm{\overrightarrow{\mu}(}c\bm{)}}$};
            \draw[edge] (eq3.east) -- (eq4.west) node[midway, above, varlabel] {$\textcolor{accblue}{q(c^{\scriptscriptstyle I})}$} node[pos=0.08, below, varlabel, scale=0.7] {$\textcolor{accorange}{\bm{\overleftarrow{\mu}(}c^{\scriptscriptstyle I}\bm{)}}$} node[pos=0.92, below, varlabel, scale=0.7] {$\textcolor{accorange}{\bm{\overrightarrow{\mu}(}c^{\scriptscriptstyle I}\bm{)}}$};
            \draw[edge] (eq4.east) -| ([xshift=-0.5cm]A2.north) node[pos=0.25, above, varlabel] {$\textcolor{accblue}{q(c^{\scriptscriptstyle III})}$} node[pos=0.045, below, varlabel, scale=0.7] {$\textcolor{accorange}{\bm{\overleftarrow{\mu}(}c^{\scriptscriptstyle III}\bm{)}}$} node[pos=0.965, right=-0.05, varlabel, scale=0.7] {$\textcolor{accorange}{\bm{{\downarrow}{\mu}(}c^{\scriptscriptstyle III}\bm{)}}$} node[pos=1, circle, draw=black, fill=white, inner sep=2pt] {};
            \draw[edge] (eq3.south) -- (eq3.south |- A0.north) node[pos=0.2, left, varlabel] {$\textcolor{accblue}{q(c^{\scriptscriptstyle II})}$} node[pos=0.07, right, varlabel, scale=0.7] {$\textcolor{accorange}{\bm{{\uparrow}{\mu}(}c^{\scriptscriptstyle II}\bm{)}}$} node[pos=0.93, right, varlabel, scale=0.7] {$\textcolor{accorange}{\bm{{\downarrow}{\mu}(}c^{\scriptscriptstyle II}\bm{)}}$} node[pos=1, circle, draw=black, fill=white, inner sep=2pt] {};
            \draw[edge] (eq4.south) -- (eq4.south |- A1.north) node[pos=0.2, left, varlabel] {$\textcolor{accblue}{q(c^{\scriptscriptstyle IV})}$} node[pos=0.07, right, varlabel, scale=0.7] {$\textcolor{accorange}{\bm{{\uparrow}{\mu}(}c^{\scriptscriptstyle IV}\bm{)}}$} node[pos=0.93, right, varlabel, scale=0.7] {$\textcolor{accorange}{\bm{{\downarrow}{\mu}(}c^{\scriptscriptstyle IV}\bm{)}}$} node[pos=1, circle, draw=black, fill=white, inner sep=2pt] {};

            \draw[edge, rounded corners=3.5pt]
            (eq0.south |- A0.north)
            -- ([yshift=-0.12cm]eq0.south |- A0.north)
            -| ([xshift=-0.1cm]A0.east)
            |- ([yshift=0.12cm]A0.south)
            -- (A0.south);

            \draw[edge, rounded corners=3.5pt]
            (eq1.south |- A1.north)
            -- ([yshift=-0.12cm]eq1.south |- A1.north)
            -| ([xshift=-0.1cm]A1.east)
            |- ([yshift=0.12cm]A1.south)
            -- (A1.south);

            \draw[edge, rounded corners=3.5pt]
            ([xshift=0.57cm]A2.north)
                -- ([xshift=0.57cm, yshift=-0.12cm]A2.north)
            -| ([xshift=-0.1cm]A2.east)
            |- ([yshift=0.12cm]A2.south)
            -- (A2.south);

        \end{tikzpicture}
    }
    \caption{Factor Graph of the context-BFE agent with variables' belief distributions \(\textcolor{accblue}{q}\) and messages \(\textcolor{accorange}{\mu}\).}
    \label{fig:35}
\end{figure}

Importantly, likelihood nodes send structured variational messages to break the cyclic dependencies introduced by the additional context state.  Similarly to the GFE agent these special messages of the context-BFE agent also requires recursive execution of the algorithm until convergence of the belief distributions. Moreover, the distributions \(q(o_0,s^{\scriptscriptstyle II}_0)\), \(q(o_1,s^{\scriptscriptstyle II}_1)\), \(q(o_2,s_2)\), \(q(c^{\scriptscriptstyle II})\), \(q(c^{\scriptscriptstyle IV})\), and \(q(c^{\scriptscriptstyle III})\) must be initialized (e.g., uniformly) before execution.
The resulting algorithm is given in \cref{alg:4}.
\begin{algorithm}[htbp]
    \caption{Message Passing Algorithm of the Context-BFE Agent}
    \label{alg:4}
    \small
    \begin{algorithmic}[1] % Die [1] sorgt für automatische Zeilennummern!
        \linespread{1.4}\selectfont
        \setlength{\leftskip}{0.3cm}
        \State \tikzmark{algtop4}$\textcolor{accorange}{\bm{\overrightarrow{\mu}(}s_0\bm{)}} \propto p(s_0)$,\quad$\textcolor{accorange}{\bm{{\uparrow}{\mu}(}o_0\bm{)}} \propto \delta_{o_0,\hat{o_0}}$,\quad$\textcolor{accorange}{\bm{{\downarrow}{\mu}(}a_0\bm{)}} \propto p(a_0)$,\quad$\textcolor{accorange}{\bm{{\uparrow}{\mu}(}o_1\bm{)}} \propto \tilde{p}(o_1)$,\quad$\textcolor{accorange}{\bm{{\downarrow}{\mu}(}a_1\bm{)}} \propto p(a_1)$,
        \Statex $\textcolor{accorange}{\bm{{\uparrow}{\mu}(}c^{\scriptscriptstyle II}\bm{)}} \propto \exp \left( \sum_{o_0,s^{\scriptscriptstyle II}_0} q(o_0,s^{\scriptscriptstyle II}_0) \ln p(o_0|s^{\scriptscriptstyle II}_0,c^{\scriptscriptstyle II}) \right)$,\quad$\textcolor{accorange}{\bm{{\downarrow}{\mu}(}o_0\bm{)}} \propto \delta_{o_0,\hat{o_0}}$,
        \Statex $\textcolor{accorange}{\bm{{\uparrow}{\mu}(}c^{\scriptscriptstyle IV}\bm{)}} \propto \exp \left( \sum_{o_1,s^{\scriptscriptstyle II}_1} q(o_1,s^{\scriptscriptstyle II}_1) \ln p(o_1|s^{\scriptscriptstyle II}_1,c^{\scriptscriptstyle IV}) \right)$,\quad$\textcolor{accorange}{\bm{{\uparrow}{\mu}(}o_2\bm{)}} \propto \tilde{p}(o_2)$,
        \Statex $\textcolor{accorange}{\bm{\overleftarrow{\mu}(}c^{\scriptscriptstyle III}\bm{)}} \propto \exp \left( \sum_{o_2,s_2} q(o_2,s_2) \ln p(o_2|s_2,c^{\scriptscriptstyle III}) \right)$,\quad$\textcolor{accorange}{\bm{\overrightarrow{\mu}(}c\bm{)}} \propto p(c)$
        \State $\textcolor{accorange}{\bm{{\uparrow}{\mu}(}s^{\scriptscriptstyle II}_0\bm{)}} \propto \sum_{o_0} \exp \left( \sum_{c^{\scriptscriptstyle II}} q(c^{\scriptscriptstyle II}) \ln p(o_0|s^{\scriptscriptstyle II}_0,c^{\scriptscriptstyle II}) \right) \textcolor{black}{\bm{{\uparrow}{\mu}(}o_0\bm{)}}$,\quad$\textcolor{accorange}{\bm{\overrightarrow{\mu}(}c^{\scriptscriptstyle I}\bm{)}} \propto \textcolor{black}{\bm{\overrightarrow{\mu}(}c\bm{)}} \textcolor{black}{\bm{{\uparrow}{\mu}(}c^{\scriptscriptstyle II}\bm{)}}$,
        \Statex $\textcolor{accorange}{\bm{{\uparrow}{\mu}(}s^{\scriptscriptstyle II}_1\bm{)}} \propto \sum_{o_1} \exp \left( \sum_{c^{\scriptscriptstyle IV}} q(c^{\scriptscriptstyle IV}) \ln p(o_1|s^{\scriptscriptstyle II}_1,c^{\scriptscriptstyle IV}) \right) \textcolor{black}{\bm{{\uparrow}{\mu}(}o_1\bm{)}}$,\quad$\textcolor{accorange}{\bm{\overleftarrow{\mu}(}c^{\scriptscriptstyle I}\bm{)}} \propto \textcolor{black}{\bm{{\uparrow}{\mu}(}c^{\scriptscriptstyle IV}\bm{)}} \textcolor{black}{\bm{\overleftarrow{\mu}(}c^{\scriptscriptstyle III}\bm{)}}$,
        \Statex $\textcolor{accorange}{\bm{\overleftarrow{\mu}(}s_2\bm{)}} \propto \sum_{o_2} \exp \left( \sum_{c^{\scriptscriptstyle III}} q(c^{\scriptscriptstyle III}) \ln p(o_2|s_2,c^{\scriptscriptstyle III}) \right) \textcolor{black}{\bm{{\uparrow}{\mu}(}o_2\bm{)}}$
        \State $\textcolor{accorange}{\bm{\overrightarrow{\mu}(}s^{\scriptscriptstyle I}_0\bm{)}} \propto \sum_{s_0,s^{\scriptscriptstyle II}_0} \delta_{s_0,s^{\scriptscriptstyle I}_0}\delta_{s_0,s^{\scriptscriptstyle II}_0} \textcolor{black}{\bm{\overrightarrow{\mu}(}s_0\bm{)}} \textcolor{black}{\bm{{\uparrow}{\mu}(}s^{\scriptscriptstyle II}_0\bm{)}}$,\quad$\textcolor{accorange}{\bm{{\downarrow}{\mu}(}c^{\scriptscriptstyle III}\bm{)}} \propto \textcolor{black}{\bm{\overrightarrow{\mu}(}c^{\scriptscriptstyle I}\bm{)}} \textcolor{black}{\bm{{\uparrow}{\mu}(}c^{\scriptscriptstyle IV}\bm{)}}$,
        \Statex $\textcolor{accorange}{\bm{\overleftarrow{\mu}(}s^{\scriptscriptstyle I}_1\bm{)}} \propto \sum_{a_1,s_2} p(s_2|s^{\scriptscriptstyle I}_1,a_1) \textcolor{black}{\bm{{\downarrow}{\mu}(}a_1\bm{)}} \textcolor{black}{\bm{\overleftarrow{\mu}(}s_2\bm{)}}$,\quad$\textcolor{accorange}{\bm{{\downarrow}{\mu}(}c^{\scriptscriptstyle IV}\bm{)}} \propto \textcolor{black}{\bm{\overrightarrow{\mu}(}c^{\scriptscriptstyle I}\bm{)}} \textcolor{black}{\bm{\overleftarrow{\mu}(}c^{\scriptscriptstyle III}\bm{)}}$,
        \Statex $\textcolor{accorange}{\bm{{\downarrow}{\mu}(}c^{\scriptscriptstyle II}\bm{)}} \propto \textcolor{black}{\bm{\overrightarrow{\mu}(}c\bm{)}} \textcolor{black}{\bm{\overleftarrow{\mu}(}c^{\scriptscriptstyle I}\bm{)}}$,\quad$\textcolor{accorange}{\bm{\overleftarrow{\mu}(}c\bm{)}} \propto \textcolor{black}{\bm{{\uparrow}{\mu}(}c^{\scriptscriptstyle II}\bm{)}} \textcolor{black}{\bm{\overleftarrow{\mu}(}c^{\scriptscriptstyle I}\bm{)}}$
        \State $\textcolor{accorange}{\bm{\overrightarrow{\mu}(}s_1\bm{)}} \propto \sum_{s^{\scriptscriptstyle I}_0,a_0} p(s_1|s^{\scriptscriptstyle I}_0,a_0) \textcolor{black}{\bm{\overrightarrow{\mu}(}s^{\scriptscriptstyle I}_0\bm{)}} \textcolor{black}{\bm{{\downarrow}{\mu}(}a_0\bm{)}}$,\quad$\textcolor{accorange}{\bm{\overleftarrow{\mu}(}s_1\bm{)}} \propto \sum_{s^{\scriptscriptstyle I}_1,s^{\scriptscriptstyle II}_1} \delta_{s_1,s^{\scriptscriptstyle I}_1}\delta_{s_1,s^{\scriptscriptstyle II}_1} \textcolor{black}{\bm{\overleftarrow{\mu}(}s^{\scriptscriptstyle I}_1\bm{)}} \textcolor{black}{\bm{{\uparrow}{\mu}(}s^{\scriptscriptstyle II}_1\bm{)}}$
        \State $\textcolor{accorange}{\bm{\overrightarrow{\mu}(}s^{\scriptscriptstyle I}_1\bm{)}} \propto \sum_{s_1,s^{\scriptscriptstyle II}_1} \delta_{s_1,s^{\scriptscriptstyle I}_1}\delta_{s_1,s^{\scriptscriptstyle II}_1} \textcolor{black}{\bm{\overrightarrow{\mu}(}s_1\bm{)}} \textcolor{black}{\bm{{\uparrow}{\mu}(}s^{\scriptscriptstyle II}_1\bm{)}}$,\quad$\textcolor{accorange}{\bm{{\downarrow}{\mu}(}s^{\scriptscriptstyle II}_1\bm{)}} \propto \sum_{s_1,s^{\scriptscriptstyle I}_1} \delta_{s_1,s^{\scriptscriptstyle I}_1}\delta_{s_1,s^{\scriptscriptstyle II}_1} \textcolor{black}{\bm{\overrightarrow{\mu}(}s_1\bm{)}} \textcolor{black}{\bm{\overleftarrow{\mu}(}s^{\scriptscriptstyle I}_1\bm{)}}$,
        \Statex $\textcolor{accorange}{\bm{{\uparrow}{\mu}(}a_0\bm{)}} \propto \sum_{s_1,s^{\scriptscriptstyle I}_0} p(s_1|s^{\scriptscriptstyle I}_0,a_0) \textcolor{black}{\bm{\overrightarrow{\mu}(}s^{\scriptscriptstyle I}_0\bm{)}} \textcolor{black}{\bm{\overleftarrow{\mu}(}s_1\bm{)}}$,\quad$\textcolor{accorange}{\bm{\overleftarrow{\mu}(}s^{\scriptscriptstyle I}_0\bm{)}} \propto \sum_{s_1,a_0} p(s_1|s^{\scriptscriptstyle I}_0,a_0) \textcolor{black}{\bm{{\downarrow}{\mu}(}a_0\bm{)}} \textcolor{black}{\bm{\overleftarrow{\mu}(}s_1\bm{)}}$
        \State $\textcolor{accorange}{\bm{{\uparrow}{\mu}(}a_1\bm{)}} \propto \sum_{s^{\scriptscriptstyle I}_1,s_2} p(s_2|s^{\scriptscriptstyle I}_1,a_1) \textcolor{black}{\bm{\overrightarrow{\mu}(}s^{\scriptscriptstyle I}_1\bm{)}} \textcolor{black}{\bm{\overleftarrow{\mu}(}s_2\bm{)}}$,\quad$\textcolor{accorange}{\bm{{\downarrow}{\mu}(}s_2\bm{)}} \propto \sum_{s^{\scriptscriptstyle I}_1,a_1} p(s_2|s^{\scriptscriptstyle I}_1,a_1) \textcolor{black}{\bm{\overrightarrow{\mu}(}s^{\scriptscriptstyle I}_1\bm{)}} \textcolor{black}{\bm{{\downarrow}{\mu}(}a_1\bm{)}}$,
        \Statex $\textcolor{accorange}{\bm{{\downarrow}{\mu}(}o_1\bm{)}} \propto \sum_{s^{\scriptscriptstyle II}_1} \exp \left( \sum_{c^{\scriptscriptstyle IV}} q(c^{\scriptscriptstyle IV}) \ln p(o_1|s^{\scriptscriptstyle II}_1,c^{\scriptscriptstyle IV}) \right) \textcolor{black}{\bm{{\downarrow}{\mu}(}s^{\scriptscriptstyle II}_1\bm{)}}$,
        \Statex $\textcolor{accorange}{\bm{\overleftarrow{\mu}(}s_0\bm{)}} \propto \sum_{s^{\scriptscriptstyle I}_0,s^{\scriptscriptstyle II}_0} \delta_{s_0,s^{\scriptscriptstyle I}_0}\delta_{s_0,s^{\scriptscriptstyle II}_0} \textcolor{black}{\bm{\overleftarrow{\mu}(}s^{\scriptscriptstyle I}_0\bm{)}} \textcolor{black}{\bm{{\uparrow}{\mu}(}s^{\scriptscriptstyle II}_0\bm{)}}$,\quad$\textcolor{accorange}{\bm{{\downarrow}{\mu}(}s^{\scriptscriptstyle II}_0\bm{)}} \propto \sum_{s_0,s^{\scriptscriptstyle I}_0} \delta_{s_0,s^{\scriptscriptstyle I}_0}\delta_{s_0,s^{\scriptscriptstyle II}_0} \textcolor{black}{\bm{\overrightarrow{\mu}(}s_0\bm{)}} \textcolor{black}{\bm{\overleftarrow{\mu}(}s^{\scriptscriptstyle I}_0\bm{)}}$
        \State $\textcolor{accorange}{\bm{{\downarrow}{\mu}(}o_2\bm{)}} \propto \sum_{s_2} \exp \left( \sum_{c^{\scriptscriptstyle III}} q(c^{\scriptscriptstyle III}) \ln p(o_2|s_2,c^{\scriptscriptstyle III}) \right) \textcolor{black}{\bm{{\downarrow}{\mu}(}s_2\bm{)}}$
        \State \tikzmark{algbottom4}$\textcolor{accblue}{q(s_0)} \propto \textcolor{black}{\bm{\overleftarrow{\mu}(}s_0\bm{)}} \textcolor{black}{\bm{\overrightarrow{\mu}(}s_0\bm{)}}$,\quad$\textcolor{accblue}{q(s^{\scriptscriptstyle I}_0)} \propto \textcolor{black}{\bm{\overleftarrow{\mu}(}s^{\scriptscriptstyle I}_0\bm{)}} \textcolor{black}{\bm{\overrightarrow{\mu}(}s^{\scriptscriptstyle I}_0\bm{)}}$,\quad$\textcolor{accblue}{q(s^{\scriptscriptstyle II}_0)} \propto \textcolor{black}{\bm{{\uparrow}{\mu}(}s^{\scriptscriptstyle II}_0\bm{)}} \textcolor{black}{\bm{{\downarrow}{\mu}(}s^{\scriptscriptstyle II}_0\bm{)}}$,
        \Statex $\textcolor{accblue}{q(s_1)} \propto \textcolor{black}{\bm{\overleftarrow{\mu}(}s_1\bm{)}} \textcolor{black}{\bm{\overrightarrow{\mu}(}s_1\bm{)}}$,\quad$\textcolor{accblue}{q(s^{\scriptscriptstyle I}_1)} \propto \textcolor{black}{\bm{\overleftarrow{\mu}(}s^{\scriptscriptstyle I}_1\bm{)}} \textcolor{black}{\bm{\overrightarrow{\mu}(}s^{\scriptscriptstyle I}_1\bm{)}}$,\quad$\textcolor{accblue}{q(s^{\scriptscriptstyle II}_1)} \propto \textcolor{black}{\bm{{\uparrow}{\mu}(}s^{\scriptscriptstyle II}_1\bm{)}} \textcolor{black}{\bm{{\downarrow}{\mu}(}s^{\scriptscriptstyle II}_1\bm{)}}$,
        \Statex $\textcolor{accblue}{q(s_2)} \propto \textcolor{black}{\bm{\overleftarrow{\mu}(}s_2\bm{)}} \textcolor{black}{\bm{{\downarrow}{\mu}(}s_2\bm{)}}$,\quad$\textcolor{accblue}{q(a_0)} \propto \textcolor{black}{\bm{{\uparrow}{\mu}(}a_0\bm{)}} \textcolor{black}{\bm{{\downarrow}{\mu}(}a_0\bm{)}}$,\quad$\textcolor{accblue}{q(a_1)} \propto \textcolor{black}{\bm{{\uparrow}{\mu}(}a_1\bm{)}} \textcolor{black}{\bm{{\downarrow}{\mu}(}a_1\bm{)}}$,
        \Statex $\textcolor{accblue}{q(o_0)} \propto \textcolor{black}{\bm{{\uparrow}{\mu}(}o_0\bm{)}} \textcolor{black}{\bm{{\downarrow}{\mu}(}o_0\bm{)}}$,\quad$\textcolor{accblue}{q(o_1)} \propto \textcolor{black}{\bm{{\uparrow}{\mu}(}o_1\bm{)}} \textcolor{black}{\bm{{\downarrow}{\mu}(}o_1\bm{)}}$,\quad$\textcolor{accblue}{q(o_2)} \propto \textcolor{black}{\bm{{\uparrow}{\mu}(}o_2\bm{)}} \textcolor{black}{\bm{{\downarrow}{\mu}(}o_2\bm{)}}$,
        \Statex $\textcolor{accblue}{q(c)} \propto \textcolor{black}{\bm{\overleftarrow{\mu}(}c\bm{)}} \textcolor{black}{\bm{\overrightarrow{\mu}(}c\bm{)}}$,\quad$\textcolor{accblue}{q(c^{\scriptscriptstyle I})} \propto \textcolor{black}{\bm{\overleftarrow{\mu}(}c^{\scriptscriptstyle I}\bm{)}} \textcolor{black}{\bm{\overrightarrow{\mu}(}c^{\scriptscriptstyle I}\bm{)}}$,\quad$\textcolor{accblue}{q(c^{\scriptscriptstyle II})} \propto \textcolor{black}{\bm{{\uparrow}{\mu}(}c^{\scriptscriptstyle II}\bm{)}} \textcolor{black}{\bm{{\downarrow}{\mu}(}c^{\scriptscriptstyle II}\bm{)}}$,
        \Statex $\textcolor{accblue}{q(c^{\scriptscriptstyle III})} \propto \textcolor{black}{\bm{\overleftarrow{\mu}(}c^{\scriptscriptstyle III}\bm{)}} \textcolor{black}{\bm{\overrightarrow{\mu}(}c^{\scriptscriptstyle III}\bm{)}}$,\quad$\textcolor{accblue}{q(c^{\scriptscriptstyle IV})} \propto \textcolor{black}{\bm{{\uparrow}{\mu}(}c^{\scriptscriptstyle IV}\bm{)}} \textcolor{black}{\bm{{\downarrow}{\mu}(}c^{\scriptscriptstyle IV}\bm{)}}$,
        \Statex $\textcolor{black}{q(o_0,s^{\scriptscriptstyle II}_0)} \propto \exp \left( \sum_{o_0,s^{\scriptscriptstyle II}_0} q(o_0,s^{\scriptscriptstyle II}_0) \ln p(o_0|s^{\scriptscriptstyle II}_0,c^{\scriptscriptstyle II}) \right) \textcolor{black}{\bm{{\downarrow}{\mu}(}s^{\scriptscriptstyle II}_0\bm{)}} \textcolor{black}{\bm{{\uparrow}{\mu}(}o_0\bm{)}}$,
        \Statex $\textcolor{black}{q(o_1,s^{\scriptscriptstyle II}_1)} \propto \exp \left( \sum_{o_1,s^{\scriptscriptstyle II}_1} q(o_1,s^{\scriptscriptstyle II}_1) \ln p(o_1|s^{\scriptscriptstyle II}_1,c^{\scriptscriptstyle IV}) \right) \textcolor{black}{\bm{{\downarrow}{\mu}(}s^{\scriptscriptstyle II}_1\bm{)}} \textcolor{black}{\bm{{\uparrow}{\mu}(}o_1\bm{)}}$,
        \Statex $\textcolor{black}{q(o_2,s_2)} \propto \exp \left( \sum_{o_2,s_2} q(o_2,s_2) \ln p(o_2|s_2,c^{\scriptscriptstyle III}) \right) \textcolor{black}{\bm{{\downarrow}{\mu}(}s_2\bm{)}} \textcolor{black}{\bm{{\uparrow}{\mu}(}o_2\bm{)}}$
    \end{algorithmic}

    \begin{tikzpicture}[remember picture, overlay]
        \draw[-{Stealth[length=2.2mm, width=1.6mm]}, thick, gray!75!black, rounded corners=5pt]
        ([yshift=0.5ex, xshift=-2.5ex]pic cs:algbottom4) -- ++(-0.4cm, 0) |- ([xshift=-2.5ex, yshift=0.5ex]pic cs:algtop4);
    \end{tikzpicture}

\end{algorithm}


\subsection{Context-GFE Agent}

In this section I present the derivation of GFE-based messages for the context state \(c\). Following this, the resulting Context-GFE agent will be able to reduce its uncertainty regarding the context state rather than only the hidden state \(s\), which is the central goal of this thesis. The agent is based on the same generative model as the Context-BFE agent, i.e., \cref{eq:35}, and consequently also utilizes the same foundational factor graph (\cref{fig:28}). Interestingly, the derivation also leads to GFE-based messages for the hidden state \(s\), although this was not the central goal. It can be suspected that this circumstance is based on the independence assumption between hidden state \(s\) and the other two variables \(c\) and \(o\) that is made during the derivation.

\subsubsection{GFE-Based Messages for the Context}

In this section I derive the GFE-based message (\(\textcolor{accorange}{\bm{{\uparrow}{\mu}(}c\bm{)}}\) in \cref{fig:36.2}) that is sent by future likelihood nodes upwards to transition nodes and enables the agent with contextual information seeking (also including the GFE-based message \(\textcolor{black}{\bm{{\uparrow}{\mu}(}s\bm{)}}\)).

\begin{figure}[htbp]
    \centering
    \begin{subfigure}[t]{0.48\textwidth}
        \centering
        \begin{tikzpicture}[
                factor/.style={draw=black, thick, rectangle, minimum size=0.8cm, font=\normalsize},
                eqfactor/.style={draw=black, thick, rectangle, minimum size=0.5cm, font=\normalsize},
                edge/.style={thick},
                varlabel/.style={font=\normalsize},
                greycircle/.style={draw=black, thick, circle, fill=gray!30, minimum size=0.5cm, font=\normalsize}
            ]

            % Knoten
            \node[factor, minimum width=2.0cm, label={[varlabel]left:$q_a(o,s,c)$}] (A1) {$p(o|s,c)$};
            \node[factor, below=1.5cm of A1, label={[varlabel]left:$q_b(o)$}] (c1) {$\tilde{p}(o)$};

            % Kanten zwischen Knoten
            \draw[edge] ([yshift=1.5cm,xshift=0.7cm]A1.north) -- ([xshift=0.7cm]A1.north) node[midway, left, varlabel] {$q_i(s)$};
            \draw[edge] (A1.south) -- (c1.north) node[midway, left, varlabel] {$q_j(o)$};
            \draw[edge] ([yshift=1.5cm,xshift=-0.7cm]A1.north) -- ([xshift=-0.7cm]A1.north) node[midway, left, varlabel] {$q_k(c)$};

        \end{tikzpicture}
        \caption{Perception-related factor graph sector with Bethe free energy.}
        \label{fig:36.1}
    \end{subfigure}
    \hfill
    \begin{subfigure}[t]{0.48\textwidth}
        \centering
        \begin{tikzpicture}[
                factor/.style={draw=black, thick, rectangle, minimum size=0.8cm, font=\normalsize},
                eqfactor/.style={draw=black, thick, rectangle, minimum size=0.5cm, font=\normalsize},
                edge/.style={thick},
                varlabel/.style={font=\normalsize},
                greycircle/.style={draw=black, thick, circle, fill=gray!30, minimum size=0.5cm, font=\normalsize}
            ]

            % Knoten
            \node[factor, minimum width=3.0cm, label={[varlabel]left:$q_a(o)q_a(s)q_a(c)$}] (A1) {$p(o|s,c)$};
            \node[factor, below=1.5cm of A1, label={[varlabel]left:$q_b(o)$}] (c1) {$\tilde{p}(o)$};

            % Kanten zwischen Knoten
            \draw[edge] ([yshift=1.5cm,xshift=1.2cm]A1.north) -- ([xshift=1.2cm]A1.north) node[midway, left, varlabel] {$q_i(s)$} node[pos=0.25, right, varlabel] {$\textcolor{black}{\bm{{\uparrow}{\mu}(}s\bm{)}}$} node[pos=0.75, right, varlabel] {$\bm{{\downarrow}{\mu}(}s\bm{)}$} {} node[pos=1, circle, draw=black, fill=white, inner sep=2pt] {};
            \draw[edge] (A1.south) -- (c1.north) node[midway, left, varlabel] {$q_j(o)$} node[pos=0, rectangle, draw=black, fill=white, inner sep=2.5pt] {};
            \draw[edge] ([yshift=1.5cm,xshift=-1.2cm]A1.north) -- ([xshift=-1.2cm]A1.north) node[midway, left, varlabel] {$q_k(c)$} node[pos=0.25, right, varlabel] {$\textcolor{accorange}{\bm{{\uparrow}{\mu}(}c\bm{)}}$} node[pos=0.75, right, varlabel] {$\bm{{\downarrow}{\mu}(}c\bm{)}$} node[pos=1, circle, draw=black, fill=white, inner sep=2pt] {};

        \end{tikzpicture}
        \caption{Perception-related factor graph sector with generalized free energy and corresponding GFE-based message (\(\textcolor{accorange}{\bm{{\uparrow}{\mu}(}c\bm{)}}\)) for  contextual information seeking.}
        \label{fig:36.2}
    \end{subfigure}
    \caption{Perception-related factor graph sectors.}
    \label{fig:36}
\end{figure}

Starting point for this derivation is the Lagrangian \(L\) of this perception-related factor graph sector with Bethe free energy (\cref{fig:36.1}):
\begin{alignat}{2}
    L = & \sum_{o,s,c} q_a(o,s,c) \ln \frac{q_a(o,s,c)}{p(o|s,c)} + \sum_o q_b(o) \ln \frac{q_b(o)}{\tilde{p}(o)} + \sum_o q_j(o) \ln \frac{1}{q_j(o)} \nonumber \\
        & + \sum_c \lambda_{ka}(c) \left( q_k(c) - \sum_{o,s} q_a(o,s,c) \right) + \sum_s \lambda_{ia}(s) \left( q_i(s) - \sum_{o,c} q_a(o,s,c) \right) \nonumber \\
        & + \sum_o \lambda_{ja}(o) \left( q_j(o) - \sum_{s,c} q_a(o,s,c) \right) + \sum_o \lambda_{jb}(o) \left( q_j(o) - \sum_{s,c} q_a(o,s,c) \right)
\end{alignat}

Similar to the GFE-based messages (\cref{sec:gfe}) we also introduce a mutual information term to the Lagrangian. In this case the term accounts for the variables \(o\) and \(c\), which means the agent is inclined to reduce its contextual uncertainty from observations. The resulting Lagrangain \(L_{\text{MI}}\) is:
\begin{equation}
    L_{\text{MI}} = L + \sum_{o,s} q_a(o,s) \ln \frac{q_a(o)q_a(s)}{q_a(o,s)}
\end{equation}

Next, we add \(0 = \sum_s q_a(s) \ln \frac{q_a(s)}{q_a(s)}\) to \(L_{\text{MI}}\) and also assume independence between \(s\) and the two variables \(o\) and \(c\), i.e., \(q_a(o,s,c) = q_a(o,c)q_a(s)\), which allows for the following rearrangement of \(L_{\text{MI}}\) (compare with \cref{sec:gfe}):
\begin{alignat}{2}
    L_{\text{MI}} = & \sum_{o,s,c} q_a(o,s,c) \ln \frac{q_a(o)q_a(c)q_a(s)}{p(o|s,c)\tilde{p}(o)} \nonumber \\
    & + \sum_c \lambda_{ka}(c) \left( q_k(c) - \sum_{o,s} q_a(o,s,c) \right) + \sum_s \lambda_{ia}(s) \left( q_i(s) - \sum_{o,c} q_a(o,s,c) \right)
\end{alignat}

Next, we perform the p-substitution from \cref{sec:gfe}. Together with \(q_a(o,s,c) = q_a(o,c)q_a(s) \Rightarrow q_a(s,c) = q_a(s)q_a(c)\) it follows (\(L_{\text{GFE}}\)):
\begin{alignat}{2}
    L_{\text{GFE}} = & \sum_{o,s,c} p(o|s,c)q_a(s)q_a(c) \ln \frac{q_a(o)q_a(c)q_a(s)}{p(o|s,c)\tilde{p}(o)} \nonumber \\
    & + \sum_c \lambda_{ka}(c) \left( q_k(c) - \sum_{o,s} q_a(o,s,c) \right) + \sum_s \lambda_{ia}(s) \left( q_i(s) - \sum_{o,c} q_a(o,s,c) \right)
\end{alignat}

Next, the Lagrangian is differentiated with respect to \(q_a(c)\), set to \(0\), and finally solved for \(q^*_a(c)\):
\begin{alignat}{3}
     & \frac{\delta L_{\text{GFE}}}{\delta q_a(c)} &  & = \sum_{o,s} p(o|s,c)q_a(s) \ln \frac{q_a(o)}{p(o|s,c)\tilde{p}(o)} + \ln q_a(c) + 1 - \lambda_{ka}(c) \overset{!}{=} 0                                \\
     &                                             &  & \Rightarrow q^*_a(c) \propto \exp \left( \sum_{o,s} p(o|s,c)q_a(s) \ln \frac{p(o|s,c)\tilde{p}(o)}{q_a(o)} \right) \exp \left( \lambda_{ka}(c) \right)
\end{alignat}

Finally, we use the result of \cref{sec:1} for the local stationary point of \(q_k(c)\)
\begin{equation}
    q^*_k(c) \propto \exp \left( \lambda_{ka}(c) \right) \exp \left( \lambda_{k\bullet}(c) \right),
\end{equation}
and use the marginalization constraint \(q_k(c) \overset{!}{=} q_a(c)\) to determine the message \(\textcolor{accorange}{\bm{{\uparrow}{\mu}(}c\bm{)}}\) that enables contextual information seeking:
\begin{alignat}{2}
    & q^*_k(c)                                                                      &  & \overset{!}{=} q^*_a(c)                                                                                               \\
    & \exp \left( \lambda_{ka}(c) \right) \exp \left( \lambda_{k\bullet}(s) \right) &  & \propto \exp \left( \sum_{o,s} p(o|s,c)q_a(s) \ln \frac{p(o|s,c)\tilde{p}(o)}{q_a(o)} \right) \exp \left( \lambda_{ka}(c) \right) \\
    & \exp \left( \lambda_{k\bullet}(c) \right)                                     &  & \propto \exp \left( \sum_{o,s} p(o|s,c)q_a(s) \ln \frac{p(o|s,c)\tilde{p}(o)}{q_a(o)} \right)                                     \\
    & \textcolor{accorange}{\bm{{\uparrow}{\mu}(}c\bm{)}}                                                  &  & \propto \exp \left( \sum_{o,s} p(o|s,c)q_a(s) \ln \frac{p(o|s,c)\tilde{p}(o)}{q_a(o)} \right)
\end{alignat}

As mentioned in \cref{sec:gfe}, using this message can lead to unstable message passing algorithms. Consequently, we also use an alternative approach that locally optimizes (via iteration index \(l\)) the message instead (details in \cref{sec:gfe}):
\begin{alignat}{2}
    &q^{(l)}_j(o) &&\propto \sum_{s,c} p(o|s,c) q_i(s) q^{(l-1)}_k(c). \label{eq:46} \\
    & \textcolor{accorange}{\bm{{\uparrow}{\mu}}^{(l)}\bm{(}c\bm{)}}                                                  &  & \propto \exp \left( \sum_{o,s} p(o|s,c)q_a(s) \ln \frac{p(o|s,c)\tilde{p}(o)}{q^{(l)}_a(o)} \right) \label{eq:47}\\
    &q^{(l)}_k(c) &&\propto \textcolor{accorange}{\bm{{\uparrow}{\mu}}^{(l)}\bm{(}c\bm{)}} \bm{{\downarrow}{\mu}(}c\bm{)} \label{eq:48}\\
    &\textcolor{accorange}{\bm{{\uparrow}{\mu}(}c\bm{)}} &&\propto \frac{q^*_k(c)}{\bm{{\downarrow}{\mu}(}c\bm{)}} \label{eq:49}
\end{alignat}
Equations \cref{eq:46,eq:47,eq:48} are iterated until convergence of \(q_k(c)\) (with \(q_j(o) = q_a(o)\)), which is then used in \cref{eq:49}.

Moreover, the Lagrangian \(L_{\text{GFE}}\) also leads to a GFE-based message \(\bm{{\uparrow}{\mu}(}s\bm{)}\), i.e., for the hidden state \(s\). The derivation is analog to the prior derivation for the GFE-based message over the context. The resulting optimization loop (analog to \cref{eq:46,eq:47,eq:48})) that leads to the message \(\bm{{\uparrow}{\mu}(}s\bm{)}\) is given in the following:
\begin{alignat}{2}
    &q^{(l)}_j(o) &&\propto \sum_{s,c} p(o|s,c) q^{(l-1)}_i(s) q_k(c). \\
    & \bm{{\uparrow}{\mu}}^{(l)}\bm{(}s\bm{)}                                                  &  & \propto \exp \left( \sum_{o,c} p(o|s,c)q_a(c) \ln \frac{p(o|s,c)\tilde{p}(o)}{q^{(l)}_a(o)} \right) \\
    &q^{(l)}_i(s) &&\propto \bm{{\uparrow}{\mu}}^{(l)}\bm{(}s\bm{)} \bm{{\downarrow}{\mu}(}s\bm{)}\\
    &\bm{{\uparrow}{\mu}(}s\bm{)} &&\propto \frac{q^*_i(s)}{\bm{{\downarrow}{\mu}(}s\bm{)}}
\end{alignat}


\subsubsection{Message Passing Algorithm}
In this section I give the message passing algorithm of the Context-GFE agent, which is based on its factor graph (\cref{fig:38}) that again exemplary includes an observation \(\hat {o_0}\). Analog to the algorithm of the GFE agent, the message passing algorithm of the Context-GFE agent also requires recursive execution and the initialization of the belief distributions \(q(o_0)\), \(q(s^{\scriptscriptstyle II}_0)\), \(q(s^{\scriptscriptstyle II}_1)\), \(q(c^{\scriptscriptstyle IV})\), \(q(s_2)\), \(q(c^{\scriptscriptstyle III})\), \(q(c^{\scriptscriptstyle II})\) and messages \(\bm{{\downarrow}{\mu}(}c^{\scriptscriptstyle IV}\bm{)}\), \(\bm{{\downarrow}{\mu}(}c^{\scriptscriptstyle III}\bm{)}\), \(\bm{{\downarrow}{\mu}(}s^{\scriptscriptstyle II}_1\bm{)}\), \(\bm{{\downarrow}{\mu}(}s_2\bm{)}\) (e.g., uniformly) before it is started. The resulting message passing algorithm of the Context-GFE agent is given in \cref{alg:5}.
\begin{figure}[htbp]
    \centering

    \resizebox{1.0\textwidth}{!}{%
        \begin{tikzpicture}[
                factor/.style={draw=black, thick, rectangle, minimum size=0.8cm, font=\normalsize},
                eqfactor/.style={draw=black, thick, rectangle, minimum size=0.5cm, font=\normalsize},
                edge/.style={thick},
                varlabel/.style={font=\normalsize},
                greycircle/.style={draw=black, thick, circle, fill=gray!30, minimum size=0.5cm, font=\normalsize}
            ]

            % Knoten
            \node[factor] (d) {$p(s_0)$};
            \node[eqfactor, right=1.8cm of d] (eq0) {$=$};
            \node[factor] (A0) at ([xshift=1.5cm, yshift=-2.25cm] d.east) {$p(o_0|s^{\scriptscriptstyle II}_0,c^{\scriptscriptstyle II})$};
            \node[factor, below=1.6cm of A0] (c0) {$\tilde{p}(o_0)$};
            \node[factor, right=1.8cm of eq0] (B0) {$p(s_1|s^{\scriptscriptstyle I}_0,a_0)$};
            \node[eqfactor, right=1.8cm of B0] (eq1) {$=$};
            \node[factor] (A1) at ([xshift=1.5cm, yshift=-2.25cm] B0.east) {$p(o_1|s^{\scriptscriptstyle II}_1,c^{\scriptscriptstyle IV})$};
            \node[factor, below=1.6cm of A1] (c1) {$\tilde{p}(o_1)$};
            \node[factor, right=1.8cm of eq1] (B1) {$p(s_2|s^{\scriptscriptstyle I}_1,a_1)$};
            \node[factor] (A2) at ([xshift=1.5cm] B1.east |- A1) {$p(o_2|s_2,c^{\scriptscriptstyle III})$};
            \node[factor, below=1.6cm of A2] (c2) {$\tilde{p}(o_2)$};
            \node[factor, above=1.6cm of B0] (u0) {$p(a_0)$};
            \node[factor, above=1.6cm of B1] (u1) {$p(a_1)$};
            \node[eqfactor, above right=3cm and 0.65cm of d] (eq3) {$=$};
            \node[eqfactor, above right=3cm and 0.65cm of B0] (eq4) {$=$};
            \node[factor, left=1.8cm of eq3] (dc) {$p(c)$};

            % Kanten zwischen Knoten
            \draw[edge] (d.east) -- (eq0.west) node[pos=0.77, above, varlabel] {$\textcolor{accblue}{q(s_0)}$} node[near start, below, varlabel, scale=0.7] {$\textcolor{accorange}{\bm{\overleftarrow{\mu}(}s_0\bm{)}}$} node[near end, below, varlabel, scale=0.7] {$\textcolor{accorange}{\bm{\overrightarrow{\mu}(}s_0\bm{)}}$};
            \draw[edge] (eq0.south) -- (eq0.south |- A0.north) node[midway, left=-0.05, varlabel] {$\textcolor{accblue}{q(s^{\scriptscriptstyle II}_0)}$} node[pos=0.25, right, varlabel, scale=0.7] {$\textcolor{accorange}{\bm{{\uparrow}{\mu}(}s^{\scriptscriptstyle II}_0\bm{)}}$} node[pos=0.75, right, varlabel, scale=0.7] {$\textcolor{accorange}{\bm{{\downarrow}{\mu}(}s^{\scriptscriptstyle II}_0\bm{)}}$}  node[pos=1, circle, draw=black, fill=white, inner sep=2pt] {};
            \draw[edge] (A0.south) -- (c0.north) node[midway, greycircle] {$\delta$} node[midway, left=0.3, varlabel] {$\textcolor{accblue}{q(o_0)}$} node[pos=0.25, right=0.25, varlabel, scale=0.7] {$\textcolor{accorange}{\bm{{\uparrow}{\mu}(}o_0\bm{)}}$} node[pos=0.75, right=0.25, varlabel, scale=0.7] {$\textcolor{accorange}{\bm{{\downarrow}{\mu}(}o_0\bm{)}}$} node[pos=0, circle, draw=black, fill=white, inner sep=2pt] {};
            \draw[edge] (eq0.east) -- (B0.west) node[midway, above, varlabel] {$\textcolor{accblue}{q(s^{\scriptscriptstyle I}_0)}$} node[near start, below, varlabel, scale=0.7] {$\textcolor{accorange}{\bm{\overleftarrow{\mu}(}s^{\scriptscriptstyle I}_0\bm{)}}$} node[near end, below, varlabel, scale=0.7] {$\textcolor{accorange}{\bm{\overrightarrow{\mu}(}s^{\scriptscriptstyle I}_0\bm{)}}$};
            \draw[edge] (B0.north) -- (u0.south) node[midway, left, varlabel] {$\textcolor{accblue}{q(a_0)}$} node[pos=0.75, right, varlabel, scale=0.7] {$\textcolor{accorange}{\bm{{\uparrow}{\mu}(}a_0\bm{)}}$} node[pos=0.25, right, varlabel, scale=0.7] {$\textcolor{accorange}{\bm{{\downarrow}{\mu}(}a_0\bm{)}}$};
            \draw[edge] (B0.east) -- (eq1.west) node[pos=0.77, above, varlabel] {$\textcolor{accblue}{q(s_1)}$} node[near start, below, varlabel, scale=0.7] {$\textcolor{accorange}{\bm{\overleftarrow{\mu}(}s_1\bm{)}}$} node[near end, below, varlabel, scale=0.7] {$\textcolor{accorange}{\bm{\overrightarrow{\mu}(}s_1\bm{)}}$};
            \draw[edge] (eq1.south) -- (eq1.south |- A1.north) node[midway, left, varlabel] {$\textcolor{accblue}{q(s^{\scriptscriptstyle II}_1)}$} node[pos=0.25, right, varlabel, scale=0.7] {$\textcolor{accorange}{\bm{{\uparrow}{\mu}(}s^{\scriptscriptstyle II}_1\bm{)}}$} node[pos=0.75, right, varlabel, scale=0.7] {$\textcolor{accorange}{\bm{{\downarrow}{\mu}(}s^{\scriptscriptstyle II}_1\bm{)}}$}  node[pos=1, circle, draw=black, fill=white, inner sep=2pt] {};
            \draw[edge] (A1.south) -- (c1.north) node[midway, left, varlabel] {$\textcolor{accblue}{q(o_1)}$} node[pos=0, rectangle, draw=black, fill=white, inner sep=2.5pt] {};
            \draw[edge] (eq1.east) -- (B1.west) node[midway, above, varlabel] {$\textcolor{accblue}{q(s^{\scriptscriptstyle I}_1)}$} node[near start, below, varlabel, scale=0.7] {$\textcolor{accorange}{\bm{\overleftarrow{\mu}(}s^{\scriptscriptstyle I}_1\bm{)}}$} node[near end, below, varlabel, scale=0.7] {$\textcolor{accorange}{\bm{\overrightarrow{\mu}(}s^{\scriptscriptstyle I}_1\bm{)}}$};
            \draw[edge] (B1.east) -| ([xshift=0.57cm]A2.north) node[pos=0.37, above, varlabel] {$\textcolor{accblue}{q(s_2)}$} node[pos=0.11, below, varlabel, scale=0.7] {$\textcolor{accorange}{\bm{\overleftarrow{\mu}(}s_2\bm{)}}$} node[pos=0.89, right, varlabel, scale=0.7] {$\textcolor{accorange}{\bm{{\downarrow}{\mu}(}s_2\bm{)}}$}  node[pos=1, circle, draw=black, fill=white, inner sep=2pt] {};
            \draw[edge] (A2.south) -- (c2.north) node[midway, left, varlabel] {$\textcolor{accblue}{q(o_2)}$} node[pos=0, rectangle, draw=black, fill=white, inner sep=2.5pt] {};
            \draw[edge] (B1.north) -- (u1.south) node[midway, left, varlabel] {$\textcolor{accblue}{q(a_1)}$} node[pos=0.75, right, varlabel, scale=0.7] {$\textcolor{accorange}{\bm{{\uparrow}{\mu}(}a_1\bm{)}}$} node[pos=0.25, right, varlabel, scale=0.7] {$\textcolor{accorange}{\bm{{\downarrow}{\mu}(}a_1\bm{)}}$};
            \draw[edge] (dc.east) -- (eq3.west) node[midway, above, varlabel] {$\textcolor{accblue}{q(c)}$} node[near start, below, varlabel, scale=0.7] {$\textcolor{accorange}{\bm{\overleftarrow{\mu}(}c\bm{)}}$} node[near end, below, varlabel, scale=0.7] {$\textcolor{accorange}{\bm{\overrightarrow{\mu}(}c\bm{)}}$};
            \draw[edge] (eq3.east) -- (eq4.west) node[midway, above, varlabel] {$\textcolor{accblue}{q(c^{\scriptscriptstyle I})}$} node[pos=0.08, below, varlabel, scale=0.7] {$\textcolor{accorange}{\bm{\overleftarrow{\mu}(}c^{\scriptscriptstyle I}\bm{)}}$} node[pos=0.92, below, varlabel, scale=0.7] {$\textcolor{accorange}{\bm{\overrightarrow{\mu}(}c^{\scriptscriptstyle I}\bm{)}}$};
            \draw[edge] (eq4.east) -| ([xshift=-0.5cm]A2.north) node[pos=0.25, above, varlabel] {$\textcolor{accblue}{q(c^{\scriptscriptstyle III})}$} node[pos=0.045, below, varlabel, scale=0.7] {$\textcolor{accorange}{\bm{\overleftarrow{\mu}(}c^{\scriptscriptstyle III}\bm{)}}$} node[pos=0.965, right=-0.05, varlabel, scale=0.7] {$\textcolor{accorange}{\bm{{\downarrow}{\mu}(}c^{\scriptscriptstyle III}\bm{)}}$} node[pos=1, circle, draw=black, fill=white, inner sep=2pt] {};
            \draw[edge] (eq3.south) -- (eq3.south |- A0.north) node[pos=0.2, left, varlabel] {$\textcolor{accblue}{q(c^{\scriptscriptstyle II})}$} node[pos=0.07, right, varlabel, scale=0.7] {$\textcolor{accorange}{\bm{{\uparrow}{\mu}(}c^{\scriptscriptstyle II}\bm{)}}$} node[pos=0.93, right, varlabel, scale=0.7] {$\textcolor{accorange}{\bm{{\downarrow}{\mu}(}c^{\scriptscriptstyle II}\bm{)}}$} node[pos=1, circle, draw=black, fill=white, inner sep=2pt] {};
            \draw[edge] (eq4.south) -- (eq4.south |- A1.north) node[pos=0.2, left, varlabel] {$\textcolor{accblue}{q(c^{\scriptscriptstyle IV})}$} node[pos=0.07, right, varlabel, scale=0.7] {$\textcolor{accorange}{\bm{{\uparrow}{\mu}(}c^{\scriptscriptstyle IV}\bm{)}}$} node[pos=0.93, right, varlabel, scale=0.7] {$\textcolor{accorange}{\bm{{\downarrow}{\mu}(}c^{\scriptscriptstyle IV}\bm{)}}$} node[pos=1, circle, draw=black, fill=white, inner sep=2pt] {};

        \end{tikzpicture}
    }
    \caption{Factor Graph of the Context-GFE agent with variables' belief distributions \(\textcolor{accblue}{q}\) and messages \(\textcolor{accorange}{\mu}\).}
    \label{fig:38}
\end{figure}

\begin{algorithm}[htbp]
    \caption{Message Passing Algorithm of the Context-GFE Agent}
    \label{alg:5}
    \small
    \begin{algorithmic}[1]
        \linespread{1.4}\selectfont
        \setlength{\leftskip}{0.3cm}
        \State \tikzmark{algtop5}$\textcolor{accorange}{\bm{\overrightarrow{\mu}(}s_0\bm{)}} \propto p(s_0)$,\quad$\textcolor{accorange}{\bm{{\uparrow}{\mu}(}o_0\bm{)}} \propto \delta_{o_0,\hat{o_0}}$,\quad$\textcolor{accorange}{\bm{{\downarrow}{\mu}(}a_0\bm{)}} \propto p(a_0)$,\quad$\textcolor{accorange}{\bm{{\downarrow}{\mu}(}a_1\bm{)}} \propto p(a_1)$,\quad$\textcolor{accorange}{\bm{{\downarrow}{\mu}(}o_0\bm{)}} \propto \delta_{o_0,\hat{o_0}}$,
        \Statex $\textcolor{accorange}{\bm{\overrightarrow{\mu}(}c\bm{)}} \propto p(c)$,\quad$\textcolor{accorange}{\bm{{\uparrow}{\mu}(}c^{\scriptscriptstyle II}\bm{)}} \propto \exp \left( \sum_{o_0,s^{\scriptscriptstyle II}_0} q(o_0)q(s^{\scriptscriptstyle II}_0) \ln p(o_0|s^{\scriptscriptstyle II}_0,c^{\scriptscriptstyle II}) \right)$,
        \Statex \quad\quad\tikzmark{onetop5}$q^{(l)}(o_1) \propto \sum_{s^{\scriptscriptstyle II}_1,c^{\scriptscriptstyle IV}} p(o_1|s^{\scriptscriptstyle II}_1,c^{\scriptscriptstyle IV}) q(s^{\scriptscriptstyle II}_1) q^{(l-1)}(c^{\scriptscriptstyle IV})$
        \Statex \quad\quad$\bm{{\uparrow}{\mu}}^{(l)}\bm{(}c^{\scriptscriptstyle IV}\bm{)} \propto \exp \left( \sum_{o_1,s^{\scriptscriptstyle II}_1} p(o_1|s^{\scriptscriptstyle II}_1,c^{\scriptscriptstyle IV}) q(s^{\scriptscriptstyle II}_1) \ln \frac{p(o_1|s^{\scriptscriptstyle II}_1,c^{\scriptscriptstyle IV})\tilde{p}(o_1)}{q^{(l)}(o_1)} \right)$
        \Statex \quad\quad\tikzmark{onebottom5}$q^{(l)}(c^{\scriptscriptstyle IV}) \propto \bm{{\uparrow}{\mu}}^{(l)}\bm{(}c^{\scriptscriptstyle IV}\bm{)} \bm{{\downarrow}{\mu}(}c^{\scriptscriptstyle IV}\bm{)}$
        \Statex $\textcolor{accorange}{\bm{{\uparrow}{\mu}(}c^{\scriptscriptstyle IV}\bm{)}} \propto {q^*(c^{\scriptscriptstyle IV})}/{\bm{{\downarrow}{\mu}(}c^{\scriptscriptstyle IV}\bm{)}}$,
        \Statex \quad\quad\tikzmark{twotop5}$q^{(l)}(o_2) \propto \sum_{s_2,c^{\scriptscriptstyle III}} p(o_2|s_2,c^{\scriptscriptstyle III}) q(s_2) q^{(l-1)}(c^{\scriptscriptstyle III})$
        \Statex \quad\quad$\textcolor{black}{\bm{\overleftarrow{\mu}}^{(l)}\bm{(}c^{\scriptscriptstyle III}\bm{)}} \propto \exp \left( \sum_{o_2,s_2} p(o_2|s_2,c^{\scriptscriptstyle III}) q(s_2) \ln \frac{p(o_2|s_2,c^{\scriptscriptstyle III})\tilde{p}(o_2)}{q^{(l)}(o_2)} \right)$
        \Statex \quad\quad\tikzmark{twobottom5}$q^{(l)}(c^{\scriptscriptstyle III}) \propto \textcolor{black}{\bm{\overleftarrow{\mu}}^{(l)}\bm{(}c^{\scriptscriptstyle III}\bm{)}} \bm{{\downarrow}{\mu}(}c^{\scriptscriptstyle III}\bm{)}$
        \Statex $\textcolor{accorange}{\bm{\overleftarrow{\mu}(}c^{\scriptscriptstyle III}\bm{)}} \propto {q^*(c^{\scriptscriptstyle III})}/{\bm{{\downarrow}{\mu}(}c^{\scriptscriptstyle III}\bm{)}}$,
        \Statex $\textcolor{accorange}{\bm{{\uparrow}{\mu}(}s^{\scriptscriptstyle II}_0\bm{)}} \propto \exp \left( \sum_{c^{\scriptscriptstyle II},o_0} q(c^{\scriptscriptstyle II}) q(o_0) \ln p(o_0|s^{\scriptscriptstyle II}_0,c^{\scriptscriptstyle II}) \right)$,
        \Statex \quad\quad\tikzmark{threetop5}$q^{(l)}(o_1) \propto \sum_{s^{\scriptscriptstyle II}_1,c^{\scriptscriptstyle IV}} p(o_1|s^{\scriptscriptstyle II}_1,c^{\scriptscriptstyle IV}) q^{(l-1)}(s^{\scriptscriptstyle II}_1) q(c^{\scriptscriptstyle IV})$
        \Statex \quad\quad$\bm{{\uparrow}{\mu}}^{(l)}\bm{(}s^{\scriptscriptstyle II}_1\bm{)} \propto \exp \left( \sum_{o_1,c^{\scriptscriptstyle IV}} p(o_1|s^{\scriptscriptstyle II}_1,c^{\scriptscriptstyle IV}) q(c^{\scriptscriptstyle IV}) \ln \frac{p(o_1|s^{\scriptscriptstyle II}_1,c^{\scriptscriptstyle IV})\tilde{p}(o_1)}{q^{(l)}(o_1)} \right)$
        \Statex \quad\quad\tikzmark{threebottom5}$q^{(l)}(s^{\scriptscriptstyle II}_1) \propto \bm{{\uparrow}{\mu}}^{(l)}\bm{(}s^{\scriptscriptstyle II}_1\bm{)} \bm{{\downarrow}{\mu}(}s^{\scriptscriptstyle II}_1\bm{)}$
        \Statex $\textcolor{accorange}{\bm{{\uparrow}{\mu}(}s^{\scriptscriptstyle II}_1\bm{)}} \propto {q^*(s^{\scriptscriptstyle II}_1)}/{\bm{{\downarrow}{\mu}(}s^{\scriptscriptstyle II}_1\bm{)}}$,
        \Statex \quad\quad\tikzmark{fourtop5}$q^{(l)}(o_2) \propto \sum_{s_2,c^{\scriptscriptstyle III}} p(o_2|s_2,c^{\scriptscriptstyle III}) q^{(l-1)}(s_2) q(c^{\scriptscriptstyle III})$
        \Statex \quad\quad$\textcolor{black}{\bm{\overleftarrow{\mu}}^{(l)}\bm{(}s_2\bm{)}} \propto \exp \left( \sum_{o_2,c^{\scriptscriptstyle III}} p(o_2|s_2,c^{\scriptscriptstyle III}) q(c^{\scriptscriptstyle III}) \ln \frac{p(o_2|s_2,c^{\scriptscriptstyle III})\tilde{p}(o_2)}{q^{(l)}(o_2)} \right)$
        \Statex \quad\quad\tikzmark{fourbottom5}$q^{(l)}(s_2) \propto \textcolor{black}{\bm{\overleftarrow{\mu}}^{(l)}\bm{(}s_2\bm{)}} \bm{{\downarrow}{\mu}(}s_2\bm{)}$
        \Statex $\textcolor{accorange}{\bm{\overleftarrow{\mu}(}s_2\bm{)}} \propto {q^*(s_2)}/{\bm{{\downarrow}{\mu}(}s_2\bm{)}}$
        \State $\textcolor{accorange}{\bm{\overrightarrow{\mu}(}c^{\scriptscriptstyle I}\bm{)}} \propto \textcolor{black}{\bm{\overrightarrow{\mu}(}c\bm{)}} \textcolor{black}{\bm{{\uparrow}{\mu}(}c^{\scriptscriptstyle II}\bm{)}}$,\quad$\textcolor{accorange}{\bm{\overrightarrow{\mu}(}s^{\scriptscriptstyle I}_0\bm{)}} \propto \sum_{s_0,s^{\scriptscriptstyle II}_0} \delta_{s_0,s^{\scriptscriptstyle I}_0}\delta_{s_0,s^{\scriptscriptstyle II}_0} \textcolor{black}{\bm{\overrightarrow{\mu}(}s_0\bm{)}} \textcolor{black}{\bm{{\uparrow}{\mu}(}s^{\scriptscriptstyle II}_0\bm{)}}$,
        \Statex $\textcolor{accorange}{\bm{\overleftarrow{\mu}(}c^{\scriptscriptstyle I}\bm{)}} \propto \textcolor{black}{\bm{{\uparrow}{\mu}(}c^{\scriptscriptstyle IV}\bm{)}} \textcolor{black}{\bm{\overleftarrow{\mu}(}c^{\scriptscriptstyle III}\bm{)}}$,\quad$\textcolor{accorange}{\bm{\overleftarrow{\mu}(}s^{\scriptscriptstyle I}_1\bm{)}} \propto \sum_{a_1,s_2} p(s_2|s^{\scriptscriptstyle I}_1,a_1) \textcolor{black}{\bm{{\downarrow}{\mu}(}a_1\bm{)}} \textcolor{black}{\bm{\overleftarrow{\mu}(}s_2\bm{)}}$
        \State $\textcolor{accorange}{\bm{{\downarrow}{\mu}(}c^{\scriptscriptstyle III}\bm{)}} \propto \textcolor{black}{\bm{\overrightarrow{\mu}(}c^{\scriptscriptstyle I}\bm{)}} \textcolor{black}{\bm{{\uparrow}{\mu}(}c^{\scriptscriptstyle IV}\bm{)}}$,\quad$\textcolor{accorange}{\bm{{\downarrow}{\mu}(}c^{\scriptscriptstyle IV}\bm{)}} \propto \textcolor{black}{\bm{\overrightarrow{\mu}(}c^{\scriptscriptstyle I}\bm{)}} \textcolor{black}{\bm{\overleftarrow{\mu}(}c^{\scriptscriptstyle III}\bm{)}}$,\quad$\textcolor{accorange}{\bm{{\downarrow}{\mu}(}c^{\scriptscriptstyle II}\bm{)}} \propto \textcolor{black}{\bm{\overrightarrow{\mu}(}c\bm{)}} \textcolor{black}{\bm{\overleftarrow{\mu}(}c^{\scriptscriptstyle I}\bm{)}}$,
        \Statex $\textcolor{accorange}{\bm{\overleftarrow{\mu}(}c\bm{)}} \propto \textcolor{black}{\bm{{\uparrow}{\mu}(}c^{\scriptscriptstyle II}\bm{)}} \textcolor{black}{\bm{\overleftarrow{\mu}(}c^{\scriptscriptstyle I}\bm{)}}$,\quad$\textcolor{accorange}{\bm{\overrightarrow{\mu}(}s_1\bm{)}} \propto \sum_{s^{\scriptscriptstyle I}_0,a_0} p(s_1|s^{\scriptscriptstyle I}_0,a_0) \textcolor{black}{\bm{\overrightarrow{\mu}(}s^{\scriptscriptstyle I}_0\bm{)}} \textcolor{black}{\bm{{\downarrow}{\mu}(}a_0\bm{)}}$,
        \Statex $\textcolor{accorange}{\bm{\overleftarrow{\mu}(}s_1\bm{)}} \propto \sum_{s^{\scriptscriptstyle I}_1,s^{\scriptscriptstyle II}_1} \delta_{s_1,s^{\scriptscriptstyle I}_1}\delta_{s_1,s^{\scriptscriptstyle II}_1} \textcolor{black}{\bm{\overleftarrow{\mu}(}s^{\scriptscriptstyle I}_1\bm{)}} \textcolor{black}{\bm{{\uparrow}{\mu}(}s^{\scriptscriptstyle II}_1\bm{)}}$
        \State $\textcolor{accorange}{\bm{\overrightarrow{\mu}(}s^{\scriptscriptstyle I}_1\bm{)}} \propto \sum_{s_1,s^{\scriptscriptstyle II}_1} \delta_{s_1,s^{\scriptscriptstyle I}_1}\delta_{s_1,s^{\scriptscriptstyle II}_1} \textcolor{black}{\bm{\overrightarrow{\mu}(}s_1\bm{)}} \textcolor{black}{\bm{{\uparrow}{\mu}(}s^{\scriptscriptstyle II}_1\bm{)}}$,\quad$\textcolor{accorange}{\bm{\overleftarrow{\mu}(}s^{\scriptscriptstyle I}_0\bm{)}} \propto \sum_{s_1,a_0} p(s_1|s^{\scriptscriptstyle I}_0,a_0) \textcolor{black}{\bm{{\downarrow}{\mu}(}a_0\bm{)}} \textcolor{black}{\bm{\overleftarrow{\mu}(}s_1\bm{)}}$,
        \Statex $\textcolor{accorange}{\bm{{\downarrow}{\mu}(}s^{\scriptscriptstyle II}_1\bm{)}} \propto \sum_{s_1,s^{\scriptscriptstyle I}_1} \delta_{s_1,s^{\scriptscriptstyle I}_1}\delta_{s_1,s^{\scriptscriptstyle II}_1} \textcolor{black}{\bm{\overrightarrow{\mu}(}s_1\bm{)}} \textcolor{black}{\bm{\overleftarrow{\mu}(}s^{\scriptscriptstyle I}_1\bm{)}}$,\quad$\textcolor{accorange}{\bm{{\uparrow}{\mu}(}a_0\bm{)}} \propto \sum_{s_1,s^{\scriptscriptstyle I}_0} p(s_1|s^{\scriptscriptstyle I}_0,a_0) \textcolor{black}{\bm{\overrightarrow{\mu}(}s^{\scriptscriptstyle I}_0\bm{)}} \textcolor{black}{\bm{\overleftarrow{\mu}(}s_1\bm{)}}$
        \State $\textcolor{accorange}{\bm{{\uparrow}{\mu}(}a_1\bm{)}} \propto \sum_{s^{\scriptscriptstyle I}_1,s_2} p(s_2|s^{\scriptscriptstyle I}_1,a_1) \textcolor{black}{\bm{\overrightarrow{\mu}(}s^{\scriptscriptstyle I}_1\bm{)}} \textcolor{black}{\bm{\overleftarrow{\mu}(}s_2\bm{)}}$,\quad$\textcolor{accorange}{\bm{{\downarrow}{\mu}(}s_2\bm{)}} \propto \sum_{s^{\scriptscriptstyle I}_1,a_1} p(s_2|s^{\scriptscriptstyle I}_1,a_1) \textcolor{black}{\bm{\overrightarrow{\mu}(}s^{\scriptscriptstyle I}_1\bm{)}} \textcolor{black}{\bm{{\downarrow}{\mu}(}a_1\bm{)}}$,
        \Statex $\textcolor{accorange}{\bm{\overleftarrow{\mu}(}s_0\bm{)}} \propto \sum_{s^{\scriptscriptstyle I}_0,s^{\scriptscriptstyle II}_0} \delta_{s_0,s^{\scriptscriptstyle I}_0}\delta_{s_0,s^{\scriptscriptstyle II}_0} \textcolor{black}{\bm{\overleftarrow{\mu}(}s^{\scriptscriptstyle I}_0\bm{)}} \textcolor{black}{\bm{{\uparrow}{\mu}(}s^{\scriptscriptstyle II}_0\bm{)}}$,\quad$\textcolor{accorange}{\bm{{\downarrow}{\mu}(}s^{\scriptscriptstyle II}_0\bm{)}} \propto \sum_{s_0,s^{\scriptscriptstyle I}_0} \delta_{s_0,s^{\scriptscriptstyle I}_0}\delta_{s_0,s^{\scriptscriptstyle II}_0} \textcolor{black}{\bm{\overrightarrow{\mu}(}s_0\bm{)}} \textcolor{black}{\bm{\overleftarrow{\mu}(}s^{\scriptscriptstyle I}_0\bm{)}}$
        \State \tikzmark{algbottom5} $\textcolor{accblue}{q(s_0)} \propto \textcolor{black}{\bm{\overleftarrow{\mu}(}s_0\bm{)}} \textcolor{black}{\bm{\overrightarrow{\mu}(}s_0\bm{)}}$,\quad$\textcolor{accblue}{q(s^{\scriptscriptstyle I}_0)} \propto \textcolor{black}{\bm{\overleftarrow{\mu}(}s^{\scriptscriptstyle I}_0\bm{)}} \textcolor{black}{\bm{\overrightarrow{\mu}(}s^{\scriptscriptstyle I}_0\bm{)}}$,\quad$\textcolor{accblue}{q(s^{\scriptscriptstyle II}_0)} \propto \textcolor{black}{\bm{{\uparrow}{\mu}(}s^{\scriptscriptstyle II}_0\bm{)}} \textcolor{black}{\bm{{\downarrow}{\mu}(}s^{\scriptscriptstyle II}_0\bm{)}}$,
        \Statex $\textcolor{accblue}{q(s_1)} \propto \textcolor{black}{\bm{\overleftarrow{\mu}(}s_1\bm{)}} \textcolor{black}{\bm{\overrightarrow{\mu}(}s_1\bm{)}}$,\quad$\textcolor{accblue}{q(s^{\scriptscriptstyle I}_1)} \propto \textcolor{black}{\bm{\overleftarrow{\mu}(}s^{\scriptscriptstyle I}_1\bm{)}} \textcolor{black}{\bm{\overrightarrow{\mu}(}s^{\scriptscriptstyle I}_1\bm{)}}$,\quad$\textcolor{accblue}{q(s^{\scriptscriptstyle II}_1)} \propto \textcolor{black}{\bm{{\uparrow}{\mu}(}s^{\scriptscriptstyle II}_1\bm{)}} \textcolor{black}{\bm{{\downarrow}{\mu}(}s^{\scriptscriptstyle II}_1\bm{)}}$,
        \Statex $\textcolor{accblue}{q(s_2)} \propto \textcolor{black}{\bm{\overleftarrow{\mu}(}s_2\bm{)}} \textcolor{black}{\bm{{\downarrow}{\mu}(}s_2\bm{)}}$,\quad$\textcolor{accblue}{q(a_0)} \propto \textcolor{black}{\bm{{\uparrow}{\mu}(}a_0\bm{)}} \textcolor{black}{\bm{{\downarrow}{\mu}(}a_0\bm{)}}$,\quad$\textcolor{accblue}{q(a_1)} \propto \textcolor{black}{\bm{{\uparrow}{\mu}(}a_1\bm{)}} \textcolor{black}{\bm{{\downarrow}{\mu}(}a_1\bm{)}}$,
        \Statex $\textcolor{accblue}{q(o_0)} \propto \textcolor{black}{\bm{{\uparrow}{\mu}(}o_0\bm{)}} \textcolor{black}{\bm{{\downarrow}{\mu}(}o_0\bm{)}}$,\quad$\textcolor{accblue}{q(o_1)} \propto \sum_{s^{\scriptscriptstyle II}_1,c^{\scriptscriptstyle IV}} p(o_1|s^{\scriptscriptstyle II}_1,c^{\scriptscriptstyle IV}) q(s^{\scriptscriptstyle II}_1) q(c^{\scriptscriptstyle IV})$,\quad$\textcolor{accblue}{q(c)} \propto \textcolor{black}{\bm{\overleftarrow{\mu}(}c\bm{)}} \textcolor{black}{\bm{\overrightarrow{\mu}(}c\bm{)}}$,
        \Statex $\textcolor{accblue}{q(o_2)} \propto \sum_{s_2,c^{\scriptscriptstyle III}} p(o_2|s_2,c^{\scriptscriptstyle III}) q(s_2) q(c^{\scriptscriptstyle III})$,\quad$\textcolor{accblue}{q(c^{\scriptscriptstyle I})} \propto \textcolor{black}{\bm{\overleftarrow{\mu}(}c^{\scriptscriptstyle I}\bm{)}} \textcolor{black}{\bm{\overrightarrow{\mu}(}c^{\scriptscriptstyle I}\bm{)}}$,
        \Statex $\textcolor{accblue}{q(c^{\scriptscriptstyle II})} \propto \textcolor{black}{\bm{{\uparrow}{\mu}(}c^{\scriptscriptstyle II}\bm{)}} \textcolor{black}{\bm{{\downarrow}{\mu}(}c^{\scriptscriptstyle II}\bm{)}}$,\quad$\textcolor{accblue}{q(c^{\scriptscriptstyle III})} \propto \textcolor{black}{\bm{\overleftarrow{\mu}(}c^{\scriptscriptstyle III}\bm{)}} \textcolor{black}{\bm{\overrightarrow{\mu}(}c^{\scriptscriptstyle III}\bm{)}}$,\quad$\textcolor{accblue}{q(c^{\scriptscriptstyle IV})} \propto \textcolor{black}{\bm{{\uparrow}{\mu}(}c^{\scriptscriptstyle IV}\bm{)}} \textcolor{black}{\bm{{\downarrow}{\mu}(}c^{\scriptscriptstyle IV}\bm{)}}$
    \end{algorithmic}

    \begin{tikzpicture}[remember picture, overlay]
        \draw[-{Stealth[length=2.2mm, width=1.6mm]}, thick, gray!75!black, rounded corners=5pt]
        ([yshift=0.5ex, xshift=-2.5ex]pic cs:algbottom5) -- ++(-0.4cm, 0) |- ([xshift=-2.5ex, yshift=0.5ex]pic cs:algtop5);
    \end{tikzpicture}
    \begin{tikzpicture}[remember picture, overlay]
        \draw[-{Stealth[length=2.2mm, width=1.6mm]}, thick, gray!75!black, rounded corners=5pt]
        ([yshift=0.5ex, xshift=-0.5ex]pic cs:onebottom5) -- ++(-0.4cm, 0) |- ([xshift=-0.5ex, yshift=0.5ex]pic cs:onetop5);
    \end{tikzpicture}
    \begin{tikzpicture}[remember picture, overlay]
        \draw[-{Stealth[length=2.2mm, width=1.6mm]}, thick, gray!75!black, rounded corners=5pt]
        ([yshift=0.5ex, xshift=-0.5ex]pic cs:twobottom5) -- ++(-0.4cm, 0) |- ([xshift=-0.5ex, yshift=0.5ex]pic cs:twotop5);
    \end{tikzpicture}
    \begin{tikzpicture}[remember picture, overlay]
        \draw[-{Stealth[length=2.2mm, width=1.6mm]}, thick, gray!75!black, rounded corners=5pt]
        ([yshift=0.5ex, xshift=-0.5ex]pic cs:threebottom5) -- ++(-0.4cm, 0) |- ([xshift=-0.5ex, yshift=0.5ex]pic cs:threetop5);
    \end{tikzpicture}
    \begin{tikzpicture}[remember picture, overlay]
        \draw[-{Stealth[length=2.2mm, width=1.6mm]}, thick, gray!75!black, rounded corners=5pt]
        ([yshift=0.5ex, xshift=-0.5ex]pic cs:fourbottom5) -- ++(-0.4cm, 0) |- ([xshift=-0.5ex, yshift=0.5ex]pic cs:fourtop5);
    \end{tikzpicture}

\end{algorithm}

\section{Simulations}

In this section I present the simulation procederes that reveal the information seeking behavior of the Context-GFE agent. During the simulations the agent interacts with a T-Maze environment that is frequently used for showing these kinds of behaviors. The simulations consist of three experiments: In the first, the agent is evaluated with respect to its tendendency to show information seeking behavior. The second experiment shows the agent's ability to exploit the acquired information to fulfill a goal. In the third experiment the agent's overall success rate of fulfilling a goal (for which information seeking may help) is evaluated.
The experiments are not only performed for the Context-GFE agent but also for the other three agents. This is necessary because the GFE agent in particular serves as a comparison instance (also shows information seeking behavior) for the Context-GFE agent. Moreover, the remaining two agents (BFE- and Context-BFE) serve as baselines for the kind of behavior that results from the absence of the information seeking.
In the following, first I detail the setup of the T-Maze environment, second give the concrete parameter values for the underlying generative process of the T-Maze as well as of the generative models of the agents, and third present the procederes of the three experiments.

\subsection{The T-Maze Environment}

In this section I present the T-Maze environment, which all agents interact with during the simulations. Intuitively, the T-Maze is a maze in which an agent can move (through actions) and similarly make observations at the different locations in the maze. A schematic of the T-Maze is given in \cref{fig:27}. The maze consists of four fields (Left Arm, Right Arm, Start, and Cue) and at every timestep the agent can be located at one of them (here the agent is at the start field). The agent has four different actions that can move it to the four possible fields in the maze. The usual interaction between the agent and the T-Maze during a simulation is as follows.

The agent always begins at the start field. Moreover, the agent has a (strong) preference for observing the reward, which is located either at the left arm or the right arm. Unfortunately, the agent does initially not know at which arm the reward is. Following this, it has two options to find the reward. First, it could randomly move to one of the two arms and receive the reward with \(50\,\%\) chance. Second, it could first move to the cue, which gives away information about the reward location (including uncertainty). Using this knowledge the agent can then make a better decision when moving to one of the two arms. In summary, the simulation lasts at maximum two timesteps during which the agent can either directly move to one of the arms (which ends the simulation after the first timestep) or it can first move to the cue and then move to one of the arms.

\begin{figure}[htbp]
    \centering

    \begin{tikzpicture}[scale=1.5]
        \tikzset{
            agent/.pic={
                    \node[circle, draw=black, thick, fill=white, minimum size=0.5cm] at (0,0) {};
                    \fill[black] (-0.08, 0.08) circle (0.03);
                    \fill[black] (0.08, 0.08) circle (0.03);
                    \draw[thick] (-0.1, -0.05) arc (180:360:0.1);
                },
            field label/.style={
                    font=\tiny\sffamily, text=gray!70!black, anchor=south, scale=0.9, transform shape
                },
            field symbol/.style={
                    font=\large\bfseries,
                    text=gray!70,
                    anchor=center,
                    scale=1.0,
                    transform shape
                }
        }

        \draw (-0.5, 0) -- (-0.5, 2.0) -- (-1.5, 2.0) -- (-1.5, 3.0) -- (1.5, 3.0) -- (1.5, 2.0) -- (0.5, 2.0) -- (0.5, 0) -- cycle;
        \draw (-0.5, 1.0) -- (0.5, 1.0);
        \draw (-0.5, 2.0) -- (0.5, 2.0);
        \draw (-0.5, 2.0) -- (-0.5, 3.0);
        \draw (0.5, 2.0) -- (0.5, 3.0);

        \pic[scale=0.8, transform shape] at (0, 1.5) {agent};
        \node[field label] at (0, 0.97) {Start};
        \node[field label] at (0, -0.03) {Cue};
        \node[field label] at (-1.0, 1.97) {Left Arm};
        \node[field label] at (1.0, 1.93) {Right Arm};
        \node[field symbol] at (0, 0.5) {!};
        \node[field symbol] at (-1.0, 2.5) {?};
        \node[field symbol] at (1.0, 2.5) {?};
    \end{tikzpicture}

    \caption{The figure shows the T-Maze environment with its four fields (Left Arm, Right Arm, Start, and Cue) and the agent (Smiley) that is currently located at the start field. The empty space above the start field is not considered a field because the agent cannot move there.}
    \label{fig:27}
\end{figure}

The T-Maze environment is precisely discribed by a corresponding generative process that relates actions from the agent, hidden environmental states and observations generated by the environment probabilistically. Specifically, observations \(o\) are generated by (depend on) hidden environmental states \(s\) via the likelihood distribution \(p(o|s)\). Hidden environmental states evolve over time (from \(t\) to \(t+1\)) based on actions \(a\) via \(p(s_{t+1}|s_t, a_t)\). At timestep \(t=0\) hidden environmental states are given by the prior \(p(s_0)\). Before giving the probabilities of these relationships, I first define the observations, hidden states, and actions. Specifically, observations and hidden states can be defined as a cartesian product of three other sets. These sets are


\subsection{Parametrization}

Next, I assign concrete probabilities to the distributions of the generative process (of the T-Maze) and the generative models of the agents. There are large overlaps between the generative process and the generative models, which is why I give the concrete values here in a single place. However, it is important to realize that the generative process is in general not equal to the generative models of the agents. In the following I will first present the concrete events for which the distributions give probabilities. 
The events can be categorized into observations, (hidden) states, and actions. Regarding observations and hidden states, these can be defined as a cartesian product of three other sets. These sets are
\begin{equation}
    \text{Field} = \{\text{L}, \text{R}, \text{S}, \text{C}\},
\end{equation}
that defines the four fields (L (Left Arm), R (Right Arm), S (Start), C (Cue)) in the T-Maze,
\begin{equation}
    \text{MazeType} = \{\text{A}, \text{B}\},
\end{equation}
that defines the two possible types of the T-Maze, and
\begin{equation}
    \text{Outcome} = \{\text{RW}, \text{NR}, \text{CL}, \text{CR}\},
\end{equation}
that defines the four possible outcomes (not to be confused with observations) at a field (RW (Reward), NR (No Reward), CL (Cue hints at Left arm), CR (Cue hints at Right arm)).
Observations \(o \in O\) are defined as
\begin{alignat}{2}
    O =\; & \text{Field} \times \text{Outcome}                                                                                                                                             \nonumber \\
    =\;   & \{(\text{L},\text{RW}),(\text{L},\text{NR}),(\text{L},\text{CL}),(\text{L},\text{CR}),(\text{R},\text{RW}),(\text{R},\text{NR}),(\text{R},\text{CL}),(\text{R},\text{CR}),     \nonumber \\
          & \;\;(\text{S},\text{RW}),(\text{S},\text{NR}),(\text{S},\text{CL}),(\text{S},\text{CR}),(\text{C},\text{RW}),(\text{C},\text{NR}),(\text{C},\text{CL}),(\text{C},\text{CR})\},
\end{alignat}
hidden states \(s \in S\) (except for the context agents) are given by
\begin{alignat}{2}
    S =\; & \text{Field} \times \text{MazeType}                                                                                                                                  \nonumber \\
    =\;   & \{(\text{L},\text{A}),(\text{L},\text{B}),(\text{R},\text{A}),(\text{R},\text{B}),(\text{S},\text{A}),(\text{S},\text{B}),(\text{C},\text{A}),(\text{C},\text{B})\},
\end{alignat}
and actions \(a \in A\) are
\begin{equation}
    A = \{\text{GtL},\text{GtR},\text{GtS},\text{GtC}\},
\end{equation}
meaning GtL (Go to Left arm), GtR (Go to Right arm), GtS (Go to Start), and GtC (Go to Cue).

Next, it is important to notice that some parametrizations of distributions are shared among the generative process, the generative models of non-context agents (BFE- and GFE) and the generative models of context agents (Context-BFE- and Context-GFE). Following this, some definitions below refer to the generative process and/or multiple generative models (from non-context- and context agents). 
The initial hidden states \(s \in S\) at timestep \(k=0\) of the generative process (T-Maze) are randomly defined as
\begin{alignat}{1}
    p\bigl(s_0=(\text{S},\text{A})\bigr)=1
\end{alignat}
or
\begin{alignat}{1}
    p\bigl(s_0=(\text{S},\text{B})\bigr)=1,
\end{alignat}
i.e., at the beginning of a simulation the T-Maze type either assumes type A- or B. This is central to the information seeking mechanism of the GFE- and Context-GFE agents because they initially do not know the type and may seek information (about this) to find out.

Following this, the initial hidden states \(s \in S\) at timestep \(k=0\) of both non-context agents are defined as
\begin{alignat}{1}
    p\bigl(s_0=(\text{S},\text{A})\bigr)=p\bigl(s_0=(\text{S},\text{B})\bigr)=0.5.
\end{alignat}

The two context agents are equally uninformed but outsource the T-Maze type to their context states \(c \in \text{MazeType}\) with
\begin{alignat}{1}
    p\bigl(c=\text{A}\bigr)=p\bigl(c=\text{B}\bigr)=0.5.
\end{alignat}

Similarly to the non-context agents the context agents know that they are initially located at the start field via their hidden state \(s \in \text{Field}\) and
\begin{alignat}{1}
    p\bigl(s=\text{S}\bigr)=1.
\end{alignat}

The generative process and the non-context agents model the generation of observations \(o\) based on hidden states \(s \in S\) via
\begin{alignat}{3}
    & p\bigl(o=(\text{S},\text{CL})\bigm|s=(\text{S},\text{A})\bigr) &  & = p\bigl(o=(\text{S},\text{CL})\bigm|s=(\text{S},\text{B})\bigr) &  & = 0.5      \nonumber \\
    & p\bigl(o=(\text{S},\text{CR})\bigm|s=(\text{S},\text{A})\bigr) &  & = p\bigl(o=(\text{S},\text{CR})\bigm|s=(\text{S},\text{B})\bigr) &  & = 0.5      \nonumber \\
    & p\bigl(o=(\text{C},\text{CL})\bigm|s=(\text{C},\text{A})\bigr) &  & = p\bigl(o=(\text{C},\text{CR})\bigm|s=(\text{C},\text{B})\bigr) &  & = 1        \nonumber \\
    & p\bigl(o=(\text{L},\text{RW})\bigm|s=(\text{L},\text{A})\bigr) &  & = p\bigl(o=(\text{L},\text{NR})\bigm|s=(\text{L},\text{B})\bigr) &  & = \alpha   \nonumber \\
    & p\bigl(o=(\text{R},\text{RW})\bigm|s=(\text{R},\text{B})\bigr) &  & = p\bigl(o=(\text{R},\text{NR})\bigm|s=(\text{R},\text{A})\bigr) &  & = \alpha   \nonumber \\
    & p\bigl(o=(\text{L},\text{NR})\bigm|s=(\text{L},\text{B})\bigr) &  & = p\bigl(o=(\text{L},\text{RW})\bigm|s=(\text{L},\text{B})\bigr) &  & = 1-\alpha \nonumber \\
    & p\bigl(o=(\text{R},\text{NR})\bigm|s=(\text{R},\text{B})\bigr) &  & = p\bigl(o=(\text{R},\text{RW})\bigm|s=(\text{R},\text{A})\bigr) &  & = 1-\alpha,
\end{alignat}
where \(\alpha\) is a parameter that gets modulated in the simulations (explained below) and other probabilities \(p(o|s)\) are zero. 

In the case of the context agents the observations depend on both the hidden state \(s \in \text{Field}\) and the context state \(c \in \text{MazeType}\) via
\begin{alignat}{3}
    & p\bigl(o=(\text{S},\text{CL})\bigm|s=\text{S}, c=\text{A}\bigr) &  & = p\bigl(o=(\text{S},\text{CL})\bigm|s=\text{S}, c=\text{B}\bigr) &  & = 0.5      \nonumber \\
    & p\bigl(o=(\text{S},\text{CR})\bigm|s=\text{S}, c=\text{A}\bigr) &  & = p\bigl(o=(\text{S},\text{CR})\bigm|s=\text{S}, c=\text{B}\bigr) &  & = 0.5      \nonumber \\
    & p\bigl(o=(\text{C},\text{CL})\bigm|s=\text{C}, c=\text{A}\bigr) &  & = p\bigl(o=(\text{C},\text{CR})\bigm|s=\text{C}, c=\text{B}\bigr) &  & = 1        \nonumber \\
    & p\bigl(o=(\text{L},\text{RW})\bigm|s=\text{L}, c=\text{A}\bigr) &  & = p\bigl(o=(\text{L},\text{NR})\bigm|s=\text{L}, c=\text{B}\bigr) &  & = \alpha   \nonumber \\
    & p\bigl(o=(\text{R},\text{RW})\bigm|s=\text{R}, c=\text{B}\bigr) &  & = p\bigl(o=(\text{R},\text{NR})\bigm|s=\text{R}, c=\text{A}\bigr) &  & = \alpha   \nonumber \\
    & p\bigl(o=(\text{L},\text{NR})\bigm|s=\text{L}, c=\text{B}\bigr) &  & = p\bigl(o=(\text{L},\text{RW})\bigm|s=\text{L}, c=\text{B}\bigr) &  & = 1-\alpha \nonumber \\
    & p\bigl(o=(\text{R},\text{NR})\bigm|s=\text{R}, c=\text{B}\bigr) &  & = p\bigl(o=(\text{R},\text{RW})\bigm|s=\text{R}, c=\text{A}\bigr) &  & = 1-\alpha,
\end{alignat}
where other probabilities \(p(o|s,c)\) are zero. The parameter \(\alpha\) indicates how strongly the maze type is correlated with the reward location in the T-Maze. In the simulations the modulation of this parameter allows for a detailed analysis of the information seeking behavior of the GFE- and Context-GFE agents.

For both the generative process and the generative models of the non-context agents hidden states \(s_{k+1} \in S\) at time \(k+1\) depend on \(s_k \in S\) and actions \(a_k \in A\) at timestep \(k\) with probabilities
\begin{alignat}{3}
     & p\bigl(s_{k+1}=(\text{S},\text{A})\bigm|s_k=(\text{S},\text{A}),a_k=\text{GtS}\bigr) &  & =p\bigl(s_{k+1}=(\text{S},\text{B})\bigm|s_k=(\text{S},\text{B}),a_k=\text{GtS}\bigr) &  & = 1 \nonumber \\
     & p\bigl(s_{k+1}=(\text{C},\text{A})\bigm|s_k=(\text{S},\text{A}),a_k=\text{GtC}\bigr) &  & =p\bigl(s_{k+1}=(\text{C},\text{B})\bigm|s_k=(\text{S},\text{B}),a_k=\text{GtC}\bigr) &  & = 1 \nonumber \\
     & p\bigl(s_{k+1}=(\text{L},\text{A})\bigm|s_k=(\text{S},\text{A}),a_k=\text{GtL}\bigr) &  & =p\bigl(s_{k+1}=(\text{L},\text{B})\bigm|s_k=(\text{S},\text{B}),a_k=\text{GtL}\bigr) &  & = 1 \nonumber \\
     & p\bigl(s_{k+1}=(\text{R},\text{A})\bigm|s_k=(\text{S},\text{A}),a_k=\text{GtR}\bigr) &  & =p\bigl(s_{k+1}=(\text{R},\text{B})\bigm|s_k=(\text{S},\text{B}),a_k=\text{GtR}\bigr) &  & = 1 \nonumber \\[1.5em]
     & p\bigl(s_{k+1}=(\text{S},\text{A})\bigm|s_k=(\text{C},\text{A}),a_k=\text{GtS}\bigr) &  & =p\bigl(s_{k+1}=(\text{S},\text{B})\bigm|s_k=(\text{C},\text{B}),a_k=\text{GtS}\bigr) &  & = 1 \nonumber \\
     & p\bigl(s_{k+1}=(\text{C},\text{A})\bigm|s_k=(\text{C},\text{A}),a_k=\text{GtC}\bigr) &  & =p\bigl(s_{k+1}=(\text{C},\text{B})\bigm|s_k=(\text{C},\text{B}),a_k=\text{GtC}\bigr) &  & = 1 \nonumber \\
     & p\bigl(s_{k+1}=(\text{L},\text{A})\bigm|s_k=(\text{C},\text{A}),a_k=\text{GtL}\bigr) &  & =p\bigl(s_{k+1}=(\text{L},\text{B})\bigm|s_k=(\text{C},\text{B}),a_k=\text{GtL}\bigr) &  & = 1 \nonumber \\
     & p\bigl(s_{k+1}=(\text{R},\text{A})\bigm|s_k=(\text{C},\text{A}),a_k=\text{GtR}\bigr) &  & =p\bigl(s_{k+1}=(\text{R},\text{B})\bigm|s_k=(\text{C},\text{B}),a_k=\text{GtR}\bigr) &  & = 1 \nonumber \\[1.5em]
     & p\bigl(s_{k+1}=(\text{S},\text{A})\bigm|s_k=(\text{L},\text{A}),a_k=\text{GtS}\bigr) &  & =p\bigl(s_{k+1}=(\text{S},\text{B})\bigm|s_k=(\text{L},\text{B}),a_k=\text{GtS}\bigr) &  & = 1 \nonumber \\
     & p\bigl(s_{k+1}=(\text{S},\text{A})\bigm|s_k=(\text{L},\text{A}),a_k=\text{GtC}\bigr) &  & =p\bigl(s_{k+1}=(\text{S},\text{B})\bigm|s_k=(\text{L},\text{B}),a_k=\text{GtC}\bigr) &  & = 1 \nonumber \\
     & p\bigl(s_{k+1}=(\text{S},\text{A})\bigm|s_k=(\text{L},\text{A}),a_k=\text{GtL}\bigr) &  & =p\bigl(s_{k+1}=(\text{S},\text{B})\bigm|s_k=(\text{L},\text{B}),a_k=\text{GtL}\bigr) &  & = 1 \nonumber \\
     & p\bigl(s_{k+1}=(\text{S},\text{A})\bigm|s_k=(\text{L},\text{A}),a_k=\text{GtR}\bigr) &  & =p\bigl(s_{k+1}=(\text{S},\text{B})\bigm|s_k=(\text{L},\text{B}),a_k=\text{GtR}\bigr) &  & = 1 \nonumber \\[1.5em]
     & p\bigl(s_{k+1}=(\text{S},\text{A})\bigm|s_k=(\text{R},\text{A}),a_k=\text{GtS}\bigr) &  & =p\bigl(s_{k+1}=(\text{S},\text{B})\bigm|s_k=(\text{R},\text{B}),a_k=\text{GtS}\bigr) &  & = 1 \nonumber \\
     & p\bigl(s_{k+1}=(\text{S},\text{A})\bigm|s_k=(\text{R},\text{A}),a_k=\text{GtC}\bigr) &  & =p\bigl(s_{k+1}=(\text{S},\text{B})\bigm|s_k=(\text{R},\text{B}),a_k=\text{GtC}\bigr) &  & = 1 \nonumber \\
     & p\bigl(s_{k+1}=(\text{S},\text{A})\bigm|s_k=(\text{R},\text{A}),a_k=\text{GtL}\bigr) &  & =p\bigl(s_{k+1}=(\text{S},\text{B})\bigm|s_k=(\text{R},\text{B}),a_k=\text{GtL}\bigr) &  & = 1 \nonumber \\
     & p\bigl(s_{k+1}=(\text{S},\text{A})\bigm|s_k=(\text{R},\text{A}),a_k=\text{GtR}\bigr) &  & =p\bigl(s_{k+1}=(\text{S},\text{B})\bigm|s_k=(\text{R},\text{B}),a_k=\text{GtR}\bigr) &  & = 1
\end{alignat}
and other probabilities \(p(s_{k+1}|s_k,a_k)\) are zero.

Context agents have qualitatively equal transition probabilities for the hidden state \(s \in \text{Field}\) with
\begin{alignat}{3}
    & p\bigl(s_{k+1}=\text{S}\bigm|s_k=\text{S},a_k=\text{GtS}\bigr) &  & =p\bigl(s_{k+1}=\text{L}\bigm|s_k=\text{S},a_k=\text{GtL}\bigr) &  & = 1 \nonumber \\
    & p\bigl(s_{k+1}=\text{C}\bigm|s_k=\text{S},a_k=\text{GtC}\bigr) &  & =p\bigl(s_{k+1}=\text{R}\bigm|s_k=\text{S},a_k=\text{GtR}\bigr) &  & = 1 \nonumber \\
    & p\bigl(s_{k+1}=\text{S}\bigm|s_k=\text{C},a_k=\text{GtS}\bigr) &  & =p\bigl(s_{k+1}=\text{L}\bigm|s_k=\text{C},a_k=\text{GtL}\bigr) &  & = 1 \nonumber \\
    & p\bigl(s_{k+1}=\text{C}\bigm|s_k=\text{C},a_k=\text{GtC}\bigr) &  & =p\bigl(s_{k+1}=\text{R}\bigm|s_k=\text{C},a_k=\text{GtR}\bigr) &  & = 1 \nonumber \\
    & p\bigl(s_{k+1}=\text{S}\bigm|s_k=\text{L},a_k=\text{GtS}\bigr) &  & =p\bigl(s_{k+1}=\text{S}\bigm|s_k=\text{L},a_k=\text{GtL}\bigr) &  & = 1 \nonumber \\
    & p\bigl(s_{k+1}=\text{S}\bigm|s_k=\text{L},a_k=\text{GtC}\bigr) &  & =p\bigl(s_{k+1}=\text{S}\bigm|s_k=\text{L},a_k=\text{GtR}\bigr) &  & = 1 \nonumber \\
    & p\bigl(s_{k+1}=\text{S}\bigm|s_k=\text{R},a_k=\text{GtS}\bigr) &  & =p\bigl(s_{k+1}=\text{S}\bigm|s_k=\text{R},a_k=\text{GtL}\bigr) &  & = 1 \nonumber \\
    & p\bigl(s_{k+1}=\text{S}\bigm|s_k=\text{R},a_k=\text{GtC}\bigr) &  & =p\bigl(s_{k+1}=\text{S}\bigm|s_k=\text{R},a_k=\text{GtR}\bigr) &  & = 1
\end{alignat}
and other probabilities \(p(s_{k+1}|s_k,a_k)\) are zero.

The context state of context agents do not have a transition model like the hidden state \(s\), i.e., there is only a single context during a trial in the simulations.
All four agents have preferences for observations that are independent of the timestep. These preferences are the main drive in all agents that lead to purposeful behavior, i.e., acting to receive prefered observation. In the agents with information seeking these preferences compete with their curiousity. The preferences are defined via the parameter \(g\) in the goal distribution \(\tilde{p}(o)\):
\begin{alignat}{4}
    & \tilde{p}\bigl(o=(\text{L},\text{RW})\bigr) &  & =\tilde{p}\bigl(o=(\text{R},\text{RW})\bigr) &  & =\tilde{p}\bigl(o=(\text{S},\text{RW})\bigr) &  & =\tilde{p}\bigl(o=(\text{C},\text{RW})\bigr) \propto \exp(g) \nonumber \\
    & \tilde{p}\bigl(o=(\text{L},\text{NR})\bigr) &  & =\tilde{p}\bigl(o=(\text{R},\text{NR})\bigr) &  & =\tilde{p}\bigl(o=(\text{S},\text{NR})\bigr) &  & =\tilde{p}\bigl(o=(\text{C},\text{NR})\bigr) \propto \exp(-g)
\end{alignat}
After normalization, higher \(g\) values lead to stronger preferences for observing the reward (RW), weaker preferences for observing no reward (NR), and in between preferences for the remaining observations.

All agents also have a prior over actions that is also independent of the timestep. This prior is always defined as uniform distribution over all four actions:
\begin{alignat}{1}
    p(a=\text{GtL}) = p(a=\text{GtR}) = p(a=\text{GtS}) = p(a=\text{GtC}) = 0.25
\end{alignat}
In this way, the agents do not have biases for certain actions.


\subsection{Experiments}
\label{sec:simulations}

This section presents the schedule of the three experiments ("Initial Curiousity", "Exploiting Information", and "Long-Term Success"). Each simulation is run for each agent type ("BFE", "Context-BFE", "GFE", and "Context-GFE") and parameters \(\alpha \in [0.0, 0.1, \dots, 1.0]\) and \(c \in [0, 1, \dots, 10]\) of their generative models. The simulations reveal different aspects of the agents' behavior.
All simulations have in common that the agent begins in the start field of the T-Maze where it also initially receives an observation. This circumstance requires further explanation. First, the agent has a prior over its initial position in the T-Maze. This prior already entails the information that the agent begins at the start field, i.e., the agent already knows that it is located at this start field. Following this, the observation at this field does not yield any new information for the agent--the observation reveals to the agents its current location which the agent already "knows". Furthermore, this initial observation then has no effect on the agent, i.e., inferred actions are not dependent on this observation.

\subsubsection{Experiment 1: Initial Curiousity}

This experiment is designed to show the agents' initial drive to visit the cue location in the T-Maze. This is not to be understood as an empirical test, i.e., we do not have a cohort of the same agent type and average their behaviors. Instead, for each agent type the simulation is run once and their inclination of initially visiting the cue can directly be read of from their action specific belief distribution in the current timestep. Concretely, the respective agent first gets exposed to the environment (T-Maze) where it is initially located at the start position and the environment generates an observation that the agent perceives. Next, the agent performs inference, i.e., its message passing algorithm is executed until convergence. As a result, the agent has infered all of its belief distributions including the one that prescribes action in the current timestep. This action specific belief distribution contains probabilities over the four possible actions (go to (stay at) start, go to the left arm, go to the right arm, go to the cue). The agent's inclination of going to the cue is then exactly the respective probability of performing this action. This probability is the central result of this simulation. The procedere is exemplary detailed using the BFE agent in \cref{fig:24}.

\begin{figure}[htbp]
    \centering

    % ================= TIMESTEP 1 =================
    \begin{subfigure}[b]{\textwidth}
        \centering
        \fbox{%
            \begin{minipage}{0.97\linewidth}
                \vspace{0.15cm} % 1. Fester Abstand: OBEN
                % ==========================================
                % BEREICH 1: Das T-Maze (Flexibel zentriert)
                % ==========================================
                \hfill % Drückt das Maze von links weg
                \begin{tabular}{@{}c@{}}
                    \begin{tikzpicture}[scale=1.2]
                        \tikzset{
                            agent/.pic={
                                    \node[circle, draw=black, thick, fill=white, minimum size=0.5cm] at (0,0) {};
                                    \fill[black] (-0.08, 0.08) circle (0.03);
                                    \fill[black] (0.08, 0.08) circle (0.03);
                                    \draw[thick] (-0.1, -0.05) arc (180:360:0.1);
                                },
                            field label/.style={
                                    font=\tiny\sffamily, text=gray!70!black, anchor=south, scale=0.9, transform shape
                                },
                            field symbol/.style={
                                    font=\large\bfseries,
                                    text=gray!70,
                                    anchor=center,
                                    scale=1.0,
                                    transform shape
                                }
                        }

                        \draw (-0.5, 0) -- (-0.5, 2.0) -- (-1.5, 2.0) -- (-1.5, 3.0) -- (1.5, 3.0) -- (1.5, 2.0) -- (0.5, 2.0) -- (0.5, 0) -- cycle;
                        \draw (-0.5, 1.0) -- (0.5, 1.0);
                        \draw (-0.5, 2.0) -- (0.5, 2.0);
                        \draw (-0.5, 2.0) -- (-0.5, 3.0);
                        \draw (0.5, 2.0) -- (0.5, 3.0);

                        \pic[scale=0.8, transform shape] at (0, 1.5) {agent};
                        % Observation Arrow
                        \draw[
                            -{Stealth[scale=1.0]},
                            line width=1.6pt,
                            color=gray!80!black
                        ]
                        (0.51, 1.3) .. controls (1.0, 1.1) and (1.3, 2.0) .. (0.3, 1.6);
                        \node[font=\footnotesize, text=gray!80!black] at (1.2, 1.4) {$\textcolor{orange}{\hat {o_0}}$};
                        \node[field label] at (0, 0.97) {Start};
                        \node[field label] at (0, -0.03) {Cue};
                        \node[field label] at (-1.0, 1.97) {Left Arm};
                        \node[field label] at (1.0, 1.93) {Right Arm};
                        \node[field symbol] at (0, 0.5) {!};
                        \node[field symbol] at (-1.0, 2.5) {?};
                        \node[field symbol] at (1.0, 2.5) {?};
                    \end{tikzpicture}
                \end{tabular}%
                \hfill % Drückt das Maze von der gestrichelten Linie weg (zentriert es somit perfekt im linken Rest-Raum)
                % ==========================================
                % BEREICH 2: Die Trennlinie
                % ==========================================
                \begin{tabular}{@{}c@{}}
                    \tikz \draw[dashed, thin, gray] (0,-2.1) -- (0,2.1);
                \end{tabular}%
                \hspace{0.3cm}% 2. Fester Abstand: LINKS VOM FFG
                % ==========================================
                % BEREICH 3: Der FFG (gibt die Breite vor)
                % ==========================================
                \begin{tabular}{@{}c@{}}
                    % Der FFG wird jetzt in seiner absolut natürlichen Größe gerendert.
                    \resizebox{9cm}{!}{\begin{tikzpicture}[
    factor/.style={draw=black, thick, rectangle, minimum size=0.8cm, font=\normalsize},
    eqfactor/.style={draw=black, thick, rectangle, minimum size=0.5cm, font=\normalsize},
    edge/.style={thick},
    varlabel/.style={font=\normalsize},
    greycircle/.style={draw=black, thick, circle, fill=gray!30, minimum size=0.5cm, font=\normalsize}
]

% Knoten
\node[factor] (d) {$p(s_0)$};
\node[eqfactor, right=1.8cm of d] (eq0) {$=$};
\node[factor, below=1.5cm of eq0] (A0) {$p(o_0|s^{\scriptscriptstyle II}_0)$};
\node[factor, below=1.5cm of A0] (c0) {$\tilde{p}(o_0)$};
\node[factor, right=1.8cm of eq0] (B0) {$p(s_1|s^{\scriptscriptstyle I}_0,u_0)$};
\node[eqfactor, right=1.8cm of B0] (eq1) {$=$};
\node[factor, below=1.5cm of eq1] (A1) {$p(o_1|s^{\scriptscriptstyle II}_1)$};
\node[factor, below=1.5cm of A1] (c1) {$\tilde{p}(o_1)$};
\node[factor, right=1.8cm of eq1] (B1) {$p(s_2|s^{\scriptscriptstyle I}_1,u_1)$};
\node[factor] (A2) at ([xshift=2cm] B1.east |- A1) {$p(o_2|s_2)$};
\node[factor, below=1.5cm of A2] (c2) {$\tilde{p}(o_2)$};
\node[factor, above=1.5cm of B0] (u0) {$p(u_0)$};
\node[factor, above=1.5cm of B1] (u1) {$p(u_1)$};

% Kanten zwischen Knoten
\draw[edge] (d.east) -- (eq0.west) node[midway, above, varlabel] {$\textcolor{black}{q(s_0)}$};
\draw[edge] (eq0.south) -- (A0.north) node[midway, left, varlabel] {$\textcolor{black}{q(s^{\scriptscriptstyle II}_0)}$};
\draw[edge] (A0.south) -- (c0.north) node[midway, greycircle] {$\delta$} node[midway, left=0.3, varlabel] {$q(\textcolor{orange}{o_0})$};
\draw[edge] (eq0.east) -- (B0.west) node[midway, above, varlabel] {$\textcolor{black}{q(s^{\scriptscriptstyle I}_0)}$};
\draw[edge] (B0.north) -- (u0.south) node[midway, left, varlabel] {$\textcolor{black}{q(u_0)}$};
\draw[edge] (B0.east) -- (eq1.west) node[midway, above, varlabel] {$\textcolor{black}{q(s_1)}$};
\draw[edge] (eq1.south) -- (A1.north) node[midway, left, varlabel] {$\textcolor{black}{q(s^{\scriptscriptstyle II}_1)}$};
\draw[edge] (A1.south) -- (c1.north) node[midway, left, varlabel] {$\textcolor{black}{q(o_1)}$};
\draw[edge] (eq1.east) -- (B1.west) node[midway, above, varlabel] {$\textcolor{black}{q(s^{\scriptscriptstyle I}_1)}$};
\draw[edge] (B1.east) -| (A2.north) node[pos=0.25, above, varlabel] {$\textcolor{black}{q(s_2)}$};
\draw[edge] (A2.south) -- (c2.north) node[midway, left, varlabel] {$\textcolor{black}{q(o_2)}$};
\draw[edge] (B1.north) -- (u1.south) node[midway, left, varlabel] {$\textcolor{black}{q(u_1)}$};

\end{tikzpicture}}
                \end{tabular}%
                \llap{\smash{\raisebox{2.0cm}[0pt][0pt]{\sffamily\bfseries\small (a)\hspace*{-0.0cm}}}}%
                \hspace{0.1cm}% 3. Fester Abstand: RECHTS VOM FFG (Zum Box-Rand)
                \vspace{0.02cm} % 4. Fester Abstand: UNTEN
            \end{minipage}%
        }
        \phantomcaption
        \label{fig:24.1}
    \end{subfigure}

    \vspace{0.2cm} % Abstand zwischen den Kästen

    % ================= TIMESTEP 1 =================
    \begin{subfigure}[b]{\textwidth}
        \centering
        \fbox{%
            \begin{minipage}{0.97\linewidth}
                \vspace{0.15cm} % 1. Fester Abstand: OBEN
                % ==========================================
                % BEREICH 1: Das T-Maze (Flexibel zentriert)
                % ==========================================
                \hfill % Drückt das Maze von links weg
                \begin{tabular}{@{}c@{}}
                    \begin{tikzpicture}[scale=1.2]
                        \tikzset{
                            agent/.pic={
                                    \node[circle, draw=black, thick, fill=white, minimum size=0.5cm] at (0,0) {};
                                    \fill[black] (-0.08, 0.08) circle (0.03);
                                    \fill[black] (0.08, 0.08) circle (0.03);
                                    \draw[thick] (-0.1, -0.05) arc (180:360:0.1);
                                },
                            field label/.style={
                                    font=\tiny\sffamily, text=gray!70!black, anchor=south, scale=0.9, transform shape
                                },
                            field symbol/.style={
                                    font=\large\bfseries,
                                    text=gray!70,
                                    anchor=center,
                                    scale=1.0,
                                    transform shape
                                }
                        }

                        \draw (-0.5, 0) -- (-0.5, 2.0) -- (-1.5, 2.0) -- (-1.5, 3.0) -- (1.5, 3.0) -- (1.5, 2.0) -- (0.5, 2.0) -- (0.5, 0) -- cycle;
                        \draw (-0.5, 1.0) -- (0.5, 1.0);
                        \draw (-0.5, 2.0) -- (0.5, 2.0);
                        \draw (-0.5, 2.0) -- (-0.5, 3.0);
                        \draw (0.5, 2.0) -- (0.5, 3.0);

                        \pic[scale=0.8, transform shape] at (0, 1.5) {agent};
                        \node[field label] at (0, 0.97) {Start};
                        \node[field label] at (0, -0.03) {Cue};
                        \node[field label] at (-1.0, 1.97) {Left Arm};
                        \node[field label] at (1.0, 1.93) {Right Arm};
                        \node[field symbol] at (0, 0.5) {!};
                        \node[field symbol] at (-1.0, 2.5) {?};
                        \node[field symbol] at (1.0, 2.5) {?};

                        % --- NEU: Gedankenblase für Variational Free Energy Minimierung ---
                        % Kleine Bläschen, die vom Kopf (ca. 0.3, 1.7) nach rechts oben steigen
                        \draw[thin, gray!60, fill=white] (0.35, 1.55) circle (0.04);
                        \draw[thin, gray!60, fill=white] (0.52, 1.65) circle (0.07);

                        % Die Haupt-Gedankenblase mit dem Minimierungs-Symbol / Free Energy F
                        \node[
                            ellipse,
                            draw=gray!40,
                            thick,
                            fill=white,
                            minimum width=1.0cm,
                            minimum height=0.55cm,
                            inner sep=2pt,
                            font=\scriptsize\sffamily,
                            text=black!100!black
                        ] at (1.2, 1.7) {$\min F$};
                    \end{tikzpicture}
                \end{tabular}%
                \hfill % Drückt das Maze von der gestrichelten Linie weg (zentriert es somit perfekt im linken Rest-Raum)
                % ==========================================
                % BEREICH 2: Die Trennlinie
                % ==========================================
                \begin{tabular}{@{}c@{}}
                    \tikz \draw[dashed, thin, gray] (0,-2.1) -- (0,2.1);
                \end{tabular}%
                \hspace{0.3cm}% 2. Fester Abstand: LINKS VOM FFG
                % ==========================================
                % BEREICH 3: Der FFG (gibt die Breite vor)
                % ==========================================
                \begin{tabular}{@{}c@{}}
                    % Der FFG wird jetzt in seiner absolut natürlichen Größe gerendert.
                    \resizebox{9cm}{!}{\begin{tikzpicture}[
    factor/.style={draw=black, thick, rectangle, minimum size=0.8cm, font=\normalsize},
    eqfactor/.style={draw=black, thick, rectangle, minimum size=0.5cm, font=\normalsize},
    edge/.style={thick},
    varlabel/.style={font=\normalsize},
    greycircle/.style={draw=black, thick, circle, fill=gray!30, minimum size=0.5cm, font=\normalsize}
]

% Knoten
\node[factor] (d) {$p(s_0)$};
\node[eqfactor, right=1.8cm of d] (eq0) {$=$};
\node[factor, below=1.5cm of eq0] (A0) {$p(o_0|s^{\scriptscriptstyle II}_0)$};
\node[factor, below=1.5cm of A0] (c0) {$\tilde{p}(o_0)$};
\node[factor, right=1.8cm of eq0] (B0) {$p(s_1|s^{\scriptscriptstyle I}_0,u_0)$};
\node[eqfactor, right=1.8cm of B0] (eq1) {$=$};
\node[factor, below=1.5cm of eq1] (A1) {$p(o_1|s^{\scriptscriptstyle II}_1)$};
\node[factor, below=1.5cm of A1] (c1) {$\tilde{p}(o_1)$};
\node[factor, right=1.8cm of eq1] (B1) {$p(s_2|s^{\scriptscriptstyle I}_1,u_1)$};
\node[factor] (A2) at ([xshift=2cm] B1.east |- A1) {$p(o_2|s_2)$};
\node[factor, below=1.5cm of A2] (c2) {$\tilde{p}(o_2)$};
\node[factor, above=1.5cm of B0] (u0) {$p(u_0)$};
\node[factor, above=1.5cm of B1] (u1) {$p(u_1)$};

% Kanten zwischen Knoten
\draw[edge] (d.east) -- (eq0.west) node[midway, above, varlabel] {$\textcolor{black}{q(s_0)}$} node[near end, below, varlabel, scale=0.7] {$\textcolor{orange}{\bm{\overrightarrow{\mu}(}s_0\bm{)}}$} node[near start, below, varlabel, scale=0.7] {$\textcolor{orange}{\bm{\overleftarrow{\mu}(}s_0\bm{)}}$};
\draw[edge] (eq0.south) -- (A0.north) node[midway, left, varlabel] {$\textcolor{black}{q(s^{\scriptscriptstyle II}_0)}$} node[pos=0.25, right=0.05, varlabel, scale=0.7] {$\textcolor{orange}{\bm{{\uparrow}{\mu}(}s^{\scriptscriptstyle II}_0\bm{)}}$} node[pos=0.75, right=0.05, varlabel, scale=0.7] {$\textcolor{orange}{\bm{{\downarrow}{\mu}(}s^{\scriptscriptstyle II}_0\bm{)}}$};
\draw[edge] (A0.south) -- (c0.north) node[midway, greycircle] {$\delta$} node[midway, left=0.3, varlabel] {$\textcolor{black}{q(o_0)}$} node[pos=0.2, right=0.05, varlabel, scale=0.7] {$\textcolor{orange}{\bm{{\uparrow}{\mu}(}o_0\bm{)}}$} node[pos=0.8, right=0.05, varlabel, scale=0.7] {$\textcolor{orange}{\bm{{\downarrow}{\mu}(}o_0\bm{)}}$};
\draw[edge] (eq0.east) -- (B0.west) node[midway, above, varlabel] {$\textcolor{black}{q(s^{\scriptscriptstyle I}_0)}$} node[near start, below, varlabel, scale=0.7] {$\textcolor{orange}{\bm{\overleftarrow{\mu}(}s^{\scriptscriptstyle I}_0\bm{)}}$} node[near end, below, varlabel, scale=0.7] {$\textcolor{orange}{\bm{\overrightarrow{\mu}(}s^{\scriptscriptstyle I}_0\bm{)}}$};
\draw[edge] (B0.north) -- (u0.south) node[midway, left, varlabel] {$\textcolor{black}{q(u_0)}$} node[pos=0.2, right=0.05, varlabel, scale=0.7] {$\textcolor{orange}{\bm{{\downarrow}{\mu}(}u_0\bm{)}}$} node[pos=0.8, right=0.05, varlabel, scale=0.7] {$\textcolor{orange}{\bm{{\uparrow}{\mu}(}u_0\bm{)}}$};
\draw[edge] (B0.east) -- (eq1.west) node[midway, above, varlabel] {$\textcolor{black}{q(s_1)}$} node[near end, below, varlabel, scale=0.7] {$\textcolor{orange}{\bm{\overrightarrow{\mu}(}s_1\bm{)}}$} node[near start, below, varlabel, scale=0.7] {$\textcolor{orange}{\bm{\overleftarrow{\mu}(}s_1\bm{)}}$};
\draw[edge] (eq1.south) -- (A1.north) node[midway, left, varlabel] {$\textcolor{black}{q(s^{\scriptscriptstyle II}_1)}$} node[pos=0.25, right=0.05, varlabel, scale=0.7] {$\textcolor{orange}{\bm{{\uparrow}{\mu}(}s^{\scriptscriptstyle II}_1\bm{)}}$} node[pos=0.75, right=0.05, varlabel, scale=0.7] {$\textcolor{orange}{\bm{{\downarrow}{\mu}(}s^{\scriptscriptstyle II}_1\bm{)}}$};
\draw[edge] (A1.south) -- (c1.north) node[midway, left, varlabel] {$\textcolor{black}{q(o_1)}$} node[pos=0.2, right=0.05, varlabel, scale=0.7] {$\textcolor{orange}{\bm{{\uparrow}{\mu}(}o_1\bm{)}}$} node[pos=0.8, right=0.05, varlabel, scale=0.7] {$\textcolor{orange}{\bm{{\downarrow}{\mu}(}o_1\bm{)}}$};
\draw[edge] (eq1.east) -- (B1.west) node[midway, above, varlabel] {$\textcolor{black}{q(s^{\scriptscriptstyle I}_1)}$} node[near start, below, varlabel, scale=0.7] {$\textcolor{orange}{\bm{\overleftarrow{\mu}(}s^{\scriptscriptstyle I}_1\bm{)}}$} node[near end, below, varlabel, scale=0.7] {$\textcolor{orange}{\bm{\overrightarrow{\mu}(}s^{\scriptscriptstyle I}_1\bm{)}}$};
\draw[edge] (B1.east) -| (A2.north) node[pos=0.25, above, varlabel] {$\textcolor{black}{q(s_2)}$} node[pos=0.89, right=0.05, varlabel, scale=0.7] {$\textcolor{orange}{\bm{{\downarrow}{\mu}(}s_2\bm{)}}$} node[pos=0.11, below, varlabel, scale=0.7] {$\textcolor{orange}{\bm{\overleftarrow{\mu}(}s_2\bm{)}}$};
\draw[edge] (A2.south) -- (c2.north) node[midway, left, varlabel] {$\textcolor{black}{q(o_2)}$} node[pos=0.2, right=0.05, varlabel, scale=0.7] {$\textcolor{orange}{\bm{{\uparrow}{\mu}(}o_2\bm{)}}$} node[pos=0.8, right=0.05, varlabel, scale=0.7] {$\textcolor{orange}{\bm{{\downarrow}{\mu}(}o_2\bm{)}}$};
\draw[edge] (B1.north) -- (u1.south) node[midway, left, varlabel] {$\textcolor{black}{q(u_1)}$} node[pos=0.2, right=0.05, varlabel, scale=0.7] {$\textcolor{orange}{\bm{{\downarrow}{\mu}(}u_1\bm{)}}$} node[pos=0.8, right=0.05, varlabel, scale=0.7] {$\textcolor{orange}{\bm{{\uparrow}{\mu}(}u_1\bm{)}}$};

\end{tikzpicture}}
                \end{tabular}%
                \llap{\smash{\raisebox{2.0cm}[0pt][0pt]{\sffamily\bfseries\small (b)\hspace*{-0.0cm}}}}%
                \hspace{0.1cm}% 3. Fester Abstand: RECHTS VOM FFG (Zum Box-Rand)
                \vspace{0.02cm} % 4. Fester Abstand: UNTEN
            \end{minipage}%
        }
        \phantomcaption
        \label{fig:24.2}
    \end{subfigure}

    \caption{The figure presents the two crucial steps of the "Initial Curiousity" simulation using the example of the BFE agent: First, the agent receives an observation \(\hat {o_0}\) from the environment (\cref{fig:24.1}). The observation has an effect on the agent's temporal generative model, i.e., in its factor graph representation, the belief distribution \(q(o_0)\) that describes its belief about observations of timestep \(t=0\), becomes a delta distribution with all probability mass at \(\hat {o_0}\). Second, the agent performs inference, i.e., it runs its message passing algorithm (messages sent along the edges in \cref{fig:24.2}). In the case of the BFE agent this refers to \cref{alg:2}. After message passing has finished, the beliefs of executing a certain action in the current timestep are then directly given by the belief \(q(u_0)\).}
    \label{fig:24}
\end{figure}


\subsubsection{Experiment 2: Exploiting Information}

In this experiment, the agent receives the cue which reveals the location of the reward (left arm or right arm) in the T-Maze. The result of this experiment is then the agent's inclination of going to the respective arm in the T-Maze that the cue indicates as given by the probabilities in its action-specific belief distribution. The simulation spans over timesteps \(t=0\) and \(t=1\). First, the agent receives an observation \(o_0\) generated be the environment (at \(t=0\)). Second, it executes the "go to cue" action and thus moves to the cue field in the environment, i.e., the action is not sampled from the action belief distribution \(q(u_0)\) but is externally triggered (at \(t=0\)). Third, it receives the second observation \(o_1\) while being at the cue field (at \(t=1\)). Fourth, the agent runs its messages passing algorithm to infer its belief distributions. After inference, the result of the experiment is given as the agent's belief (a probability) of executing the specific action that moves it to the arm that is indicated by the cue. A detailed explanation of the procedere is given in \cref{fig:25}.

\begin{figure}[bp!]
    \centering
    \resizebox{1.0\textwidth}{!}{%
        \begin{tikzpicture}[
                panel/.style={
                    draw=gray!80, thick, fill=blue!5,
                    rounded corners=4pt,
                    minimum width=8.0cm,
                    minimum height=4.5cm,
                    inner sep=0pt,
                    align=center
                },
                sublabel/.style={
                    anchor=north east,
                    font=\sffamily\bfseries\Large,
                    text=black
                },
                divider/.style={
                    draw=gray!70,
                    dashed,
                    line width=1.0pt
                }
            ]

            % --- 2x2 Grid (Feste Positionierung) ---
            % Obere Reihe: (a) und (b)
            \node[panel] (boxA) at (0,0) {
                % \usebox{\boxA}
            };

            \node[panel, right=0.3cm of boxA] (boxB) {
                % \usebox{\boxB}
            };

            % Untere Reihe: (c) und (d)
            \node[panel, below=0.3cm of boxA] (boxC) {
                % \usebox{\boxC}
            };

            \node[panel, right=0.3cm of boxC] (boxD) {
                % \usebox{\boxD}
            };

            \foreach \b in {boxA, boxB, boxC, boxD} {
                % yshift sorgt für einen kleinen Abstand zum oberen/unteren Rand
                \draw[divider] ([yshift=-8pt]\b.north) -- ([yshift=8pt]\b.south);
            }

            % --- Labels (a), (b), (c), (d) fix oben rechts verankert ---
            %\node[sublabel] at ([xshift=-10pt, yshift=-8pt]boxA.north east) {(a)};
            %\node[sublabel] at ([xshift=-10pt, yshift=-8pt]boxB.north east) {(b)};
            %\node[sublabel] at ([xshift=-10pt, yshift=-8pt]boxC.north east) {(c)};
            %\node[sublabel] at ([xshift=-10pt, yshift=-8pt]boxD.north east) {(d)};

        \end{tikzpicture}
    }
    \caption{Overview of the four experimental setups: (a) initial curiosity, (b) belief state, (c) free energy minimization during planning, and (d) action execution.}
    \label{fig:four_panels}
\end{figure}

\begin{figure}[b!]
    \centering

    % ================= TIMESTEP 1 =================
    \begin{subfigure}[b]{\textwidth}
        \centering
        \fbox{%
            \begin{minipage}{0.97\linewidth}
                \vspace{0.15cm} % 1. Fester Abstand: OBEN
                % ==========================================
                % BEREICH 1: Das T-Maze (Flexibel zentriert)
                % ==========================================
                \hfill % Drückt das Maze von links weg
                \begin{tabular}{@{}c@{}}
                    \begin{tikzpicture}[scale=1.2]
                        \tikzset{
                            agent/.pic={
                                    \node[circle, draw=black, thick, fill=white, minimum size=0.5cm] at (0,0) {};
                                    \fill[black] (-0.08, 0.08) circle (0.03);
                                    \fill[black] (0.08, 0.08) circle (0.03);
                                    \draw[thick] (-0.1, -0.05) arc (180:360:0.1);
                                },
                            field label/.style={
                                    font=\tiny\sffamily, text=gray!70!black, anchor=south, scale=0.9, transform shape
                                },
                            field symbol/.style={
                                    font=\large\bfseries,
                                    text=gray!70,
                                    anchor=center,
                                    scale=1.0,
                                    transform shape
                                }
                        }

                        \draw (-0.5, 0) -- (-0.5, 2.0) -- (-1.5, 2.0) -- (-1.5, 3.0) -- (1.5, 3.0) -- (1.5, 2.0) -- (0.5, 2.0) -- (0.5, 0) -- cycle;
                        \draw (-0.5, 1.0) -- (0.5, 1.0);
                        \draw (-0.5, 2.0) -- (0.5, 2.0);
                        \draw (-0.5, 2.0) -- (-0.5, 3.0);
                        \draw (0.5, 2.0) -- (0.5, 3.0);

                        \pic[scale=0.8, transform shape] at (0, 1.5) {agent};
                        % Observation Arrow
                        \draw[
                            -{Stealth[scale=1.0]},
                            line width=1.6pt,
                            color=gray!80!black
                        ]
                        (0.51, 1.3) .. controls (1.0, 1.1) and (1.3, 2.0) .. (0.3, 1.6);
                        \node[font=\footnotesize, text=gray!80!black] at (1.2, 1.4) {$\textcolor{orange}{\hat {o_0}}$};
                        \node[field label] at (0, 0.97) {Start};
                        \node[field label] at (0, -0.03) {Cue};
                        \node[field label] at (-1.0, 1.97) {Left Arm};
                        \node[field label] at (1.0, 1.93) {Right Arm};
                        \node[field symbol] at (0, 0.5) {!};
                        \node[field symbol] at (-1.0, 2.5) {?};
                        \node[field symbol] at (1.0, 2.5) {?};
                    \end{tikzpicture}
                \end{tabular}%
                \hfill % Drückt das Maze von der gestrichelten Linie weg (zentriert es somit perfekt im linken Rest-Raum)
                % ==========================================
                % BEREICH 2: Die Trennlinie
                % ==========================================
                \begin{tabular}{@{}c@{}}
                    \tikz \draw[dashed, thin, gray] (0,-2.1) -- (0,2.1);
                \end{tabular}%
                \hspace{0.3cm}% 2. Fester Abstand: LINKS VOM FFG
                % ==========================================
                % BEREICH 3: Der FFG (gibt die Breite vor)
                % ==========================================
                \begin{tabular}{@{}c@{}}
                    % Der FFG wird jetzt in seiner absolut natürlichen Größe gerendert.
                    \resizebox{9cm}{!}{\begin{tikzpicture}[
    factor/.style={draw=black, thick, rectangle, minimum size=0.8cm, font=\normalsize},
    eqfactor/.style={draw=black, thick, rectangle, minimum size=0.5cm, font=\normalsize},
    edge/.style={thick},
    varlabel/.style={font=\normalsize},
    greycircle/.style={draw=black, thick, circle, fill=gray!30, minimum size=0.5cm, font=\normalsize}
]

% Knoten
\node[factor] (d) {$p(s_0)$};
\node[eqfactor, right=1.8cm of d] (eq0) {$=$};
\node[factor, below=1.5cm of eq0] (A0) {$p(o_0|s^{\scriptscriptstyle II}_0)$};
\node[factor, below=1.5cm of A0] (c0) {$\tilde{p}(o_0)$};
\node[factor, right=1.8cm of eq0] (B0) {$p(s_1|s^{\scriptscriptstyle I}_0,u_0)$};
\node[eqfactor, right=1.8cm of B0] (eq1) {$=$};
\node[factor, below=1.5cm of eq1] (A1) {$p(o_1|s^{\scriptscriptstyle II}_1)$};
\node[factor, below=1.5cm of A1] (c1) {$\tilde{p}(o_1)$};
\node[factor, right=1.8cm of eq1] (B1) {$p(s_2|s^{\scriptscriptstyle I}_1,u_1)$};
\node[factor] (A2) at ([xshift=2cm] B1.east |- A1) {$p(o_2|s_2)$};
\node[factor, below=1.5cm of A2] (c2) {$\tilde{p}(o_2)$};
\node[factor, above=1.5cm of B0] (u0) {$p(u_0)$};
\node[factor, above=1.5cm of B1] (u1) {$p(u_1)$};

% Kanten zwischen Knoten
\draw[edge] (d.east) -- (eq0.west) node[midway, above, varlabel] {$\textcolor{black}{q(s_0)}$};
\draw[edge] (eq0.south) -- (A0.north) node[midway, left, varlabel] {$\textcolor{black}{q(s^{\scriptscriptstyle II}_0)}$};
\draw[edge] (A0.south) -- (c0.north) node[midway, greycircle] {$\delta$} node[midway, left=0.3, varlabel] {$q(\textcolor{orange}{o_0})$};
\draw[edge] (eq0.east) -- (B0.west) node[midway, above, varlabel] {$\textcolor{black}{q(s^{\scriptscriptstyle I}_0)}$};
\draw[edge] (B0.north) -- (u0.south) node[midway, left, varlabel] {$\textcolor{black}{q(u_0)}$};
\draw[edge] (B0.east) -- (eq1.west) node[midway, above, varlabel] {$\textcolor{black}{q(s_1)}$};
\draw[edge] (eq1.south) -- (A1.north) node[midway, left, varlabel] {$\textcolor{black}{q(s^{\scriptscriptstyle II}_1)}$};
\draw[edge] (A1.south) -- (c1.north) node[midway, left, varlabel] {$\textcolor{black}{q(o_1)}$};
\draw[edge] (eq1.east) -- (B1.west) node[midway, above, varlabel] {$\textcolor{black}{q(s^{\scriptscriptstyle I}_1)}$};
\draw[edge] (B1.east) -| (A2.north) node[pos=0.25, above, varlabel] {$\textcolor{black}{q(s_2)}$};
\draw[edge] (A2.south) -- (c2.north) node[midway, left, varlabel] {$\textcolor{black}{q(o_2)}$};
\draw[edge] (B1.north) -- (u1.south) node[midway, left, varlabel] {$\textcolor{black}{q(u_1)}$};

\end{tikzpicture}}
                \end{tabular}%
                \llap{\smash{\raisebox{2.0cm}[0pt][0pt]{\sffamily\bfseries\small (a)\hspace*{-0.0cm}}}}%
                \hspace{0.1cm}% 3. Fester Abstand: RECHTS VOM FFG (Zum Box-Rand)
                \vspace{0.02cm} % 4. Fester Abstand: UNTEN
            \end{minipage}%
        }
        \phantomcaption
        \label{fig:25.1}
    \end{subfigure}

    \vspace{0.2cm} % Abstand zwischen den Kästen

    % ================= TIMESTEP 1 =================
    \begin{subfigure}[b]{\textwidth}
        \centering
        \fbox{%
            \begin{minipage}{0.97\linewidth}
                \vspace{0.15cm} % 1. Fester Abstand: OBEN
                % ==========================================
                % BEREICH 1: Das T-Maze (Flexibel zentriert)
                % ==========================================
                \hfill % Drückt das Maze von links weg
                \begin{tabular}{@{}c@{}}
                    \begin{tikzpicture}[scale=1.2]
                        \tikzset{
                            agent/.pic={
                                    \node[circle, draw=black, thick, fill=white, minimum size=0.5cm] at (0,0) {};
                                    \fill[black] (-0.08, 0.08) circle (0.03);
                                    \fill[black] (0.08, 0.08) circle (0.03);
                                    \draw[thick] (-0.1, -0.05) arc (180:360:0.1);
                                },
                            field label/.style={
                                    font=\tiny\sffamily, text=gray!70!black, anchor=south, scale=0.9, transform shape
                                },
                            field symbol/.style={
                                    font=\large\bfseries,
                                    text=gray!70,
                                    anchor=center,
                                    scale=1.0,
                                    transform shape
                                }
                        }

                        \draw (-0.5, 0) -- (-0.5, 2.0) -- (-1.5, 2.0) -- (-1.5, 3.0) -- (1.5, 3.0) -- (1.5, 2.0) -- (0.5, 2.0) -- (0.5, 0) -- cycle;
                        \draw (-0.5, 1.0) -- (0.5, 1.0);
                        \draw (-0.5, 2.0) -- (0.5, 2.0);
                        \draw (-0.5, 2.0) -- (-0.5, 3.0);
                        \draw (0.5, 2.0) -- (0.5, 3.0);

                        \pic[scale=0.8, transform shape] at (0, 1.5) {agent};
                        % Observation Arrow
                        \draw[
                            {Stealth[scale=1.0]}-,
                            line width=1.6pt,
                            color=gray!80!black
                        ]
                        (0.51, 1.3) .. controls (1.0, 1.1) and (1.3, 2.0) .. (0.3, 1.6);
                        \node[font=\footnotesize, text=gray!80!black] at (1.2, 1.4) {$\textcolor{orange}{\hat {a_0}}$};
                        \node[field label] at (0, 0.97) {Start};
                        \node[field label] at (0, -0.03) {Cue};
                        \node[field label] at (-1.0, 1.97) {Left Arm};
                        \node[field label] at (1.0, 1.93) {Right Arm};
                        \node[field symbol] at (0, 0.5) {!};
                        \node[field symbol] at (-1.0, 2.5) {?};
                        \node[field symbol] at (1.0, 2.5) {?};
                    \end{tikzpicture}
                \end{tabular}%
                \hfill % Drückt das Maze von der gestrichelten Linie weg (zentriert es somit perfekt im linken Rest-Raum)
                % ==========================================
                % BEREICH 2: Die Trennlinie
                % ==========================================
                \begin{tabular}{@{}c@{}}
                    \tikz \draw[dashed, thin, gray] (0,-2.1) -- (0,2.1);
                \end{tabular}%
                \hspace{0.3cm}% 2. Fester Abstand: LINKS VOM FFG
                % ==========================================
                % BEREICH 3: Der FFG (gibt die Breite vor)
                % ==========================================
                \begin{tabular}{@{}c@{}}
                    % Der FFG wird jetzt in seiner absolut natürlichen Größe gerendert.
                    \resizebox{9cm}{!}{\begin{tikzpicture}[
    factor/.style={draw=black, thick, rectangle, minimum size=0.8cm, font=\normalsize},
    eqfactor/.style={draw=black, thick, rectangle, minimum size=0.5cm, font=\normalsize},
    edge/.style={thick},
    varlabel/.style={font=\normalsize},
    greycircle/.style={draw=black, thick, circle, fill=gray!30, minimum size=0.5cm, font=\normalsize}
]

% Knoten
\node[factor] (d) {$p(s_0)$};
\node[eqfactor, right=1.8cm of d] (eq0) {$=$};
\node[factor, below=1.5cm of eq0] (A0) {$p(o_0|s^{\scriptscriptstyle II}_0)$};
\node[factor, below=1.5cm of A0] (c0) {$\tilde{p}(o_0)$};
\node[factor, right=1.8cm of eq0] (B0) {$p(s_1|s^{\scriptscriptstyle I}_0,u_0)$};
\node[eqfactor, right=1.8cm of B0] (eq1) {$=$};
\node[factor, below=1.5cm of eq1] (A1) {$p(o_1|s^{\scriptscriptstyle II}_1)$};
\node[factor, below=1.5cm of A1] (c1) {$\tilde{p}(o_1)$};
\node[factor, right=1.8cm of eq1] (B1) {$p(s_2|s^{\scriptscriptstyle I}_1,u_1)$};
\node[factor] (A2) at ([xshift=2cm] B1.east |- A1) {$p(o_2|s_2)$};
\node[factor, below=1.5cm of A2] (c2) {$\tilde{p}(o_2)$};
\node[factor, above=1.5cm of B0] (u0) {$p(u_0)$};
\node[factor, above=1.5cm of B1] (u1) {$p(u_1)$};

% Kanten zwischen Knoten
\draw[edge] (d.east) -- (eq0.west) node[midway, above, varlabel] {$\textcolor{black}{q(s_0)}$};
\draw[edge] (eq0.south) -- (A0.north) node[midway, left, varlabel] {$\textcolor{black}{q(s^{\scriptscriptstyle II}_0)}$};
\draw[edge] (A0.south) -- (c0.north) node[midway, greycircle] {$\delta$} node[midway, left=0.3, varlabel] {$q(\textcolor{black}{o_0})$};
\draw[edge] (eq0.east) -- (B0.west) node[midway, above, varlabel] {$\textcolor{black}{q(s^{\scriptscriptstyle I}_0)}$};
\draw[edge] (B0.north) -- (u0.south) node[midway, greycircle] {$\delta$} node[midway, left=0.3, varlabel] {$q(\textcolor{orange}{u_0})$};
\draw[edge] (B0.east) -- (eq1.west) node[midway, above, varlabel] {$\textcolor{black}{q(s_1)}$};
\draw[edge] (eq1.south) -- (A1.north) node[midway, left, varlabel] {$\textcolor{black}{q(s^{\scriptscriptstyle II}_1)}$};
\draw[edge] (A1.south) -- (c1.north) node[midway, left, varlabel] {$\textcolor{black}{q(o_1)}$};
\draw[edge] (eq1.east) -- (B1.west) node[midway, above, varlabel] {$\textcolor{black}{q(s^{\scriptscriptstyle I}_1)}$};
\draw[edge] (B1.east) -| (A2.north) node[pos=0.25, above, varlabel] {$\textcolor{black}{q(s_2)}$};
\draw[edge] (A2.south) -- (c2.north) node[midway, left, varlabel] {$\textcolor{black}{q(o_2)}$};
\draw[edge] (B1.north) -- (u1.south) node[midway, left, varlabel] {$\textcolor{black}{q(u_1)}$};

\end{tikzpicture}}
                \end{tabular}%
                \llap{\smash{\raisebox{2.0cm}[0pt][0pt]{\sffamily\bfseries\small (b)\hspace*{-0.0cm}}}}%
                \hspace{0.1cm}% 3. Fester Abstand: RECHTS VOM FFG (Zum Box-Rand)
                \vspace{0.02cm} % 4. Fester Abstand: UNTEN
            \end{minipage}%
        }
        \phantomcaption
        \label{fig:25.2}
    \end{subfigure}

    \vspace{0.2cm} % Abstand zwischen den Kästen

    % ================= TIMESTEP 1 =================
    \begin{subfigure}[b]{\textwidth}
        \centering
        \fbox{%
            \begin{minipage}{0.97\linewidth}
                \vspace{0.15cm} % 1. Fester Abstand: OBEN
                % ==========================================
                % BEREICH 1: Das T-Maze (Flexibel zentriert)
                % ==========================================
                \hfill % Drückt das Maze von links weg
                \begin{tabular}{@{}c@{}}
                    \begin{tikzpicture}[scale=1.2]
                        \tikzset{
                            agent/.pic={
                                    \node[circle, draw=black, thick, fill=white, minimum size=0.5cm] at (0,0) {};
                                    \fill[black] (-0.08, 0.08) circle (0.03);
                                    \fill[black] (0.08, 0.08) circle (0.03);
                                    \draw[thick] (-0.1, -0.05) arc (180:360:0.1);
                                },
                            field label/.style={
                                    font=\tiny\sffamily, text=gray!70!black, anchor=south, scale=0.9, transform shape
                                },
                            field symbol/.style={
                                    font=\large\bfseries,
                                    text=gray!70,
                                    anchor=center,
                                    scale=1.0,
                                    transform shape
                                }
                        }

                        \draw (-0.5, 0) -- (-0.5, 2.0) -- (-1.5, 2.0) -- (-1.5, 3.0) -- (1.5, 3.0) -- (1.5, 2.0) -- (0.5, 2.0) -- (0.5, 0) -- cycle;
                        \draw (-0.5, 1.0) -- (0.5, 1.0);
                        \draw (-0.5, 2.0) -- (0.5, 2.0);
                        \draw (-0.5, 2.0) -- (-0.5, 3.0);
                        \draw (0.5, 2.0) -- (0.5, 3.0);

                        \pic[scale=0.8, transform shape] at (0, 0.5) {agent};
                        % Observation Arrow
                        \draw[
                            -{Stealth[scale=1.0]},
                            line width=1.6pt,
                            color=gray!80!black
                        ]
                        (0.51, 0.3) .. controls (1.0, 0.1) and (1.3, 1.0) .. (0.3, 0.6);
                        \node[font=\footnotesize, text=gray!80!black] at (1.2, 0.4) {$\textcolor{orange}{\hat {o_1}}$};
                        \node[field label] at (0, 0.97) {Start};
                        \node[field label] at (0, -0.03) {Cue};
                        \node[field label] at (-1.0, 1.97) {Left Arm};
                        \node[field label] at (1.0, 1.93) {Right Arm};
                        %\node[field symbol] at (0, 0.5) {!};
                        \node[field symbol] at (-1.0, 2.5) {?};
                        \node[field symbol] at (1.0, 2.5) {?};
                    \end{tikzpicture}
                \end{tabular}%
                \hfill % Drückt das Maze von der gestrichelten Linie weg (zentriert es somit perfekt im linken Rest-Raum)
                % ==========================================
                % BEREICH 2: Die Trennlinie
                % ==========================================
                \begin{tabular}{@{}c@{}}
                    \tikz \draw[dashed, thin, gray] (0,-2.1) -- (0,2.1);
                \end{tabular}%
                \hspace{0.3cm}% 2. Fester Abstand: LINKS VOM FFG
                % ==========================================
                % BEREICH 3: Der FFG (gibt die Breite vor)
                % ==========================================
                \begin{tabular}{@{}c@{}}
                    % Der FFG wird jetzt in seiner absolut natürlichen Größe gerendert.
                    \resizebox{9cm}{!}{\begin{tikzpicture}[
    factor/.style={draw=black, thick, rectangle, minimum size=0.8cm, font=\normalsize},
    eqfactor/.style={draw=black, thick, rectangle, minimum size=0.5cm, font=\normalsize},
    edge/.style={thick},
    varlabel/.style={font=\normalsize},
    greycircle/.style={draw=black, thick, circle, fill=gray!30, minimum size=0.5cm, font=\normalsize}
]

% Knoten
\node[factor] (d) {$p(s_0)$};
\node[eqfactor, right=1.8cm of d] (eq0) {$=$};
\node[factor, below=1.5cm of eq0] (A0) {$p(o_0|s^{\scriptscriptstyle II}_0)$};
\node[factor, below=1.5cm of A0] (c0) {$\tilde{p}(o_0)$};
\node[factor, right=1.8cm of eq0] (B0) {$p(s_1|s^{\scriptscriptstyle I}_0,u_0)$};
\node[eqfactor, right=1.8cm of B0] (eq1) {$=$};
\node[factor, below=1.5cm of eq1] (A1) {$p(o_1|s^{\scriptscriptstyle II}_1)$};
\node[factor, below=1.5cm of A1] (c1) {$\tilde{p}(o_1)$};
\node[factor, right=1.8cm of eq1] (B1) {$p(s_2|s^{\scriptscriptstyle I}_1,u_1)$};
\node[factor] (A2) at ([xshift=2cm] B1.east |- A1) {$p(o_2|s_2)$};
\node[factor, below=1.5cm of A2] (c2) {$\tilde{p}(o_2)$};
\node[factor, above=1.5cm of B0] (u0) {$p(u_0)$};
\node[factor, above=1.5cm of B1] (u1) {$p(u_1)$};

% Kanten zwischen Knoten
\draw[edge] (d.east) -- (eq0.west) node[midway, above, varlabel] {$\textcolor{black}{q(s_0)}$};
\draw[edge] (eq0.south) -- (A0.north) node[midway, left, varlabel] {$\textcolor{black}{q(s^{\scriptscriptstyle II}_0)}$};
\draw[edge] (A0.south) -- (c0.north) node[midway, greycircle] {$\delta$} node[midway, left=0.3, varlabel] {$q(\textcolor{black}{o_0})$};
\draw[edge] (eq0.east) -- (B0.west) node[midway, above, varlabel] {$\textcolor{black}{q(s^{\scriptscriptstyle I}_0)}$};
\draw[edge] (B0.north) -- (u0.south) node[midway, greycircle] {$\delta$} node[midway, left=0.3, varlabel] {$q(\textcolor{black}{u_0})$};
\draw[edge] (B0.east) -- (eq1.west) node[midway, above, varlabel] {$\textcolor{black}{q(s_1)}$};
\draw[edge] (eq1.south) -- (A1.north) node[midway, left, varlabel] {$\textcolor{black}{q(s^{\scriptscriptstyle II}_1)}$};
\draw[edge] (A1.south) -- (c1.north) node[midway, greycircle] {$\delta$} node[midway, left=0.3, varlabel] {$q(\textcolor{orange}{o_1})$};
\draw[edge] (eq1.east) -- (B1.west) node[midway, above, varlabel] {$\textcolor{black}{q(s^{\scriptscriptstyle I}_1)}$};
\draw[edge] (B1.east) -| (A2.north) node[pos=0.25, above, varlabel] {$\textcolor{black}{q(s_2)}$};
\draw[edge] (A2.south) -- (c2.north) node[midway, left, varlabel] {$\textcolor{black}{q(o_2)}$};
\draw[edge] (B1.north) -- (u1.south) node[midway, left, varlabel] {$\textcolor{black}{q(u_1)}$};

\end{tikzpicture}}
                \end{tabular}%
                \llap{\smash{\raisebox{2.0cm}[0pt][0pt]{\sffamily\bfseries\small (c)\hspace*{-0.0cm}}}}%
                \hspace{0.1cm}% 3. Fester Abstand: RECHTS VOM FFG (Zum Box-Rand)
                \vspace{0.02cm} % 4. Fester Abstand: UNTEN
            \end{minipage}%
        }
        \phantomcaption
        \label{fig:25.3}
    \end{subfigure}
    %\vspace{0.2cm} % Abstand zwischen den Kästen
\end{figure}

\begin{figure}[htbp]
    \ContinuedFloat

    % ================= TIMESTEP 1 =================
    \begin{subfigure}[b]{\textwidth}
        \centering
        \fbox{%
            \begin{minipage}{0.97\linewidth}
                \vspace{0.15cm} % 1. Fester Abstand: OBEN
                % ==========================================
                % BEREICH 1: Das T-Maze (Flexibel zentriert)
                % ==========================================
                \hfill % Drückt das Maze von links weg
                \begin{tabular}{@{}c@{}}
                    \begin{tikzpicture}[scale=1.2]
                        \tikzset{
                            agent/.pic={
                                    \node[circle, draw=black, thick, fill=white, minimum size=0.5cm] at (0,0) {};
                                    \fill[black] (-0.08, 0.08) circle (0.03);
                                    \fill[black] (0.08, 0.08) circle (0.03);
                                    \draw[thick] (-0.1, -0.05) arc (180:360:0.1);
                                },
                            field label/.style={
                                    font=\tiny\sffamily, text=gray!70!black, anchor=south, scale=0.9, transform shape
                                },
                            field symbol/.style={
                                    font=\large\bfseries,
                                    text=gray!70,
                                    anchor=center,
                                    scale=1.0,
                                    transform shape
                                }
                        }

                        \draw (-0.5, 0) -- (-0.5, 2.0) -- (-1.5, 2.0) -- (-1.5, 3.0) -- (1.5, 3.0) -- (1.5, 2.0) -- (0.5, 2.0) -- (0.5, 0) -- cycle;
                        \draw (-0.5, 1.0) -- (0.5, 1.0);
                        \draw (-0.5, 2.0) -- (0.5, 2.0);
                        \draw (-0.5, 2.0) -- (-0.5, 3.0);
                        \draw (0.5, 2.0) -- (0.5, 3.0);

                        \pic[scale=0.8, transform shape] at (0, 0.5) {agent};
                        \node[field label] at (0, 0.97) {Start};
                        \node[field label] at (0, -0.03) {Cue};
                        \node[field label] at (-1.0, 1.97) {Left Arm};
                        \node[field label] at (1.0, 1.93) {Right Arm};
                        %\node[field symbol] at (0, 0.5) {!};
                        \node[field symbol] at (-1.0, 2.5) {?};
                        \node[field symbol] at (1.0, 2.5) {?};

                        % --- NEU: Gedankenblase für Variational Free Energy Minimierung ---
                        % Kleine Bläschen, die vom Kopf (ca. 0.3, 1.7) nach rechts oben steigen
                        \draw[thin, gray!60, fill=white] (0.35, 0.55) circle (0.04);
                        \draw[thin, gray!60, fill=white] (0.52, 0.65) circle (0.07);

                        % Die Haupt-Gedankenblase mit dem Minimierungs-Symbol / Free Energy F
                        \node[
                            ellipse,
                            draw=gray!40,
                            thick,
                            fill=white,
                            minimum width=1.0cm,
                            minimum height=0.55cm,
                            inner sep=2pt,
                            font=\scriptsize\sffamily,
                            text=black!100!black
                        ] at (1.2, 0.7) {$\min F$};
                    \end{tikzpicture}
                \end{tabular}%
                \hfill % Drückt das Maze von der gestrichelten Linie weg (zentriert es somit perfekt im linken Rest-Raum)
                % ==========================================
                % BEREICH 2: Die Trennlinie
                % ==========================================
                \begin{tabular}{@{}c@{}}
                    \tikz \draw[dashed, thin, gray] (0,-2.1) -- (0,2.1);
                \end{tabular}%
                \hspace{0.3cm}% 2. Fester Abstand: LINKS VOM FFG
                % ==========================================
                % BEREICH 3: Der FFG (gibt die Breite vor)
                % ==========================================
                \begin{tabular}{@{}c@{}}
                    % Der FFG wird jetzt in seiner absolut natürlichen Größe gerendert.
                    \resizebox{9cm}{!}{\begin{tikzpicture}[
    factor/.style={draw=black, thick, rectangle, minimum size=0.8cm, font=\normalsize},
    eqfactor/.style={draw=black, thick, rectangle, minimum size=0.5cm, font=\normalsize},
    edge/.style={thick},
    varlabel/.style={font=\normalsize},
    greycircle/.style={draw=black, thick, circle, fill=gray!30, minimum size=0.5cm, font=\normalsize}
]

% Knoten
\node[factor] (d) {$p(s_0)$};
\node[eqfactor, right=1.8cm of d] (eq0) {$=$};
\node[factor, below=1.5cm of eq0] (A0) {$p(o_0|s^{\scriptscriptstyle II}_0)$};
\node[factor, below=1.5cm of A0] (c0) {$\tilde{p}(o_0)$};
\node[factor, right=1.8cm of eq0] (B0) {$p(s_1|s^{\scriptscriptstyle I}_0,u_0)$};
\node[eqfactor, right=1.8cm of B0] (eq1) {$=$};
\node[factor, below=1.5cm of eq1] (A1) {$p(o_1|s^{\scriptscriptstyle II}_1)$};
\node[factor, below=1.5cm of A1] (c1) {$\tilde{p}(o_1)$};
\node[factor, right=1.8cm of eq1] (B1) {$p(s_2|s^{\scriptscriptstyle I}_1,u_1)$};
\node[factor] (A2) at ([xshift=2cm] B1.east |- A1) {$p(o_2|s_2)$};
\node[factor, below=1.5cm of A2] (c2) {$\tilde{p}(o_2)$};
\node[factor, above=1.5cm of B0] (u0) {$p(u_0)$};
\node[factor, above=1.5cm of B1] (u1) {$p(u_1)$};

% Kanten zwischen Knoten
\draw[edge] (d.east) -- (eq0.west) node[midway, above, varlabel] {$\textcolor{black}{q(s_0)}$} node[near end, below, varlabel, scale=0.7] {$\textcolor{orange}{\bm{\overrightarrow{\mu}(}s_0\bm{)}}$} node[near start, below, varlabel, scale=0.7] {$\textcolor{orange}{\bm{\overleftarrow{\mu}(}s_0\bm{)}}$};
\draw[edge] (eq0.south) -- (A0.north) node[midway, left, varlabel] {$\textcolor{black}{q(s^{\scriptscriptstyle II}_0)}$} node[pos=0.25, right=0.05, varlabel, scale=0.7] {$\textcolor{orange}{\bm{{\uparrow}{\mu}(}s^{\scriptscriptstyle II}_0\bm{)}}$} node[pos=0.75, right=0.05, varlabel, scale=0.7] {$\textcolor{orange}{\bm{{\downarrow}{\mu}(}s^{\scriptscriptstyle II}_0\bm{)}}$};
\draw[edge] (A0.south) -- (c0.north) node[midway, greycircle] {$\delta$} node[midway, left=0.3, varlabel] {$\textcolor{black}{q(o_0)}$} node[pos=0.2, right=0.05, varlabel, scale=0.7] {$\textcolor{orange}{\bm{{\uparrow}{\mu}(}o_0\bm{)}}$} node[pos=0.8, right=0.05, varlabel, scale=0.7] {$\textcolor{orange}{\bm{{\downarrow}{\mu}(}o_0\bm{)}}$};
\draw[edge] (eq0.east) -- (B0.west) node[midway, above, varlabel] {$\textcolor{black}{q(s^{\scriptscriptstyle I}_0)}$} node[near start, below, varlabel, scale=0.7] {$\textcolor{orange}{\bm{\overleftarrow{\mu}(}s^{\scriptscriptstyle I}_0\bm{)}}$} node[near end, below, varlabel, scale=0.7] {$\textcolor{orange}{\bm{\overrightarrow{\mu}(}s^{\scriptscriptstyle I}_0\bm{)}}$};
\draw[edge] (B0.north) -- (u0.south) node[midway, greycircle] {$\delta$} node[midway, left=0.3, varlabel] {$\textcolor{black}{q(u_0)}$} node[pos=0.2, right=0.05, varlabel, scale=0.7] {$\textcolor{orange}{\bm{{\downarrow}{\mu}(}u_0\bm{)}}$} node[pos=0.8, right=0.05, varlabel, scale=0.7] {$\textcolor{orange}{\bm{{\uparrow}{\mu}(}u_0\bm{)}}$};
\draw[edge] (B0.east) -- (eq1.west) node[midway, above, varlabel] {$\textcolor{black}{q(s_1)}$} node[near end, below, varlabel, scale=0.7] {$\textcolor{orange}{\bm{\overrightarrow{\mu}(}s_1\bm{)}}$} node[near start, below, varlabel, scale=0.7] {$\textcolor{orange}{\bm{\overleftarrow{\mu}(}s_1\bm{)}}$};
\draw[edge] (eq1.south) -- (A1.north) node[midway, left, varlabel] {$\textcolor{black}{q(s^{\scriptscriptstyle II}_1)}$} node[pos=0.25, right=0.05, varlabel, scale=0.7] {$\textcolor{orange}{\bm{{\uparrow}{\mu}(}s^{\scriptscriptstyle II}_1\bm{)}}$} node[pos=0.75, right=0.05, varlabel, scale=0.7] {$\textcolor{orange}{\bm{{\downarrow}{\mu}(}s^{\scriptscriptstyle II}_1\bm{)}}$};
\draw[edge] (A1.south) -- (c1.north) node[midway, greycircle] {$\delta$} node[midway, left=0.3, varlabel] {$\textcolor{black}{q(o_1)}$} node[pos=0.2, right=0.05, varlabel, scale=0.7] {$\textcolor{orange}{\bm{{\uparrow}{\mu}(}o_1\bm{)}}$} node[pos=0.8, right=0.05, varlabel, scale=0.7] {$\textcolor{orange}{\bm{{\downarrow}{\mu}(}o_1\bm{)}}$};
\draw[edge] (eq1.east) -- (B1.west) node[midway, above, varlabel] {$\textcolor{black}{q(s^{\scriptscriptstyle I}_1)}$} node[near start, below, varlabel, scale=0.7] {$\textcolor{orange}{\bm{\overleftarrow{\mu}(}s^{\scriptscriptstyle I}_1\bm{)}}$} node[near end, below, varlabel, scale=0.7] {$\textcolor{orange}{\bm{\overrightarrow{\mu}(}s^{\scriptscriptstyle I}_1\bm{)}}$};
\draw[edge] (B1.east) -| (A2.north) node[pos=0.25, above, varlabel] {$\textcolor{black}{q(s_2)}$} node[pos=0.89, right=0.05, varlabel, scale=0.7] {$\textcolor{orange}{\bm{{\downarrow}{\mu}(}s_2\bm{)}}$} node[pos=0.11, below, varlabel, scale=0.7] {$\textcolor{orange}{\bm{\overleftarrow{\mu}(}s_2\bm{)}}$};
\draw[edge] (A2.south) -- (c2.north) node[midway, left, varlabel] {$\textcolor{black}{q(o_2)}$} node[pos=0.2, right=0.05, varlabel, scale=0.7] {$\textcolor{orange}{\bm{{\uparrow}{\mu}(}o_2\bm{)}}$} node[pos=0.8, right=0.05, varlabel, scale=0.7] {$\textcolor{orange}{\bm{{\downarrow}{\mu}(}o_2\bm{)}}$};
\draw[edge] (B1.north) -- (u1.south) node[midway, left, varlabel] {$\textcolor{black}{q(u_1)}$} node[pos=0.2, right=0.05, varlabel, scale=0.7] {$\textcolor{orange}{\bm{{\downarrow}{\mu}(}u_1\bm{)}}$} node[pos=0.8, right=0.05, varlabel, scale=0.7] {$\textcolor{orange}{\bm{{\uparrow}{\mu}(}u_1\bm{)}}$};

\end{tikzpicture}}
                \end{tabular}%
                \llap{\smash{\raisebox{2.0cm}[0pt][0pt]{\sffamily\bfseries\small (d)\hspace*{-0.0cm}}}}%
                \hspace{0.1cm}% 3. Fester Abstand: RECHTS VOM FFG (Zum Box-Rand)
                \vspace{0.02cm} % 4. Fester Abstand: UNTEN
            \end{minipage}%
        }
        \phantomcaption
        \label{fig:25.4}
    \end{subfigure}

    \caption{The figure presents the four steps of the simulation "Exploiting Information" via the example of the BFE agent. First, the agents receives an observation \(\hat {o_0}\) which makes \(q(o_0\) a delta distribution with all its probability mass at \(\hat {o_0}\) (\cref{fig:25.1}). Second, the action "go to cue" is externally triggered and the respective belief distribution \(q(u_0)\) also becomes a delta distribution with its mass at this action (\cref{fig:25.2}). Third, being at the cue triggeres an observation from the environment that is informative regarding the reward location. Similar to the first observation, it leads to a delta distribution at \(q(o_1)\) (\cref{fig:25.3}). Lastly, the agent executes its message passing algorithm to infer all its belief distributions (\cref{fig:25.4}). The inclination of moving to the cued arm that promises the reward is then given as a probability in the action-specific belief distribution \(q(u_1)\).}
    \label{fig:25}
\end{figure}

\subsubsection{Experiment 3: Long-Term Success}

This experiment is about empirically measuring the agents' succes in receiving the reward during a trial. Specifically, there are two ways how an agent can achieve this: First, it initially "decides" against going to the cue and instead moves to one of the two arms. It is then either lucky and observes the reward or not. Second, the agent initially moves to the cue and only then moves to one of the two arms. Again, it then either observes the reward or not. Generally, a trial consists of three timesteps \(t \in {0,1,2}\) and the agent can either observe the reward in timestep \(t=1\) (when directly moving to an arm) or in timestep \(t=2\) (when first moving to the cue in \(t=1\)). A trial is succesfull when the agent observes the reward in timesteps \(t=1\) or \(t=2\) (a trial is prematurely terminated when it receives in reward in timestep \(t=1\)). There are \(1000\) trials for each agent (and \(\alpha\)- and \(c\) value). The result is the percentage of succesfull trials. The procedere is detailed in \cref{fig:26}.

\begin{figure}[b!]
    \centering

    % ================= TIMESTEP 1.1 =================
    \begin{subfigure}[b]{0.49\textwidth}
        \centering
        \fbox{%
            \begin{minipage}{0.97\linewidth}
                \vspace{0.15cm} % 1. Fester Abstand: OBEN
                % ==========================================
                % BEREICH 1: Das T-Maze (Flexibel zentriert)
                % ==========================================
                \hfill % Drückt das Maze von links weg
                \begin{tabular}{@{}c@{}}
                    \begin{tikzpicture}[scale=0.6]
                        \tikzset{
                            agent/.pic={
                                    \node[circle, draw=black, thick, fill=white, minimum size=0.5cm] at (0,0) {};
                                    \fill[black] (-0.08, 0.08) circle (0.03);
                                    \fill[black] (0.08, 0.08) circle (0.03);
                                    \draw[thick] (-0.1, -0.05) arc (180:360:0.1);
                                },
                            field label/.style={
                                    font=\tiny\sffamily, text=gray!70!black, anchor=south, scale=0.9, transform shape
                                },
                            field symbol/.style={
                                    font=\large\bfseries,
                                    text=gray!70,
                                    anchor=center,
                                    scale=1.0,
                                    transform shape
                                }
                        }

                        \draw (-0.5, 0) -- (-0.5, 2.0) -- (-1.5, 2.0) -- (-1.5, 3.0) -- (1.5, 3.0) -- (1.5, 2.0) -- (0.5, 2.0) -- (0.5, 0) -- cycle;
                        \draw (-0.5, 1.0) -- (0.5, 1.0);
                        \draw (-0.5, 2.0) -- (0.5, 2.0);
                        \draw (-0.5, 2.0) -- (-0.5, 3.0);
                        \draw (0.5, 2.0) -- (0.5, 3.0);

                        \pic[scale=0.8, transform shape] at (0, 1.5) {agent};
                        % Observation Arrow
                        \draw[
                            -{Stealth[scale=1.0]},
                            line width=1.6pt,
                            color=gray!80!black
                        ]
                        (0.51, 1.3) .. controls (1.0, 1.1) and (1.3, 2.0) .. (0.3, 1.6);
                        \node[font=\footnotesize, text=gray!80!black] at (1.2, 1.2) {$\textcolor{orange}{\hat {o_0}}$};
                        \node[field label] at (0, 0.97) {Start};
                        \node[field label] at (0, -0.03) {Cue};
                        \node[field label] at (-1.0, 1.97) {Left Arm};
                        \node[field label] at (1.0, 1.93) {Right Arm};
                        \node[field symbol] at (0, 0.5) {!};
                        \node[field symbol] at (-1.0, 2.5) {?};
                        \node[field symbol] at (1.0, 2.5) {?};
                    \end{tikzpicture}
                \end{tabular}%
                \hfill % Drückt das Maze von der gestrichelten Linie weg (zentriert es somit perfekt im linken Rest-Raum)
                % ==========================================
                % BEREICH 2: Die Trennlinie
                % ==========================================
                \begin{tabular}{@{}c@{}}
                    \tikz \draw[dashed, thin, gray] (0,-1.05) -- (0,1.05);
                \end{tabular}%
                \hspace{0.3cm}% 2. Fester Abstand: LINKS VOM FFG
                % ==========================================
                % BEREICH 3: Der FFG (gibt die Breite vor)
                % ==========================================
                \begin{tabular}{@{}c@{}}
                    % Der FFG wird jetzt in seiner absolut natürlichen Größe gerendert.
                    \resizebox{4.5cm}{!}{\begin{tikzpicture}[
    factor/.style={draw=black, thick, rectangle, minimum size=0.8cm, font=\normalsize},
    eqfactor/.style={draw=black, thick, rectangle, minimum size=0.5cm, font=\normalsize},
    edge/.style={thick},
    varlabel/.style={font=\normalsize},
    greycircle/.style={draw=black, thick, circle, fill=gray!30, minimum size=0.5cm, font=\normalsize}
]

% Knoten
\node[factor] (d) {$p(s_0)$};
\node[eqfactor, right=1.8cm of d] (eq0) {$=$};
\node[factor, below=1.5cm of eq0] (A0) {$p(o_0|s^{\scriptscriptstyle II}_0)$};
\node[factor, below=1.5cm of A0] (c0) {$\tilde{p}(o_0)$};
\node[factor, right=1.8cm of eq0] (B0) {$p(s_1|s^{\scriptscriptstyle I}_0,u_0)$};
\node[eqfactor, right=1.8cm of B0] (eq1) {$=$};
\node[factor, below=1.5cm of eq1] (A1) {$p(o_1|s^{\scriptscriptstyle II}_1)$};
\node[factor, below=1.5cm of A1] (c1) {$\tilde{p}(o_1)$};
\node[factor, right=1.8cm of eq1] (B1) {$p(s_2|s^{\scriptscriptstyle I}_1,u_1)$};
\node[factor] (A2) at ([xshift=2cm] B1.east |- A1) {$p(o_2|s_2)$};
\node[factor, below=1.5cm of A2] (c2) {$\tilde{p}(o_2)$};
\node[factor, above=1.5cm of B0] (u0) {$p(u_0)$};
\node[factor, above=1.5cm of B1] (u1) {$p(u_1)$};

% Kanten zwischen Knoten
\draw[edge] (d.east) -- (eq0.west) node[midway, above, varlabel] {$\textcolor{black}{q(s_0)}$};
\draw[edge] (eq0.south) -- (A0.north) node[midway, left, varlabel] {$\textcolor{black}{q(s^{\scriptscriptstyle II}_0)}$};
\draw[edge] (A0.south) -- (c0.north) node[midway, greycircle] {$\delta$} node[midway, left=0.3, varlabel] {$q(\textcolor{orange}{o_0})$};
\draw[edge] (eq0.east) -- (B0.west) node[midway, above, varlabel] {$\textcolor{black}{q(s^{\scriptscriptstyle I}_0)}$};
\draw[edge] (B0.north) -- (u0.south) node[midway, left, varlabel] {$\textcolor{black}{q(u_0)}$};
\draw[edge] (B0.east) -- (eq1.west) node[midway, above, varlabel] {$\textcolor{black}{q(s_1)}$};
\draw[edge] (eq1.south) -- (A1.north) node[midway, left, varlabel] {$\textcolor{black}{q(s^{\scriptscriptstyle II}_1)}$};
\draw[edge] (A1.south) -- (c1.north) node[midway, left, varlabel] {$\textcolor{black}{q(o_1)}$};
\draw[edge] (eq1.east) -- (B1.west) node[midway, above, varlabel] {$\textcolor{black}{q(s^{\scriptscriptstyle I}_1)}$};
\draw[edge] (B1.east) -| (A2.north) node[pos=0.25, above, varlabel] {$\textcolor{black}{q(s_2)}$};
\draw[edge] (A2.south) -- (c2.north) node[midway, left, varlabel] {$\textcolor{black}{q(o_2)}$};
\draw[edge] (B1.north) -- (u1.south) node[midway, left, varlabel] {$\textcolor{black}{q(u_1)}$};

\end{tikzpicture}}
                \end{tabular}%
                \llap{\smash{\raisebox{1.0cm}[0pt][0pt]{\sffamily\bfseries\small (a)\hspace*{-0.05cm}}}}%
                \hspace{0.1cm}% 3. Fester Abstand: RECHTS VOM FFG (Zum Box-Rand)
                \vspace{0.02cm} % 4. Fester Abstand: UNTEN
            \end{minipage}%
        }
        \phantomcaption
        \label{fig:26.1a}
    \end{subfigure}
    \hfill
    % ================= TIMESTEP 1.2 =================
    \begin{subfigure}[b]{0.49\textwidth}
        \centering
        \fbox{%
            \begin{minipage}{0.97\linewidth}
                \vspace{0.15cm} % 1. Fester Abstand: OBEN
                % ==========================================
                % BEREICH 1: Das T-Maze (Flexibel zentriert)
                % ==========================================
                \hfill % Drückt das Maze von links weg
                \begin{tabular}{@{}c@{}}
                    \begin{tikzpicture}[scale=0.6]
                        \tikzset{
                            agent/.pic={
                                    \node[circle, draw=black, thick, fill=white, minimum size=0.5cm] at (0,0) {};
                                    \fill[black] (-0.08, 0.08) circle (0.03);
                                    \fill[black] (0.08, 0.08) circle (0.03);
                                    \draw[thick] (-0.1, -0.05) arc (180:360:0.1);
                                },
                            field label/.style={
                                    font=\tiny\sffamily, text=gray!70!black, anchor=south, scale=0.9, transform shape
                                },
                            field symbol/.style={
                                    font=\large\bfseries,
                                    text=gray!70,
                                    anchor=center,
                                    scale=1.0,
                                    transform shape
                                }
                        }

                        \draw (-0.5, 0) -- (-0.5, 2.0) -- (-1.5, 2.0) -- (-1.5, 3.0) -- (1.5, 3.0) -- (1.5, 2.0) -- (0.5, 2.0) -- (0.5, 0) -- cycle;
                        \draw (-0.5, 1.0) -- (0.5, 1.0);
                        \draw (-0.5, 2.0) -- (0.5, 2.0);
                        \draw (-0.5, 2.0) -- (-0.5, 3.0);
                        \draw (0.5, 2.0) -- (0.5, 3.0);

                        \pic[scale=0.8, transform shape] at (0, 1.5) {agent};
                        % Observation Arrow
                        \draw[
                            -{Stealth[scale=1.0]},
                            line width=1.6pt,
                            color=gray!80!black
                        ]
                        (0.51, 1.3) .. controls (1.0, 1.1) and (1.3, 2.0) .. (0.3, 1.6);
                        \node[font=\footnotesize, text=gray!80!black] at (1.2, 1.2) {$\textcolor{orange}{\hat {o_0}}$};
                        \node[field label] at (0, 0.97) {Start};
                        \node[field label] at (0, -0.03) {Cue};
                        \node[field label] at (-1.0, 1.97) {Left Arm};
                        \node[field label] at (1.0, 1.93) {Right Arm};
                        \node[field symbol] at (0, 0.5) {!};
                        \node[field symbol] at (-1.0, 2.5) {?};
                        \node[field symbol] at (1.0, 2.5) {?};
                    \end{tikzpicture}
                \end{tabular}%
                \hfill % Drückt das Maze von der gestrichelten Linie weg (zentriert es somit perfekt im linken Rest-Raum)
                % ==========================================
                % BEREICH 2: Die Trennlinie
                % ==========================================
                \begin{tabular}{@{}c@{}}
                    \tikz \draw[dashed, thin, gray] (0,-1.05) -- (0,1.05);
                \end{tabular}%
                \hspace{0.3cm}% 2. Fester Abstand: LINKS VOM FFG
                % ==========================================
                % BEREICH 3: Der FFG (gibt die Breite vor)
                % ==========================================
                \begin{tabular}{@{}c@{}}
                    % Der FFG wird jetzt in seiner absolut natürlichen Größe gerendert.
                    \resizebox{4.5cm}{!}{\begin{tikzpicture}[
    factor/.style={draw=black, thick, rectangle, minimum size=0.8cm, font=\normalsize},
    eqfactor/.style={draw=black, thick, rectangle, minimum size=0.5cm, font=\normalsize},
    edge/.style={thick},
    varlabel/.style={font=\normalsize},
    greycircle/.style={draw=black, thick, circle, fill=gray!30, minimum size=0.5cm, font=\normalsize}
]

% Knoten
\node[factor] (d) {$p(s_0)$};
\node[eqfactor, right=1.8cm of d] (eq0) {$=$};
\node[factor, below=1.5cm of eq0] (A0) {$p(o_0|s^{\scriptscriptstyle II}_0)$};
\node[factor, below=1.5cm of A0] (c0) {$\tilde{p}(o_0)$};
\node[factor, right=1.8cm of eq0] (B0) {$p(s_1|s^{\scriptscriptstyle I}_0,u_0)$};
\node[eqfactor, right=1.8cm of B0] (eq1) {$=$};
\node[factor, below=1.5cm of eq1] (A1) {$p(o_1|s^{\scriptscriptstyle II}_1)$};
\node[factor, below=1.5cm of A1] (c1) {$\tilde{p}(o_1)$};
\node[factor, right=1.8cm of eq1] (B1) {$p(s_2|s^{\scriptscriptstyle I}_1,u_1)$};
\node[factor] (A2) at ([xshift=2cm] B1.east |- A1) {$p(o_2|s_2)$};
\node[factor, below=1.5cm of A2] (c2) {$\tilde{p}(o_2)$};
\node[factor, above=1.5cm of B0] (u0) {$p(u_0)$};
\node[factor, above=1.5cm of B1] (u1) {$p(u_1)$};

% Kanten zwischen Knoten
\draw[edge] (d.east) -- (eq0.west) node[midway, above, varlabel] {$\textcolor{black}{q(s_0)}$};
\draw[edge] (eq0.south) -- (A0.north) node[midway, left, varlabel] {$\textcolor{black}{q(s^{\scriptscriptstyle II}_0)}$};
\draw[edge] (A0.south) -- (c0.north) node[midway, greycircle] {$\delta$} node[midway, left=0.3, varlabel] {$q(\textcolor{orange}{o_0})$};
\draw[edge] (eq0.east) -- (B0.west) node[midway, above, varlabel] {$\textcolor{black}{q(s^{\scriptscriptstyle I}_0)}$};
\draw[edge] (B0.north) -- (u0.south) node[midway, left, varlabel] {$\textcolor{black}{q(u_0)}$};
\draw[edge] (B0.east) -- (eq1.west) node[midway, above, varlabel] {$\textcolor{black}{q(s_1)}$};
\draw[edge] (eq1.south) -- (A1.north) node[midway, left, varlabel] {$\textcolor{black}{q(s^{\scriptscriptstyle II}_1)}$};
\draw[edge] (A1.south) -- (c1.north) node[midway, left, varlabel] {$\textcolor{black}{q(o_1)}$};
\draw[edge] (eq1.east) -- (B1.west) node[midway, above, varlabel] {$\textcolor{black}{q(s^{\scriptscriptstyle I}_1)}$};
\draw[edge] (B1.east) -| (A2.north) node[pos=0.25, above, varlabel] {$\textcolor{black}{q(s_2)}$};
\draw[edge] (A2.south) -- (c2.north) node[midway, left, varlabel] {$\textcolor{black}{q(o_2)}$};
\draw[edge] (B1.north) -- (u1.south) node[midway, left, varlabel] {$\textcolor{black}{q(u_1)}$};

\end{tikzpicture}}
                \end{tabular}%
                \llap{\smash{\raisebox{1.0cm}[0pt][0pt]{\sffamily\bfseries\small (b)\hspace*{-0.05cm}}}}%
                \hspace{0.1cm}% 3. Fester Abstand: RECHTS VOM FFG (Zum Box-Rand)
                \vspace{0.02cm} % 4. Fester Abstand: UNTEN
            \end{minipage}%
        }
        \phantomcaption
        \label{fig:26.1b}
    \end{subfigure}

    \vspace{0.2cm} % Abstand zwischen den Kästen

    % ================= TIMESTEP 2.1 =================
    \begin{subfigure}[b]{0.49\textwidth}
        \centering
        \fbox{%
            \begin{minipage}{0.97\linewidth}
                \vspace{0.15cm} % 1. Fester Abstand: OBEN
                % ==========================================
                % BEREICH 1: Das T-Maze (Flexibel zentriert)
                % ==========================================
                \hfill % Drückt das Maze von links weg
                \begin{tabular}{@{}c@{}}
                    \begin{tikzpicture}[scale=0.6]
                        \tikzset{
                            agent/.pic={
                                    \node[circle, draw=black, thick, fill=white, minimum size=0.5cm] at (0,0) {};
                                    \fill[black] (-0.08, 0.08) circle (0.03);
                                    \fill[black] (0.08, 0.08) circle (0.03);
                                    \draw[thick] (-0.1, -0.05) arc (180:360:0.1);
                                },
                            field label/.style={
                                    font=\tiny\sffamily, text=gray!70!black, anchor=south, scale=0.9, transform shape
                                },
                            field symbol/.style={
                                    font=\large\bfseries,
                                    text=gray!70,
                                    anchor=center,
                                    scale=1.0,
                                    transform shape
                                }
                        }

                        \draw (-0.5, 0) -- (-0.5, 2.0) -- (-1.5, 2.0) -- (-1.5, 3.0) -- (1.5, 3.0) -- (1.5, 2.0) -- (0.5, 2.0) -- (0.5, 0) -- cycle;
                        \draw (-0.5, 1.0) -- (0.5, 1.0);
                        \draw (-0.5, 2.0) -- (0.5, 2.0);
                        \draw (-0.5, 2.0) -- (-0.5, 3.0);
                        \draw (0.5, 2.0) -- (0.5, 3.0);

                        \pic[scale=0.8, transform shape] at (0, 1.5) {agent};
                        \node[field label] at (0, 0.97) {Start};
                        \node[field label] at (0, -0.03) {Cue};
                        \node[field label] at (-1.0, 1.97) {Left Arm};
                        \node[field label] at (1.0, 1.93) {Right Arm};
                        \node[field symbol] at (0, 0.5) {!};
                        \node[field symbol] at (-1.0, 2.5) {?};
                        \node[field symbol] at (1.0, 2.5) {?};

                        % --- NEU: Gedankenblase für Variational Free Energy Minimierung ---
                        % Kleine Bläschen, die vom Kopf (ca. 0.3, 1.7) nach rechts oben steigen
                        \draw[thin, gray!60, fill=white] (0.35, 1.55) circle (0.04);
                        \draw[thin, gray!60, fill=white] (0.52, 1.65) circle (0.07);

                        % Die Haupt-Gedankenblase mit dem Minimierungs-Symbol / Free Energy F
                        \node[
                            ellipse,
                            draw=gray!40,
                            thick,
                            fill=white,
                            minimum width=1.0cm,
                            minimum height=0.55cm,
                            inner sep=2pt,
                            font=\scriptsize\sffamily,
                            text=black!100!black,
                            transform shape
                        ] at (1.2, 1.7) {$\min F$};
                    \end{tikzpicture}
                \end{tabular}%
                \hfill % Drückt das Maze von der gestrichelten Linie weg (zentriert es somit perfekt im linken Rest-Raum)
                % ==========================================
                % BEREICH 2: Die Trennlinie
                % ==========================================
                \begin{tabular}{@{}c@{}}
                    \tikz \draw[dashed, thin, gray] (0,-1.05) -- (0,1.05);
                \end{tabular}%
                \hspace{0.3cm}% 2. Fester Abstand: LINKS VOM FFG
                % ==========================================
                % BEREICH 3: Der FFG (gibt die Breite vor)
                % ==========================================
                \begin{tabular}{@{}c@{}}
                    % Der FFG wird jetzt in seiner absolut natürlichen Größe gerendert.
                    \resizebox{4.4cm}{!}{\begin{tikzpicture}[
    factor/.style={draw=black, thick, rectangle, minimum size=0.8cm, font=\normalsize},
    eqfactor/.style={draw=black, thick, rectangle, minimum size=0.5cm, font=\normalsize},
    edge/.style={thick},
    varlabel/.style={font=\normalsize},
    greycircle/.style={draw=black, thick, circle, fill=gray!30, minimum size=0.5cm, font=\normalsize}
]

% Knoten
\node[factor] (d) {$p(s_0)$};
\node[eqfactor, right=1.8cm of d] (eq0) {$=$};
\node[factor, below=1.5cm of eq0] (A0) {$p(o_0|s^{\scriptscriptstyle II}_0)$};
\node[factor, below=1.5cm of A0] (c0) {$\tilde{p}(o_0)$};
\node[factor, right=1.8cm of eq0] (B0) {$p(s_1|s^{\scriptscriptstyle I}_0,u_0)$};
\node[eqfactor, right=1.8cm of B0] (eq1) {$=$};
\node[factor, below=1.5cm of eq1] (A1) {$p(o_1|s^{\scriptscriptstyle II}_1)$};
\node[factor, below=1.5cm of A1] (c1) {$\tilde{p}(o_1)$};
\node[factor, right=1.8cm of eq1] (B1) {$p(s_2|s^{\scriptscriptstyle I}_1,u_1)$};
\node[factor] (A2) at ([xshift=2cm] B1.east |- A1) {$p(o_2|s_2)$};
\node[factor, below=1.5cm of A2] (c2) {$\tilde{p}(o_2)$};
\node[factor, above=1.5cm of B0] (u0) {$p(u_0)$};
\node[factor, above=1.5cm of B1] (u1) {$p(u_1)$};

% Kanten zwischen Knoten
\draw[edge] (d.east) -- (eq0.west) node[midway, above, varlabel] {$\textcolor{black}{q(s_0)}$} node[near end, below, varlabel, scale=0.7] {$\textcolor{orange}{\bm{\overrightarrow{\mu}(}s_0\bm{)}}$} node[near start, below, varlabel, scale=0.7] {$\textcolor{orange}{\bm{\overleftarrow{\mu}(}s_0\bm{)}}$};
\draw[edge] (eq0.south) -- (A0.north) node[midway, left, varlabel] {$\textcolor{black}{q(s^{\scriptscriptstyle II}_0)}$} node[pos=0.25, right=0.05, varlabel, scale=0.7] {$\textcolor{orange}{\bm{{\uparrow}{\mu}(}s^{\scriptscriptstyle II}_0\bm{)}}$} node[pos=0.75, right=0.05, varlabel, scale=0.7] {$\textcolor{orange}{\bm{{\downarrow}{\mu}(}s^{\scriptscriptstyle II}_0\bm{)}}$};
\draw[edge] (A0.south) -- (c0.north) node[midway, greycircle] {$\delta$} node[midway, left=0.3, varlabel] {$\textcolor{black}{q(o_0)}$} node[pos=0.2, right=0.05, varlabel, scale=0.7] {$\textcolor{orange}{\bm{{\uparrow}{\mu}(}o_0\bm{)}}$} node[pos=0.8, right=0.05, varlabel, scale=0.7] {$\textcolor{orange}{\bm{{\downarrow}{\mu}(}o_0\bm{)}}$};
\draw[edge] (eq0.east) -- (B0.west) node[midway, above, varlabel] {$\textcolor{black}{q(s^{\scriptscriptstyle I}_0)}$} node[near start, below, varlabel, scale=0.7] {$\textcolor{orange}{\bm{\overleftarrow{\mu}(}s^{\scriptscriptstyle I}_0\bm{)}}$} node[near end, below, varlabel, scale=0.7] {$\textcolor{orange}{\bm{\overrightarrow{\mu}(}s^{\scriptscriptstyle I}_0\bm{)}}$};
\draw[edge] (B0.north) -- (u0.south) node[midway, left, varlabel] {$\textcolor{black}{q(u_0)}$} node[pos=0.2, right=0.05, varlabel, scale=0.7] {$\textcolor{orange}{\bm{{\downarrow}{\mu}(}u_0\bm{)}}$} node[pos=0.8, right=0.05, varlabel, scale=0.7] {$\textcolor{orange}{\bm{{\uparrow}{\mu}(}u_0\bm{)}}$};
\draw[edge] (B0.east) -- (eq1.west) node[midway, above, varlabel] {$\textcolor{black}{q(s_1)}$} node[near end, below, varlabel, scale=0.7] {$\textcolor{orange}{\bm{\overrightarrow{\mu}(}s_1\bm{)}}$} node[near start, below, varlabel, scale=0.7] {$\textcolor{orange}{\bm{\overleftarrow{\mu}(}s_1\bm{)}}$};
\draw[edge] (eq1.south) -- (A1.north) node[midway, left, varlabel] {$\textcolor{black}{q(s^{\scriptscriptstyle II}_1)}$} node[pos=0.25, right=0.05, varlabel, scale=0.7] {$\textcolor{orange}{\bm{{\uparrow}{\mu}(}s^{\scriptscriptstyle II}_1\bm{)}}$} node[pos=0.75, right=0.05, varlabel, scale=0.7] {$\textcolor{orange}{\bm{{\downarrow}{\mu}(}s^{\scriptscriptstyle II}_1\bm{)}}$};
\draw[edge] (A1.south) -- (c1.north) node[midway, left, varlabel] {$\textcolor{black}{q(o_1)}$} node[pos=0.2, right=0.05, varlabel, scale=0.7] {$\textcolor{orange}{\bm{{\uparrow}{\mu}(}o_1\bm{)}}$} node[pos=0.8, right=0.05, varlabel, scale=0.7] {$\textcolor{orange}{\bm{{\downarrow}{\mu}(}o_1\bm{)}}$};
\draw[edge] (eq1.east) -- (B1.west) node[midway, above, varlabel] {$\textcolor{black}{q(s^{\scriptscriptstyle I}_1)}$} node[near start, below, varlabel, scale=0.7] {$\textcolor{orange}{\bm{\overleftarrow{\mu}(}s^{\scriptscriptstyle I}_1\bm{)}}$} node[near end, below, varlabel, scale=0.7] {$\textcolor{orange}{\bm{\overrightarrow{\mu}(}s^{\scriptscriptstyle I}_1\bm{)}}$};
\draw[edge] (B1.east) -| (A2.north) node[pos=0.25, above, varlabel] {$\textcolor{black}{q(s_2)}$} node[pos=0.89, right=0.05, varlabel, scale=0.7] {$\textcolor{orange}{\bm{{\downarrow}{\mu}(}s_2\bm{)}}$} node[pos=0.11, below, varlabel, scale=0.7] {$\textcolor{orange}{\bm{\overleftarrow{\mu}(}s_2\bm{)}}$};
\draw[edge] (A2.south) -- (c2.north) node[midway, left, varlabel] {$\textcolor{black}{q(o_2)}$} node[pos=0.2, right=0.05, varlabel, scale=0.7] {$\textcolor{orange}{\bm{{\uparrow}{\mu}(}o_2\bm{)}}$} node[pos=0.8, right=0.05, varlabel, scale=0.7] {$\textcolor{orange}{\bm{{\downarrow}{\mu}(}o_2\bm{)}}$};
\draw[edge] (B1.north) -- (u1.south) node[midway, left, varlabel] {$\textcolor{black}{q(u_1)}$} node[pos=0.2, right=0.05, varlabel, scale=0.7] {$\textcolor{orange}{\bm{{\downarrow}{\mu}(}u_1\bm{)}}$} node[pos=0.8, right=0.05, varlabel, scale=0.7] {$\textcolor{orange}{\bm{{\uparrow}{\mu}(}u_1\bm{)}}$};

\end{tikzpicture}}
                \end{tabular}%
                \llap{\smash{\raisebox{1.0cm}[0pt][0pt]{\sffamily\bfseries\small (c)\hspace*{-0.05cm}}}}%
                \hspace{0.1cm}% 3. Fester Abstand: RECHTS VOM FFG (Zum Box-Rand)
                \vspace{0.02cm} % 4. Fester Abstand: UNTEN
            \end{minipage}%
        }
        \phantomcaption
        \label{fig:26.2a}
    \end{subfigure}
    \hfill
    % ================= TIMESTEP 2.2 =================
    \begin{subfigure}[b]{0.49\textwidth}
        \centering
        \fbox{%
            \begin{minipage}{0.97\linewidth}
                \vspace{0.15cm} % 1. Fester Abstand: OBEN
                % ==========================================
                % BEREICH 1: Das T-Maze (Flexibel zentriert)
                % ==========================================
                \hfill % Drückt das Maze von links weg
                \begin{tabular}{@{}c@{}}
                    \begin{tikzpicture}[scale=0.6]
                        \tikzset{
                            agent/.pic={
                                    \node[circle, draw=black, thick, fill=white, minimum size=0.5cm] at (0,0) {};
                                    \fill[black] (-0.08, 0.08) circle (0.03);
                                    \fill[black] (0.08, 0.08) circle (0.03);
                                    \draw[thick] (-0.1, -0.05) arc (180:360:0.1);
                                },
                            field label/.style={
                                    font=\tiny\sffamily, text=gray!70!black, anchor=south, scale=0.9, transform shape
                                },
                            field symbol/.style={
                                    font=\large\bfseries,
                                    text=gray!70,
                                    anchor=center,
                                    scale=1.0,
                                    transform shape
                                }
                        }

                        \draw (-0.5, 0) -- (-0.5, 2.0) -- (-1.5, 2.0) -- (-1.5, 3.0) -- (1.5, 3.0) -- (1.5, 2.0) -- (0.5, 2.0) -- (0.5, 0) -- cycle;
                        \draw (-0.5, 1.0) -- (0.5, 1.0);
                        \draw (-0.5, 2.0) -- (0.5, 2.0);
                        \draw (-0.5, 2.0) -- (-0.5, 3.0);
                        \draw (0.5, 2.0) -- (0.5, 3.0);

                        \pic[scale=0.8, transform shape] at (0, 1.5) {agent};
                        \node[field label] at (0, 0.97) {Start};
                        \node[field label] at (0, -0.03) {Cue};
                        \node[field label] at (-1.0, 1.97) {Left Arm};
                        \node[field label] at (1.0, 1.93) {Right Arm};
                        \node[field symbol] at (0, 0.5) {!};
                        \node[field symbol] at (-1.0, 2.5) {?};
                        \node[field symbol] at (1.0, 2.5) {?};

                        % --- NEU: Gedankenblase für Variational Free Energy Minimierung ---
                        % Kleine Bläschen, die vom Kopf (ca. 0.3, 1.7) nach rechts oben steigen
                        \draw[thin, gray!60, fill=white] (0.35, 1.55) circle (0.04);
                        \draw[thin, gray!60, fill=white] (0.52, 1.65) circle (0.07);

                        % Die Haupt-Gedankenblase mit dem Minimierungs-Symbol / Free Energy F
                        \node[
                            ellipse,
                            draw=gray!40,
                            thick,
                            fill=white,
                            minimum width=1.0cm,
                            minimum height=0.55cm,
                            inner sep=2pt,
                            font=\scriptsize\sffamily,
                            text=black!100!black,
                            transform shape
                        ] at (1.2, 1.7) {$\min F$};
                    \end{tikzpicture}
                \end{tabular}%
                \hfill % Drückt das Maze von der gestrichelten Linie weg (zentriert es somit perfekt im linken Rest-Raum)
                % ==========================================
                % BEREICH 2: Die Trennlinie
                % ==========================================
                \begin{tabular}{@{}c@{}}
                    \tikz \draw[dashed, thin, gray] (0,-1.05) -- (0,1.05);
                \end{tabular}%
                \hspace{0.3cm}% 2. Fester Abstand: LINKS VOM FFG
                % ==========================================
                % BEREICH 3: Der FFG (gibt die Breite vor)
                % ==========================================
                \begin{tabular}{@{}c@{}}
                    % Der FFG wird jetzt in seiner absolut natürlichen Größe gerendert.
                    \resizebox{4.4cm}{!}{\begin{tikzpicture}[
    factor/.style={draw=black, thick, rectangle, minimum size=0.8cm, font=\normalsize},
    eqfactor/.style={draw=black, thick, rectangle, minimum size=0.5cm, font=\normalsize},
    edge/.style={thick},
    varlabel/.style={font=\normalsize},
    greycircle/.style={draw=black, thick, circle, fill=gray!30, minimum size=0.5cm, font=\normalsize}
]

% Knoten
\node[factor] (d) {$p(s_0)$};
\node[eqfactor, right=1.8cm of d] (eq0) {$=$};
\node[factor, below=1.5cm of eq0] (A0) {$p(o_0|s^{\scriptscriptstyle II}_0)$};
\node[factor, below=1.5cm of A0] (c0) {$\tilde{p}(o_0)$};
\node[factor, right=1.8cm of eq0] (B0) {$p(s_1|s^{\scriptscriptstyle I}_0,u_0)$};
\node[eqfactor, right=1.8cm of B0] (eq1) {$=$};
\node[factor, below=1.5cm of eq1] (A1) {$p(o_1|s^{\scriptscriptstyle II}_1)$};
\node[factor, below=1.5cm of A1] (c1) {$\tilde{p}(o_1)$};
\node[factor, right=1.8cm of eq1] (B1) {$p(s_2|s^{\scriptscriptstyle I}_1,u_1)$};
\node[factor] (A2) at ([xshift=2cm] B1.east |- A1) {$p(o_2|s_2)$};
\node[factor, below=1.5cm of A2] (c2) {$\tilde{p}(o_2)$};
\node[factor, above=1.5cm of B0] (u0) {$p(u_0)$};
\node[factor, above=1.5cm of B1] (u1) {$p(u_1)$};

% Kanten zwischen Knoten
\draw[edge] (d.east) -- (eq0.west) node[midway, above, varlabel] {$\textcolor{black}{q(s_0)}$} node[near end, below, varlabel, scale=0.7] {$\textcolor{orange}{\bm{\overrightarrow{\mu}(}s_0\bm{)}}$} node[near start, below, varlabel, scale=0.7] {$\textcolor{orange}{\bm{\overleftarrow{\mu}(}s_0\bm{)}}$};
\draw[edge] (eq0.south) -- (A0.north) node[midway, left, varlabel] {$\textcolor{black}{q(s^{\scriptscriptstyle II}_0)}$} node[pos=0.25, right=0.05, varlabel, scale=0.7] {$\textcolor{orange}{\bm{{\uparrow}{\mu}(}s^{\scriptscriptstyle II}_0\bm{)}}$} node[pos=0.75, right=0.05, varlabel, scale=0.7] {$\textcolor{orange}{\bm{{\downarrow}{\mu}(}s^{\scriptscriptstyle II}_0\bm{)}}$};
\draw[edge] (A0.south) -- (c0.north) node[midway, greycircle] {$\delta$} node[midway, left=0.3, varlabel] {$\textcolor{black}{q(o_0)}$} node[pos=0.2, right=0.05, varlabel, scale=0.7] {$\textcolor{orange}{\bm{{\uparrow}{\mu}(}o_0\bm{)}}$} node[pos=0.8, right=0.05, varlabel, scale=0.7] {$\textcolor{orange}{\bm{{\downarrow}{\mu}(}o_0\bm{)}}$};
\draw[edge] (eq0.east) -- (B0.west) node[midway, above, varlabel] {$\textcolor{black}{q(s^{\scriptscriptstyle I}_0)}$} node[near start, below, varlabel, scale=0.7] {$\textcolor{orange}{\bm{\overleftarrow{\mu}(}s^{\scriptscriptstyle I}_0\bm{)}}$} node[near end, below, varlabel, scale=0.7] {$\textcolor{orange}{\bm{\overrightarrow{\mu}(}s^{\scriptscriptstyle I}_0\bm{)}}$};
\draw[edge] (B0.north) -- (u0.south) node[midway, left, varlabel] {$\textcolor{black}{q(u_0)}$} node[pos=0.2, right=0.05, varlabel, scale=0.7] {$\textcolor{orange}{\bm{{\downarrow}{\mu}(}u_0\bm{)}}$} node[pos=0.8, right=0.05, varlabel, scale=0.7] {$\textcolor{orange}{\bm{{\uparrow}{\mu}(}u_0\bm{)}}$};
\draw[edge] (B0.east) -- (eq1.west) node[midway, above, varlabel] {$\textcolor{black}{q(s_1)}$} node[near end, below, varlabel, scale=0.7] {$\textcolor{orange}{\bm{\overrightarrow{\mu}(}s_1\bm{)}}$} node[near start, below, varlabel, scale=0.7] {$\textcolor{orange}{\bm{\overleftarrow{\mu}(}s_1\bm{)}}$};
\draw[edge] (eq1.south) -- (A1.north) node[midway, left, varlabel] {$\textcolor{black}{q(s^{\scriptscriptstyle II}_1)}$} node[pos=0.25, right=0.05, varlabel, scale=0.7] {$\textcolor{orange}{\bm{{\uparrow}{\mu}(}s^{\scriptscriptstyle II}_1\bm{)}}$} node[pos=0.75, right=0.05, varlabel, scale=0.7] {$\textcolor{orange}{\bm{{\downarrow}{\mu}(}s^{\scriptscriptstyle II}_1\bm{)}}$};
\draw[edge] (A1.south) -- (c1.north) node[midway, left, varlabel] {$\textcolor{black}{q(o_1)}$} node[pos=0.2, right=0.05, varlabel, scale=0.7] {$\textcolor{orange}{\bm{{\uparrow}{\mu}(}o_1\bm{)}}$} node[pos=0.8, right=0.05, varlabel, scale=0.7] {$\textcolor{orange}{\bm{{\downarrow}{\mu}(}o_1\bm{)}}$};
\draw[edge] (eq1.east) -- (B1.west) node[midway, above, varlabel] {$\textcolor{black}{q(s^{\scriptscriptstyle I}_1)}$} node[near start, below, varlabel, scale=0.7] {$\textcolor{orange}{\bm{\overleftarrow{\mu}(}s^{\scriptscriptstyle I}_1\bm{)}}$} node[near end, below, varlabel, scale=0.7] {$\textcolor{orange}{\bm{\overrightarrow{\mu}(}s^{\scriptscriptstyle I}_1\bm{)}}$};
\draw[edge] (B1.east) -| (A2.north) node[pos=0.25, above, varlabel] {$\textcolor{black}{q(s_2)}$} node[pos=0.89, right=0.05, varlabel, scale=0.7] {$\textcolor{orange}{\bm{{\downarrow}{\mu}(}s_2\bm{)}}$} node[pos=0.11, below, varlabel, scale=0.7] {$\textcolor{orange}{\bm{\overleftarrow{\mu}(}s_2\bm{)}}$};
\draw[edge] (A2.south) -- (c2.north) node[midway, left, varlabel] {$\textcolor{black}{q(o_2)}$} node[pos=0.2, right=0.05, varlabel, scale=0.7] {$\textcolor{orange}{\bm{{\uparrow}{\mu}(}o_2\bm{)}}$} node[pos=0.8, right=0.05, varlabel, scale=0.7] {$\textcolor{orange}{\bm{{\downarrow}{\mu}(}o_2\bm{)}}$};
\draw[edge] (B1.north) -- (u1.south) node[midway, left, varlabel] {$\textcolor{black}{q(u_1)}$} node[pos=0.2, right=0.05, varlabel, scale=0.7] {$\textcolor{orange}{\bm{{\downarrow}{\mu}(}u_1\bm{)}}$} node[pos=0.8, right=0.05, varlabel, scale=0.7] {$\textcolor{orange}{\bm{{\uparrow}{\mu}(}u_1\bm{)}}$};

\end{tikzpicture}}
                \end{tabular}%
                \llap{\smash{\raisebox{1.0cm}[0pt][0pt]{\sffamily\bfseries\small (d)\hspace*{-0.05cm}}}}%
                \hspace{0.1cm}% 3. Fester Abstand: RECHTS VOM FFG (Zum Box-Rand)
                \vspace{0.02cm} % 4. Fester Abstand: UNTEN
            \end{minipage}%
        }
        \phantomcaption
        \label{fig:26.2b}
    \end{subfigure}

    \vspace{0.2cm} % Abstand zwischen den Kästen

    % ================= TIMESTEP 3.1 =================
    \begin{subfigure}[b]{0.49\textwidth}
        \centering
        \fbox{%
            \begin{minipage}{0.97\linewidth}
                \vspace{0.15cm} % 1. Fester Abstand: OBEN
                % ==========================================
                % BEREICH 1: Das T-Maze (Flexibel zentriert)
                % ==========================================
                \hfill % Drückt das Maze von links weg
                \begin{tabular}{@{}c@{}}
                    \begin{tikzpicture}[scale=0.6]
                        \tikzset{
                            agent/.pic={
                                    \node[circle, draw=black, thick, fill=white, minimum size=0.5cm] at (0,0) {};
                                    \fill[black] (-0.08, 0.08) circle (0.03);
                                    \fill[black] (0.08, 0.08) circle (0.03);
                                    \draw[thick] (-0.1, -0.05) arc (180:360:0.1);
                                },
                            field label/.style={
                                    font=\tiny\sffamily, text=gray!70!black, anchor=south, scale=0.9, transform shape
                                },
                            field symbol/.style={
                                    font=\large\bfseries,
                                    text=gray!70,
                                    anchor=center,
                                    scale=1.0,
                                    transform shape
                                }
                        }

                        \draw (-0.5, 0) -- (-0.5, 2.0) -- (-1.5, 2.0) -- (-1.5, 3.0) -- (1.5, 3.0) -- (1.5, 2.0) -- (0.5, 2.0) -- (0.5, 0) -- cycle;
                        \draw (-0.5, 1.0) -- (0.5, 1.0);
                        \draw (-0.5, 2.0) -- (0.5, 2.0);
                        \draw (-0.5, 2.0) -- (-0.5, 3.0);
                        \draw (0.5, 2.0) -- (0.5, 3.0);

                        \pic[scale=0.8, transform shape] at (0, 1.5) {agent};
                        % Observation Arrow
                        \draw[
                            {Stealth[scale=1.0]}-,
                            line width=1.6pt,
                            color=gray!80!black,
                        ]
                        (0.51, 1.3) .. controls (1.0, 1.1) and (1.3, 2.0) .. (0.3, 1.6);
                        \node[font=\footnotesize, text=gray!80!black] at (1.2, 0.9) {$\textcolor{orange}{\hat {a_0}}$};
                        \node[field label] at (0, 0.97) {Start};
                        \node[field label] at (0, -0.03) {Cue};
                        \node[field label] at (-1.0, 1.97) {Left Arm};
                        \node[field label] at (1.0, 1.93) {Right Arm};
                        \node[field symbol] at (0, 0.5) {!};
                        \node[field symbol] at (-1.0, 2.5) {?};
                        \node[field symbol] at (1.0, 2.5) {?};
                    \end{tikzpicture}
                \end{tabular}%
                \hfill % Drückt das Maze von der gestrichelten Linie weg (zentriert es somit perfekt im linken Rest-Raum)
                % ==========================================
                % BEREICH 2: Die Trennlinie
                % ==========================================
                \begin{tabular}{@{}c@{}}
                    \tikz \draw[dashed, thin, gray] (0,-1.05) -- (0,1.05);
                \end{tabular}%
                \hspace{0.3cm}% 2. Fester Abstand: LINKS VOM FFG
                % ==========================================
                % BEREICH 3: Der FFG (gibt die Breite vor)
                % ==========================================
                \begin{tabular}{@{}c@{}}
                    % Der FFG wird jetzt in seiner absolut natürlichen Größe gerendert.
                    \resizebox{4.5cm}{!}{\begin{tikzpicture}[
    factor/.style={draw=black, thick, rectangle, minimum size=0.8cm, font=\normalsize},
    eqfactor/.style={draw=black, thick, rectangle, minimum size=0.5cm, font=\normalsize},
    edge/.style={thick},
    varlabel/.style={font=\normalsize},
    greycircle/.style={draw=black, thick, circle, fill=gray!30, minimum size=0.5cm, font=\normalsize}
]

% Knoten
\node[factor] (d) {$p(s_0)$};
\node[eqfactor, right=1.8cm of d] (eq0) {$=$};
\node[factor, below=1.5cm of eq0] (A0) {$p(o_0|s^{\scriptscriptstyle II}_0)$};
\node[factor, below=1.5cm of A0] (c0) {$\tilde{p}(o_0)$};
\node[factor, right=1.8cm of eq0] (B0) {$p(s_1|s^{\scriptscriptstyle I}_0,u_0)$};
\node[eqfactor, right=1.8cm of B0] (eq1) {$=$};
\node[factor, below=1.5cm of eq1] (A1) {$p(o_1|s^{\scriptscriptstyle II}_1)$};
\node[factor, below=1.5cm of A1] (c1) {$\tilde{p}(o_1)$};
\node[factor, right=1.8cm of eq1] (B1) {$p(s_2|s^{\scriptscriptstyle I}_1,u_1)$};
\node[factor] (A2) at ([xshift=2cm] B1.east |- A1) {$p(o_2|s_2)$};
\node[factor, below=1.5cm of A2] (c2) {$\tilde{p}(o_2)$};
\node[factor, above=1.5cm of B0] (u0) {$p(u_0)$};
\node[factor, above=1.5cm of B1] (u1) {$p(u_1)$};

% Kanten zwischen Knoten
\draw[edge] (d.east) -- (eq0.west) node[midway, above, varlabel] {$\textcolor{black}{q(s_0)}$};
\draw[edge] (eq0.south) -- (A0.north) node[midway, left, varlabel] {$\textcolor{black}{q(s^{\scriptscriptstyle II}_0)}$};
\draw[edge] (A0.south) -- (c0.north) node[midway, greycircle] {$\delta$} node[midway, left=0.3, varlabel] {$q(\textcolor{black}{o_0})$};
\draw[edge] (eq0.east) -- (B0.west) node[midway, above, varlabel] {$\textcolor{black}{q(s^{\scriptscriptstyle I}_0)}$};
\draw[edge] (B0.north) -- (u0.south) node[midway, greycircle] {$\delta$} node[midway, left=0.3, varlabel] {$q(\textcolor{orange}{u_0})$};
\draw[edge] (B0.east) -- (eq1.west) node[midway, above, varlabel] {$\textcolor{black}{q(s_1)}$};
\draw[edge] (eq1.south) -- (A1.north) node[midway, left, varlabel] {$\textcolor{black}{q(s^{\scriptscriptstyle II}_1)}$};
\draw[edge] (A1.south) -- (c1.north) node[midway, left, varlabel] {$\textcolor{black}{q(o_1)}$};
\draw[edge] (eq1.east) -- (B1.west) node[midway, above, varlabel] {$\textcolor{black}{q(s^{\scriptscriptstyle I}_1)}$};
\draw[edge] (B1.east) -| (A2.north) node[pos=0.25, above, varlabel] {$\textcolor{black}{q(s_2)}$};
\draw[edge] (A2.south) -- (c2.north) node[midway, left, varlabel] {$\textcolor{black}{q(o_2)}$};
\draw[edge] (B1.north) -- (u1.south) node[midway, left, varlabel] {$\textcolor{black}{q(u_1)}$};

\end{tikzpicture}}
                \end{tabular}%
                \llap{\smash{\raisebox{1.0cm}[0pt][0pt]{\sffamily\bfseries\small (e)\hspace*{-0.05cm}}}}%
                \hspace{0.1cm}% 3. Fester Abstand: RECHTS VOM FFG (Zum Box-Rand)
                \vspace{0.02cm} % 4. Fester Abstand: UNTEN
            \end{minipage}%
        }
        \phantomcaption
        \label{fig:26.3a}
    \end{subfigure}
    \hfill
    % ================= TIMESTEP 3.2 =================
    \begin{subfigure}[b]{0.49\textwidth}
        \centering
        \fbox{%
            \begin{minipage}{0.97\linewidth}
                \vspace{0.15cm} % 1. Fester Abstand: OBEN
                % ==========================================
                % BEREICH 1: Das T-Maze (Flexibel zentriert)
                % ==========================================
                \hfill % Drückt das Maze von links weg
                \begin{tabular}{@{}c@{}}
                    \begin{tikzpicture}[scale=0.6]
                        \tikzset{
                            agent/.pic={
                                    \node[circle, draw=black, thick, fill=white, minimum size=0.5cm] at (0,0) {};
                                    \fill[black] (-0.08, 0.08) circle (0.03);
                                    \fill[black] (0.08, 0.08) circle (0.03);
                                    \draw[thick] (-0.1, -0.05) arc (180:360:0.1);
                                },
                            field label/.style={
                                    font=\tiny\sffamily, text=gray!70!black, anchor=south, scale=0.9, transform shape
                                },
                            field symbol/.style={
                                    font=\large\bfseries,
                                    text=gray!70,
                                    anchor=center,
                                    scale=1.0,
                                    transform shape
                                }
                        }

                        \draw (-0.5, 0) -- (-0.5, 2.0) -- (-1.5, 2.0) -- (-1.5, 3.0) -- (1.5, 3.0) -- (1.5, 2.0) -- (0.5, 2.0) -- (0.5, 0) -- cycle;
                        \draw (-0.5, 1.0) -- (0.5, 1.0);
                        \draw (-0.5, 2.0) -- (0.5, 2.0);
                        \draw (-0.5, 2.0) -- (-0.5, 3.0);
                        \draw (0.5, 2.0) -- (0.5, 3.0);

                        \pic[scale=0.8, transform shape] at (0, 1.5) {agent};
                        % Observation Arrow
                        \draw[
                            {Stealth[scale=1.0]}-,
                            line width=1.6pt,
                            color=gray!80!black
                        ]
                        (0.51, 1.3) .. controls (1.0, 1.1) and (1.3, 2.0) .. (0.3, 1.6);
                        \node[font=\footnotesize, text=gray!80!black] at (1.2, 0.9) {$\textcolor{orange}{\hat {a_0}}$};
                        \node[field label] at (0, 0.97) {Start};
                        \node[field label] at (0, -0.03) {Cue};
                        \node[field label] at (-1.0, 1.97) {Left Arm};
                        \node[field label] at (1.0, 1.93) {Right Arm};
                        \node[field symbol] at (0, 0.5) {!};
                        \node[field symbol] at (-1.0, 2.5) {?};
                        \node[field symbol] at (1.0, 2.5) {?};
                    \end{tikzpicture}
                \end{tabular}%
                \hfill % Drückt das Maze von der gestrichelten Linie weg (zentriert es somit perfekt im linken Rest-Raum)
                % ==========================================
                % BEREICH 2: Die Trennlinie
                % ==========================================
                \begin{tabular}{@{}c@{}}
                    \tikz \draw[dashed, thin, gray] (0,-1.05) -- (0,1.05);
                \end{tabular}%
                \hspace{0.3cm}% 2. Fester Abstand: LINKS VOM FFG
                % ==========================================
                % BEREICH 3: Der FFG (gibt die Breite vor)
                % ==========================================
                \begin{tabular}{@{}c@{}}
                    % Der FFG wird jetzt in seiner absolut natürlichen Größe gerendert.
                    \resizebox{4.5cm}{!}{\begin{tikzpicture}[
    factor/.style={draw=black, thick, rectangle, minimum size=0.8cm, font=\normalsize},
    eqfactor/.style={draw=black, thick, rectangle, minimum size=0.5cm, font=\normalsize},
    edge/.style={thick},
    varlabel/.style={font=\normalsize},
    greycircle/.style={draw=black, thick, circle, fill=gray!30, minimum size=0.5cm, font=\normalsize}
]

% Knoten
\node[factor] (d) {$p(s_0)$};
\node[eqfactor, right=1.8cm of d] (eq0) {$=$};
\node[factor, below=1.5cm of eq0] (A0) {$p(o_0|s^{\scriptscriptstyle II}_0)$};
\node[factor, below=1.5cm of A0] (c0) {$\tilde{p}(o_0)$};
\node[factor, right=1.8cm of eq0] (B0) {$p(s_1|s^{\scriptscriptstyle I}_0,u_0)$};
\node[eqfactor, right=1.8cm of B0] (eq1) {$=$};
\node[factor, below=1.5cm of eq1] (A1) {$p(o_1|s^{\scriptscriptstyle II}_1)$};
\node[factor, below=1.5cm of A1] (c1) {$\tilde{p}(o_1)$};
\node[factor, right=1.8cm of eq1] (B1) {$p(s_2|s^{\scriptscriptstyle I}_1,u_1)$};
\node[factor] (A2) at ([xshift=2cm] B1.east |- A1) {$p(o_2|s_2)$};
\node[factor, below=1.5cm of A2] (c2) {$\tilde{p}(o_2)$};
\node[factor, above=1.5cm of B0] (u0) {$p(u_0)$};
\node[factor, above=1.5cm of B1] (u1) {$p(u_1)$};

% Kanten zwischen Knoten
\draw[edge] (d.east) -- (eq0.west) node[midway, above, varlabel] {$\textcolor{black}{q(s_0)}$};
\draw[edge] (eq0.south) -- (A0.north) node[midway, left, varlabel] {$\textcolor{black}{q(s^{\scriptscriptstyle II}_0)}$};
\draw[edge] (A0.south) -- (c0.north) node[midway, greycircle] {$\delta$} node[midway, left=0.3, varlabel] {$q(\textcolor{black}{o_0})$};
\draw[edge] (eq0.east) -- (B0.west) node[midway, above, varlabel] {$\textcolor{black}{q(s^{\scriptscriptstyle I}_0)}$};
\draw[edge] (B0.north) -- (u0.south) node[midway, greycircle] {$\delta$} node[midway, left=0.3, varlabel] {$q(\textcolor{orange}{u_0})$};
\draw[edge] (B0.east) -- (eq1.west) node[midway, above, varlabel] {$\textcolor{black}{q(s_1)}$};
\draw[edge] (eq1.south) -- (A1.north) node[midway, left, varlabel] {$\textcolor{black}{q(s^{\scriptscriptstyle II}_1)}$};
\draw[edge] (A1.south) -- (c1.north) node[midway, left, varlabel] {$\textcolor{black}{q(o_1)}$};
\draw[edge] (eq1.east) -- (B1.west) node[midway, above, varlabel] {$\textcolor{black}{q(s^{\scriptscriptstyle I}_1)}$};
\draw[edge] (B1.east) -| (A2.north) node[pos=0.25, above, varlabel] {$\textcolor{black}{q(s_2)}$};
\draw[edge] (A2.south) -- (c2.north) node[midway, left, varlabel] {$\textcolor{black}{q(o_2)}$};
\draw[edge] (B1.north) -- (u1.south) node[midway, left, varlabel] {$\textcolor{black}{q(u_1)}$};

\end{tikzpicture}}
                \end{tabular}%
                \llap{\smash{\raisebox{1.0cm}[0pt][0pt]{\sffamily\bfseries\small (f)\hspace*{-0.05cm}}}}%
                \hspace{0.1cm}% 3. Fester Abstand: RECHTS VOM FFG (Zum Box-Rand)
                \vspace{0.02cm} % 4. Fester Abstand: UNTEN
            \end{minipage}%
        }
        \phantomcaption
        \label{fig:26.3b}
    \end{subfigure}
\end{figure}

\begin{figure}[htbp]
    \ContinuedFloat

    \vspace{0.2cm}
    % ================= TIMESTEP 4.1 =================
    \begin{subfigure}[b]{0.49\textwidth}
        \centering
        \fbox{%
            \begin{minipage}{0.97\linewidth}
                \vspace{0.075cm} % 1. Fester Abstand: OBEN
                % ==========================================
                % BEREICH 1: Das T-Maze (angepasst für 0.48\textwidth)
                % ==========================================
                \hfill
                \begin{tabular}{@{}c@{}}
                    \begin{tikzpicture}[scale=0.6]
                        \tikzset{
                            agent/.pic={
                                    \node[circle, draw=black, thick, fill=white, minimum size=0.5cm] at (0,0) {};
                                    \fill[black] (-0.08, 0.08) circle (0.03);
                                    \fill[black] (0.08, 0.08) circle (0.03);
                                    \draw[thick] (-0.1, -0.05) arc (180:360:0.1);
                                },
                            field label/.style={
                                    font=\tiny\sffamily, text=gray!70!black, anchor=south, scale=0.9, transform shape
                                },
                            field symbol/.style={
                                    font=\large\bfseries,
                                    text=gray!70,
                                    anchor=center,
                                    scale=1.0,
                                    transform shape
                                }
                        }

                        \draw (-0.5, 0) -- (-0.5, 2.0) -- (-1.5, 2.0) -- (-1.5, 3.0) -- (1.5, 3.0) -- (1.5, 2.0) -- (0.5, 2.0) -- (0.5, 0) -- cycle;
                        \draw (-0.5, 1.0) -- (0.5, 1.0);
                        \draw (-0.5, 2.0) -- (0.5, 2.0);
                        \draw (-0.5, 2.0) -- (-0.5, 3.0);
                        \draw (0.5, 2.0) -- (0.5, 3.0);

                        \pic[scale=0.8, transform shape] at (-1, 2.5) {agent};
                        \node[field label] at (0, 0.97) {Start};
                        \node[field label] at (0, -0.03) {Cue};
                        \node[field label] at (-1.0, 1.97) {Left Arm};
                        \node[field label] at (1.0, 1.93) {Right Arm};
                        %\node[field symbol] at (-1.0, 2.5) {?};
                        \node[field symbol] at (1.0, 2.5) {?};
                        \node[field symbol] at (0.0, 0.5) {!};

                        % Observation Arrow
                        \draw[
                            -{Stealth[scale=1.0]},
                            line width=1.4pt,
                            color=gray!80!black
                        ]
                        (-0.49, 2.3) .. controls (0.0, 2.1) and (0.3, 3.0) .. (-0.7, 2.6);
                        \node[font=\footnotesize, text=gray!80!black] at (0.2, 2.3) {$\textcolor{orange}{\hat {o_1}}$};
                    \end{tikzpicture}
                \end{tabular}%
                \hfill
                % ==========================================
                % BEREICH 2: Die Trennlinie
                % ==========================================
                \begin{tabular}{@{}c@{}}
                    \tikz \draw[dashed, thin, gray] (0,-1.05) -- (0,1.05);
                \end{tabular}%
                \hspace{0.15cm}
                % ==========================================
                % BEREICH 3: Der FFG
                % ==========================================
                \begin{tabular}{@{}c@{}}
                    \resizebox{4.5cm}{!}{\begin{tikzpicture}[
    factor/.style={draw=black, thick, rectangle, minimum size=0.8cm, font=\normalsize},
    eqfactor/.style={draw=black, thick, rectangle, minimum size=0.5cm, font=\normalsize},
    edge/.style={thick},
    varlabel/.style={font=\normalsize},
    greycircle/.style={draw=black, thick, circle, fill=gray!30, minimum size=0.5cm, font=\normalsize}
]

% Knoten
\node[factor] (d) {$p(s_0)$};
\node[eqfactor, right=1.8cm of d] (eq0) {$=$};
\node[factor, below=1.5cm of eq0] (A0) {$p(o_0|s^{\scriptscriptstyle II}_0)$};
\node[factor, below=1.5cm of A0] (c0) {$\tilde{p}(o_0)$};
\node[factor, right=1.8cm of eq0] (B0) {$p(s_1|s^{\scriptscriptstyle I}_0,u_0)$};
\node[eqfactor, right=1.8cm of B0] (eq1) {$=$};
\node[factor, below=1.5cm of eq1] (A1) {$p(o_1|s^{\scriptscriptstyle II}_1)$};
\node[factor, below=1.5cm of A1] (c1) {$\tilde{p}(o_1)$};
\node[factor, right=1.8cm of eq1] (B1) {$p(s_2|s^{\scriptscriptstyle I}_1,u_1)$};
\node[factor] (A2) at ([xshift=2cm] B1.east |- A1) {$p(o_2|s_2)$};
\node[factor, below=1.5cm of A2] (c2) {$\tilde{p}(o_2)$};
\node[factor, above=1.5cm of B0] (u0) {$p(u_0)$};
\node[factor, above=1.5cm of B1] (u1) {$p(u_1)$};

% Kanten zwischen Knoten
\draw[edge] (d.east) -- (eq0.west) node[midway, above, varlabel] {$\textcolor{black}{q(s_0)}$};
\draw[edge] (eq0.south) -- (A0.north) node[midway, left, varlabel] {$\textcolor{black}{q(s^{\scriptscriptstyle II}_0)}$};
\draw[edge] (A0.south) -- (c0.north) node[midway, greycircle] {$\delta$} node[midway, left=0.3, varlabel] {$q(\textcolor{black}{o_0})$};
\draw[edge] (eq0.east) -- (B0.west) node[midway, above, varlabel] {$\textcolor{black}{q(s^{\scriptscriptstyle I}_0)}$};
\draw[edge] (B0.north) -- (u0.south) node[midway, greycircle] {$\delta$} node[midway, left=0.3, varlabel] {$q(\textcolor{black}{u_0})$};
\draw[edge] (B0.east) -- (eq1.west) node[midway, above, varlabel] {$\textcolor{black}{q(s_1)}$};
\draw[edge] (eq1.south) -- (A1.north) node[midway, left, varlabel] {$\textcolor{black}{q(s^{\scriptscriptstyle II}_1)}$};
\draw[edge] (A1.south) -- (c1.north) node[midway, greycircle] {$\delta$} node[midway, left=0.3, varlabel] {$q(\textcolor{orange}{o_1})$};
\draw[edge] (eq1.east) -- (B1.west) node[midway, above, varlabel] {$\textcolor{black}{q(s^{\scriptscriptstyle I}_1)}$};
\draw[edge] (B1.east) -| (A2.north) node[pos=0.25, above, varlabel] {$\textcolor{black}{q(s_2)}$};
\draw[edge] (A2.south) -- (c2.north) node[midway, left, varlabel] {$\textcolor{black}{q(o_2)}$};
\draw[edge] (B1.north) -- (u1.south) node[midway, left, varlabel] {$\textcolor{black}{q(u_1)}$};

\end{tikzpicture}}
                \end{tabular}%
                \llap{\smash{\raisebox{1.0cm}[0pt][0pt]{\sffamily\bfseries\small (g)\hspace*{-0.05cm}}}}%
                \hspace{0.05cm}%
                \vspace{0.01cm}
            \end{minipage}%
        }
        \phantomcaption
        \label{fig:26.4a}
    \end{subfigure}
    \hfill % Schiebt die beiden Subfigures an die äußeren Ränder
    % ================= TIMESTEP 4.2 =================
    \begin{subfigure}[b]{0.49\textwidth}
        \centering
        \fbox{%
            \begin{minipage}{0.97\linewidth}
                \vspace{0.075cm} % 1. Fester Abstand: OBEN
                % ==========================================
                % BEREICH 1: Das T-Maze (angepasst für 0.48\textwidth)
                % ==========================================
                \hfill
                \begin{tabular}{@{}c@{}}
                    \begin{tikzpicture}[scale=0.6]
                        \tikzset{
                            agent/.pic={
                                    \node[circle, draw=black, thick, fill=white, minimum size=0.5cm] at (0,0) {};
                                    \fill[black] (-0.08, 0.08) circle (0.03);
                                    \fill[black] (0.08, 0.08) circle (0.03);
                                    \draw[thick] (-0.1, -0.05) arc (180:360:0.1);
                                },
                            field label/.style={
                                    font=\tiny\sffamily, text=gray!70!black, anchor=south, scale=0.9, transform shape
                                },
                            field symbol/.style={
                                    font=\large\bfseries,
                                    text=gray!70,
                                    anchor=center,
                                    scale=1.0,
                                    transform shape
                                }
                        }

                        \draw (-0.5, 0) -- (-0.5, 2.0) -- (-1.5, 2.0) -- (-1.5, 3.0) -- (1.5, 3.0) -- (1.5, 2.0) -- (0.5, 2.0) -- (0.5, 0) -- cycle;
                        \draw (-0.5, 1.0) -- (0.5, 1.0);
                        \draw (-0.5, 2.0) -- (0.5, 2.0);
                        \draw (-0.5, 2.0) -- (-0.5, 3.0);
                        \draw (0.5, 2.0) -- (0.5, 3.0);

                        \pic[scale=0.8, transform shape] at (0, 0.5) {agent};
                        \node[field label] at (0, 0.97) {Start};
                        \node[field label] at (0, -0.03) {Cue};
                        \node[field label] at (-1.0, 1.97) {Left Arm};
                        \node[field label] at (1.0, 1.93) {Right Arm};
                        \node[field symbol] at (-1.0, 2.5) {?};
                        \node[field symbol] at (1.0, 2.5) {?};

                        % Observation Arrow
                        \draw[
                            -{Stealth[scale=1.0]},
                            line width=1.4pt,
                            color=gray!80!black
                        ]
                        (0.51, 0.3) .. controls (1.0, 0.1) and (1.3, 1.0) .. (0.3, 0.6);
                        \node[font=\footnotesize, text=gray!80!black] at (1.2, 0.3) {$\textcolor{orange}{\hat {o_1}}$};
                    \end{tikzpicture}
                \end{tabular}%
                \hfill
                % ==========================================
                % BEREICH 2: Die Trennlinie
                % ==========================================
                \begin{tabular}{@{}c@{}}
                    \tikz \draw[dashed, thin, gray] (0,-1.05) -- (0,1.05);
                \end{tabular}%
                \hspace{0.15cm}%
                % ==========================================
                % BEREICH 3: Der FFG
                % ==========================================
                \begin{tabular}{@{}c@{}}
                    \resizebox{4.5cm}{!}{\begin{tikzpicture}[
    factor/.style={draw=black, thick, rectangle, minimum size=0.8cm, font=\normalsize},
    eqfactor/.style={draw=black, thick, rectangle, minimum size=0.5cm, font=\normalsize},
    edge/.style={thick},
    varlabel/.style={font=\normalsize},
    greycircle/.style={draw=black, thick, circle, fill=gray!30, minimum size=0.5cm, font=\normalsize}
]

% Knoten
\node[factor] (d) {$p(s_0)$};
\node[eqfactor, right=1.8cm of d] (eq0) {$=$};
\node[factor, below=1.5cm of eq0] (A0) {$p(o_0|s^{\scriptscriptstyle II}_0)$};
\node[factor, below=1.5cm of A0] (c0) {$\tilde{p}(o_0)$};
\node[factor, right=1.8cm of eq0] (B0) {$p(s_1|s^{\scriptscriptstyle I}_0,u_0)$};
\node[eqfactor, right=1.8cm of B0] (eq1) {$=$};
\node[factor, below=1.5cm of eq1] (A1) {$p(o_1|s^{\scriptscriptstyle II}_1)$};
\node[factor, below=1.5cm of A1] (c1) {$\tilde{p}(o_1)$};
\node[factor, right=1.8cm of eq1] (B1) {$p(s_2|s^{\scriptscriptstyle I}_1,u_1)$};
\node[factor] (A2) at ([xshift=2cm] B1.east |- A1) {$p(o_2|s_2)$};
\node[factor, below=1.5cm of A2] (c2) {$\tilde{p}(o_2)$};
\node[factor, above=1.5cm of B0] (u0) {$p(u_0)$};
\node[factor, above=1.5cm of B1] (u1) {$p(u_1)$};

% Kanten zwischen Knoten
\draw[edge] (d.east) -- (eq0.west) node[midway, above, varlabel] {$\textcolor{black}{q(s_0)}$};
\draw[edge] (eq0.south) -- (A0.north) node[midway, left, varlabel] {$\textcolor{black}{q(s^{\scriptscriptstyle II}_0)}$};
\draw[edge] (A0.south) -- (c0.north) node[midway, greycircle] {$\delta$} node[midway, left=0.3, varlabel] {$q(\textcolor{black}{o_0})$};
\draw[edge] (eq0.east) -- (B0.west) node[midway, above, varlabel] {$\textcolor{black}{q(s^{\scriptscriptstyle I}_0)}$};
\draw[edge] (B0.north) -- (u0.south) node[midway, greycircle] {$\delta$} node[midway, left=0.3, varlabel] {$q(\textcolor{black}{u_0})$};
\draw[edge] (B0.east) -- (eq1.west) node[midway, above, varlabel] {$\textcolor{black}{q(s_1)}$};
\draw[edge] (eq1.south) -- (A1.north) node[midway, left, varlabel] {$\textcolor{black}{q(s^{\scriptscriptstyle II}_1)}$};
\draw[edge] (A1.south) -- (c1.north) node[midway, greycircle] {$\delta$} node[midway, left=0.3, varlabel] {$q(\textcolor{orange}{o_1})$};
\draw[edge] (eq1.east) -- (B1.west) node[midway, above, varlabel] {$\textcolor{black}{q(s^{\scriptscriptstyle I}_1)}$};
\draw[edge] (B1.east) -| (A2.north) node[pos=0.25, above, varlabel] {$\textcolor{black}{q(s_2)}$};
\draw[edge] (A2.south) -- (c2.north) node[midway, left, varlabel] {$\textcolor{black}{q(o_2)}$};
\draw[edge] (B1.north) -- (u1.south) node[midway, left, varlabel] {$\textcolor{black}{q(u_1)}$};

\end{tikzpicture}}
                \end{tabular}%
                \llap{\smash{\raisebox{1.0cm}[0pt][0pt]{\sffamily\bfseries\small (h)\hspace*{-0.05cm}}}}%
                \hspace{0.05cm}%
                \vspace{0.01cm}
            \end{minipage}%
        }
        \phantomcaption
        \label{fig:26.4b}
    \end{subfigure}

    \vspace{0.2cm} % Abstand zwischen den Kästen

    \hfill % Schiebt die beiden Subfigures an die äußeren Ränder
    % ================= TIMESTEP 5.2 =================
    \begin{subfigure}[b]{0.49\textwidth}
        \centering
        \fbox{%
            \begin{minipage}{0.97\linewidth}
                \vspace{0.15cm} % 1. Fester Abstand: OBEN
                % ==========================================
                % BEREICH 1: Das T-Maze (Flexibel zentriert)
                % ==========================================
                \hfill % Drückt das Maze von links weg
                \begin{tabular}{@{}c@{}}
                    \begin{tikzpicture}[scale=0.6]
                        \tikzset{
                            agent/.pic={
                                    \node[circle, draw=black, thick, fill=white, minimum size=0.5cm] at (0,0) {};
                                    \fill[black] (-0.08, 0.08) circle (0.03);
                                    \fill[black] (0.08, 0.08) circle (0.03);
                                    \draw[thick] (-0.1, -0.05) arc (180:360:0.1);
                                },
                            field label/.style={
                                    font=\tiny\sffamily, text=gray!70!black, anchor=south, scale=0.9, transform shape
                                },
                            field symbol/.style={
                                    font=\large\bfseries,
                                    text=gray!70,
                                    anchor=center,
                                    scale=1.0,
                                    transform shape
                                }
                        }

                        \draw (-0.5, 0) -- (-0.5, 2.0) -- (-1.5, 2.0) -- (-1.5, 3.0) -- (1.5, 3.0) -- (1.5, 2.0) -- (0.5, 2.0) -- (0.5, 0) -- cycle;
                        \draw (-0.5, 1.0) -- (0.5, 1.0);
                        \draw (-0.5, 2.0) -- (0.5, 2.0);
                        \draw (-0.5, 2.0) -- (-0.5, 3.0);
                        \draw (0.5, 2.0) -- (0.5, 3.0);

                        \pic[scale=0.8, transform shape] at (0, 0.5) {agent};
                        \node[field label] at (0, 0.97) {Start};
                        \node[field label] at (0, -0.03) {Cue};
                        \node[field label] at (-1.0, 1.97) {Left Arm};
                        \node[field label] at (1.0, 1.93) {Right Arm};
                        %\node[field symbol] at (0, 0.5) {!};
                        \node[field symbol] at (-1.0, 2.5) {?};
                        \node[field symbol] at (1.0, 2.5) {?};

                        % --- NEU: Gedankenblase für Variational Free Energy Minimierung ---
                        % Kleine Bläschen, die vom Kopf (ca. 0.3, 1.7) nach rechts oben steigen
                        \draw[thin, gray!60, fill=white] (0.35, 0.55) circle (0.04);
                        \draw[thin, gray!60, fill=white] (0.52, 0.65) circle (0.07);

                        % Die Haupt-Gedankenblase mit dem Minimierungs-Symbol / Free Energy F
                        \node[
                            ellipse,
                            draw=gray!40,
                            thick,
                            fill=white,
                            minimum width=1.0cm,
                            minimum height=0.55cm,
                            inner sep=2pt,
                            font=\scriptsize\sffamily,
                            text=black!100!black,
                            transform shape
                        ] at (1.2, 0.7) {$\min F$};
                    \end{tikzpicture}
                \end{tabular}%
                \hfill % Drückt das Maze von der gestrichelten Linie weg (zentriert es somit perfekt im linken Rest-Raum)
                % ==========================================
                % BEREICH 2: Die Trennlinie
                % ==========================================
                \begin{tabular}{@{}c@{}}
                    \tikz \draw[dashed, thin, gray] (0,-1.05) -- (0,1.05);
                \end{tabular}%
                \hspace{0.3cm}% 2. Fester Abstand: LINKS VOM FFG
                % ==========================================
                % BEREICH 3: Der FFG (gibt die Breite vor)
                % ==========================================
                \begin{tabular}{@{}c@{}}
                    % Der FFG wird jetzt in seiner absolut natürlichen Größe gerendert.
                    \resizebox{4.4cm}{!}{\begin{tikzpicture}[
    factor/.style={draw=black, thick, rectangle, minimum size=0.8cm, font=\normalsize},
    eqfactor/.style={draw=black, thick, rectangle, minimum size=0.5cm, font=\normalsize},
    edge/.style={thick},
    varlabel/.style={font=\normalsize},
    greycircle/.style={draw=black, thick, circle, fill=gray!30, minimum size=0.5cm, font=\normalsize}
]

% Knoten
\node[factor] (d) {$p(s_0)$};
\node[eqfactor, right=1.8cm of d] (eq0) {$=$};
\node[factor, below=1.5cm of eq0] (A0) {$p(o_0|s^{\scriptscriptstyle II}_0)$};
\node[factor, below=1.5cm of A0] (c0) {$\tilde{p}(o_0)$};
\node[factor, right=1.8cm of eq0] (B0) {$p(s_1|s^{\scriptscriptstyle I}_0,u_0)$};
\node[eqfactor, right=1.8cm of B0] (eq1) {$=$};
\node[factor, below=1.5cm of eq1] (A1) {$p(o_1|s^{\scriptscriptstyle II}_1)$};
\node[factor, below=1.5cm of A1] (c1) {$\tilde{p}(o_1)$};
\node[factor, right=1.8cm of eq1] (B1) {$p(s_2|s^{\scriptscriptstyle I}_1,u_1)$};
\node[factor] (A2) at ([xshift=2cm] B1.east |- A1) {$p(o_2|s_2)$};
\node[factor, below=1.5cm of A2] (c2) {$\tilde{p}(o_2)$};
\node[factor, above=1.5cm of B0] (u0) {$p(u_0)$};
\node[factor, above=1.5cm of B1] (u1) {$p(u_1)$};

% Kanten zwischen Knoten
\draw[edge] (d.east) -- (eq0.west) node[midway, above, varlabel] {$\textcolor{black}{q(s_0)}$} node[near end, below, varlabel, scale=0.7] {$\textcolor{orange}{\bm{\overrightarrow{\mu}(}s_0\bm{)}}$} node[near start, below, varlabel, scale=0.7] {$\textcolor{orange}{\bm{\overleftarrow{\mu}(}s_0\bm{)}}$};
\draw[edge] (eq0.south) -- (A0.north) node[midway, left, varlabel] {$\textcolor{black}{q(s^{\scriptscriptstyle II}_0)}$} node[pos=0.25, right=0.05, varlabel, scale=0.7] {$\textcolor{orange}{\bm{{\uparrow}{\mu}(}s^{\scriptscriptstyle II}_0\bm{)}}$} node[pos=0.75, right=0.05, varlabel, scale=0.7] {$\textcolor{orange}{\bm{{\downarrow}{\mu}(}s^{\scriptscriptstyle II}_0\bm{)}}$};
\draw[edge] (A0.south) -- (c0.north) node[midway, greycircle] {$\delta$} node[midway, left=0.3, varlabel] {$\textcolor{black}{q(o_0)}$} node[pos=0.2, right=0.05, varlabel, scale=0.7] {$\textcolor{orange}{\bm{{\uparrow}{\mu}(}o_0\bm{)}}$} node[pos=0.8, right=0.05, varlabel, scale=0.7] {$\textcolor{orange}{\bm{{\downarrow}{\mu}(}o_0\bm{)}}$};
\draw[edge] (eq0.east) -- (B0.west) node[midway, above, varlabel] {$\textcolor{black}{q(s^{\scriptscriptstyle I}_0)}$} node[near start, below, varlabel, scale=0.7] {$\textcolor{orange}{\bm{\overleftarrow{\mu}(}s^{\scriptscriptstyle I}_0\bm{)}}$} node[near end, below, varlabel, scale=0.7] {$\textcolor{orange}{\bm{\overrightarrow{\mu}(}s^{\scriptscriptstyle I}_0\bm{)}}$};
\draw[edge] (B0.north) -- (u0.south) node[midway, greycircle] {$\delta$} node[midway, left=0.3, varlabel] {$\textcolor{black}{q(u_0)}$} node[pos=0.2, right=0.05, varlabel, scale=0.7] {$\textcolor{orange}{\bm{{\downarrow}{\mu}(}u_0\bm{)}}$} node[pos=0.8, right=0.05, varlabel, scale=0.7] {$\textcolor{orange}{\bm{{\uparrow}{\mu}(}u_0\bm{)}}$};
\draw[edge] (B0.east) -- (eq1.west) node[midway, above, varlabel] {$\textcolor{black}{q(s_1)}$} node[near end, below, varlabel, scale=0.7] {$\textcolor{orange}{\bm{\overrightarrow{\mu}(}s_1\bm{)}}$} node[near start, below, varlabel, scale=0.7] {$\textcolor{orange}{\bm{\overleftarrow{\mu}(}s_1\bm{)}}$};
\draw[edge] (eq1.south) -- (A1.north) node[midway, left, varlabel] {$\textcolor{black}{q(s^{\scriptscriptstyle II}_1)}$} node[pos=0.25, right=0.05, varlabel, scale=0.7] {$\textcolor{orange}{\bm{{\uparrow}{\mu}(}s^{\scriptscriptstyle II}_1\bm{)}}$} node[pos=0.75, right=0.05, varlabel, scale=0.7] {$\textcolor{orange}{\bm{{\downarrow}{\mu}(}s^{\scriptscriptstyle II}_1\bm{)}}$};
\draw[edge] (A1.south) -- (c1.north) node[midway, greycircle] {$\delta$} node[midway, left=0.3, varlabel] {$\textcolor{black}{q(o_1)}$} node[pos=0.2, right=0.05, varlabel, scale=0.7] {$\textcolor{orange}{\bm{{\uparrow}{\mu}(}o_1\bm{)}}$} node[pos=0.8, right=0.05, varlabel, scale=0.7] {$\textcolor{orange}{\bm{{\downarrow}{\mu}(}o_1\bm{)}}$};
\draw[edge] (eq1.east) -- (B1.west) node[midway, above, varlabel] {$\textcolor{black}{q(s^{\scriptscriptstyle I}_1)}$} node[near start, below, varlabel, scale=0.7] {$\textcolor{orange}{\bm{\overleftarrow{\mu}(}s^{\scriptscriptstyle I}_1\bm{)}}$} node[near end, below, varlabel, scale=0.7] {$\textcolor{orange}{\bm{\overrightarrow{\mu}(}s^{\scriptscriptstyle I}_1\bm{)}}$};
\draw[edge] (B1.east) -| (A2.north) node[pos=0.25, above, varlabel] {$\textcolor{black}{q(s_2)}$} node[pos=0.89, right=0.05, varlabel, scale=0.7] {$\textcolor{orange}{\bm{{\downarrow}{\mu}(}s_2\bm{)}}$} node[pos=0.11, below, varlabel, scale=0.7] {$\textcolor{orange}{\bm{\overleftarrow{\mu}(}s_2\bm{)}}$};
\draw[edge] (A2.south) -- (c2.north) node[midway, left, varlabel] {$\textcolor{black}{q(o_2)}$} node[pos=0.2, right=0.05, varlabel, scale=0.7] {$\textcolor{orange}{\bm{{\uparrow}{\mu}(}o_2\bm{)}}$} node[pos=0.8, right=0.05, varlabel, scale=0.7] {$\textcolor{orange}{\bm{{\downarrow}{\mu}(}o_2\bm{)}}$};
\draw[edge] (B1.north) -- (u1.south) node[midway, left, varlabel] {$\textcolor{black}{q(u_1)}$} node[pos=0.2, right=0.05, varlabel, scale=0.7] {$\textcolor{orange}{\bm{{\downarrow}{\mu}(}u_1\bm{)}}$} node[pos=0.8, right=0.05, varlabel, scale=0.7] {$\textcolor{orange}{\bm{{\uparrow}{\mu}(}u_1\bm{)}}$};

\end{tikzpicture}}
                \end{tabular}%
                \llap{\smash{\raisebox{1.0cm}[0pt][0pt]{\sffamily\bfseries\small (i)\hspace*{-0.05cm}}}}%
                \hspace{0.1cm}% 3. Fester Abstand: RECHTS VOM FFG (Zum Box-Rand)
                \vspace{0.02cm} % 4. Fester Abstand: UNTEN
            \end{minipage}%
        }
        \phantomcaption
        \label{fig:26.5b}
    \end{subfigure}

    \vspace{0.2cm}

    \hfill
    % ================= TIMESTEP 6.2 =================
    \begin{subfigure}[b]{0.49\textwidth}
        \centering
        \fbox{%
            \begin{minipage}{0.97\linewidth}
                \vspace{0.15cm} % 1. Fester Abstand: OBEN
                % ==========================================
                % BEREICH 1: Das T-Maze (Flexibel zentriert)
                % ==========================================
                \hfill % Drückt das Maze von links weg
                \begin{tabular}{@{}c@{}}
                    \begin{tikzpicture}[scale=0.6]
                        \tikzset{
                            agent/.pic={
                                    \node[circle, draw=black, thick, fill=white, minimum size=0.5cm] at (0,0) {};
                                    \fill[black] (-0.08, 0.08) circle (0.03);
                                    \fill[black] (0.08, 0.08) circle (0.03);
                                    \draw[thick] (-0.1, -0.05) arc (180:360:0.1);
                                },
                            field label/.style={
                                    font=\tiny\sffamily, text=gray!70!black, anchor=south, scale=0.9, transform shape
                                },
                            field symbol/.style={
                                    font=\large\bfseries,
                                    text=gray!70,
                                    anchor=center,
                                    scale=1.0,
                                    transform shape
                                }
                        }

                        \draw (-0.5, 0) -- (-0.5, 2.0) -- (-1.5, 2.0) -- (-1.5, 3.0) -- (1.5, 3.0) -- (1.5, 2.0) -- (0.5, 2.0) -- (0.5, 0) -- cycle;
                        \draw (-0.5, 1.0) -- (0.5, 1.0);
                        \draw (-0.5, 2.0) -- (0.5, 2.0);
                        \draw (-0.5, 2.0) -- (-0.5, 3.0);
                        \draw (0.5, 2.0) -- (0.5, 3.0);

                        \pic[scale=0.8, transform shape] at (0, 0.5) {agent};
                        % Observation Arrow
                        \draw[
                            {Stealth[scale=1.0]}-,
                            line width=1.6pt,
                            color=gray!80!black
                        ]
                        (0.51, 0.3) .. controls (1.0, 0.1) and (1.3, 1.0) .. (0.3, 0.6);
                        \node[font=\footnotesize, text=gray!80!black] at (1.2, 0.9) {$\textcolor{orange}{\hat {a_1}}$};
                        \node[field label] at (0, 0.97) {Start};
                        \node[field label] at (0, -0.03) {Cue};
                        \node[field label] at (-1.0, 1.97) {Left Arm};
                        \node[field label] at (1.0, 1.93) {Right Arm};
                        %\node[field symbol] at (0, 0.5) {!};
                        \node[field symbol] at (-1.0, 2.5) {?};
                        \node[field symbol] at (1.0, 2.5) {?};
                    \end{tikzpicture}
                \end{tabular}%
                \hfill % Drückt das Maze von der gestrichelten Linie weg (zentriert es somit perfekt im linken Rest-Raum)
                % ==========================================
                % BEREICH 2: Die Trennlinie
                % ==========================================
                \begin{tabular}{@{}c@{}}
                    \tikz \draw[dashed, thin, gray] (0,-1.05) -- (0,1.05);
                \end{tabular}%
                \hspace{0.3cm}% 2. Fester Abstand: LINKS VOM FFG
                % ==========================================
                % BEREICH 3: Der FFG (gibt die Breite vor)
                % ==========================================
                \begin{tabular}{@{}c@{}}
                    % Der FFG wird jetzt in seiner absolut natürlichen Größe gerendert.
                    \resizebox{4.5cm}{!}{\begin{tikzpicture}[
    factor/.style={draw=black, thick, rectangle, minimum size=0.8cm, font=\normalsize},
    eqfactor/.style={draw=black, thick, rectangle, minimum size=0.5cm, font=\normalsize},
    edge/.style={thick},
    varlabel/.style={font=\normalsize},
    greycircle/.style={draw=black, thick, circle, fill=gray!30, minimum size=0.5cm, font=\normalsize}
]

% Knoten
\node[factor] (d) {$p(s_0)$};
\node[eqfactor, right=1.8cm of d] (eq0) {$=$};
\node[factor, below=1.5cm of eq0] (A0) {$p(o_0|s^{\scriptscriptstyle II}_0)$};
\node[factor, below=1.5cm of A0] (c0) {$\tilde{p}(o_0)$};
\node[factor, right=1.8cm of eq0] (B0) {$p(s_1|s^{\scriptscriptstyle I}_0,u_0)$};
\node[eqfactor, right=1.8cm of B0] (eq1) {$=$};
\node[factor, below=1.5cm of eq1] (A1) {$p(o_1|s^{\scriptscriptstyle II}_1)$};
\node[factor, below=1.5cm of A1] (c1) {$\tilde{p}(o_1)$};
\node[factor, right=1.8cm of eq1] (B1) {$p(s_2|s^{\scriptscriptstyle I}_1,u_1)$};
\node[factor] (A2) at ([xshift=2cm] B1.east |- A1) {$p(o_2|s_2)$};
\node[factor, below=1.5cm of A2] (c2) {$\tilde{p}(o_2)$};
\node[factor, above=1.5cm of B0] (u0) {$p(u_0)$};
\node[factor, above=1.5cm of B1] (u1) {$p(u_1)$};

% Kanten zwischen Knoten
\draw[edge] (d.east) -- (eq0.west) node[midway, above, varlabel] {$\textcolor{black}{q(s_0)}$};
\draw[edge] (eq0.south) -- (A0.north) node[midway, left, varlabel] {$\textcolor{black}{q(s^{\scriptscriptstyle II}_0)}$};
\draw[edge] (A0.south) -- (c0.north) node[midway, greycircle] {$\delta$} node[midway, left=0.3, varlabel] {$q(\textcolor{black}{o_0})$};
\draw[edge] (eq0.east) -- (B0.west) node[midway, above, varlabel] {$\textcolor{black}{q(s^{\scriptscriptstyle I}_0)}$};
\draw[edge] (B0.north) -- (u0.south) node[midway, greycircle] {$\delta$} node[midway, left=0.3, varlabel] {$q(\textcolor{black}{u_0})$};
\draw[edge] (B0.east) -- (eq1.west) node[midway, above, varlabel] {$\textcolor{black}{q(s_1)}$};
\draw[edge] (eq1.south) -- (A1.north) node[midway, left, varlabel] {$\textcolor{black}{q(s^{\scriptscriptstyle II}_1)}$};
\draw[edge] (A1.south) -- (c1.north) node[midway, greycircle] {$\delta$} node[midway, left=0.3, varlabel] {$q(\textcolor{black}{o_1})$};
\draw[edge] (eq1.east) -- (B1.west) node[midway, above, varlabel] {$\textcolor{black}{q(s^{\scriptscriptstyle I}_1)}$};
\draw[edge] (B1.east) -| (A2.north) node[pos=0.25, above, varlabel] {$\textcolor{black}{q(s_2)}$};
\draw[edge] (A2.south) -- (c2.north) node[midway, left, varlabel] {$\textcolor{black}{q(o_2)}$};
\draw[edge] (B1.north) -- (u1.south) node[midway, greycircle] {$\delta$} node[midway, left=0.3, varlabel] {$q(\textcolor{orange}{u_1})$};

\end{tikzpicture}}
                \end{tabular}%
                \llap{\smash{\raisebox{1.0cm}[0pt][0pt]{\sffamily\bfseries\small (j)\hspace*{-0.05cm}}}}%
                \hspace{0.1cm}% 3. Fester Abstand: RECHTS VOM FFG (Zum Box-Rand)
                \vspace{0.02cm} % 4. Fester Abstand: UNTEN
            \end{minipage}%
        }
        \phantomcaption
        \label{fig:26.6b}
    \end{subfigure}

    \vspace{0.2cm} % Abstand zwischen den Kästen

    \hfill
    % ================= TIMESTEP 7.2 =================
    \begin{subfigure}[b]{0.49\textwidth}
        \centering
        \fbox{%
            \begin{minipage}{0.97\linewidth}
                \vspace{0.075cm} % 1. Fester Abstand: OBEN
                % ==========================================
                % BEREICH 1: Das T-Maze (angepasst für 0.48\textwidth)
                % ==========================================
                \hfill
                \begin{tabular}{@{}c@{}}
                    \begin{tikzpicture}[scale=0.6]
                        \tikzset{
                            agent/.pic={
                                    \node[circle, draw=black, thick, fill=white, minimum size=0.5cm] at (0,0) {};
                                    \fill[black] (-0.08, 0.08) circle (0.03);
                                    \fill[black] (0.08, 0.08) circle (0.03);
                                    \draw[thick] (-0.1, -0.05) arc (180:360:0.1);
                                },
                            field label/.style={
                                    font=\tiny\sffamily, text=gray!70!black, anchor=south, scale=0.9, transform shape
                                },
                            field symbol/.style={
                                    font=\large\bfseries,
                                    text=gray!70,
                                    anchor=center,
                                    scale=1.0,
                                    transform shape
                                }
                        }


                        \draw (-0.5, 0) -- (-0.5, 2.0) -- (-1.5, 2.0) -- (-1.5, 3.0) -- (1.5, 3.0) -- (1.5, 2.0) -- (0.5, 2.0) -- (0.5, 0) -- cycle;
                        \draw (-0.5, 1.0) -- (0.5, 1.0);
                        \draw (-0.5, 2.0) -- (0.5, 2.0);
                        \draw (-0.5, 2.0) -- (-0.5, 3.0);
                        \draw (0.5, 2.0) -- (0.5, 3.0);

                        \pic[scale=0.8, transform shape] at (-1, 2.5) {agent};
                        \node[field label] at (0, 0.97) {Start};
                        \node[field label] at (0, -0.03) {Cue};
                        \node[field label] at (-1.0, 1.97) {Left Arm};
                        \node[field label] at (1.0, 1.93) {Right Arm};
                        %\node[field symbol] at (-1.0, 2.5) {?};
                        \node[field symbol] at (1.0, 2.5) {?};
                        \node[field symbol] at (0.0, 0.5) {!};

                        % Observation Arrow
                        \draw[
                            -{Stealth[scale=1.0]},
                            line width=1.4pt,
                            color=gray!80!black
                        ]
                        (-0.49, 2.3) .. controls (0.0, 2.1) and (0.3, 3.0) .. (-0.7, 2.6);
                        \node[font=\footnotesize, text=gray!80!black] at (0.2, 2.3) {$\textcolor{orange}{\hat {o_2}}$};
                    \end{tikzpicture}
                \end{tabular}%
                \hfill
                % ==========================================
                % BEREICH 2: Die Trennlinie
                % ==========================================
                \begin{tabular}{@{}c@{}}
                    \tikz \draw[dashed, thin, gray] (0,-1.05) -- (0,1.05);
                \end{tabular}%
                \hspace{0.15cm}
                % ==========================================
                % BEREICH 3: Der FFG
                % ==========================================
                \begin{tabular}{@{}c@{}}
                    \resizebox{4.5cm}{!}{\begin{tikzpicture}[
    factor/.style={draw=black, thick, rectangle, minimum size=0.8cm, font=\normalsize},
    eqfactor/.style={draw=black, thick, rectangle, minimum size=0.5cm, font=\normalsize},
    edge/.style={thick},
    varlabel/.style={font=\normalsize},
    greycircle/.style={draw=black, thick, circle, fill=gray!30, minimum size=0.5cm, font=\normalsize}
]

% Knoten
\node[factor] (d) {$p(s_0)$};
\node[eqfactor, right=1.8cm of d] (eq0) {$=$};
\node[factor, below=1.5cm of eq0] (A0) {$p(o_0|s^{\scriptscriptstyle II}_0)$};
\node[factor, below=1.5cm of A0] (c0) {$\tilde{p}(o_0)$};
\node[factor, right=1.8cm of eq0] (B0) {$p(s_1|s^{\scriptscriptstyle I}_0,u_0)$};
\node[eqfactor, right=1.8cm of B0] (eq1) {$=$};
\node[factor, below=1.5cm of eq1] (A1) {$p(o_1|s^{\scriptscriptstyle II}_1)$};
\node[factor, below=1.5cm of A1] (c1) {$\tilde{p}(o_1)$};
\node[factor, right=1.8cm of eq1] (B1) {$p(s_2|s^{\scriptscriptstyle I}_1,u_1)$};
\node[factor] (A2) at ([xshift=2cm] B1.east |- A1) {$p(o_2|s_2)$};
\node[factor, below=1.5cm of A2] (c2) {$\tilde{p}(o_2)$};
\node[factor, above=1.5cm of B0] (u0) {$p(u_0)$};
\node[factor, above=1.5cm of B1] (u1) {$p(u_1)$};

% Kanten zwischen Knoten
\draw[edge] (d.east) -- (eq0.west) node[midway, above, varlabel] {$\textcolor{black}{q(s_0)}$};
\draw[edge] (eq0.south) -- (A0.north) node[midway, left, varlabel] {$\textcolor{black}{q(s^{\scriptscriptstyle II}_0)}$};
\draw[edge] (A0.south) -- (c0.north) node[midway, greycircle] {$\delta$} node[midway, left=0.3, varlabel] {$q(\textcolor{black}{o_0})$};
\draw[edge] (eq0.east) -- (B0.west) node[midway, above, varlabel] {$\textcolor{black}{q(s^{\scriptscriptstyle I}_0)}$};
\draw[edge] (B0.north) -- (u0.south) node[midway, greycircle] {$\delta$} node[midway, left=0.3, varlabel] {$q(\textcolor{black}{u_0})$};
\draw[edge] (B0.east) -- (eq1.west) node[midway, above, varlabel] {$\textcolor{black}{q(s_1)}$};
\draw[edge] (eq1.south) -- (A1.north) node[midway, left, varlabel] {$\textcolor{black}{q(s^{\scriptscriptstyle II}_1)}$};
\draw[edge] (A1.south) -- (c1.north) node[midway, greycircle] {$\delta$} node[midway, left=0.3, varlabel] {$q(\textcolor{black}{o_1})$};
\draw[edge] (eq1.east) -- (B1.west) node[midway, above, varlabel] {$\textcolor{black}{q(s^{\scriptscriptstyle I}_1)}$};
\draw[edge] (B1.east) -| (A2.north) node[pos=0.25, above, varlabel] {$\textcolor{black}{q(s_2)}$};
\draw[edge] (A2.south) -- (c2.north) node[midway, greycircle] {$\delta$} node[midway, left=0.3, varlabel] {$q(\textcolor{orange}{o_2})$};
\draw[edge] (B1.north) -- (u1.south) node[midway, greycircle] {$\delta$} node[midway, left=0.3, varlabel] {$q(\textcolor{black}{u_1})$};

\end{tikzpicture}}
                \end{tabular}%
                \llap{\smash{\raisebox{1.0cm}[0pt][0pt]{\sffamily\bfseries\small (k)\hspace*{-0.05cm}}}}%
                \hspace{0.05cm}%
                \vspace{0.01cm}
            \end{minipage}%
        }
        \phantomcaption
        \label{fig:26.7b}
    \end{subfigure}

    \caption{The figure shows the two behavioral paths the agent can take during a trial (left column and right column). In the left column the agent decides against moving to the cue and instead directly moves to one of the two arms. In the right column the agent first moves to the cue and then decides for one of the two arms. Both paths start similarly. Specifically, the agent always begins at the start field in the T-Maze and receives an observation (\cref{fig:26.1a,fig:26.1b}). Next, the agent infers its belief distributions (\cref{fig:26.2a,fig:26.2b}). Based on the inferred action-specifc belief distribution the agent executes an action (\cref{fig:26.3a,fig:26.3b}) where this action either moves it to the cue (\cref{fig:26.4b}) or directly to one of the arms (\cref{fig:26.4a}). Directly moving to one of the arms ends the the trial (left column). Next, the agent in the right column again infers its belief distributions also taking into account the information received by the cue (\cref{fig:26.5b}). Following this second inference, the agent can again act again ((\cref{fig:26.6b})) to finally move to one of the two arm ((\cref{fig:26.7b})). Note that in principle there are more behavioral paths than the two shown here but these other paths do never lead to the reward, e.g., always taking actions that keeps the agent at the start field. The two paths shown here are the only paths that have non-zero probabilities of leading to the reward, i.e., also with the two paths the agent might not receive the reward when moving to the wrong arm.}
    \label{fig:26}
\end{figure}